\documentclass[11pt,letterpaper]{report}
\newcommand{\Empresa}{Synaptix}
\newcommand{\EmpresaArchivo}{SYNAPTIX}
\newcommand{\RazonSocial}{Synaptix Ingeniería de Software SpA}
\newcommand{\RUTEmpresa}{76.913.482-4}
\newcommand{\CodigoLicitacion}{TFEP-01/2026}
\newcommand{\RepresentanteLegal}{Ignacio Silva}
\newcommand{\NombreOferta}{Puerto Deportivo Bahía Panitao}
\newcommand{\VersionDocumento}{0.6}
\newcommand{\FechaDocumento}{9 de octubre de 2026}
\newcommand{\ClasificacionDocumento}{Confidencial}

\usepackage{estilo/propuesta-tecnica}
% Portada vectorial reutilizable. No requiere una imagen de fondo.
\NewDocumentCommand{\PortadaDocumento}{m m}{%
  % La portada consume un folio sin reiniciar el contador de páginas.
  \clearpage
  \begingroup
    \thispagestyle{empty}
    \begin{tikzpicture}[remember picture,overlay]
      \fill[SynNavy] (current page.south west) rectangle (current page.north east);
      \fill[SynNavyLight,opacity=.52]
        ($(current page.north east)+(-4.2cm,0)$) rectangle (current page.south east);

      % Red de conexiones inspirada en la identidad digital de Synaptix.
      \draw[SynCyan,line width=2.2pt,opacity=.92]
        ($(current page.south west)+(8.2cm,2.2cm)$)
        .. controls +(3.4cm,4.0cm) and +(-4.4cm,-3.0cm) ..
        ($(current page.north east)+(-1.0cm,-4.2cm)$);
      \draw[SynBlue,line width=1.2pt,opacity=.78]
        ($(current page.south west)+(10.0cm,0.4cm)$)
        .. controls +(1.0cm,5.7cm) and +(-3.3cm,-1.0cm) ..
        ($(current page.east)+(0cm,3.0cm)$);
      \draw[SynTeal,line width=1.25pt,opacity=.70]
        ($(current page.south west)+(6.4cm,5.0cm)$)
        .. controls +(4.2cm,2.4cm) and +(-2.2cm,-4.0cm) ..
        ($(current page.north east)+(-2.0cm,-1.0cm)$);
      \draw[SynCyan!65,line width=.8pt,opacity=.62]
        ($(current page.south west)+(11.4cm,3.0cm)$)
        .. controls +(2.6cm,2.8cm) and +(-1.5cm,-3.2cm) ..
        ($(current page.north east)+(-.4cm,-7.8cm)$);

      \foreach \p/\s in {
        {($(current page.south west)+(8.2cm,2.2cm)$)}/3.2pt,
        {($(current page.south west)+(11.0cm,6.2cm)$)}/4.2pt,
        {($(current page.east)+(-1.3cm,3.0cm)$)}/3.2pt,
        {($(current page.north east)+(-2.0cm,-1.0cm)$)}/4.4pt,
        {($(current page.north east)+(-1.0cm,-4.2cm)$)}/5.0pt,
        {($(current page.north east)+(-.4cm,-7.8cm)$)}/3.1pt
      }{
        \fill[SynCyan] \p circle (\s);
        \draw[SynCyan!45,line width=.7pt] \p circle ({\s+3.2pt});
      }

      \node[anchor=north west] at ($(current page.north west)+(1.55cm,-1.35cm)$)
        {\fontsize{18}{20}\selectfont\textbf{\textcolor{SynWhite}{SYNAP}\textcolor{SynCyan}{TIX}}};
      \node[anchor=north east,text=SynCyan] at ($(current page.north east)+(-1.45cm,-1.46cm)$)
        {\fontsize{9}{11}\selectfont\bfseries INTEGRACIÓN · DATOS · DECISIONES};

      \node[anchor=west,align=left,text width=.66\paperwidth]
        at ($(current page.west)+(1.55cm,2.7cm)$) {%
          {\fontsize{13}{16}\selectfont\bfseries\textcolor{SynCyan}{PROPUESTA}}\\[2mm]
          {\fontsize{34}{38}\selectfont\bfseries\textcolor{SynWhite}{TÉCNICA}}\\[7mm]
          {\fontsize{12}{15}\selectfont\textcolor{SynWhite!82}{\NombreOferta}}
        };

      \node[anchor=west,align=left,text width=.68\paperwidth]
        at ($(current page.west)+(1.55cm,-1.0cm)$) {%
          {\fontsize{11}{14}\selectfont\bfseries\textcolor{SynCyan}{#1}}\\[2mm]
          {\fontsize{19}{23}\selectfont\bfseries\textcolor{SynWhite}{#2}}
        };

      \node[anchor=west,align=left,text width=.68\paperwidth]
        at ($(current page.west)+(1.55cm,-4.5cm)$) {%
          {\fontsize{11}{14}\selectfont\bfseries\textcolor{SynWhite}{\RazonSocial}}\\[2mm]
          {\fontsize{11}{14}\selectfont\textcolor{SynWhite!82}{RUT \RUTEmpresa\quad Licitación N\textsuperscript{o} \CodigoLicitacion}}
        };

      \draw[SynCyan,line width=1.5pt]
        ($(current page.south west)+(1.55cm,3.25cm)$) -- ++(3.2cm,0);
      \node[anchor=south west,align=left,text width=.54\paperwidth,text=SynWhite!74]
        at ($(current page.south west)+(1.55cm,1.35cm)$) {%
          \fontsize{11}{14}\selectfont
          \textbf{Versión} \VersionDocumento\quad
          \textbf{Fecha} \FechaDocumento\\[1mm]
          \textbf{Clasificación} \ClasificacionDocumento
        };

      \node[anchor=south east,text=SynWhite!84]
        at ($(current page.south east)+(-1.45cm,2.65cm)$) {%
          \FirmaCompletaRepresentante
        };

      % Media firma y folio propios de la portada sobre el fondo oscuro.
      \node[anchor=south east,text=SynWhite!78]
        at ($(current page.south east)+(-1.45cm,0.72cm)$) {%
          \fontsize{8}{10}\selectfont
          \MediaFirmaRepresentante\hspace{5mm}%
          Página\enspace\textbf{\textcolor{SynCyan}{\thepage}}
        };
    \end{tikzpicture}
    \null
  \clearpage
  \endgroup
}

\NewDocumentCommand{\PortadaPropuesta}{m m}{%
  \PortadaDocumento{SUBDOCUMENTO #1}{#2}%
}

% Cada formulario abre con una portada independiente de su índice.
\NewDocumentCommand{\PortadaFormulario}{m m}{%
  \PortadaDocumento{FORMULARIO #1}{#2}%
}

% Actualice estas filas en cada subdocumento para reflejar el uso real.
\NewDocumentCommand{\FilaDeclaracionIA}{m m m m m m}{%
  #1 & #2 & #3 & #4 & #5 & #6 \\
  \hline
}

\providecommand{\DeclaracionesIA}{%
  \FilaDeclaracionIA
    {Plantilla base}
    {OpenAI Codex}
    {Diseño y generación inicial de la plantilla LaTeX}
    {Medio}
    {Alto}
    {Equipo de Synaptix: control de tipografía, tablas y paginación pendiente antes de emisión}%
  \FilaDeclaracionIA{Formato del SD7, 2026-10-07}{OpenAI Codex}{Adaptación de portada, metadatos, índice, encabezados, firmas, foliación y cierre al formato común de Synaptix.}{Bajo}{Ninguno}{José Lara: responsable del control editorial y visual previo a emisión; revisión pendiente.}
}

\newenvironment{declaracionia}{%
  \begingroup
  \setlength{\tabcolsep}{3pt}
  \begin{longtable}{L{\dimexpr.105\textwidth-2\tabcolsep\relax}L{\dimexpr.14\textwidth-2\tabcolsep\relax}L{\dimexpr.21\textwidth-2\tabcolsep\relax}L{\dimexpr.115\textwidth-2\tabcolsep\relax}L{\dimexpr.14\textwidth-2\tabcolsep\relax}L{\dimexpr.29\textwidth-2\tabcolsep\relax}}
    \caption{Declaración de uso de IA por sección y archivo asociado}
    \label{tab:declaracion-ia} \\
    \hline
    \EncabezadoTabla{Sección} &
    \EncabezadoTabla{Herramienta} &
    \EncabezadoTabla{Finalidad del uso} &
    \EncabezadoTabla{Nivel en texto} &
    \EncabezadoTabla{Nivel en diagramas} &
    \EncabezadoTabla{Revisión humana: responsable y verificación} \\
    \hline
    \endfirsthead
    \hline
    \EncabezadoTabla{Sección} &
    \EncabezadoTabla{Herramienta} &
    \EncabezadoTabla{Finalidad del uso} &
    \EncabezadoTabla{Nivel en texto} &
    \EncabezadoTabla{Nivel en diagramas} &
    \EncabezadoTabla{Revisión humana: responsable y verificación} \\
    \hline
    \endhead
}{%
  \end{longtable}
  \endgroup
}

\newcommand{\CierreSubdocumento}{%
  \clearpage
  \printbibliography[heading=referenciassynaptix,title={Referencias}]
  \clearpage
  \chapter*{Declaración de uso de IA}
  \addcontentsline{toc}{chapter}{Declaración de uso de IA}
  \markboth{Declaración de uso de IA}{Declaración de uso de IA}
  Esta declaración identifica las herramientas de inteligencia artificial utilizadas, su finalidad y los controles humanos consignados. Las filas distinguen responsabilidades asignadas de verificaciones realizadas; los controles pendientes se completan antes de emisión. Los niveles admitidos son Ninguno, Bajo, Medio y Alto. Synaptix mantiene el registro conforme al uso efectivo y lo consolida en el Formulario A-6.

  La Tabla~\ref{tab:declaracion-ia} registra el uso declarado en este documento y sus archivos asociados.
  \begin{declaracionia}
    \DeclaracionesIA
  \end{declaracionia}
  \FuenteTabla{Registro de uso de IA de este subdocumento.}
}

\graphicspath{{figuras/}}
\newcommand{\EnFormularioTCatorce}{}

\begin{document}
\pagenumbering{arabic}
\PortadaFormulario{T-14}{Plan de trabajo, EDT y carta Gantt}
\tableofcontents
\clearpage
\setcounter{chapter}{6}
\chapter{Plan de trabajo, EDT y carta Gantt}
Synaptix presenta la estructura de entregables y el diccionario de trabajo asociado al Subdocumento 7. Los códigos de cuentas identifican los mismos componentes en este formulario, en el capítulo y en los registros de esfuerzo. La planificación cubre el contrato de 56 meses y conserva las ventanas de marcha blanca, los pasos a producción y el inicio de operación del artículo 17 \parencite[art.~17, pp.~12--13; Formulario~T-14]{bases-admin}.

\section{Estructura de entregables y trazabilidad}
La EDT organiza los resultados del contrato por entregables. El inventario distingue componentes superiores y paquetes de ejecución; el diccionario desarrolla cada resultado y presenta mediante familias los lotes y servicios repetidos, con sus instancias identificadas en los registros técnicos.
\subsection{Criterio de descomposición y código de cuentas}
Synaptix adopta una descomposición por entregables. El código raíz \texttt{1} identifica el contrato; los códigos \texttt{1.n} agrupan entregables principales; y los niveles inferiores identifican componentes y paquetes. La etapa contractual será un atributo del trabajo, de modo que una capacidad desarrollada progresivamente conserve su relación con el mismo entregable. La EDT expresa pertenencia y cobertura; las fechas y dependencias se establecerán en la red de actividades.

La descomposición se detendrá cuando el paquete tenga un resultado medible, una condición de aceptación, un responsable único y una estimación defendible. Las instancias actualmente estimadas tienen entre 8 y 80 horas-persona; este rango describe la estimación y no constituye una exigencia de las Bases ni un límite para definir nuevos paquetes. Se mantiene un período quincenal de control. La dedicación parcial y las esperas externas se declararán por separado; el esfuerzo del paquete no equivale automáticamente a su duración. Los servicios recurrentes se presentan mediante familias con resultado, responsable, frecuencia y aceptación comunes; cada período conserva su instancia estimada y asignada.

Cada rama incluirá todo el trabajo de su componente superior y excluirá el trabajo imputado a otra rama. La verificación de cobertura cruzará requisitos, entregables y paquetes; la suma de porcentajes de un gráfico no sustituirá esa comprobación. Las pruebas integradas, los servicios compartidos y las innovaciones tendrán una imputación única; el uso de un componente por varios dominios se registrará como dependencia.

\subsection{Inventario de entregables principales}
El inventario organiza los productos y servicios establecidos en los Subdocumentos 3, 4, 5 y 13. Las Tablas~\ref{tab:sd7-inventario}, \ref{tab:sd7-inventario-datos} y \ref{tab:sd7-inventario-servicio} identifican las ramas y su frontera de trabajo. Sus códigos agrupan entregables y no representan todavía paquetes terminales estimados.

\begin{table}[htbp]
\centering
\begingroup
\setlength{\tabcolsep}{4pt}
\begin{tabularx}{\linewidth}{L{12mm}L{36mm}Y}
\hline
\EncabezadoTabla{Código} & \EncabezadoTabla{Entregable} & \EncabezadoTabla{Cobertura y frontera}\\\hline
1.1 & Expediente de dirección y control & Línea base, trazabilidad, actas, cambios, riesgos, informes y control físico y financiero. Incluye el plan de reversibilidad de los primeros 90 días y sus actualizaciones.\\\hline
1.2 & Plataforma híbrida habilitada & Cinco ambientes DEV, QA, PREPROD, PROD y DR; nube, recinto y terminales; configuración, conectividad, continuidad, despliegue y observabilidad técnica. SD4 y T-11 delimitan los activos.\\\hline
1.3 & Controles de auditoría y seguridad habilitados & Identidades, permisos, consentimiento, cifrado, gestión de claves, trazabilidad y controles transversales. Los ensayos integrales de estos controles pertenecen a 1.10.\\\hline
1.4 & Capacidades del Núcleo Operacional & Operación marítima, Escuela náutica e Infraestructura: salidas y regresos, amarras, participación y retiro, inventario físico e inspección de boyas sin conexión.\\\hline
\end{tabularx}
\caption{Entregables de gobierno, plataforma y núcleo operacional}
\label{tab:sd7-inventario}
\FuenteTabla{Síntesis de Synaptix: SD3, secciones 3.1--3.4; SD4, secciones 4.1--4.3; SD5, secciones 5.1--5.4; SD13, innovaciones 1--5.}
\endgroup
\end{table}

\begin{table}[htbp]
\centering
\begingroup
\setlength{\tabcolsep}{4pt}
\begin{tabularx}{\linewidth}{L{12mm}L{36mm}Y}
\hline
\EncabezadoTabla{Código} & \EncabezadoTabla{Entregable} & \EncabezadoTabla{Cobertura y frontera}\\\hline
1.5 & Capacidades del Núcleo de Servicios & Varadero y contratistas, Combustible y Ambiente y consumos: faenas, habilitación, despacho, medición individual, residuos e indicadores. La medición inicial corresponde a E1; las capacidades completas, a E2.\\\hline
1.6 & Capacidades del Núcleo Administrativo & Clientes y contratos y Finanzas y cobranza: maestros, asignaciones, espera, visitas, hechos facturables, saldos y expedientes. Conserva la autoridad tributaria del sistema contable externo.\\\hline
1.7 & Canales y servicios comunes integrados & Aplicación Android, portal y consola; interfaces externas, notificaciones, telemetría y sincronización. La construcción compartida se imputa aquí; las reglas propias de cada dominio, a 1.4--1.6.\\\hline
1.8 & Datos migrados y archivo histórico consultable & Diagnóstico, herramientas, saneamiento, cargas, conciliación y consulta independiente del sistema de 2014. Comprende la carga inicial de E1 y la cobertura histórica completa de E2 conforme a SD5.\\\hline
\end{tabularx}
\caption{Entregables de servicios, administración, integración y datos}
\label{tab:sd7-inventario-datos}
\FuenteTabla{Síntesis de Synaptix: SD3, secciones 3.1--3.4; SD4, secciones 4.1--4.3; SD5, secciones 5.1--5.4; SD13, innovaciones 1--5.}
\endgroup
\end{table}

\begin{table}[htbp]
\centering
\begingroup
\setlength{\tabcolsep}{4pt}
\begin{tabularx}{\linewidth}{L{12mm}L{36mm}Y}
\hline
\EncabezadoTabla{Código} & \EncabezadoTabla{Entregable} & \EncabezadoTabla{Cobertura y frontera}\\\hline
1.9 & Cinco innovaciones incorporadas y evaluadas & Solicitudes de mantenimiento, mantenimiento preventivo en invernada, lectura asistida, servicio gestionado por resultados y Pasaporte Ambiental Náutico. Incluye su trabajo incremental sin repetir capacidades base.\\\hline
1.10 & Expediente de verificación y aceptación & Protocolos y pruebas integradas, seguridad, desempeño, estrés, aislamiento, recuperación y migración; incidencias y evidencia para T-17. Incluye los dos ensayos completos de migración de SD5.\\\hline
1.11 & Usuarios preparados y etapas implantadas & Materiales por perfil, formación, olas por zona o proceso, convivencia, conciliación diaria, marcha blanca, reversión y estabilización con presencia en terreno.\\\hline
1.12 & Servicio gestionado y transferencia de salida & Operación de ambas etapas en meses 21--56: soporte, NOC/SOC, mantenimiento, respaldos, seguimiento de niveles de servicio y transferencia final de activos, datos y conocimiento.\\\hline
\end{tabularx}
\caption{Entregables de innovación, aceptación, implantación y operación}
\label{tab:sd7-inventario-servicio}
\FuenteTabla{Síntesis de Synaptix: SD3, secciones 3.1--3.4; SD4, secciones 4.1--4.3; SD5, secciones 5.1--5.4; SD13, innovaciones 1--5.}
\endgroup
\end{table}

El inventario separa la construcción de capacidades de su comprobación integral. Cada paquete de construcción incluirá sus pruebas unitarias y documentación; 1.10 reunirá los ensayos entre componentes y la evidencia de recepción. Las pruebas internas específicas de una innovación permanecerán en 1.9 y sus resultados alimentarán el expediente general, sin volver a estimar la misma ejecución.

Las compras y obras físicas atribuidas al cliente por el Subdocumento 2 se registrarán como habilitaciones externas con responsable, fecha requerida y condiciones de recepción. Synaptix mantendrá en sus paquetes las especificaciones, la coordinación y la verificación de integración que le corresponden. Estas dependencias no se tratarán como instalaciones ya ejecutadas ni como ampliaciones automáticas del plazo \parencite[cap.~11; cap.~17, secc.~17.5]{caso07}.

Los códigos \texttt{IN1} a \texttt{IN5} del Subdocumento 13 conservarán su identidad como referencias de origen. El trabajo de diseño, construcción, pruebas y pilotos de las innovaciones se imputará a 1.9; el servicio recurrente, a 1.12. En particular, los informes \texttt{IN1.7.m.q} y la operación \texttt{IN2.7} y \texttt{IN3.7} conservarán su referencia a SD13 dentro de la rama de servicio. Cada paquete de ejecución tendrá una sola imputación de esfuerzo; si un componente de SD13 requiere subdivisión, la matriz de correspondencia identificará todos sus paquetes hijos sin contar también el esfuerzo del padre.

El calendario de la innovación 2 conservará $M_{\mathrm{inv}}$ como el primer junio disponible con IN1 y sus interfaces aceptadas, no anterior al mes 21, y mantendrá la secuencia de invernada y botaduras definida en SD13. Los componentes \texttt{IN4.4} e \texttt{IN5.4} distinguirán sus resultados de puesta en servicio de sus trabajos recurrentes. La correspondencia detallada y las cargas se consolidarán antes de cerrar T-14 y T-15.

\subsection{Expediente de diagnóstico de las fuentes de migración}
El componente \texttt{1.8.1} entregará un expediente que permita decidir cómo extraer, interpretar y conciliar los antecedentes antes de construir las cargas definitivas. Su cobertura comprende los seis conjuntos obligatorios de SD5, sección 5.3.1: socios y armadores; contratos vigentes y documentos; seis años de pagos y saldos; embarcaciones e invernadas; los 14 expedientes de abandono; y tres años de matrícula escolar. El inventario abarcará las fuentes de los doce años de historia para delimitar su tratamiento, sin extender automáticamente a todo ese período la carga operacional \parencite[cap.~15, RT-05.15]{caso07}.

La Figura~\ref{fig:sd7-rama-diagnostico} muestra la pertenencia de los cuatro paquetes al componente 1.8.1. La Tabla~\ref{tab:sd7-paquetes-diagnostico} resume sus resultados y esfuerzos. Las horas son una estimación ascendente inicial de la oferta, basada en las tareas enumeradas en el diccionario; no son mediciones de productividad ni duración comprometida.

\begin{figure}[htbp]
\centering
\begin{tikzpicture}[
 paquete/.style={draw=SynLine,fill=SynSurface,rounded corners=1mm,
 text width=82mm,minimum height=13mm,inner sep=3mm,align=left,
 font=\fontsize{9}{12}\selectfont},
 rama/.style={draw=SynNavy,fill=SynNavy,text=SynWhite,
 text width=32mm,inner sep=3mm,align=left,font=\fontsize{9}{12}\selectfont}]
\node[rama] (raiz) at (-5,0) {\textbf{1.8.1}\\Expediente de diagnóstico de las fuentes de migración};
\node[paquete] (p1) at (2,2.7) {\textbf{1.8.1.1}\\Catálogo de fuentes y accesos};
\node[paquete] (p2) at (2,.9) {\textbf{1.8.1.2}\\Procedimiento de extracción diagnóstica};
\node[paquete] (p3) at (2,-.9) {\textbf{1.8.1.3}\\Informe reproducible de calidad del origen};
\node[paquete] (p4) at (2,-2.7) {\textbf{1.8.1.4}\\Matriz de cobertura y decisiones de migración};
\draw[SynBlue,line width=.8pt] (raiz.east) -- (-2.6,0);
\draw[SynBlue,line width=.8pt] (-2.6,2.7) -- (-2.6,-2.7);
\foreach \paquete/\altura in {p1/2.7,p2/.9,p3/-.9,p4/-2.7}{
 \draw[SynBlue,line width=.8pt] (-2.6,\altura) -- (\paquete.west);
}
\end{tikzpicture}
\caption{Descomposición del expediente de diagnóstico en paquetes}
\label{fig:sd7-rama-diagnostico}
\FuenteFigura{EDT propuesta por Synaptix para el diagnóstico previo a la migración.}
\end{figure}

La jerarquía asigna cada paquete a un único componente superior. Las líneas expresan esa pertenencia; las precedencias de ejecución se definirán separadamente.

\begin{table}[htbp]
\centering
\begin{tabularx}{\linewidth}{L{25mm}YR{13mm}L{31mm}}
\hline
\EncabezadoTabla{Código} & \EncabezadoTabla{Resultado} & \EncabezadoTabla{Horas} & \EncabezadoTabla{Responsable}\\\hline
1.8.1.1 & Catálogo de fuentes y accesos & 24 & Líder de Datos\\\hline
1.8.1.2 & Procedimiento de extracción diagnóstica & 40 & Líder de Datos\\\hline
1.8.1.3 & Informe reproducible de calidad del origen & 32 & Líder de Datos\\\hline
1.8.1.4 & Matriz de cobertura y decisiones de migración & 24 & Líder de Datos\\\hline
\textbf{Total} & \textbf{Expediente de diagnóstico} & \textbf{120} & \\\hline
\end{tabularx}
\caption{Paquetes del expediente de diagnóstico de fuentes}
\label{tab:sd7-paquetes-diagnostico}
\FuenteTabla{Estimación ascendente propuesta por Synaptix; alcance definido en SD5, sección 5.3.}
\end{table}

Las 120 horas comprenden 24 de Análisis Funcional, 72 de Datos, 8 de Seguridad y 16 de Calidad. Incluyen documentación y comprobación de los productos de diagnóstico. No incluyen saneamiento masivo, carga definitiva, consulta histórica implementada, los dos ensayos completos de migración ni la operación del servicio; esos trabajos requieren otros componentes de 1.8, 1.10 y 1.12. La aceptación de este expediente no equivale a la aceptación de la migración.

\ifdefined\EnFormularioTCatorce
El diccionario detallado de los cuatro paquetes se presenta en el apartado siguiente de este formulario.
\else
El diccionario detallado de los cuatro paquetes se presenta en el archivo independiente \path{SYNAPTIX-Formulario-T-14.pdf}.
\fi
Sus fichas conservan estos códigos, esfuerzos, responsables y fronteras; el resumen no sustituye sus condiciones de aceptación.

\section{Plataforma, seguridad y servicios comunes}
Las cuentas 1.2, 1.3 y 1.7 reúnen la habilitación de plataforma, los controles de seguridad y los servicios compartidos. Las fichas siguientes identifican sus resultados y condiciones de recepción, conservando el inventario de equipos en T-11.
\subsection{Plataforma, seguridad y servicios comunes}
Las cuentas 1.2, 1.3 y 1.7 se descomponen en 41 paquetes técnicos: 18 de plataforma, 10 de seguridad y 13 de servicios comunes. Su estimación ascendente suma 2.024 HH: 776, 512 y 736, respectivamente. El T-14 presenta las fichas con contenido, exclusiones, aceptación, responsable último, esfuerzo por perfil, costo, recursos, supuestos, dependencias e hitos. El T-15 registra las actividades derivadas y compara su carga con las cuentas ya cuantificadas.

La plataforma incluye DEV, QA, PREPROD, PROD y DR: los cinco deberán estar habilitados y operativos para H3, conforme a RT-04.01. Cada ambiente recibe una ficha y evidencia propia; el expediente de entrega consolida esas evidencias sin volver a estimar su ejecución. La creación reproducible y el gobierno del repositorio del cliente se mantienen separados de la construcción de módulos, mientras las pruebas integrales de continuidad y rendimiento pertenecen a 1.10 \parencite[RT-03.03, RT-04.01]{bases-transversales}.

La seguridad comprende identidad central y autoridad local, permisos, claves, consentimiento, auditoría, amenazas, datos de prueba y respuesta técnica. Servicios comunes comprende contratos/eventos, mensajería, sincronización, telemetría, canales y adaptadores externos. Los nueve dominios funcionales consumirán esos componentes sin duplicar su construcción. Las reglas, pantallas y decisiones de cada dominio se estimarán en 1.4--1.6, y las extensiones incrementales conservarán sus paquetes de 1.9.

\subsection{Habilitaciones externas y lotes de terminales}
Los identificadores EXT representan hitos de dependencia, sin esfuerzo de Synaptix adicional al ya incluido en las fichas: EXT-RECINTO, obra y condiciones de instalación recibidas; EXT-EQUIPOS, activos nuevos y conformes a T-11; EXT-RED, tendidos, accesos y enlaces habilitados; EXT-TELEMETRIA, buses y concentradores instalados y accesibles; EXT-LICENCIAS, licencias efectivas y soporte autorizados; EXT-IDENTIDAD, fuente laboral y responsables del cliente; EXT-CONTABLE y EXT-PAGOS, acceso y contratos de integración de terceros. El responsable de coordinación será Fernando Guerrero, con Proyecto y la contraparte designada por el cliente para compras/obras. La coordinación de terceros no convierte sus trabajos en instalaciones ya ejecutadas ni autoriza modificar el alcance o el plazo.

Las habilitaciones de recinto, activos, red y telemetría deben permitir completar las fichas técnicas antes de H3; la fuente de identidad debe estar disponible antes de configurar su integración. Las interfaces contable y de pagos se requieren antes de construir y validar sus adaptadores E2. Proyecto registrará fecha requerida, fecha ofrecida, responsable externo, condición de recepción y efecto sobre dependencias. Si una condición falla, el expediente registrará el bloqueo y se gestionará su resolución dentro de los hitos contractuales.

\paragraph{1.2.3.6.t.b: parque de terminales enrolado y comprobado}
La ficha de referencia 1.2.3.4 define y prueba políticas; este paquete aplica esas políticas al parque. $t$ identifica Windows o Android y $b$ un lote de hasta cuatro equipos, con activo y usuario o función asignados. Bruno Toro responde por el resultado. Incluye enrolamiento, cifrado, versión, actualización, protección, retiro y registro de acceso; excluye adquisición, obras físicas, capacitación de procesos e implementación de la aplicación.

Se recibirá con un registro por activo y comprobación de política, cobertura de protección, bloqueo y recuperación autorizada; un equipo sin licencia o agente efectivo no se considerará protegido. Esfuerzo por lote: Operación 8, Seguridad 4 y Calidad 4 HH, total 16. Recursos: terminales recibidos de T-11, UEM/licencias efectivas, política recibida, red y usuarios autorizados. Se supone hardware compatible y una configuración estándar por tipo; incompatibilidades se dimensionan aparte antes de enrolar.

Para cantidades efectivas $N_W,N_A$, se genera $L=\lceil N_W/4\rceil+\lceil N_A/4\rceil$ lotes, con cero lotes sólo para un tipo confirmado sin equipos. El esfuerzo adicional es $16L$ HH; no se cuenta otra vez la prueba de referencia. El inventario actualizado deberá instanciar los lotes sin reemplazar el detalle de T-11. Cada lote se controla dentro de una quincena y antes de la habilitación del grupo de usuarios que lo necesita; las fechas se fijan con la nivelación y las ventanas autorizadas. La ausencia de inventario confirmado no representa cero horas. Predecesoras: 1.2.3.4, terminales y licencias recibidos; hito: preparación técnica H3 y habilitaciones E1/E2 correspondientes. Referencias: SD4, puestos y gestión de terminales; T-11; SD2, SUP-04; BT capítulo 8.

\section{Diccionario de plataforma, seguridad y servicios comunes}
Las fichas descomponen las cuentas 1.2, 1.3 y 1.7 por resultados técnicos, con responsable último único y esfuerzo entre 8 y 80 HH. Las comprobaciones propias de configuración o construcción se incluyen en cada ficha; la validación conjunta de carga, estrés, resiliencia, seguridad ofensiva, accesibilidad y operación desconectada pertenece a 1.10. La recepción de un componente no sustituye la certificación de la solución.

Se supone disponibilidad de especificaciones SD4, modelo SD5 aprobado en H2, recursos/licencias seleccionados y habilitaciones del cliente conformes a SD2. Las horas corresponden a configurar recursos o construir componentes acotados sobre las tecnologías seleccionadas; no son productividad medida ni horas de compras/obras del cliente. Un fallo de recepción o requisito técnico pendiente bloqueará el trabajo dependiente. Los períodos indicados son ventanas para nivelación y no acreditan factibilidad del calendario.

La compra, obra y suministro de medidores se conserva en las responsabilidades del SD2. Synaptix especifica, coordina y verifica la integración que le corresponde. El inventario detallado permanece en T-11 y los cálculos de capacidad, electricidad y autonomía en SD4. Las variantes o cantidades no verificadas se contrastarán antes de instalación; este diccionario no transforma una condición de diseño en dato del equipo existente.

El perfil I corresponde a Fernando Guerrero; O y L, a José Lara Arce. Sus funciones comparten una disponibilidad por persona. Las actividades derivadas se conservan en \path{componentes/actividades_plataforma.csv}, sin añadir esfuerzo al de las fichas.

La recepción contractual corresponde a la contraparte autorizada. Todas las fichas requieren insumos y versiones recibidos, registros y entornos autorizados, herramientas y contraparte programada. Las HH excluyen dedicación del cliente y cobertura externa continua. La falta de insumos o evidencia bloquea la recepción y exige reprogramar dentro de las condiciones contractuales. Estas condiciones comunes se aplican además de los criterios específicos de cada ficha.

Las fichas de la Tabla~\ref{tab:dic-compacto-plataforma} presentan resultado, responsable, esfuerzo y aceptación. Las condiciones particulares, recursos y ventanas por código se conservan en \path{componentes/condiciones_fichas_t14.csv}, junto con la asignación y precedencias del diccionario técnico integrado. Las actividades de T-15 desarrollan los aportes por perfil y el T-18 desarrolla la ejecución de implantación.


\begingroup\small\setlength{\tabcolsep}{3pt}
\begin{longtable}{L{52mm}L{91mm}}
\caption{Diccionario de plataforma y seguridad}\label{tab:dic-compacto-plataforma}\\\hline
Paquete, resultado y responsable & Contenido y criterio de aceptación\\\hline\endfirsthead
\hline Paquete, resultado y responsable & Contenido y criterio de aceptación\\\hline\endhead
\texttt{1.\allowbreak{}2.\allowbreak{}1.\allowbreak{}1}\newline Repositorio e infraestructura como código gobernados\newline Fernando Guerrero; 40 HH. & Repositorio del cliente, cuenta raíz documentada, módulos OpenTofu, estado cifrado y bloqueo; excluye módulos de aplicación. Titularidad, MFA, revisión independiente, estado recuperable y ausencia de secretos en Git quedan comprobados. Responsable: Fernando Guerrero (Arquitectura/Integración).\\\hline
\texttt{1.\allowbreak{}2.\allowbreak{}1.\allowbreak{}2}\newline Red cloud primaria reproducible\newline Fernando Guerrero; 48 HH. & VPC, subredes privadas en dos zonas, rutas, endpoints y acceso de administración por código. Se reproduce la red y se comprueba acceso permitido, bloqueo de exposición de datos y diferencias del plan IaC. Responsable: Fernando Guerrero (Arquitectura/Integración).\\\hline
\texttt{1.\allowbreak{}2.\allowbreak{}1.\allowbreak{}3}\newline Servicios de persistencia y objetos habilitados\newline Nicolás Torfan; 48 HH. & RDS PostgreSQL, almacén de proyecciones y S3 con políticas; excluye modelo funcional, migración y cargas de negocio. Se identifican autoridades y recursos separados; conectividad TLS, políticas, respaldo y prohibición de acceso directo entre dominios se comprueban. Nicolás Torfan, Datos, responde por el resultado y presenta la evidencia; la contraparte mantiene la recepción contractual.\\\hline
\texttt{1.\allowbreak{}2.\allowbreak{}1.\allowbreak{}4}\newline Plataforma de contenedores y exposición habilitada\newline Fernando Guerrero; 48 HH. & ECS/Fargate, balanceo privado, API Gateway, DNS y protección de entrada por IaC; excluye servicios de negocio. Imagen de comprobación se ejecuta por despliegue reproducible, sin publicar subredes de aplicación/datos ni claves estáticas. Responsable: Fernando Guerrero (Arquitectura/Integración).\\\hline
\texttt{1.\allowbreak{}2.\allowbreak{}1.\allowbreak{}5}\newline Cadena de promoción y artefactos protegidos\newline Fernando Guerrero; 48 HH. & Repositorio, registro OCI, firma, evidencias de dependencias y promoción; excluye construir módulos funcionales. El mismo artefacto identificado se promueve sin recompilar por ambiente y una versión no autorizada queda bloqueada. Responsable: Fernando Guerrero (Arquitectura/Integración).\\\hline
\texttt{1.\allowbreak{}2.\allowbreak{}2.\allowbreak{}1}\newline Ambiente DEV habilitado\newline Fernando Guerrero; 24 HH. & Configuración, acceso y despliegue de comprobación DEV con datos sintéticos; excluye horas de construcción de software. DEV se reproduce, usa sólo datos permitidos y mantiene aislamiento y horario versionado. Responsable: Fernando Guerrero (Arquitectura/Integración).\\\hline
\texttt{1.\allowbreak{}2.\allowbreak{}2.\allowbreak{}2}\newline Ambiente QA habilitado\newline Davor Serey; 24 HH. & Configuración QA, datos autorizados y ejecución de comprobación; excluye regresión de todos los módulos. QA queda aislado de producción, identifica versiones y sólo usa datos sintéticos o anonimizados aprobados. Davor Serey, Calidad, responde por el resultado y presenta la evidencia; la contraparte mantiene la recepción contractual.\\\hline
\texttt{1.\allowbreak{}2.\allowbreak{}2.\allowbreak{}3}\newline Ambiente PREPROD habilitado\newline Fernando Guerrero; 32 HH. & Topología productiva a escala de prueba y restitución controlada; excluye ejecución de carga, UAT y ensayos de migración. Se comprueban aislamiento, capacidad de restaurar configuración y control de copias excepcionales anonimizadas. Responsable: Fernando Guerrero (Arquitectura/Integración).\\\hline
\texttt{1.\allowbreak{}2.\allowbreak{}2.\allowbreak{}4}\newline Ambiente PROD habilitado\newline José Lara Arce; 40 HH. & Configuración Multi-AZ, límites, permisos y comprobación técnica sin usuarios de negocio. Producción tiene configuración identificada, aislamiento, monitoreo y acceso nominativo; habilitar ambiente no acredita paso a producción E1. José Lara Arce, Operación/SRE, responde por el resultado y presenta la evidencia; la contraparte mantiene la recepción contractual.\\\hline
\texttt{1.\allowbreak{}2.\allowbreak{}2.\allowbreak{}5}\newline Ambiente DR habilitado\newline José Lara Arce; 56 HH. & Warm standby México, réplica cifrada, imágenes, objetos y configuración mínima; excluye prueba integral de desastre. Se comprueban salud, replicación, permisos y restitución del entorno; capacidad de respaldo no se declara por una réplica sin evidencia. José Lara Arce, Operación/SRE, responde por el resultado y presenta la evidencia; la contraparte mantiene la recepción contractual.\\\hline
\texttt{1.\allowbreak{}2.\allowbreak{}3.\allowbreak{}1}\newline Expediente de habilitación del recinto y recepción de activos\newline Fernando Guerrero; 40 HH. & Contraste de plano y condiciones SD4, compras/obras del cliente y T-11; excluye ejecutar obras o comprar equipamiento. Cada habilitación tiene especificación, titular, fecha requerida y evidencia de recepción pendiente o conforme; ningún dato proyectado se atribuye al recinto existente. Responsable: Fernando Guerrero (Arquitectura/Integración).\\\hline
\texttt{1.\allowbreak{}2.\allowbreak{}3.\allowbreak{}2}\newline Configuración de nodos locales y almacenamiento redundante\newline José Lara Arce; 48 HH. & Configuración versionada de los dos nodos, almacenamiento y servicios locales; excluye construcción de reglas de dominio. Inventario y firmware corresponden a activos recibidos; restauración de configuración, salud de discos y aislamiento se comprueban. José Lara Arce, Operación/SRE, responde por el resultado y presenta la evidencia; la contraparte mantiene la recepción contractual.\\\hline
\texttt{1.\allowbreak{}2.\allowbreak{}3.\allowbreak{}3}\newline Configuración de red local, segmentación y enlaces\newline Fernando Guerrero; 56 HH. & Configuración versionada de cortafuegos, switches, VLAN, WAN y túneles; excluye obras y tendido físico del cliente. Rutas autorizadas y segmentación se verifican; fallos de un enlace no abren acceso no autorizado ni exponen bases de datos. Responsable: Fernando Guerrero (Arquitectura/Integración).\\\hline
\texttt{1.\allowbreak{}2.\allowbreak{}3.\allowbreak{}4}\newline Gestión y protección de terminales de referencia\newline Bruno Toro; 40 HH. & Políticas UEM, cifrado, actualización y protección para una estación y un Android de referencia; enrolamiento del parque se estima por lotes. Se verifican enrolamiento, bloqueo, pérdida, política, actualización y cobertura efectiva; la selección de una licencia no equivale a cobertura comprobada. Bruno Toro, Seguridad, responde por el resultado y presenta la evidencia; la contraparte mantiene la recepción contractual.\\\hline
\texttt{1.\allowbreak{}2.\allowbreak{}3.\allowbreak{}5}\newline Configuración de concentradores y broker de borde\newline Fernando Guerrero; 48 HH. & Configuración de diez concentradores y EMQX, reloj, certificados y cola; excluye software de sincronización y puesta en servicio de medidores. Cada concentrador tiene activo, identidad, UTC, muestreo por minuto, agregado cada 15 minutos y cola de 24 horas comprobados con señales de prueba. Responsable: Fernando Guerrero (Arquitectura/Integración).\\\hline
\texttt{1.\allowbreak{}2.\allowbreak{}4.\allowbreak{}1}\newline Observabilidad técnica y alertas habilitadas\newline José Lara Arce; 48 HH. & OpenTelemetry, métricas, registros, trazas y alertas de infraestructura; excluye indicadores funcionales de cada módulo. Las fallas simuladas generan alerta con activo, correlación y responsable; se comprueban cobertura y falta de agente. José Lara Arce, Operación/SRE, responde por el resultado y presenta la evidencia; la contraparte mantiene la recepción contractual.\\\hline
\texttt{1.\allowbreak{}2.\allowbreak{}4.\allowbreak{}2}\newline Procedimientos de respaldo y restauración técnica\newline José Lara Arce; 48 HH. & Políticas de copia, retención, restauración de muestra y permisos; excluye recuperación integral y retorno de migración. Una copia identificada se restaura con integridad, permisos y tiempos registrados; la existencia de backup no sustituye su restauración. José Lara Arce, Operación/SRE, responde por el resultado y presenta la evidencia; la contraparte mantiene la recepción contractual.\\\hline
\texttt{1.\allowbreak{}2.\allowbreak{}4.\allowbreak{}3}\newline Expediente técnico de entrega de cinco ambientes\newline Fernando Guerrero; 40 HH. & Consolidación de evidencia de infraestructura, observabilidad y habilitaciones; excluye repetir las comprobaciones anteriores. DEV, QA, PREPROD, PROD y DR tienen evidencia y pendientes identificados; condiciones bloqueantes impiden recepción H3. Responsable: Fernando Guerrero (Arquitectura/Integración).\\\hline
\texttt{1.\allowbreak{}3.\allowbreak{}1.\allowbreak{}1}\newline Gestión de claves, secretos y certificados\newline Bruno Toro; 40 HH. & KMS, almacén de secretos, certificados y procedimiento de rotación/revocación; excluye cifrado funcional del código de cada módulo. Credenciales temporales, alcance mínimo, revocación, rotación y lectura autorizada quedan comprobados sin secretos en repositorio. Bruno Toro, Seguridad, responde por el resultado y presenta la evidencia; la contraparte mantiene la recepción contractual.\\\hline
\texttt{1.\allowbreak{}3.\allowbreak{}1.\allowbreak{}2}\newline Identidad central y federación configuradas\newline Bruno Toro; 48 HH. & Cognito, OIDC/SAML, MFA, identidades nominativas y recuperación DR; excluye autoridad local y permisos funcionales de los dominios. Se rechazan cuentas compartidas y acceso sin factor exigido; recuperación DR no presupone replicación de contraseñas no soportada. Bruno Toro, Seguridad, responde por el resultado y presenta la evidencia; la contraparte mantiene la recepción contractual.\\\hline
\texttt{1.\allowbreak{}3.\allowbreak{}1.\allowbreak{}3}\newline Matriz de permisos y separación de funciones\newline Bruno Toro; 40 HH. & Roles, ámbitos y separación entre preparación/aprobación para los nueve dominios; excluye implementación en cada servicio. Cada operación tiene actor, ámbito y denegación; el cliente valida las decisiones y no quedan permisos globales implícitos. Bruno Toro, Seguridad, responde por el resultado y presenta la evidencia; la contraparte mantiene la recepción contractual.\\\hline
\texttt{1.\allowbreak{}3.\allowbreak{}2.\allowbreak{}1}\newline Autoridad local de acceso implementada\newline Bruno Toro; 64 HH. & Componente local de permisos firmados y cifrados y vínculo a fuente laboral; excluye reglas propias de módulos y prueba integral offline. Se comprueban identidad nominativa, ámbito y vencimiento máximo de 24 horas desde emisión, con denegación de administración global. Bruno Toro, Seguridad, responde por el resultado y presenta la evidencia; la contraparte mantiene la recepción contractual.\\\hline
\texttt{1.\allowbreak{}3.\allowbreak{}2.\allowbreak{}2}\newline Revocación y actualización de habilitación local\newline Bruno Toro; 56 HH. & Propagación de cambios, bajas, permisos y fuente laboral actual; excluye provisión de datos laborales por el cliente. Bajas y cambios se aplican con trazabilidad y fallo de fuente visible; no se mantiene acceso por caché indefinida. Bruno Toro, Seguridad, responde por el resultado y presenta la evidencia; la contraparte mantiene la recepción contractual.\\\hline
\texttt{1.\allowbreak{}3.\allowbreak{}3.\allowbreak{}1}\newline Registro de auditoría protegido\newline Bruno Toro; 64 HH. & Eventos, correlación, cuenta separada y conservación protegida; excluye generación de hechos funcionales por módulos. Actor, origen, instante y versión quedan preservados; alteración, acceso indebido y falta de envío se detectan. Bruno Toro, Seguridad, responde por el resultado y presenta la evidencia; la contraparte mantiene la recepción contractual.\\\hline
\texttt{1.\allowbreak{}3.\allowbreak{}3.\allowbreak{}2}\newline Consentimiento y restricciones de finalidad implementados\newline Bruno Toro; 64 HH. & Servicio de consentimiento y evento de cambios; excluye adhesión y reglas adicionales del pasaporte IN5. Alta, revocación, finalidad y conservación tienen decisión trazada; la revocación impide nuevos usos dependientes sin borrar evidencia obligatoria. Bruno Toro, Seguridad, responde por el resultado y presenta la evidencia; la contraparte mantiene la recepción contractual.\\\hline
\texttt{1.\allowbreak{}3.\allowbreak{}4.\allowbreak{}1}\newline Modelo de amenazas y plan de seguridad aprobables\newline Bruno Toro; 48 HH. & Amenazas por autoridad, canal, WAN, terminal y tercero, mitigación y comprobación prevista; excluye pruebas ofensivas integrales. Cada amenaza relevante tiene control, cuenta de trabajo y método de verificación; se identifican pendientes antes de H2. Bruno Toro, Seguridad, responde por el resultado y presenta la evidencia; la contraparte mantiene la recepción contractual.\\\hline
\texttt{1.\allowbreak{}3.\allowbreak{}4.\allowbreak{}2}\newline Detección, cobertura y respuesta técnica configuradas\newline Bruno Toro; 48 HH. & Controles de endpoint, runtime y servicios administrados y procedimientos de contención; excluye operación continua SOC. Se registra alcance efectivo por activo y fallo de agente; una alerta de prueba produce contención, evidencia y responsable identificados. Bruno Toro, Seguridad, responde por el resultado y presenta la evidencia; la contraparte mantiene la recepción contractual.\\\hline
\texttt{1.\allowbreak{}3.\allowbreak{}4.\allowbreak{}3}\newline Expediente de privacidad y control de datos de prueba\newline Bruno Toro; 40 HH. & Inventario de sensibilidad, reglas de anonimización y copias autorizadas; excluye anonimización masiva y migración. Cada conjunto tiene finalidad, permiso y conservación; se impide uso productivo en DEV y uso de copia no aprobada en QA/PREPROD. Bruno Toro, Seguridad, responde por el resultado y presenta la evidencia; la contraparte mantiene la recepción contractual.\\\hline
\texttt{1.\allowbreak{}7.\allowbreak{}1.\allowbreak{}1}\newline Contratos de interfaces y eventos comunes\newline Fernando Guerrero; 40 HH. & OpenAPI, sobre de eventos, autoridades, idempotencia, correlación y errores; excluye especificar reglas funcionales completas. Contratos tienen versiones y ejemplos de rechazo y reintento y no conceden acceso directo a tablas de otro dominio. Responsable: Fernando Guerrero (Arquitectura/Integración).\\\hline
\texttt{1.\allowbreak{}7.\allowbreak{}1.\allowbreak{}2}\newline Infraestructura de mensajería y DLQ configurada\newline Fernando Guerrero; 48 HH. & EventBridge, SQS/SNS, rutas y dead-letter queues por IaC; excluye código productor/consumidor de cada módulo. Permisos, entrega, reintento, DLQ y monitoreo se comprueban con mensajes de prueba sin duplicar hechos de negocio. Responsable: Fernando Guerrero (Arquitectura/Integración).\\\hline
\texttt{1.\allowbreak{}7.\allowbreak{}1.\allowbreak{}3}\newline Biblioteca de outbox, idempotencia y correlación\newline Fernando Guerrero; 56 HH. & Componente compartido para emisión y recepción verificables; excluye implementar handlers de cada dominio. Duplicados, reintentos y fallo entre escritura/emisión conservan un efecto y recuperación trazable. Responsable: Fernando Guerrero (Arquitectura/Integración).\\\hline
\texttt{1.\allowbreak{}7.\allowbreak{}2.\allowbreak{}1}\newline Sincronizador de lotes y confirmaciones\newline Fernando Guerrero; 64 HH. & Agente y coordinador HTTPS/mTLS con secuencia y ACK; excluye reglas de conciliación particulares y retorno de migración. La cola sólo se elimina tras confirmación; interrupción y reintento preservan orden e integridad y muestran pendientes. Responsable: Fernando Guerrero (Arquitectura/Integración).\\\hline
\texttt{1.\allowbreak{}7.\allowbreak{}2.\allowbreak{}2}\newline Conciliación de autoridad y conflictos implementada\newline Nicolás Torfan; 64 HH. & Reglas comunes de autoridad por agregado, conflictos y revisión; excluye decidir hechos monetarios o cambios de negocio por el cliente. Conflictos se conservan y escalan sin última escritura silenciosa; decisión y replay tienen evidencia. Nicolás Torfan, Datos, responde por el resultado y presenta la evidencia; la contraparte mantiene la recepción contractual.\\\hline
\texttt{1.\allowbreak{}7.\allowbreak{}2.\allowbreak{}3}\newline Ingesta y consolidación de telemetría implementadas\newline Fernando Guerrero; 64 HH. & Consumo local MQTT, lotes cloud y series con calidad; excluye facturación, medición física e IN5. Se validan secuencia, UTC, unidad, calidad y origen; no se suman totalizadores ni se habilita consumo directo del broker desde cloud. Responsable: Fernando Guerrero (Arquitectura/Integración).\\\hline
\texttt{1.\allowbreak{}7.\allowbreak{}3.\allowbreak{}1}\newline Base Android de operación en terreno\newline Nikolai Navea; 56 HH. & Navegación común, sesión y almacenamiento cifrado durable; excluye pantallas y reglas de cada dominio. La sesión limita rol y ámbito; reinicio conserva pendientes protegidos y el canal no accede directo a la base. Nikolai Navea, Desarrollo, responde por el resultado y presenta la evidencia; la contraparte mantiene la recepción contractual.\\\hline
\texttt{1.\allowbreak{}7.\allowbreak{}3.\allowbreak{}2}\newline Cola móvil, evidencia y sincronización de canal\newline Nikolai Navea; 64 HH. & Cola local, fotos, revisión de pendientes y adaptación a sincronizador; excluye registrar inspecciones funcionales. Desconexión, reinicio y ACK tardío no pierden fotos/hechos ni crean duplicados; permisos vencidos bloquean nuevas acciones. Nikolai Navea, Desarrollo, responde por el resultado y presenta la evidencia; la contraparte mantiene la recepción contractual.\\\hline
\texttt{1.\allowbreak{}7.\allowbreak{}3.\allowbreak{}3}\newline Base de portal y consola accesibles\newline Nikolai Navea; 56 HH. & Estructura React/PWA, sesión, navegación, errores y componentes comunes; excluye vistas funcionales de módulos e innovaciones. Se comprueban teclado, foco, mensajes, aislamiento de sesión e idiomas previstos; la base no declara WCAG integral de toda aplicación. Nikolai Navea, Desarrollo, responde por el resultado y presenta la evidencia; la contraparte mantiene la recepción contractual.\\\hline
\texttt{1.\allowbreak{}7.\allowbreak{}4.\allowbreak{}1}\newline Servicio común de notificaciones y preferencias\newline Fernando Guerrero; 56 HH. & Canales, preferencias, evidencias de envío y reintento; excluye decidir destinatarios por las reglas de cada dominio. Un fallo conserva estado y responsable; mensajes repetidos no duplican efectos y cada envío respeta autorización y finalidad. Responsable: Fernando Guerrero (Arquitectura/Integración).\\\hline
\texttt{1.\allowbreak{}7.\allowbreak{}4.\allowbreak{}2}\newline Adaptador con sistema laboral y fuente de identidad\newline Fernando Guerrero; 56 HH. & Intercambio documentado de altas/bajas y habilitación; excluye administrar el sistema laboral externo. Cambios y fallos quedan visibles, se conserva persona-identidad y se rechaza tratar una respuesta vencida como habilitación actual. Responsable: Fernando Guerrero (Arquitectura/Integración).\\\hline
\texttt{1.\allowbreak{}7.\allowbreak{}4.\allowbreak{}3}\newline Adaptador contable y conciliación técnica\newline Fernando Guerrero; 56 HH. & Intercambio de hechos/saldos con autoridad contable externa; excluye reglas de cobranza, emisión tributaria y conciliación de negocio. Acuse, rechazo, reintentos e identificadores se conservan; no se emite un documento tributario desde un canal no autorizado. Responsable: Fernando Guerrero (Arquitectura/Integración).\\\hline
\texttt{1.\allowbreak{}7.\allowbreak{}4.\allowbreak{}4}\newline Adaptador de pagos y webhooks protegidos\newline Fernando Guerrero; 56 HH. & Webhook, autenticidad, idempotencia y estados técnicos; excluye reglas monetarias y reversión de cobros por cliente. Eventos repetidos o fuera de orden no crean otro pago; rechazo, conciliación pendiente y origen quedan trazados. Responsable: Fernando Guerrero (Arquitectura/Integración).\\\hline
\end{longtable}
\endgroup
\FuenteTabla{Elaboración propia. Las horas y los códigos corresponden al catálogo técnico integrado.}


\section{Núcleos funcionales}
Las cuentas 1.4--1.6 desarrollan las capacidades operacionales, de servicios y administrativas. Cada componente distingue especificación, construcción y vistas de tarea, con fronteras y aceptación propias.
\subsection{Componentes funcionales y separación de responsabilidades}
Las cuentas 1.4, 1.5 y 1.6 contienen 30 componentes funcionales y 306 paquetes terminales. La cuenta 1.4 desarrolla Operación marítima, Escuela náutica e Infraestructura; 1.5, Varadero y contratistas, Combustible y Ambiente y consumos; 1.6, Clientes y contratos y Finanzas y cobranza. Auditoría y seguridad permanece en 1.3. Esta organización conserva los nueve dominios de SD3/SD4/SD5: un componente EDT representa trabajo para entregar una capacidad y no crea otro dominio o autoridad de datos.

Cada componente conserva sus tres paquetes originales y añade lotes estimables de análisis/diseño, construcción, verificación y coordinación. Los 90 códigos originales y sus 4.200 HH se acreditan una vez; la selección UCP suma 12.708 HH en 306 paquetes, con 7.766 HH en E1 y 4.942 HH en E2. Los componentes usan las bases comunes 1.2/1.3/1.7 y los datos migrados 1.8. El diccionario de conciliación identifica las HH adicionales y la matriz de actividades muestra el crédito del trabajo original, sin reconstruir servicios comunes.

Los resultados se vinculan con los ámbitos RF de SD3 y las entidades/autoridades de SD5. La comprobación exhaustiva de cobertura deberá relacionar cada ID aplicable de T-12 con el paquete y su prueba, incluidos requisitos transversales y complementos del caso. Una referencia al ámbito NAV, ESC o FAC no equivale a demostrar por sí sola todos sus requisitos; la matriz conservará IDs y enunciados originales \parencite[caps.~19--20]{bases-transversales}.

E1 mantiene movimientos, planes, alertas, amarra compatible, escuela, inspección de boyas, captura de combustible, medición e historia ambiental inicial, junto con los maestros y contratos que requieren esos procesos. E2 completa agenda y faenas, contratistas, residuos, imputación comercial, espera, cobranza y expedientes. La medición temprana conserva su dependencia de instalación y serie anterior a 2029. Las innovaciones conservan su alcance incremental y sus pilotos de SD13, sin sustituir capacidades obligatorias.

\section{Diccionario de los núcleos funcionales}
Estas fichas conservan los 90 paquetes originales y sus 4.200 HH como crédito del trabajo de cada componente. El esfuerzo seleccionado es 12.708 HH y se completa con los lotes del diccionario de conciliación siguiente; las estimaciones originales no representan por sí solas el presupuesto vigente del núcleo. Las fechas individuales y predecesoras vigentes se consultan en el diccionario integrado.

Cada componente tiene tres paquetes terminales: especificación y protocolo; lógica, persistencia y contratos; vistas de tarea y comprobación. El componente agrupa el resultado de negocio y no recibe horas adicionales. La especificación contiene decisiones, datos, reglas, actores, excepciones, ejemplos y protocolo; la lógica implementa esas reglas y su persistencia; las vistas utilizan los canales comunes para la tarea. Se conservan los nueve dominios canónicos y sus autoridades de SD3/SD4/SD5, sin convertir cada componente en un dominio nuevo.

Las estimaciones iniciales de 32, 64 y 44 HH suponen una ampliación acotada sobre identidad, auditoría, persistencia, contratos y canales ya construidos. Son presupuestos de ingeniería por resultado, no mediciones ni tamaños obtenidos mediante UCP. El contraste independiente se presenta en la sección de estimación funcional. Su diferencia permanece abierta; no se atribuye validación de estas HH. Se revisarán en refinamiento y se ajustará la descomposición cuando un resultado no pueda estimarse, asignarse o verificarse de forma independiente. No se aplican porcentajes generales de análisis/programación/pruebas que vuelvan a sumar el mismo trabajo. La especificación reserva 16 HH Funcional para reglas y ejemplos, 8 Datos para atributos y autoridad, 4 Integración para contratos y 4 Seguridad para acceso y sensibilidad. La lógica reserva 48 Desarrollo, 8 Datos, 4 Integración y 4 Calidad. La vista reserva 32 Desarrollo, 4 Funcional, 4 Seguridad y 4 Calidad; incluye unidades y comprobación de tarea, no UAT o carga integral.

 Infraestructura, licencias, procesamiento, reservas y costos de terceros se presupuestan aparte y una sola vez. La migración 1.8 carga los datos existentes; estas fichas construyen el modelo de aplicación y su uso funcional. Las innovaciones 1.9 amplían capacidades adicionales sin duplicar estos resultados. Pruebas integradas, formación general, implantación y servicio permanecen en 1.10--1.12.

La recepción contractual corresponde a la contraparte autorizada. Todas las fichas requieren insumos y versiones recibidos, registros y entornos autorizados, herramientas y contraparte programada. Las HH excluyen dedicación del cliente y cobertura externa continua. La falta de insumos o evidencia bloquea la recepción y exige reprogramar dentro de las condiciones contractuales. Estas condiciones comunes se aplican además de los criterios específicos de cada ficha.

Las fichas de la Tabla~\ref{tab:dic-compacto-nucleos} presentan resultado, responsable, esfuerzo y aceptación. Las condiciones particulares, recursos y ventanas por código se conservan en \path{componentes/condiciones_fichas_t14.csv}, junto con la asignación y precedencias del diccionario técnico integrado. Las actividades de T-15 desarrollan los aportes por perfil y el T-18 desarrolla la ejecución de implantación.


\begingroup\small\setlength{\tabcolsep}{3pt}
\begin{longtable}{L{52mm}L{91mm}}
\caption{Diccionario de núcleos funcionales}\label{tab:dic-compacto-nucleos}\\\hline
 \EncabezadoTabla{Paquete, resultado y responsable} & \EncabezadoTabla{Contenido y criterio de aceptación}\\\hline\endfirsthead
\hline \EncabezadoTabla{Paquete, resultado y responsable} & \EncabezadoTabla{Contenido y criterio de aceptación}\\\hline\endhead
\texttt{1.\allowbreak{}4.\allowbreak{}1.\allowbreak{}1.\allowbreak{}1}\newline Especificación y protocolo: Movimientos de salida y regreso identificados\newline Ulysses Barreto; 32 HH. & Reglas, datos, actores, excepciones, contrato y protocolo de recepción, sin construir software ni copiar todo el modelo de SD5. Exclusiones del componente: Otorgar zarpe de autoridad marítima y búsqueda o salvamento. La especificación permite derivar casos de prueba positivos y de rechazo; datos, autoridad y decisiones pendientes están identificados. Regreso seguro exige confirmación del marinero u observación radial atribuida; declaración del armador no confirma ocupación física ni elimina discrepancias. Ulysses Barreto, Funcional, responde por este paquete; la contraparte valida decisiones de negocio y conserva la recepción contractual. El cumplimiento de tiempos o disponibilidad se verifica después en la prueba integral, sin acreditarlo por la entrega del código.\\\hline
\texttt{1.\allowbreak{}4.\allowbreak{}1.\allowbreak{}1.\allowbreak{}2}\newline Lógica, persistencia y contratos: Movimientos de salida y regreso identificados\newline Nikolai Navea; 64 HH. & Servicio, persistencia y contrato propios del resultado, pruebas unitarias y documentación; excluye vistas y componentes de canal compartidos. Exclusiones del componente: Otorgar zarpe de autoridad marítima y búsqueda o salvamento. La implementación pasa sus pruebas unitarias, restricciones de datos, duplicados y contratos para los escenarios definidos. Regreso seguro exige confirmación del marinero u observación radial atribuida; declaración del armador no confirma ocupación física ni elimina discrepancias. Nikolai Navea, Desarrollo, responde por este paquete; la contraparte valida decisiones de negocio y conserva la recepción contractual. El cumplimiento de tiempos o disponibilidad se verifica después en la prueba integral, sin acreditarlo por la entrega del código.\\\hline
\texttt{1.\allowbreak{}4.\allowbreak{}1.\allowbreak{}1.\allowbreak{}3}\newline Vistas de tarea y comprobación: Movimientos de salida y regreso identificados\newline Nikolai Navea; 44 HH. & Formulario o vista de tarea en canal móvil/portal autorizado, validación funcional y de acceso; reutiliza navegación, sesión y cola de 1.7, sin reconstruirlos. Exclusiones del componente: Otorgar zarpe de autoridad marítima y búsqueda o salvamento. La tarea se reproduce sobre versión identificada con permisos, errores y evidencia visibles; se comprueban teclado/foco y controles adecuados al canal. Regreso seguro exige confirmación del marinero u observación radial atribuida; declaración del armador no confirma ocupación física ni elimina discrepancias. Nikolai Navea, Desarrollo, responde por este paquete; la contraparte valida decisiones de negocio y conserva la recepción contractual. El cumplimiento de tiempos o disponibilidad se verifica después en la prueba integral, sin acreditarlo por la entrega del código.\\\hline
\texttt{1.\allowbreak{}4.\allowbreak{}1.\allowbreak{}2.\allowbreak{}1}\newline Especificación y protocolo: Planes voluntarios y discrepancias operacionales\newline Ulysses Barreto; 32 HH. & Reglas, datos, actores, excepciones, contrato y protocolo de recepción, sin construir software ni copiar todo el modelo de SD5. Exclusiones del componente: Notificación común y medición de ocupación por un sensor como autoridad única. La especificación permite derivar casos de prueba positivos y de rechazo; datos, autoridad y decisiones pendientes están identificados. Plan conserva autor, vencimiento, cierre y discrepancia; cierre tardío no borra una actuación ni reabre un plan cerrado. Ulysses Barreto, Funcional, responde por este paquete; la contraparte valida decisiones de negocio y conserva la recepción contractual. El cumplimiento de tiempos o disponibilidad se verifica después en la prueba integral, sin acreditarlo por la entrega del código.\\\hline
\texttt{1.\allowbreak{}4.\allowbreak{}1.\allowbreak{}2.\allowbreak{}2}\newline Lógica, persistencia y contratos: Planes voluntarios y discrepancias operacionales\newline Nikolai Navea; 64 HH. & Servicio, persistencia y contrato propios del resultado, pruebas unitarias y documentación; excluye vistas y componentes de canal compartidos. Exclusiones del componente: Notificación común y medición de ocupación por un sensor como autoridad única. La implementación pasa sus pruebas unitarias, restricciones de datos, duplicados y contratos para los escenarios definidos. Plan conserva autor, vencimiento, cierre y discrepancia; cierre tardío no borra una actuación ni reabre un plan cerrado. Nikolai Navea, Desarrollo, responde por este paquete; la contraparte valida decisiones de negocio y conserva la recepción contractual. El cumplimiento de tiempos o disponibilidad se verifica después en la prueba integral, sin acreditarlo por la entrega del código.\\\hline
\texttt{1.\allowbreak{}4.\allowbreak{}1.\allowbreak{}2.\allowbreak{}3}\newline Vistas de tarea y comprobación: Planes voluntarios y discrepancias operacionales\newline Nikolai Navea; 44 HH. & Formulario o vista de tarea en canal móvil/portal autorizado, validación funcional y de acceso; reutiliza navegación, sesión y cola de 1.7, sin reconstruirlos. Exclusiones del componente: Notificación común y medición de ocupación por un sensor como autoridad única. La tarea se reproduce sobre versión identificada con permisos, errores y evidencia visibles; se comprueban teclado/foco y controles adecuados al canal. Plan conserva autor, vencimiento, cierre y discrepancia; cierre tardío no borra una actuación ni reabre un plan cerrado. Nikolai Navea, Desarrollo, responde por este paquete; la contraparte valida decisiones de negocio y conserva la recepción contractual. El cumplimiento de tiempos o disponibilidad se verifica después en la prueba integral, sin acreditarlo por la entrega del código.\\\hline
\texttt{1.\allowbreak{}4.\allowbreak{}1.\allowbreak{}3.\allowbreak{}1}\newline Especificación y protocolo: Atención y escalamiento de planes vencidos\newline Ulysses Barreto; 32 HH. & Reglas, datos, actores, excepciones, contrato y protocolo de recepción, sin construir software ni copiar todo el modelo de SD5. Exclusiones del componente: Construcción del servicio común de notificaciones y ejecución de salvamento. La especificación permite derivar casos de prueba positivos y de rechazo; datos, autoridad y decisiones pendientes están identificados. Vencimiento genera solicitud de alerta para guardia dentro de quince minutos, sin gracia; contacto y escalamiento conservan responsable y resultado. Ulysses Barreto, Funcional, responde por este paquete; la contraparte valida decisiones de negocio y conserva la recepción contractual. El cumplimiento de tiempos o disponibilidad se verifica después en la prueba integral, sin acreditarlo por la entrega del código.\\\hline
\texttt{1.\allowbreak{}4.\allowbreak{}1.\allowbreak{}3.\allowbreak{}2}\newline Lógica, persistencia y contratos: Atención y escalamiento de planes vencidos\newline Nikolai Navea; 64 HH. & Servicio, persistencia y contrato propios del resultado, pruebas unitarias y documentación; excluye vistas y componentes de canal compartidos. Exclusiones del componente: Construcción del servicio común de notificaciones y ejecución de salvamento. La implementación pasa sus pruebas unitarias, restricciones de datos, duplicados y contratos para los escenarios definidos. Vencimiento genera solicitud de alerta para guardia dentro de quince minutos, sin gracia; contacto y escalamiento conservan responsable y resultado. Nikolai Navea, Desarrollo, responde por este paquete; la contraparte valida decisiones de negocio y conserva la recepción contractual. El cumplimiento de tiempos o disponibilidad se verifica después en la prueba integral, sin acreditarlo por la entrega del código.\\\hline
\texttt{1.\allowbreak{}4.\allowbreak{}1.\allowbreak{}3.\allowbreak{}3}\newline Vistas de tarea y comprobación: Atención y escalamiento de planes vencidos\newline Nikolai Navea; 44 HH. & Formulario o vista de tarea en canal móvil/portal autorizado, validación funcional y de acceso; reutiliza navegación, sesión y cola de 1.7, sin reconstruirlos. Exclusiones del componente: Construcción del servicio común de notificaciones y ejecución de salvamento. La tarea se reproduce sobre versión identificada con permisos, errores y evidencia visibles; se comprueban teclado/foco y controles adecuados al canal. Vencimiento genera solicitud de alerta para guardia dentro de quince minutos, sin gracia; contacto y escalamiento conservan responsable y resultado. Nikolai Navea, Desarrollo, responde por este paquete; la contraparte valida decisiones de negocio y conserva la recepción contractual. El cumplimiento de tiempos o disponibilidad se verifica después en la prueba integral, sin acreditarlo por la entrega del código.\\\hline
\texttt{1.\allowbreak{}4.\allowbreak{}1.\allowbreak{}4.\allowbreak{}1}\newline Especificación y protocolo: Autorización operacional de amarra y compatibilidad\newline Ulysses Barreto; 32 HH. & Reglas, datos, actores, excepciones, contrato y protocolo de recepción, sin construir software ni copiar todo el modelo de SD5. Exclusiones del componente: Gestión comercial del contrato y considerar reserva como ocupación comprobada. La especificación permite derivar casos de prueba positivos y de rechazo; datos, autoridad y decisiones pendientes están identificados. Una autoridad local serializa puesto e intervalo y rechaza doble asignación; compatibilidad y última observación son explícitas. Ulysses Barreto, Funcional, responde por este paquete; la contraparte valida decisiones de negocio y conserva la recepción contractual. El cumplimiento de tiempos o disponibilidad se verifica después en la prueba integral, sin acreditarlo por la entrega del código.\\\hline
\texttt{1.\allowbreak{}4.\allowbreak{}1.\allowbreak{}4.\allowbreak{}2}\newline Lógica, persistencia y contratos: Autorización operacional de amarra y compatibilidad\newline Nikolai Navea; 64 HH. & Servicio, persistencia y contrato propios del resultado, pruebas unitarias y documentación; excluye vistas y componentes de canal compartidos. Exclusiones del componente: Gestión comercial del contrato y considerar reserva como ocupación comprobada. La implementación pasa sus pruebas unitarias, restricciones de datos, duplicados y contratos para los escenarios definidos. Una autoridad local serializa puesto e intervalo y rechaza doble asignación; compatibilidad y última observación son explícitas. Nikolai Navea, Desarrollo, responde por este paquete; la contraparte valida decisiones de negocio y conserva la recepción contractual. El cumplimiento de tiempos o disponibilidad se verifica después en la prueba integral, sin acreditarlo por la entrega del código.\\\hline
\texttt{1.\allowbreak{}4.\allowbreak{}1.\allowbreak{}4.\allowbreak{}3}\newline Vistas de tarea y comprobación: Autorización operacional de amarra y compatibilidad\newline Nikolai Navea; 44 HH. & Formulario o vista de tarea en canal móvil/portal autorizado, validación funcional y de acceso; reutiliza navegación, sesión y cola de 1.7, sin reconstruirlos. Exclusiones del componente: Gestión comercial del contrato y considerar reserva como ocupación comprobada. La tarea se reproduce sobre versión identificada con permisos, errores y evidencia visibles; se comprueban teclado/foco y controles adecuados al canal. Una autoridad local serializa puesto e intervalo y rechaza doble asignación; compatibilidad y última observación son explícitas. Nikolai Navea, Desarrollo, responde por este paquete; la contraparte valida decisiones de negocio y conserva la recepción contractual. El cumplimiento de tiempos o disponibilidad se verifica después en la prueba integral, sin acreditarlo por la entrega del código.\\\hline
\texttt{1.\allowbreak{}4.\allowbreak{}2.\allowbreak{}1.\allowbreak{}1}\newline Especificación y protocolo: Clases, nóminas y participación por intervalo\newline Ulysses Barreto; 32 HH. & Reglas, datos, actores, excepciones, contrato y protocolo de recepción, sin construir software ni copiar todo el modelo de SD5. Exclusiones del componente: Matrícula y autorizaciones migradas, plataforma de canal y prueba integral de carga. La especificación permite derivar casos de prueba positivos y de rechazo; datos, autoridad y decisiones pendientes están identificados. Alumno, bote y monitor se vinculan por intervalo con cambios trazados; hasta 18 embarcaciones salen en 60 segundos en protocolo previsto. Ulysses Barreto, Funcional, responde por este paquete; la contraparte valida decisiones de negocio y conserva la recepción contractual. El cumplimiento de tiempos o disponibilidad se verifica después en la prueba integral, sin acreditarlo por la entrega del código.\\\hline
\texttt{1.\allowbreak{}4.\allowbreak{}2.\allowbreak{}1.\allowbreak{}2}\newline Lógica, persistencia y contratos: Clases, nóminas y participación por intervalo\newline Nikolai Navea; 64 HH. & Servicio, persistencia y contrato propios del resultado, pruebas unitarias y documentación; excluye vistas y componentes de canal compartidos. Exclusiones del componente: Matrícula y autorizaciones migradas, plataforma de canal y prueba integral de carga. La implementación pasa sus pruebas unitarias, restricciones de datos, duplicados y contratos para los escenarios definidos. Alumno, bote y monitor se vinculan por intervalo con cambios trazados; hasta 18 embarcaciones salen en 60 segundos en protocolo previsto. Nikolai Navea, Desarrollo, responde por este paquete; la contraparte valida decisiones de negocio y conserva la recepción contractual. El cumplimiento de tiempos o disponibilidad se verifica después en la prueba integral, sin acreditarlo por la entrega del código.\\\hline
\texttt{1.\allowbreak{}4.\allowbreak{}2.\allowbreak{}1.\allowbreak{}3}\newline Vistas de tarea y comprobación: Clases, nóminas y participación por intervalo\newline Nikolai Navea; 44 HH. & Formulario o vista de tarea en canal móvil/portal autorizado, validación funcional y de acceso; reutiliza navegación, sesión y cola de 1.7, sin reconstruirlos. Exclusiones del componente: Matrícula y autorizaciones migradas, plataforma de canal y prueba integral de carga. La tarea se reproduce sobre versión identificada con permisos, errores y evidencia visibles; se comprueban teclado/foco y controles adecuados al canal. Alumno, bote y monitor se vinculan por intervalo con cambios trazados; hasta 18 embarcaciones salen en 60 segundos en protocolo previsto. Nikolai Navea, Desarrollo, responde por este paquete; la contraparte valida decisiones de negocio y conserva la recepción contractual. El cumplimiento de tiempos o disponibilidad se verifica después en la prueba integral, sin acreditarlo por la entrega del código.\\\hline
\texttt{1.\allowbreak{}4.\allowbreak{}2.\allowbreak{}2.\allowbreak{}1}\newline Especificación y protocolo: Retiros coordinados y entrega única de alumno\newline Ulysses Barreto; 32 HH. & Reglas, datos, actores, excepciones, contrato y protocolo de recepción, sin construir software ni copiar todo el modelo de SD5. Exclusiones del componente: Solicitud del apoderado tratada como entrega efectiva y reemplazar coordinación segura por radio. La especificación permite derivar casos de prueba positivos y de rechazo; datos, autoridad y decisiones pendientes están identificados. Adulto e identidad/autorización se verifican; se confirma una sola entrega efectiva y se bloquea confirmación contradictoria sin estado comprobable. Ulysses Barreto, Funcional, responde por este paquete; la contraparte valida decisiones de negocio y conserva la recepción contractual. El cumplimiento de tiempos o disponibilidad se verifica después en la prueba integral, sin acreditarlo por la entrega del código.\\\hline
\texttt{1.\allowbreak{}4.\allowbreak{}2.\allowbreak{}2.\allowbreak{}2}\newline Lógica, persistencia y contratos: Retiros coordinados y entrega única de alumno\newline Nikolai Navea; 64 HH. & Servicio, persistencia y contrato propios del resultado, pruebas unitarias y documentación; excluye vistas y componentes de canal compartidos. Exclusiones del componente: Solicitud del apoderado tratada como entrega efectiva y reemplazar coordinación segura por radio. La implementación pasa sus pruebas unitarias, restricciones de datos, duplicados y contratos para los escenarios definidos. Adulto e identidad/autorización se verifican; se confirma una sola entrega efectiva y se bloquea confirmación contradictoria sin estado comprobable. Nikolai Navea, Desarrollo, responde por este paquete; la contraparte valida decisiones de negocio y conserva la recepción contractual. El cumplimiento de tiempos o disponibilidad se verifica después en la prueba integral, sin acreditarlo por la entrega del código.\\\hline
\texttt{1.\allowbreak{}4.\allowbreak{}2.\allowbreak{}2.\allowbreak{}3}\newline Vistas de tarea y comprobación: Retiros coordinados y entrega única de alumno\newline Nikolai Navea; 44 HH. & Formulario o vista de tarea en canal móvil/portal autorizado, validación funcional y de acceso; reutiliza navegación, sesión y cola de 1.7, sin reconstruirlos. Exclusiones del componente: Solicitud del apoderado tratada como entrega efectiva y reemplazar coordinación segura por radio. La tarea se reproduce sobre versión identificada con permisos, errores y evidencia visibles; se comprueban teclado/foco y controles adecuados al canal. Adulto e identidad/autorización se verifican; se confirma una sola entrega efectiva y se bloquea confirmación contradictoria sin estado comprobable. Nikolai Navea, Desarrollo, responde por este paquete; la contraparte valida decisiones de negocio y conserva la recepción contractual. El cumplimiento de tiempos o disponibilidad se verifica después en la prueba integral, sin acreditarlo por la entrega del código.\\\hline
\texttt{1.\allowbreak{}4.\allowbreak{}2.\allowbreak{}3.\allowbreak{}1}\newline Especificación y protocolo: Retorno de participantes y vista mínima de salud\newline Ulysses Barreto; 32 HH. & Reglas, datos, actores, excepciones, contrato y protocolo de recepción, sin construir software ni copiar todo el modelo de SD5. Exclusiones del componente: Difusión completa de ficha de salud y reconstrucción de datos ausentes. La especificación permite derivar casos de prueba positivos y de rechazo; datos, autoridad y decisiones pendientes están identificados. Cada participante se confirma antes de cerrar clase; monitor ve sólo antecedentes mínimos de alumnos vinculados y se solicita aviso al apoderado. Ulysses Barreto, Funcional, responde por este paquete; la contraparte valida decisiones de negocio y conserva la recepción contractual. El cumplimiento de tiempos o disponibilidad se verifica después en la prueba integral, sin acreditarlo por la entrega del código.\\\hline
\texttt{1.\allowbreak{}4.\allowbreak{}2.\allowbreak{}3.\allowbreak{}2}\newline Lógica, persistencia y contratos: Retorno de participantes y vista mínima de salud\newline Nikolai Navea; 64 HH. & Servicio, persistencia y contrato propios del resultado, pruebas unitarias y documentación; excluye vistas y componentes de canal compartidos. Exclusiones del componente: Difusión completa de ficha de salud y reconstrucción de datos ausentes. La implementación pasa sus pruebas unitarias, restricciones de datos, duplicados y contratos para los escenarios definidos. Cada participante se confirma antes de cerrar clase; monitor ve sólo antecedentes mínimos de alumnos vinculados y se solicita aviso al apoderado. Nikolai Navea, Desarrollo, responde por este paquete; la contraparte valida decisiones de negocio y conserva la recepción contractual. El cumplimiento de tiempos o disponibilidad se verifica después en la prueba integral, sin acreditarlo por la entrega del código.\\\hline
\texttt{1.\allowbreak{}4.\allowbreak{}2.\allowbreak{}3.\allowbreak{}3}\newline Vistas de tarea y comprobación: Retorno de participantes y vista mínima de salud\newline Nikolai Navea; 44 HH. & Formulario o vista de tarea en canal móvil/portal autorizado, validación funcional y de acceso; reutiliza navegación, sesión y cola de 1.7, sin reconstruirlos. Exclusiones del componente: Difusión completa de ficha de salud y reconstrucción de datos ausentes. La tarea se reproduce sobre versión identificada con permisos, errores y evidencia visibles; se comprueban teclado/foco y controles adecuados al canal. Cada participante se confirma antes de cerrar clase; monitor ve sólo antecedentes mínimos de alumnos vinculados y se solicita aviso al apoderado. Nikolai Navea, Desarrollo, responde por este paquete; la contraparte valida decisiones de negocio y conserva la recepción contractual. El cumplimiento de tiempos o disponibilidad se verifica después en la prueba integral, sin acreditarlo por la entrega del código.\\\hline
\texttt{1.\allowbreak{}4.\allowbreak{}3.\allowbreak{}1.\allowbreak{}1}\newline Especificación y protocolo: Inventario y condiciones de puestos y fondeos\newline Ulysses Barreto; 32 HH. & Reglas, datos, actores, excepciones, contrato y protocolo de recepción, sin construir software ni copiar todo el modelo de SD5. Exclusiones del componente: Instrumentación física adicional o piloto no seleccionado. La especificación permite derivar casos de prueba positivos y de rechazo; datos, autoridad y decisiones pendientes están identificados. Puesto, boya y tren de fondeo tienen identidad, condición y última observación con antigüedad visible; una reserva no prueba presencia. Ulysses Barreto, Funcional, responde por este paquete; la contraparte valida decisiones de negocio y conserva la recepción contractual. El cumplimiento de tiempos o disponibilidad se verifica después en la prueba integral, sin acreditarlo por la entrega del código.\\\hline
\texttt{1.\allowbreak{}4.\allowbreak{}3.\allowbreak{}1.\allowbreak{}2}\newline Lógica, persistencia y contratos: Inventario y condiciones de puestos y fondeos\newline Nikolai Navea; 64 HH. & Servicio, persistencia y contrato propios del resultado, pruebas unitarias y documentación; excluye vistas y componentes de canal compartidos. Exclusiones del componente: Instrumentación física adicional o piloto no seleccionado. La implementación pasa sus pruebas unitarias, restricciones de datos, duplicados y contratos para los escenarios definidos. Puesto, boya y tren de fondeo tienen identidad, condición y última observación con antigüedad visible; una reserva no prueba presencia. Nikolai Navea, Desarrollo, responde por este paquete; la contraparte valida decisiones de negocio y conserva la recepción contractual. El cumplimiento de tiempos o disponibilidad se verifica después en la prueba integral, sin acreditarlo por la entrega del código.\\\hline
\texttt{1.\allowbreak{}4.\allowbreak{}3.\allowbreak{}1.\allowbreak{}3}\newline Vistas de tarea y comprobación: Inventario y condiciones de puestos y fondeos\newline Nikolai Navea; 44 HH. & Formulario o vista de tarea en canal móvil/portal autorizado, validación funcional y de acceso; reutiliza navegación, sesión y cola de 1.7, sin reconstruirlos. Exclusiones del componente: Instrumentación física adicional o piloto no seleccionado. La tarea se reproduce sobre versión identificada con permisos, errores y evidencia visibles; se comprueban teclado/foco y controles adecuados al canal. Puesto, boya y tren de fondeo tienen identidad, condición y última observación con antigüedad visible; una reserva no prueba presencia. Nikolai Navea, Desarrollo, responde por este paquete; la contraparte valida decisiones de negocio y conserva la recepción contractual. El cumplimiento de tiempos o disponibilidad se verifica después en la prueba integral, sin acreditarlo por la entrega del código.\\\hline
\texttt{1.\allowbreak{}4.\allowbreak{}3.\allowbreak{}2.\allowbreak{}1}\newline Especificación y protocolo: Inspecciones de boyas y evidencia desconectada\newline Ulysses Barreto; 32 HH. & Reglas, datos, actores, excepciones, contrato y protocolo de recepción, sin construir software ni copiar todo el modelo de SD5. Exclusiones del componente: Construcción de cola móvil compartida y suministro de energía/comunicaciones no seleccionadas. La especificación permite derivar casos de prueba positivos y de rechazo; datos, autoridad y decisiones pendientes están identificados. Observación y fotos conservan activo, autor, tiempo y pendientes; inspección de ocho horas sin cobertura se prepara para comprobación integral. Ulysses Barreto, Funcional, responde por este paquete; la contraparte valida decisiones de negocio y conserva la recepción contractual. El cumplimiento de tiempos o disponibilidad se verifica después en la prueba integral, sin acreditarlo por la entrega del código.\\\hline
\texttt{1.\allowbreak{}4.\allowbreak{}3.\allowbreak{}2.\allowbreak{}2}\newline Lógica, persistencia y contratos: Inspecciones de boyas y evidencia desconectada\newline Nikolai Navea; 64 HH. & Servicio, persistencia y contrato propios del resultado, pruebas unitarias y documentación; excluye vistas y componentes de canal compartidos. Exclusiones del componente: Construcción de cola móvil compartida y suministro de energía/comunicaciones no seleccionadas. La implementación pasa sus pruebas unitarias, restricciones de datos, duplicados y contratos para los escenarios definidos. Observación y fotos conservan activo, autor, tiempo y pendientes; inspección de ocho horas sin cobertura se prepara para comprobación integral. Nikolai Navea, Desarrollo, responde por este paquete; la contraparte valida decisiones de negocio y conserva la recepción contractual. El cumplimiento de tiempos o disponibilidad se verifica después en la prueba integral, sin acreditarlo por la entrega del código.\\\hline
\texttt{1.\allowbreak{}4.\allowbreak{}3.\allowbreak{}2.\allowbreak{}3}\newline Vistas de tarea y comprobación: Inspecciones de boyas y evidencia desconectada\newline Nikolai Navea; 44 HH. & Formulario o vista de tarea en canal móvil/portal autorizado, validación funcional y de acceso; reutiliza navegación, sesión y cola de 1.7, sin reconstruirlos. Exclusiones del componente: Construcción de cola móvil compartida y suministro de energía/comunicaciones no seleccionadas. La tarea se reproduce sobre versión identificada con permisos, errores y evidencia visibles; se comprueban teclado/foco y controles adecuados al canal. Observación y fotos conservan activo, autor, tiempo y pendientes; inspección de ocho horas sin cobertura se prepara para comprobación integral. Nikolai Navea, Desarrollo, responde por este paquete; la contraparte valida decisiones de negocio y conserva la recepción contractual. El cumplimiento de tiempos o disponibilidad se verifica después en la prueba integral, sin acreditarlo por la entrega del código.\\\hline
\texttt{1.\allowbreak{}4.\allowbreak{}3.\allowbreak{}3.\allowbreak{}1}\newline Especificación y protocolo: Asociaciones temporales de medidores e instalaciones\newline Ulysses Barreto; 32 HH. & Reglas, datos, actores, excepciones, contrato y protocolo de recepción, sin construir software ni copiar todo el modelo de SD5. Exclusiones del componente: Cálculo del consumo, adquisición e instalación física de medidores. La especificación permite derivar casos de prueba positivos y de rechazo; datos, autoridad y decisiones pendientes están identificados. Un medidor queda asociado a servicio/amarra por período sin superposición; sustitución y época de contador conservan historial. Ulysses Barreto, Funcional, responde por este paquete; la contraparte valida decisiones de negocio y conserva la recepción contractual. El cumplimiento de tiempos o disponibilidad se verifica después en la prueba integral, sin acreditarlo por la entrega del código.\\\hline
\texttt{1.\allowbreak{}4.\allowbreak{}3.\allowbreak{}3.\allowbreak{}2}\newline Lógica, persistencia y contratos: Asociaciones temporales de medidores e instalaciones\newline Nikolai Navea; 64 HH. & Servicio, persistencia y contrato propios del resultado, pruebas unitarias y documentación; excluye vistas y componentes de canal compartidos. Exclusiones del componente: Cálculo del consumo, adquisición e instalación física de medidores. La implementación pasa sus pruebas unitarias, restricciones de datos, duplicados y contratos para los escenarios definidos. Un medidor queda asociado a servicio/amarra por período sin superposición; sustitución y época de contador conservan historial. Nikolai Navea, Desarrollo, responde por este paquete; la contraparte valida decisiones de negocio y conserva la recepción contractual. El cumplimiento de tiempos o disponibilidad se verifica después en la prueba integral, sin acreditarlo por la entrega del código.\\\hline
\texttt{1.\allowbreak{}4.\allowbreak{}3.\allowbreak{}3.\allowbreak{}3}\newline Vistas de tarea y comprobación: Asociaciones temporales de medidores e instalaciones\newline Nikolai Navea; 44 HH. & Formulario o vista de tarea en canal móvil/portal autorizado, validación funcional y de acceso; reutiliza navegación, sesión y cola de 1.7, sin reconstruirlos. Exclusiones del componente: Cálculo del consumo, adquisición e instalación física de medidores. La tarea se reproduce sobre versión identificada con permisos, errores y evidencia visibles; se comprueban teclado/foco y controles adecuados al canal. Un medidor queda asociado a servicio/amarra por período sin superposición; sustitución y época de contador conservan historial. Nikolai Navea, Desarrollo, responde por este paquete; la contraparte valida decisiones de negocio y conserva la recepción contractual. El cumplimiento de tiempos o disponibilidad se verifica después en la prueba integral, sin acreditarlo por la entrega del código.\\\hline
\texttt{1.\allowbreak{}5.\allowbreak{}1.\allowbreak{}1.\allowbreak{}1}\newline Especificación y protocolo: Maestros de faena y habilitaciones de consulta inicial\newline Ulysses Barreto; 32 HH. & Reglas, datos, actores, excepciones, contrato y protocolo de recepción, sin construir software ni copiar todo el modelo de SD5. Exclusiones del componente: Agenda y órdenes completas E2 y fuente laboral externa. La especificación permite derivar casos de prueba positivos y de rechazo; datos, autoridad y decisiones pendientes están identificados. Activos, contratistas y vigencias necesarios para E1 tienen identidad y autoridad; una copia visible no autoriza una faena. Ulysses Barreto, Funcional, responde por este paquete; la contraparte valida decisiones de negocio y conserva la recepción contractual. El cumplimiento de tiempos o disponibilidad se verifica después en la prueba integral, sin acreditarlo por la entrega del código.\\\hline
\texttt{1.\allowbreak{}5.\allowbreak{}1.\allowbreak{}1.\allowbreak{}2}\newline Lógica, persistencia y contratos: Maestros de faena y habilitaciones de consulta inicial\newline Nikolai Navea; 64 HH. & Servicio, persistencia y contrato propios del resultado, pruebas unitarias y documentación; excluye vistas y componentes de canal compartidos. Exclusiones del componente: Agenda y órdenes completas E2 y fuente laboral externa. La implementación pasa sus pruebas unitarias, restricciones de datos, duplicados y contratos para los escenarios definidos. Activos, contratistas y vigencias necesarios para E1 tienen identidad y autoridad; una copia visible no autoriza una faena. Nikolai Navea, Desarrollo, responde por este paquete; la contraparte valida decisiones de negocio y conserva la recepción contractual. El cumplimiento de tiempos o disponibilidad se verifica después en la prueba integral, sin acreditarlo por la entrega del código.\\\hline
\texttt{1.\allowbreak{}5.\allowbreak{}1.\allowbreak{}1.\allowbreak{}3}\newline Vistas de tarea y comprobación: Maestros de faena y habilitaciones de consulta inicial\newline Nikolai Navea; 44 HH. & Formulario o vista de tarea en canal móvil/portal autorizado, validación funcional y de acceso; reutiliza navegación, sesión y cola de 1.7, sin reconstruirlos. Exclusiones del componente: Agenda y órdenes completas E2 y fuente laboral externa. La tarea se reproduce sobre versión identificada con permisos, errores y evidencia visibles; se comprueban teclado/foco y controles adecuados al canal. Activos, contratistas y vigencias necesarios para E1 tienen identidad y autoridad; una copia visible no autoriza una faena. Nikolai Navea, Desarrollo, responde por este paquete; la contraparte valida decisiones de negocio y conserva la recepción contractual. El cumplimiento de tiempos o disponibilidad se verifica después en la prueba integral, sin acreditarlo por la entrega del código.\\\hline
\texttt{1.\allowbreak{}5.\allowbreak{}1.\allowbreak{}2.\allowbreak{}1}\newline Especificación y protocolo: Agenda de varadero y posiciones compatibles\newline Ulysses Barreto; 32 HH. & Reglas, datos, actores, excepciones, contrato y protocolo de recepción, sin construir software ni copiar todo el modelo de SD5. Exclusiones del componente: Inventario de puestos de Infraestructura y asignar capacidad adicional no prevista. La especificación permite derivar casos de prueba positivos y de rechazo; datos, autoridad y decisiones pendientes están identificados. Agenda comprueba posición, recurso, intervalo y dependencias, y conserva autor de cambios y reprogramaciones. Ulysses Barreto, Funcional, responde por este paquete; la contraparte valida decisiones de negocio y conserva la recepción contractual. El cumplimiento de tiempos o disponibilidad se verifica después en la prueba integral, sin acreditarlo por la entrega del código.\\\hline
\texttt{1.\allowbreak{}5.\allowbreak{}1.\allowbreak{}2.\allowbreak{}2}\newline Lógica, persistencia y contratos: Agenda de varadero y posiciones compatibles\newline Nikolai Navea; 64 HH. & Servicio, persistencia y contrato propios del resultado, pruebas unitarias y documentación; excluye vistas y componentes de canal compartidos. Exclusiones del componente: Inventario de puestos de Infraestructura y asignar capacidad adicional no prevista. La implementación pasa sus pruebas unitarias, restricciones de datos, duplicados y contratos para los escenarios definidos. Agenda comprueba posición, recurso, intervalo y dependencias, y conserva autor de cambios y reprogramaciones. Nikolai Navea, Desarrollo, responde por este paquete; la contraparte valida decisiones de negocio y conserva la recepción contractual. El cumplimiento de tiempos o disponibilidad se verifica después en la prueba integral, sin acreditarlo por la entrega del código.\\\hline
\texttt{1.\allowbreak{}5.\allowbreak{}1.\allowbreak{}2.\allowbreak{}3}\newline Vistas de tarea y comprobación: Agenda de varadero y posiciones compatibles\newline Nikolai Navea; 44 HH. & Formulario o vista de tarea en canal móvil/portal autorizado, validación funcional y de acceso; reutiliza navegación, sesión y cola de 1.7, sin reconstruirlos. Exclusiones del componente: Inventario de puestos de Infraestructura y asignar capacidad adicional no prevista. La tarea se reproduce sobre versión identificada con permisos, errores y evidencia visibles; se comprueban teclado/foco y controles adecuados al canal. Agenda comprueba posición, recurso, intervalo y dependencias, y conserva autor de cambios y reprogramaciones. Nikolai Navea, Desarrollo, responde por este paquete; la contraparte valida decisiones de negocio y conserva la recepción contractual. El cumplimiento de tiempos o disponibilidad se verifica después en la prueba integral, sin acreditarlo por la entrega del código.\\\hline
\texttt{1.\allowbreak{}5.\allowbreak{}1.\allowbreak{}3.\allowbreak{}1}\newline Especificación y protocolo: Plan de izaje, validación y suspensión de faena\newline Ulysses Barreto; 32 HH. & Reglas, datos, actores, excepciones, contrato y protocolo de recepción, sin construir software ni copiar todo el modelo de SD5. Exclusiones del componente: Enclavamiento físico del travelift y sustituir la decisión de un operador calificado. La especificación permite derivar casos de prueba positivos y de rechazo; datos, autoridad y decisiones pendientes están identificados. Plan identifica eslingado, restricciones y validador calificado; aviso oficial suspende; no se exige interacción con carga suspendida. Ulysses Barreto, Funcional, responde por este paquete; la contraparte valida decisiones de negocio y conserva la recepción contractual. El cumplimiento de tiempos o disponibilidad se verifica después en la prueba integral, sin acreditarlo por la entrega del código.\\\hline
\texttt{1.\allowbreak{}5.\allowbreak{}1.\allowbreak{}3.\allowbreak{}2}\newline Lógica, persistencia y contratos: Plan de izaje, validación y suspensión de faena\newline Nikolai Navea; 64 HH. & Servicio, persistencia y contrato propios del resultado, pruebas unitarias y documentación; excluye vistas y componentes de canal compartidos. Exclusiones del componente: Enclavamiento físico del travelift y sustituir la decisión de un operador calificado. La implementación pasa sus pruebas unitarias, restricciones de datos, duplicados y contratos para los escenarios definidos. Plan identifica eslingado, restricciones y validador calificado; aviso oficial suspende; no se exige interacción con carga suspendida. Nikolai Navea, Desarrollo, responde por este paquete; la contraparte valida decisiones de negocio y conserva la recepción contractual. El cumplimiento de tiempos o disponibilidad se verifica después en la prueba integral, sin acreditarlo por la entrega del código.\\\hline
\texttt{1.\allowbreak{}5.\allowbreak{}1.\allowbreak{}3.\allowbreak{}3}\newline Vistas de tarea y comprobación: Plan de izaje, validación y suspensión de faena\newline Nikolai Navea; 44 HH. & Formulario o vista de tarea en canal móvil/portal autorizado, validación funcional y de acceso; reutiliza navegación, sesión y cola de 1.7, sin reconstruirlos. Exclusiones del componente: Enclavamiento físico del travelift y sustituir la decisión de un operador calificado. La tarea se reproduce sobre versión identificada con permisos, errores y evidencia visibles; se comprueban teclado/foco y controles adecuados al canal. Plan identifica eslingado, restricciones y validador calificado; aviso oficial suspende; no se exige interacción con carga suspendida. Nikolai Navea, Desarrollo, responde por este paquete; la contraparte valida decisiones de negocio y conserva la recepción contractual. El cumplimiento de tiempos o disponibilidad se verifica después en la prueba integral, sin acreditarlo por la entrega del código.\\\hline
\texttt{1.\allowbreak{}5.\allowbreak{}1.\allowbreak{}4.\allowbreak{}1}\newline Especificación y protocolo: Órdenes, ejecución y reprogramación de trabajos\newline Ulysses Barreto; 32 HH. & Reglas, datos, actores, excepciones, contrato y protocolo de recepción, sin construir software ni copiar todo el modelo de SD5. Exclusiones del componente: Solicitudes adicionales de mantenimiento IN1 y mantenimiento preventivo IN2. La especificación permite derivar casos de prueba positivos y de rechazo; datos, autoridad y decisiones pendientes están identificados. Orden conserva aceptación, actuación, evidencia y cambio de calendario; cada reprogramación comprueba posiciones y dependencias antes de notificar. Ulysses Barreto, Funcional, responde por este paquete; la contraparte valida decisiones de negocio y conserva la recepción contractual. El cumplimiento de tiempos o disponibilidad se verifica después en la prueba integral, sin acreditarlo por la entrega del código.\\\hline
\texttt{1.\allowbreak{}5.\allowbreak{}1.\allowbreak{}4.\allowbreak{}2}\newline Lógica, persistencia y contratos: Órdenes, ejecución y reprogramación de trabajos\newline Nikolai Navea; 64 HH. & Servicio, persistencia y contrato propios del resultado, pruebas unitarias y documentación; excluye vistas y componentes de canal compartidos. Exclusiones del componente: Solicitudes adicionales de mantenimiento IN1 y mantenimiento preventivo IN2. La implementación pasa sus pruebas unitarias, restricciones de datos, duplicados y contratos para los escenarios definidos. Orden conserva aceptación, actuación, evidencia y cambio de calendario; cada reprogramación comprueba posiciones y dependencias antes de notificar. Nikolai Navea, Desarrollo, responde por este paquete; la contraparte valida decisiones de negocio y conserva la recepción contractual. El cumplimiento de tiempos o disponibilidad se verifica después en la prueba integral, sin acreditarlo por la entrega del código.\\\hline
\texttt{1.\allowbreak{}5.\allowbreak{}1.\allowbreak{}4.\allowbreak{}3}\newline Vistas de tarea y comprobación: Órdenes, ejecución y reprogramación de trabajos\newline Nikolai Navea; 44 HH. & Formulario o vista de tarea en canal móvil/portal autorizado, validación funcional y de acceso; reutiliza navegación, sesión y cola de 1.7, sin reconstruirlos. Exclusiones del componente: Solicitudes adicionales de mantenimiento IN1 y mantenimiento preventivo IN2. La tarea se reproduce sobre versión identificada con permisos, errores y evidencia visibles; se comprueban teclado/foco y controles adecuados al canal. Orden conserva aceptación, actuación, evidencia y cambio de calendario; cada reprogramación comprueba posiciones y dependencias antes de notificar. Nikolai Navea, Desarrollo, responde por este paquete; la contraparte valida decisiones de negocio y conserva la recepción contractual. El cumplimiento de tiempos o disponibilidad se verifica después en la prueba integral, sin acreditarlo por la entrega del código.\\\hline
\texttt{1.\allowbreak{}5.\allowbreak{}1.\allowbreak{}5.\allowbreak{}1}\newline Especificación y protocolo: Acreditación, inducción y acceso a intervención\newline Ulysses Barreto; 32 HH. & Reglas, datos, actores, excepciones, contrato y protocolo de recepción, sin construir software ni copiar todo el modelo de SD5. Exclusiones del componente: Administrar sistema laboral externo y automatizar un control físico sin contrato compatible. La especificación permite derivar casos de prueba positivos y de rechazo; datos, autoridad y decisiones pendientes están identificados. Acceso exige identidad, vigencia, inducción e intervención autorizada; QR identifica sin conceder habilitación autónoma. Ulysses Barreto, Funcional, responde por este paquete; la contraparte valida decisiones de negocio y conserva la recepción contractual. El cumplimiento de tiempos o disponibilidad se verifica después en la prueba integral, sin acreditarlo por la entrega del código.\\\hline
\texttt{1.\allowbreak{}5.\allowbreak{}1.\allowbreak{}5.\allowbreak{}2}\newline Lógica, persistencia y contratos: Acreditación, inducción y acceso a intervención\newline Nikolai Navea; 64 HH. & Servicio, persistencia y contrato propios del resultado, pruebas unitarias y documentación; excluye vistas y componentes de canal compartidos. Exclusiones del componente: Administrar sistema laboral externo y automatizar un control físico sin contrato compatible. La implementación pasa sus pruebas unitarias, restricciones de datos, duplicados y contratos para los escenarios definidos. Acceso exige identidad, vigencia, inducción e intervención autorizada; QR identifica sin conceder habilitación autónoma. Nikolai Navea, Desarrollo, responde por este paquete; la contraparte valida decisiones de negocio y conserva la recepción contractual. El cumplimiento de tiempos o disponibilidad se verifica después en la prueba integral, sin acreditarlo por la entrega del código.\\\hline
\texttt{1.\allowbreak{}5.\allowbreak{}1.\allowbreak{}5.\allowbreak{}3}\newline Vistas de tarea y comprobación: Acreditación, inducción y acceso a intervención\newline Nikolai Navea; 44 HH. & Formulario o vista de tarea en canal móvil/portal autorizado, validación funcional y de acceso; reutiliza navegación, sesión y cola de 1.7, sin reconstruirlos. Exclusiones del componente: Administrar sistema laboral externo y automatizar un control físico sin contrato compatible. La tarea se reproduce sobre versión identificada con permisos, errores y evidencia visibles; se comprueban teclado/foco y controles adecuados al canal. Acceso exige identidad, vigencia, inducción e intervención autorizada; QR identifica sin conceder habilitación autónoma. Nikolai Navea, Desarrollo, responde por este paquete; la contraparte valida decisiones de negocio y conserva la recepción contractual. El cumplimiento de tiempos o disponibilidad se verifica después en la prueba integral, sin acreditarlo por la entrega del código.\\\hline
\texttt{1.\allowbreak{}5.\allowbreak{}2.\allowbreak{}1.\allowbreak{}1}\newline Especificación y protocolo: Despacho identificado y captura segura de combustible\newline Ulysses Barreto; 32 HH. & Reglas, datos, actores, excepciones, contrato y protocolo de recepción, sin construir software ni copiar todo el modelo de SD5. Exclusiones del componente: Construcción de broker/sincronizador y certificación o sustitución del equipo físico. La especificación permite derivar casos de prueba positivos y de rechazo; datos, autoridad y decisiones pendientes están identificados. Despacho conserva operador, nave, producto, instrumento, lecturas, unidad y evidencia en momento seguro; hecho local no depende de contabilidad. Ulysses Barreto, Funcional, responde por este paquete; la contraparte valida decisiones de negocio y conserva la recepción contractual. El cumplimiento de tiempos o disponibilidad se verifica después en la prueba integral, sin acreditarlo por la entrega del código.\\\hline
\texttt{1.\allowbreak{}5.\allowbreak{}2.\allowbreak{}1.\allowbreak{}2}\newline Lógica, persistencia y contratos: Despacho identificado y captura segura de combustible\newline Nikolai Navea; 64 HH. & Servicio, persistencia y contrato propios del resultado, pruebas unitarias y documentación; excluye vistas y componentes de canal compartidos. Exclusiones del componente: Construcción de broker/sincronizador y certificación o sustitución del equipo físico. La implementación pasa sus pruebas unitarias, restricciones de datos, duplicados y contratos para los escenarios definidos. Despacho conserva operador, nave, producto, instrumento, lecturas, unidad y evidencia en momento seguro; hecho local no depende de contabilidad. Nikolai Navea, Desarrollo, responde por este paquete; la contraparte valida decisiones de negocio y conserva la recepción contractual. El cumplimiento de tiempos o disponibilidad se verifica después en la prueba integral, sin acreditarlo por la entrega del código.\\\hline
\texttt{1.\allowbreak{}5.\allowbreak{}2.\allowbreak{}1.\allowbreak{}3}\newline Vistas de tarea y comprobación: Despacho identificado y captura segura de combustible\newline Nikolai Navea; 44 HH. & Formulario o vista de tarea en canal móvil/portal autorizado, validación funcional y de acceso; reutiliza navegación, sesión y cola de 1.7, sin reconstruirlos. Exclusiones del componente: Construcción de broker/sincronizador y certificación o sustitución del equipo físico. La tarea se reproduce sobre versión identificada con permisos, errores y evidencia visibles; se comprueban teclado/foco y controles adecuados al canal. Despacho conserva operador, nave, producto, instrumento, lecturas, unidad y evidencia en momento seguro; hecho local no depende de contabilidad. Nikolai Navea, Desarrollo, responde por este paquete; la contraparte valida decisiones de negocio y conserva la recepción contractual. El cumplimiento de tiempos o disponibilidad se verifica después en la prueba integral, sin acreditarlo por la entrega del código.\\\hline
\texttt{1.\allowbreak{}5.\allowbreak{}2.\allowbreak{}2.\allowbreak{}1}\newline Especificación y protocolo: Conciliación de tramos y rectificación de despachos\newline Ulysses Barreto; 32 HH. & Reglas, datos, actores, excepciones, contrato y protocolo de recepción, sin construir software ni copiar todo el modelo de SD5. Exclusiones del componente: Saldo contable autoritativo y convertir diferencia negativa en consumo válido. La especificación permite derivar casos de prueba positivos y de rechazo; datos, autoridad y decisiones pendientes están identificados. Reinicio/escala o lectura inválida abre tramo/anomalía; sólo volumen verificable se consolida; rectificación conserva original. Ulysses Barreto, Funcional, responde por este paquete; la contraparte valida decisiones de negocio y conserva la recepción contractual. El cumplimiento de tiempos o disponibilidad se verifica después en la prueba integral, sin acreditarlo por la entrega del código.\\\hline
\texttt{1.\allowbreak{}5.\allowbreak{}2.\allowbreak{}2.\allowbreak{}2}\newline Lógica, persistencia y contratos: Conciliación de tramos y rectificación de despachos\newline Nikolai Navea; 64 HH. & Servicio, persistencia y contrato propios del resultado, pruebas unitarias y documentación; excluye vistas y componentes de canal compartidos. Exclusiones del componente: Saldo contable autoritativo y convertir diferencia negativa en consumo válido. La implementación pasa sus pruebas unitarias, restricciones de datos, duplicados y contratos para los escenarios definidos. Reinicio/escala o lectura inválida abre tramo/anomalía; sólo volumen verificable se consolida; rectificación conserva original. Nikolai Navea, Desarrollo, responde por este paquete; la contraparte valida decisiones de negocio y conserva la recepción contractual. El cumplimiento de tiempos o disponibilidad se verifica después en la prueba integral, sin acreditarlo por la entrega del código.\\\hline
\texttt{1.\allowbreak{}5.\allowbreak{}2.\allowbreak{}2.\allowbreak{}3}\newline Vistas de tarea y comprobación: Conciliación de tramos y rectificación de despachos\newline Nikolai Navea; 44 HH. & Formulario o vista de tarea en canal móvil/portal autorizado, validación funcional y de acceso; reutiliza navegación, sesión y cola de 1.7, sin reconstruirlos. Exclusiones del componente: Saldo contable autoritativo y convertir diferencia negativa en consumo válido. La tarea se reproduce sobre versión identificada con permisos, errores y evidencia visibles; se comprueban teclado/foco y controles adecuados al canal. Reinicio/escala o lectura inválida abre tramo/anomalía; sólo volumen verificable se consolida; rectificación conserva original. Nikolai Navea, Desarrollo, responde por este paquete; la contraparte valida decisiones de negocio y conserva la recepción contractual. El cumplimiento de tiempos o disponibilidad se verifica después en la prueba integral, sin acreditarlo por la entrega del código.\\\hline
\texttt{1.\allowbreak{}5.\allowbreak{}3.\allowbreak{}1.\allowbreak{}1}\newline Especificación y protocolo: Intervalos y series conciliadas de agua y energía\newline Ulysses Barreto; 32 HH. & Reglas, datos, actores, excepciones, contrato y protocolo de recepción, sin construir software ni copiar todo el modelo de SD5. Exclusiones del componente: Captura técnica MQTT y cálculo comercial completo E2. La especificación permite derivar casos de prueba positivos y de rechazo; datos, autoridad y decisiones pendientes están identificados. Intervalo mantiene lecturas originales, unidad, época, calidad y asociación temporal; no se suman totalizadores como volumen. Ulysses Barreto, Funcional, responde por este paquete; la contraparte valida decisiones de negocio y conserva la recepción contractual. El cumplimiento de tiempos o disponibilidad se verifica después en la prueba integral, sin acreditarlo por la entrega del código.\\\hline
\texttt{1.\allowbreak{}5.\allowbreak{}3.\allowbreak{}1.\allowbreak{}2}\newline Lógica, persistencia y contratos: Intervalos y series conciliadas de agua y energía\newline Nikolai Navea; 64 HH. & Servicio, persistencia y contrato propios del resultado, pruebas unitarias y documentación; excluye vistas y componentes de canal compartidos. Exclusiones del componente: Captura técnica MQTT y cálculo comercial completo E2. La implementación pasa sus pruebas unitarias, restricciones de datos, duplicados y contratos para los escenarios definidos. Intervalo mantiene lecturas originales, unidad, época, calidad y asociación temporal; no se suman totalizadores como volumen. Nikolai Navea, Desarrollo, responde por este paquete; la contraparte valida decisiones de negocio y conserva la recepción contractual. El cumplimiento de tiempos o disponibilidad se verifica después en la prueba integral, sin acreditarlo por la entrega del código.\\\hline
\texttt{1.\allowbreak{}5.\allowbreak{}3.\allowbreak{}1.\allowbreak{}3}\newline Vistas de tarea y comprobación: Intervalos y series conciliadas de agua y energía\newline Nikolai Navea; 44 HH. & Formulario o vista de tarea en canal móvil/portal autorizado, validación funcional y de acceso; reutiliza navegación, sesión y cola de 1.7, sin reconstruirlos. Exclusiones del componente: Captura técnica MQTT y cálculo comercial completo E2. La tarea se reproduce sobre versión identificada con permisos, errores y evidencia visibles; se comprueban teclado/foco y controles adecuados al canal. Intervalo mantiene lecturas originales, unidad, época, calidad y asociación temporal; no se suman totalizadores como volumen. Nikolai Navea, Desarrollo, responde por este paquete; la contraparte valida decisiones de negocio y conserva la recepción contractual. El cumplimiento de tiempos o disponibilidad se verifica después en la prueba integral, sin acreditarlo por la entrega del código.\\\hline
\texttt{1.\allowbreak{}5.\allowbreak{}3.\allowbreak{}2.\allowbreak{}1}\newline Especificación y protocolo: Consultas de consumo y evidencia ambiental inicial\newline Ulysses Barreto; 32 HH. & Reglas, datos, actores, excepciones, contrato y protocolo de recepción, sin construir software ni copiar todo el modelo de SD5. Exclusiones del componente: Pasaporte IN5 y presumir serie previa a 2029 sin calendario real. La especificación permite derivar casos de prueba positivos y de rechazo; datos, autoridad y decisiones pendientes están identificados. Vista conserva período, cobertura, versión y fuente; evidencia incompleta impide afirmar cumplimiento ambiental o reducción. Ulysses Barreto, Funcional, responde por este paquete; la contraparte valida decisiones de negocio y conserva la recepción contractual. El cumplimiento de tiempos o disponibilidad se verifica después en la prueba integral, sin acreditarlo por la entrega del código.\\\hline
\texttt{1.\allowbreak{}5.\allowbreak{}3.\allowbreak{}2.\allowbreak{}2}\newline Lógica, persistencia y contratos: Consultas de consumo y evidencia ambiental inicial\newline Nikolai Navea; 64 HH. & Servicio, persistencia y contrato propios del resultado, pruebas unitarias y documentación; excluye vistas y componentes de canal compartidos. Exclusiones del componente: Pasaporte IN5 y presumir serie previa a 2029 sin calendario real. La implementación pasa sus pruebas unitarias, restricciones de datos, duplicados y contratos para los escenarios definidos. Vista conserva período, cobertura, versión y fuente; evidencia incompleta impide afirmar cumplimiento ambiental o reducción. Nikolai Navea, Desarrollo, responde por este paquete; la contraparte valida decisiones de negocio y conserva la recepción contractual. El cumplimiento de tiempos o disponibilidad se verifica después en la prueba integral, sin acreditarlo por la entrega del código.\\\hline
\texttt{1.\allowbreak{}5.\allowbreak{}3.\allowbreak{}2.\allowbreak{}3}\newline Vistas de tarea y comprobación: Consultas de consumo y evidencia ambiental inicial\newline Nikolai Navea; 44 HH. & Formulario o vista de tarea en canal móvil/portal autorizado, validación funcional y de acceso; reutiliza navegación, sesión y cola de 1.7, sin reconstruirlos. Exclusiones del componente: Pasaporte IN5 y presumir serie previa a 2029 sin calendario real. La tarea se reproduce sobre versión identificada con permisos, errores y evidencia visibles; se comprueban teclado/foco y controles adecuados al canal. Vista conserva período, cobertura, versión y fuente; evidencia incompleta impide afirmar cumplimiento ambiental o reducción. Nikolai Navea, Desarrollo, responde por este paquete; la contraparte valida decisiones de negocio y conserva la recepción contractual. El cumplimiento de tiempos o disponibilidad se verifica después en la prueba integral, sin acreditarlo por la entrega del código.\\\hline
\texttt{1.\allowbreak{}5.\allowbreak{}3.\allowbreak{}3.\allowbreak{}1}\newline Especificación y protocolo: Recepción de residuos y discrepancias de origen\newline Ulysses Barreto; 32 HH. & Reglas, datos, actores, excepciones, contrato y protocolo de recepción, sin construir software ni copiar todo el modelo de SD5. Exclusiones del componente: Investigación física ejecutada por software y emitir certificado de gestor. La especificación permite derivar casos de prueba positivos y de rechazo; datos, autoridad y decisiones pendientes están identificados. Registro conserva origen, cantidad, unidad y evidencia; origen desconocido permanece observado con investigación identificada. Ulysses Barreto, Funcional, responde por este paquete; la contraparte valida decisiones de negocio y conserva la recepción contractual. El cumplimiento de tiempos o disponibilidad se verifica después en la prueba integral, sin acreditarlo por la entrega del código.\\\hline
\texttt{1.\allowbreak{}5.\allowbreak{}3.\allowbreak{}3.\allowbreak{}2}\newline Lógica, persistencia y contratos: Recepción de residuos y discrepancias de origen\newline Nikolai Navea; 64 HH. & Servicio, persistencia y contrato propios del resultado, pruebas unitarias y documentación; excluye vistas y componentes de canal compartidos. Exclusiones del componente: Investigación física ejecutada por software y emitir certificado de gestor. La implementación pasa sus pruebas unitarias, restricciones de datos, duplicados y contratos para los escenarios definidos. Registro conserva origen, cantidad, unidad y evidencia; origen desconocido permanece observado con investigación identificada. Nikolai Navea, Desarrollo, responde por este paquete; la contraparte valida decisiones de negocio y conserva la recepción contractual. El cumplimiento de tiempos o disponibilidad se verifica después en la prueba integral, sin acreditarlo por la entrega del código.\\\hline
\texttt{1.\allowbreak{}5.\allowbreak{}3.\allowbreak{}3.\allowbreak{}3}\newline Vistas de tarea y comprobación: Recepción de residuos y discrepancias de origen\newline Nikolai Navea; 44 HH. & Formulario o vista de tarea en canal móvil/portal autorizado, validación funcional y de acceso; reutiliza navegación, sesión y cola de 1.7, sin reconstruirlos. Exclusiones del componente: Investigación física ejecutada por software y emitir certificado de gestor. La tarea se reproduce sobre versión identificada con permisos, errores y evidencia visibles; se comprueban teclado/foco y controles adecuados al canal. Registro conserva origen, cantidad, unidad y evidencia; origen desconocido permanece observado con investigación identificada. Nikolai Navea, Desarrollo, responde por este paquete; la contraparte valida decisiones de negocio y conserva la recepción contractual. El cumplimiento de tiempos o disponibilidad se verifica después en la prueba integral, sin acreditarlo por la entrega del código.\\\hline
\texttt{1.\allowbreak{}5.\allowbreak{}3.\allowbreak{}4.\allowbreak{}1}\newline Especificación y protocolo: Entregas RESPEL parciales y certificados vinculados\newline Ulysses Barreto; 32 HH. & Reglas, datos, actores, excepciones, contrato y protocolo de recepción, sin construir software ni copiar todo el modelo de SD5. Exclusiones del componente: Otorgar certificación externa y borrar retiro real por discrepancia de origen. La especificación permite derivar casos de prueba positivos y de rechazo; datos, autoridad y decisiones pendientes están identificados. Detalle permite dividir/agrupar residuos sin exceder disponibilidad; retiro y certificado se vinculan a entrega efectiva y diferencias quedan trazadas. Ulysses Barreto, Funcional, responde por este paquete; la contraparte valida decisiones de negocio y conserva la recepción contractual. El cumplimiento de tiempos o disponibilidad se verifica después en la prueba integral, sin acreditarlo por la entrega del código.\\\hline
\texttt{1.\allowbreak{}5.\allowbreak{}3.\allowbreak{}4.\allowbreak{}2}\newline Lógica, persistencia y contratos: Entregas RESPEL parciales y certificados vinculados\newline Nikolai Navea; 64 HH. & Servicio, persistencia y contrato propios del resultado, pruebas unitarias y documentación; excluye vistas y componentes de canal compartidos. Exclusiones del componente: Otorgar certificación externa y borrar retiro real por discrepancia de origen. La implementación pasa sus pruebas unitarias, restricciones de datos, duplicados y contratos para los escenarios definidos. Detalle permite dividir/agrupar residuos sin exceder disponibilidad; retiro y certificado se vinculan a entrega efectiva y diferencias quedan trazadas. Nikolai Navea, Desarrollo, responde por este paquete; la contraparte valida decisiones de negocio y conserva la recepción contractual. El cumplimiento de tiempos o disponibilidad se verifica después en la prueba integral, sin acreditarlo por la entrega del código.\\\hline
\texttt{1.\allowbreak{}5.\allowbreak{}3.\allowbreak{}4.\allowbreak{}3}\newline Vistas de tarea y comprobación: Entregas RESPEL parciales y certificados vinculados\newline Nikolai Navea; 44 HH. & Formulario o vista de tarea en canal móvil/portal autorizado, validación funcional y de acceso; reutiliza navegación, sesión y cola de 1.7, sin reconstruirlos. Exclusiones del componente: Otorgar certificación externa y borrar retiro real por discrepancia de origen. La tarea se reproduce sobre versión identificada con permisos, errores y evidencia visibles; se comprueban teclado/foco y controles adecuados al canal. Detalle permite dividir/agrupar residuos sin exceder disponibilidad; retiro y certificado se vinculan a entrega efectiva y diferencias quedan trazadas. Nikolai Navea, Desarrollo, responde por este paquete; la contraparte valida decisiones de negocio y conserva la recepción contractual. El cumplimiento de tiempos o disponibilidad se verifica después en la prueba integral, sin acreditarlo por la entrega del código.\\\hline
\texttt{1.\allowbreak{}5.\allowbreak{}3.\allowbreak{}5.\allowbreak{}1}\newline Especificación y protocolo: Indicadores ambientales e imputación comercial de consumo\newline Ulysses Barreto; 32 HH. & Reglas, datos, actores, excepciones, contrato y protocolo de recepción, sin construir software ni copiar todo el modelo de SD5. Exclusiones del componente: Cobro tributario, incentivo IN4 y beneficio ambiental voluntario IN5. La especificación permite derivar casos de prueba positivos y de rechazo; datos, autoridad y decisiones pendientes están identificados. Indicador tiene período, regla y registros de soporte; sólo intervalos conciliados producen hecho de consumo para Finanzas sin duplicados. Ulysses Barreto, Funcional, responde por este paquete; la contraparte valida decisiones de negocio y conserva la recepción contractual. El cumplimiento de tiempos o disponibilidad se verifica después en la prueba integral, sin acreditarlo por la entrega del código.\\\hline
\texttt{1.\allowbreak{}5.\allowbreak{}3.\allowbreak{}5.\allowbreak{}2}\newline Lógica, persistencia y contratos: Indicadores ambientales e imputación comercial de consumo\newline Nikolai Navea; 64 HH. & Servicio, persistencia y contrato propios del resultado, pruebas unitarias y documentación; excluye vistas y componentes de canal compartidos. Exclusiones del componente: Cobro tributario, incentivo IN4 y beneficio ambiental voluntario IN5. La implementación pasa sus pruebas unitarias, restricciones de datos, duplicados y contratos para los escenarios definidos. Indicador tiene período, regla y registros de soporte; sólo intervalos conciliados producen hecho de consumo para Finanzas sin duplicados. Nikolai Navea, Desarrollo, responde por este paquete; la contraparte valida decisiones de negocio y conserva la recepción contractual. El cumplimiento de tiempos o disponibilidad se verifica después en la prueba integral, sin acreditarlo por la entrega del código.\\\hline
\texttt{1.\allowbreak{}5.\allowbreak{}3.\allowbreak{}5.\allowbreak{}3}\newline Vistas de tarea y comprobación: Indicadores ambientales e imputación comercial de consumo\newline Nikolai Navea; 44 HH. & Formulario o vista de tarea en canal móvil/portal autorizado, validación funcional y de acceso; reutiliza navegación, sesión y cola de 1.7, sin reconstruirlos. Exclusiones del componente: Cobro tributario, incentivo IN4 y beneficio ambiental voluntario IN5. La tarea se reproduce sobre versión identificada con permisos, errores y evidencia visibles; se comprueban teclado/foco y controles adecuados al canal. Indicador tiene período, regla y registros de soporte; sólo intervalos conciliados producen hecho de consumo para Finanzas sin duplicados. Nikolai Navea, Desarrollo, responde por este paquete; la contraparte valida decisiones de negocio y conserva la recepción contractual. El cumplimiento de tiempos o disponibilidad se verifica después en la prueba integral, sin acreditarlo por la entrega del código.\\\hline
\texttt{1.\allowbreak{}6.\allowbreak{}1.\allowbreak{}1.\allowbreak{}1}\newline Especificación y protocolo: Maestros de personas, organizaciones y embarcaciones\newline Ulysses Barreto; 32 HH. & Reglas, datos, actores, excepciones, contrato y protocolo de recepción, sin construir software ni copiar todo el modelo de SD5. Exclusiones del componente: Extracción/saneamiento de legado y autenticación técnica Cognito. La especificación permite derivar casos de prueba positivos y de rechazo; datos, autoridad y decisiones pendientes están identificados. Identidad estable, vínculos y vigencias se consultan por contrato; persona y embarcación no se duplican por cambio de nombre o canal. Ulysses Barreto, Funcional, responde por este paquete; la contraparte valida decisiones de negocio y conserva la recepción contractual. El cumplimiento de tiempos o disponibilidad se verifica después en la prueba integral, sin acreditarlo por la entrega del código.\\\hline
\texttt{1.\allowbreak{}6.\allowbreak{}1.\allowbreak{}1.\allowbreak{}2}\newline Lógica, persistencia y contratos: Maestros de personas, organizaciones y embarcaciones\newline Nikolai Navea; 64 HH. & Servicio, persistencia y contrato propios del resultado, pruebas unitarias y documentación; excluye vistas y componentes de canal compartidos. Exclusiones del componente: Extracción/saneamiento de legado y autenticación técnica Cognito. La implementación pasa sus pruebas unitarias, restricciones de datos, duplicados y contratos para los escenarios definidos. Identidad estable, vínculos y vigencias se consultan por contrato; persona y embarcación no se duplican por cambio de nombre o canal. Nikolai Navea, Desarrollo, responde por este paquete; la contraparte valida decisiones de negocio y conserva la recepción contractual. El cumplimiento de tiempos o disponibilidad se verifica después en la prueba integral, sin acreditarlo por la entrega del código.\\\hline
\texttt{1.\allowbreak{}6.\allowbreak{}1.\allowbreak{}1.\allowbreak{}3}\newline Vistas de tarea y comprobación: Maestros de personas, organizaciones y embarcaciones\newline Nikolai Navea; 44 HH. & Formulario o vista de tarea en canal móvil/portal autorizado, validación funcional y de acceso; reutiliza navegación, sesión y cola de 1.7, sin reconstruirlos. Exclusiones del componente: Extracción/saneamiento de legado y autenticación técnica Cognito. La tarea se reproduce sobre versión identificada con permisos, errores y evidencia visibles; se comprueban teclado/foco y controles adecuados al canal. Identidad estable, vínculos y vigencias se consultan por contrato; persona y embarcación no se duplican por cambio de nombre o canal. Nikolai Navea, Desarrollo, responde por este paquete; la contraparte valida decisiones de negocio y conserva la recepción contractual. El cumplimiento de tiempos o disponibilidad se verifica después en la prueba integral, sin acreditarlo por la entrega del código.\\\hline
\texttt{1.\allowbreak{}6.\allowbreak{}1.\allowbreak{}2.\allowbreak{}1}\newline Especificación y protocolo: Contratos y documentos vigentes de consulta E1\newline Ulysses Barreto; 32 HH. & Reglas, datos, actores, excepciones, contrato y protocolo de recepción, sin construir software ni copiar todo el modelo de SD5. Exclusiones del componente: OCR IN3 y revisión documental íntegra de migración. La especificación permite derivar casos de prueba positivos y de rechazo; datos, autoridad y decisiones pendientes están identificados. Contrato, parte, documento y vigencia tienen autoridad y versión; acceso respeta ámbito y un documento pendiente no se declara vigente. Ulysses Barreto, Funcional, responde por este paquete; la contraparte valida decisiones de negocio y conserva la recepción contractual. El cumplimiento de tiempos o disponibilidad se verifica después en la prueba integral, sin acreditarlo por la entrega del código.\\\hline
\texttt{1.\allowbreak{}6.\allowbreak{}1.\allowbreak{}2.\allowbreak{}2}\newline Lógica, persistencia y contratos: Contratos y documentos vigentes de consulta E1\newline Nikolai Navea; 64 HH. & Servicio, persistencia y contrato propios del resultado, pruebas unitarias y documentación; excluye vistas y componentes de canal compartidos. Exclusiones del componente: OCR IN3 y revisión documental íntegra de migración. La implementación pasa sus pruebas unitarias, restricciones de datos, duplicados y contratos para los escenarios definidos. Contrato, parte, documento y vigencia tienen autoridad y versión; acceso respeta ámbito y un documento pendiente no se declara vigente. Nikolai Navea, Desarrollo, responde por este paquete; la contraparte valida decisiones de negocio y conserva la recepción contractual. El cumplimiento de tiempos o disponibilidad se verifica después en la prueba integral, sin acreditarlo por la entrega del código.\\\hline
\texttt{1.\allowbreak{}6.\allowbreak{}1.\allowbreak{}2.\allowbreak{}3}\newline Vistas de tarea y comprobación: Contratos y documentos vigentes de consulta E1\newline Nikolai Navea; 44 HH. & Formulario o vista de tarea en canal móvil/portal autorizado, validación funcional y de acceso; reutiliza navegación, sesión y cola de 1.7, sin reconstruirlos. Exclusiones del componente: OCR IN3 y revisión documental íntegra de migración. La tarea se reproduce sobre versión identificada con permisos, errores y evidencia visibles; se comprueban teclado/foco y controles adecuados al canal. Contrato, parte, documento y vigencia tienen autoridad y versión; acceso respeta ámbito y un documento pendiente no se declara vigente. Nikolai Navea, Desarrollo, responde por este paquete; la contraparte valida decisiones de negocio y conserva la recepción contractual. El cumplimiento de tiempos o disponibilidad se verifica después en la prueba integral, sin acreditarlo por la entrega del código.\\\hline
\texttt{1.\allowbreak{}6.\allowbreak{}1.\allowbreak{}3.\allowbreak{}1}\newline Especificación y protocolo: Atención a visitas y asignación transitoria inicial\newline Ulysses Barreto; 32 HH. & Reglas, datos, actores, excepciones, contrato y protocolo de recepción, sin construir software ni copiar todo el modelo de SD5. Exclusiones del componente: Emisión tributaria y doble autoridad sobre ocupación física. La especificación permite derivar casos de prueba positivos y de rechazo; datos, autoridad y decisiones pendientes están identificados. Solicitud conserva identidad, período y estado; confirmación requiere compatibilidad y autoridad de puesto, con información español/inglés prevista. Ulysses Barreto, Funcional, responde por este paquete; la contraparte valida decisiones de negocio y conserva la recepción contractual. El cumplimiento de tiempos o disponibilidad se verifica después en la prueba integral, sin acreditarlo por la entrega del código.\\\hline
\texttt{1.\allowbreak{}6.\allowbreak{}1.\allowbreak{}3.\allowbreak{}2}\newline Lógica, persistencia y contratos: Atención a visitas y asignación transitoria inicial\newline Nikolai Navea; 64 HH. & Servicio, persistencia y contrato propios del resultado, pruebas unitarias y documentación; excluye vistas y componentes de canal compartidos. Exclusiones del componente: Emisión tributaria y doble autoridad sobre ocupación física. La implementación pasa sus pruebas unitarias, restricciones de datos, duplicados y contratos para los escenarios definidos. Solicitud conserva identidad, período y estado; confirmación requiere compatibilidad y autoridad de puesto, con información español/inglés prevista. Nikolai Navea, Desarrollo, responde por este paquete; la contraparte valida decisiones de negocio y conserva la recepción contractual. El cumplimiento de tiempos o disponibilidad se verifica después en la prueba integral, sin acreditarlo por la entrega del código.\\\hline
\texttt{1.\allowbreak{}6.\allowbreak{}1.\allowbreak{}3.\allowbreak{}3}\newline Vistas de tarea y comprobación: Atención a visitas y asignación transitoria inicial\newline Nikolai Navea; 44 HH. & Formulario o vista de tarea en canal móvil/portal autorizado, validación funcional y de acceso; reutiliza navegación, sesión y cola de 1.7, sin reconstruirlos. Exclusiones del componente: Emisión tributaria y doble autoridad sobre ocupación física. La tarea se reproduce sobre versión identificada con permisos, errores y evidencia visibles; se comprueban teclado/foco y controles adecuados al canal. Solicitud conserva identidad, período y estado; confirmación requiere compatibilidad y autoridad de puesto, con información español/inglés prevista. Nikolai Navea, Desarrollo, responde por este paquete; la contraparte valida decisiones de negocio y conserva la recepción contractual. El cumplimiento de tiempos o disponibilidad se verifica después en la prueba integral, sin acreditarlo por la entrega del código.\\\hline
\texttt{1.\allowbreak{}6.\allowbreak{}1.\allowbreak{}4.\allowbreak{}1}\newline Especificación y protocolo: Espera y gestión contractual completa\newline Ulysses Barreto; 32 HH. & Reglas, datos, actores, excepciones, contrato y protocolo de recepción, sin construir software ni copiar todo el modelo de SD5. Exclusiones del componente: Decidir prioridades comerciales por Synaptix y rehacer compatibilidad operacional E1. La especificación permite derivar casos de prueba positivos y de rechazo; datos, autoridad y decisiones pendientes están identificados. Orden de espera, prioridad y excepción quedan autorizados; contrato y asignación se concilian sin cambiar hechos de ocupación de otro dominio. Ulysses Barreto, Funcional, responde por este paquete; la contraparte valida decisiones de negocio y conserva la recepción contractual. El cumplimiento de tiempos o disponibilidad se verifica después en la prueba integral, sin acreditarlo por la entrega del código.\\\hline
\texttt{1.\allowbreak{}6.\allowbreak{}1.\allowbreak{}4.\allowbreak{}2}\newline Lógica, persistencia y contratos: Espera y gestión contractual completa\newline Nikolai Navea; 64 HH. & Servicio, persistencia y contrato propios del resultado, pruebas unitarias y documentación; excluye vistas y componentes de canal compartidos. Exclusiones del componente: Decidir prioridades comerciales por Synaptix y rehacer compatibilidad operacional E1. La implementación pasa sus pruebas unitarias, restricciones de datos, duplicados y contratos para los escenarios definidos. Orden de espera, prioridad y excepción quedan autorizados; contrato y asignación se concilian sin cambiar hechos de ocupación de otro dominio. Nikolai Navea, Desarrollo, responde por este paquete; la contraparte valida decisiones de negocio y conserva la recepción contractual. El cumplimiento de tiempos o disponibilidad se verifica después en la prueba integral, sin acreditarlo por la entrega del código.\\\hline
\texttt{1.\allowbreak{}6.\allowbreak{}1.\allowbreak{}4.\allowbreak{}3}\newline Vistas de tarea y comprobación: Espera y gestión contractual completa\newline Nikolai Navea; 44 HH. & Formulario o vista de tarea en canal móvil/portal autorizado, validación funcional y de acceso; reutiliza navegación, sesión y cola de 1.7, sin reconstruirlos. Exclusiones del componente: Decidir prioridades comerciales por Synaptix y rehacer compatibilidad operacional E1. La tarea se reproduce sobre versión identificada con permisos, errores y evidencia visibles; se comprueban teclado/foco y controles adecuados al canal. Orden de espera, prioridad y excepción quedan autorizados; contrato y asignación se concilian sin cambiar hechos de ocupación de otro dominio. Nikolai Navea, Desarrollo, responde por este paquete; la contraparte valida decisiones de negocio y conserva la recepción contractual. El cumplimiento de tiempos o disponibilidad se verifica después en la prueba integral, sin acreditarlo por la entrega del código.\\\hline
\texttt{1.\allowbreak{}6.\allowbreak{}2.\allowbreak{}1.\allowbreak{}1}\newline Especificación y protocolo: Hechos facturables, tarifas y lotes de envío\newline Ulysses Barreto; 32 HH. & Reglas, datos, actores, excepciones, contrato y protocolo de recepción, sin construir software ni copiar todo el modelo de SD5. Exclusiones del componente: Adaptador contable técnico y emisión de documentos tributarios. La especificación permite derivar casos de prueba positivos y de rechazo; datos, autoridad y decisiones pendientes están identificados. Cada hecho conserva servicio, cantidad, unidad, moneda, tarifa aplicable y responsable de pago; lote no duplica hechos remitidos. Ulysses Barreto, Funcional, responde por este paquete; la contraparte valida decisiones de negocio y conserva la recepción contractual. El cumplimiento de tiempos o disponibilidad se verifica después en la prueba integral, sin acreditarlo por la entrega del código.\\\hline
\texttt{1.\allowbreak{}6.\allowbreak{}2.\allowbreak{}1.\allowbreak{}2}\newline Lógica, persistencia y contratos: Hechos facturables, tarifas y lotes de envío\newline Nikolai Navea; 64 HH. & Servicio, persistencia y contrato propios del resultado, pruebas unitarias y documentación; excluye vistas y componentes de canal compartidos. Exclusiones del componente: Adaptador contable técnico y emisión de documentos tributarios. La implementación pasa sus pruebas unitarias, restricciones de datos, duplicados y contratos para los escenarios definidos. Cada hecho conserva servicio, cantidad, unidad, moneda, tarifa aplicable y responsable de pago; lote no duplica hechos remitidos. Nikolai Navea, Desarrollo, responde por este paquete; la contraparte valida decisiones de negocio y conserva la recepción contractual. El cumplimiento de tiempos o disponibilidad se verifica después en la prueba integral, sin acreditarlo por la entrega del código.\\\hline
\texttt{1.\allowbreak{}6.\allowbreak{}2.\allowbreak{}1.\allowbreak{}3}\newline Vistas de tarea y comprobación: Hechos facturables, tarifas y lotes de envío\newline Nikolai Navea; 44 HH. & Formulario o vista de tarea en canal móvil/portal autorizado, validación funcional y de acceso; reutiliza navegación, sesión y cola de 1.7, sin reconstruirlos. Exclusiones del componente: Adaptador contable técnico y emisión de documentos tributarios. La tarea se reproduce sobre versión identificada con permisos, errores y evidencia visibles; se comprueban teclado/foco y controles adecuados al canal. Cada hecho conserva servicio, cantidad, unidad, moneda, tarifa aplicable y responsable de pago; lote no duplica hechos remitidos. Nikolai Navea, Desarrollo, responde por este paquete; la contraparte valida decisiones de negocio y conserva la recepción contractual. El cumplimiento de tiempos o disponibilidad se verifica después en la prueba integral, sin acreditarlo por la entrega del código.\\\hline
\texttt{1.\allowbreak{}6.\allowbreak{}2.\allowbreak{}2.\allowbreak{}1}\newline Especificación y protocolo: Respuesta contable, pagos y saldos informados\newline Ulysses Barreto; 32 HH. & Reglas, datos, actores, excepciones, contrato y protocolo de recepción, sin construir software ni copiar todo el modelo de SD5. Exclusiones del componente: Construcción de webhooks/pasarela y modificar saldo externo sin autorización. La especificación permite derivar casos de prueba positivos y de rechazo; datos, autoridad y decisiones pendientes están identificados. Aceptación/rechazo y saldo conservan fuente y fecha; un saldo local atrasado se identifica y no sustituye autoridad contable. Ulysses Barreto, Funcional, responde por este paquete; la contraparte valida decisiones de negocio y conserva la recepción contractual. El cumplimiento de tiempos o disponibilidad se verifica después en la prueba integral, sin acreditarlo por la entrega del código.\\\hline
\texttt{1.\allowbreak{}6.\allowbreak{}2.\allowbreak{}2.\allowbreak{}2}\newline Lógica, persistencia y contratos: Respuesta contable, pagos y saldos informados\newline Nikolai Navea; 64 HH. & Servicio, persistencia y contrato propios del resultado, pruebas unitarias y documentación; excluye vistas y componentes de canal compartidos. Exclusiones del componente: Construcción de webhooks/pasarela y modificar saldo externo sin autorización. La implementación pasa sus pruebas unitarias, restricciones de datos, duplicados y contratos para los escenarios definidos. Aceptación/rechazo y saldo conservan fuente y fecha; un saldo local atrasado se identifica y no sustituye autoridad contable. Nikolai Navea, Desarrollo, responde por este paquete; la contraparte valida decisiones de negocio y conserva la recepción contractual. El cumplimiento de tiempos o disponibilidad se verifica después en la prueba integral, sin acreditarlo por la entrega del código.\\\hline
\texttt{1.\allowbreak{}6.\allowbreak{}2.\allowbreak{}2.\allowbreak{}3}\newline Vistas de tarea y comprobación: Respuesta contable, pagos y saldos informados\newline Nikolai Navea; 44 HH. & Formulario o vista de tarea en canal móvil/portal autorizado, validación funcional y de acceso; reutiliza navegación, sesión y cola de 1.7, sin reconstruirlos. Exclusiones del componente: Construcción de webhooks/pasarela y modificar saldo externo sin autorización. La tarea se reproduce sobre versión identificada con permisos, errores y evidencia visibles; se comprueban teclado/foco y controles adecuados al canal. Aceptación/rechazo y saldo conservan fuente y fecha; un saldo local atrasado se identifica y no sustituye autoridad contable. Nikolai Navea, Desarrollo, responde por este paquete; la contraparte valida decisiones de negocio y conserva la recepción contractual. El cumplimiento de tiempos o disponibilidad se verifica después en la prueba integral, sin acreditarlo por la entrega del código.\\\hline
\texttt{1.\allowbreak{}6.\allowbreak{}2.\allowbreak{}3.\allowbreak{}1}\newline Especificación y protocolo: Cobranza y medidas con decisión competente\newline Ulysses Barreto; 32 HH. & Reglas, datos, actores, excepciones, contrato y protocolo de recepción, sin construir software ni copiar todo el modelo de SD5. Exclusiones del componente: Adoptar decisiones jurídicas o físicas por software. La especificación permite derivar casos de prueba positivos y de rechazo; datos, autoridad y decisiones pendientes están identificados. Actuación registra propuesta, autorización y ejecución; tramo de deuda no activa restricción física ni desposesión automática. Ulysses Barreto, Funcional, responde por este paquete; la contraparte valida decisiones de negocio y conserva la recepción contractual. El cumplimiento de tiempos o disponibilidad se verifica después en la prueba integral, sin acreditarlo por la entrega del código.\\\hline
\texttt{1.\allowbreak{}6.\allowbreak{}2.\allowbreak{}3.\allowbreak{}2}\newline Lógica, persistencia y contratos: Cobranza y medidas con decisión competente\newline Nikolai Navea; 64 HH. & Servicio, persistencia y contrato propios del resultado, pruebas unitarias y documentación; excluye vistas y componentes de canal compartidos. Exclusiones del componente: Adoptar decisiones jurídicas o físicas por software. La implementación pasa sus pruebas unitarias, restricciones de datos, duplicados y contratos para los escenarios definidos. Actuación registra propuesta, autorización y ejecución; tramo de deuda no activa restricción física ni desposesión automática. Nikolai Navea, Desarrollo, responde por este paquete; la contraparte valida decisiones de negocio y conserva la recepción contractual. El cumplimiento de tiempos o disponibilidad se verifica después en la prueba integral, sin acreditarlo por la entrega del código.\\\hline
\texttt{1.\allowbreak{}6.\allowbreak{}2.\allowbreak{}3.\allowbreak{}3}\newline Vistas de tarea y comprobación: Cobranza y medidas con decisión competente\newline Nikolai Navea; 44 HH. & Formulario o vista de tarea en canal móvil/portal autorizado, validación funcional y de acceso; reutiliza navegación, sesión y cola de 1.7, sin reconstruirlos. Exclusiones del componente: Adoptar decisiones jurídicas o físicas por software. La tarea se reproduce sobre versión identificada con permisos, errores y evidencia visibles; se comprueban teclado/foco y controles adecuados al canal. Actuación registra propuesta, autorización y ejecución; tramo de deuda no activa restricción física ni desposesión automática. Nikolai Navea, Desarrollo, responde por este paquete; la contraparte valida decisiones de negocio y conserva la recepción contractual. El cumplimiento de tiempos o disponibilidad se verifica después en la prueba integral, sin acreditarlo por la entrega del código.\\\hline
\texttt{1.\allowbreak{}6.\allowbreak{}2.\allowbreak{}4.\allowbreak{}1}\newline Especificación y protocolo: Expedientes de custodia y exportación íntegra\newline Ulysses Barreto; 32 HH. & Reglas, datos, actores, excepciones, contrato y protocolo de recepción, sin construir software ni copiar todo el modelo de SD5. Exclusiones del componente: Inventar antecedentes faltantes, repetir digitalización/revisión de migración y resolver litigios. La especificación permite derivar casos de prueba positivos y de rechazo; datos, autoridad y decisiones pendientes están identificados. Índice conserva folio, versión, documento, origen e integridad; exportación mantiene relaciones y señala faltantes para los 14 expedientes. Ulysses Barreto, Funcional, responde por este paquete; la contraparte valida decisiones de negocio y conserva la recepción contractual. El cumplimiento de tiempos o disponibilidad se verifica después en la prueba integral, sin acreditarlo por la entrega del código.\\\hline
\texttt{1.\allowbreak{}6.\allowbreak{}2.\allowbreak{}4.\allowbreak{}2}\newline Lógica, persistencia y contratos: Expedientes de custodia y exportación íntegra\newline Nikolai Navea; 64 HH. & Servicio, persistencia y contrato propios del resultado, pruebas unitarias y documentación; excluye vistas y componentes de canal compartidos. Exclusiones del componente: Inventar antecedentes faltantes, repetir digitalización/revisión de migración y resolver litigios. La implementación pasa sus pruebas unitarias, restricciones de datos, duplicados y contratos para los escenarios definidos. Índice conserva folio, versión, documento, origen e integridad; exportación mantiene relaciones y señala faltantes para los 14 expedientes. Nikolai Navea, Desarrollo, responde por este paquete; la contraparte valida decisiones de negocio y conserva la recepción contractual. El cumplimiento de tiempos o disponibilidad se verifica después en la prueba integral, sin acreditarlo por la entrega del código.\\\hline
\texttt{1.\allowbreak{}6.\allowbreak{}2.\allowbreak{}4.\allowbreak{}3}\newline Vistas de tarea y comprobación: Expedientes de custodia y exportación íntegra\newline Nikolai Navea; 44 HH. & Formulario o vista de tarea en canal móvil/portal autorizado, validación funcional y de acceso; reutiliza navegación, sesión y cola de 1.7, sin reconstruirlos. Exclusiones del componente: Inventar antecedentes faltantes, repetir digitalización/revisión de migración y resolver litigios. La tarea se reproduce sobre versión identificada con permisos, errores y evidencia visibles; se comprueban teclado/foco y controles adecuados al canal. Índice conserva folio, versión, documento, origen e integridad; exportación mantiene relaciones y señala faltantes para los 14 expedientes. Nikolai Navea, Desarrollo, responde por este paquete; la contraparte valida decisiones de negocio y conserva la recepción contractual. El cumplimiento de tiempos o disponibilidad se verifica después en la prueba integral, sin acreditarlo por la entrega del código.\\\hline
\end{longtable}
\endgroup
\FuenteTabla{Elaboración propia. Las horas y los códigos corresponden al catálogo técnico integrado.}

\subsection{Diccionario de conciliación funcional}
Los lotes completan el esfuerzo seleccionado de cada componente. Cada código es terminal y conserva aceptación individual, responsable y programación. Se excluyen identidad/canales/adaptadores compartidos, migración, UAT integral y formación; esos trabajos permanecen en sus cuentas. Las predecesoras exactas y los atributos completos se conservan en \path{componentes/diccionario_integrado.csv}.
\begingroup\small\begin{longtable}{L{29mm}L{80mm}R{12mm}L{28mm}}
\caption{Lotes de conciliación: resultado, esfuerzo y responsable}\label{tab:sd7-conciliacion-lotes}\\\hline
Código & Resultado y aceptación & HH & Responsable\\\hline\endfirsthead
\hline Código & Resultado y aceptación & HH & Responsable\\\hline\endhead
\texttt{1.\allowbreak{}4.\allowbreak{}1.\allowbreak{}1.\allowbreak{}4.\allowbreak{}1} & Análisis de reglas, excepciones y protocolo: Movimientos de salida y regreso identificados — lote 1/1\newline Aceptación: lote versionado y trazable al flujo; revisión independiente sin defectos impeditivos.\newline Perfil A; 2027-01-27 a 2027-01-28. & 51 & Ulysses Barreto\\\hline
\texttt{1.\allowbreak{}4.\allowbreak{}1.\allowbreak{}1.\allowbreak{}4.\allowbreak{}2} & Diseño de datos y restricciones del núcleo: Movimientos de salida y regreso identificados — lote 1/1\newline Aceptación: lote versionado y trazable al flujo; revisión independiente sin defectos impeditivos.\newline Perfil D; 2027-01-27 a 2027-01-29. & 44 & Nicolás Torfan\\\hline
\texttt{1.\allowbreak{}4.\allowbreak{}1.\allowbreak{}1.\allowbreak{}4.\allowbreak{}3} & Diseño de contratos internos y composición: Movimientos de salida y regreso identificados — lote 1/2\newline Aceptación: lote versionado y trazable al flujo; revisión independiente sin defectos impeditivos.\newline Perfil I; 2027-01-29 a 2027-02-01. & 33 & Fernando Guerrero\\\hline
\texttt{1.\allowbreak{}4.\allowbreak{}1.\allowbreak{}1.\allowbreak{}4.\allowbreak{}4} & Diseño de contratos internos y composición: Movimientos de salida y regreso identificados — lote 2/2\newline Aceptación: lote versionado y trazable al flujo; revisión independiente sin defectos impeditivos.\newline Perfil I; 2027-01-27 a 2027-01-29. & 32 & Fernando Guerrero\\\hline
\texttt{1.\allowbreak{}4.\allowbreak{}1.\allowbreak{}1.\allowbreak{}5.\allowbreak{}1} & Construcción funcional y pruebas unitarias: Movimientos de salida y regreso identificados — lote 1/4\newline Aceptación: lote versionado y trazable al flujo; revisión independiente sin defectos impeditivos.\newline Perfil V; 2027-07-21 a 2027-07-23. & 51 & Nikolai Navea\\\hline
\texttt{1.\allowbreak{}4.\allowbreak{}1.\allowbreak{}1.\allowbreak{}5.\allowbreak{}2} & Construcción funcional y pruebas unitarias: Movimientos de salida y regreso identificados — lote 2/4\newline Aceptación: lote versionado y trazable al flujo; revisión independiente sin defectos impeditivos.\newline Perfil V; 2027-07-23 a 2027-07-27. & 51 & Nikolai Navea\\\hline
\texttt{1.\allowbreak{}4.\allowbreak{}1.\allowbreak{}1.\allowbreak{}5.\allowbreak{}3} & Construcción funcional y pruebas unitarias: Movimientos de salida y regreso identificados — lote 3/4\newline Aceptación: lote versionado y trazable al flujo; revisión independiente sin defectos impeditivos.\newline Perfil V; 2027-07-21 a 2027-07-23. & 51 & Nikolai Navea\\\hline
\texttt{1.\allowbreak{}4.\allowbreak{}1.\allowbreak{}1.\allowbreak{}5.\allowbreak{}4} & Construcción funcional y pruebas unitarias: Movimientos de salida y regreso identificados — lote 4/4\newline Aceptación: lote versionado y trazable al flujo; revisión independiente sin defectos impeditivos.\newline Perfil V; 2027-07-23 a 2027-07-27. & 50 & Nikolai Navea\\\hline
\texttt{1.\allowbreak{}4.\allowbreak{}1.\allowbreak{}1.\allowbreak{}6.\allowbreak{}1} & Verificación del núcleo, regresión y evidencia: Movimientos de salida y regreso identificados — lote 1/2\newline Aceptación: lote versionado y trazable al flujo; revisión independiente sin defectos impeditivos.\newline Perfil Q; 2027-07-30 a 2027-08-02. & 49 & Davor Serey\\\hline
\texttt{1.\allowbreak{}4.\allowbreak{}1.\allowbreak{}1.\allowbreak{}6.\allowbreak{}2} & Verificación del núcleo, regresión y evidencia: Movimientos de salida y regreso identificados — lote 2/2\newline Aceptación: lote versionado y trazable al flujo; revisión independiente sin defectos impeditivos.\newline Perfil Q; 2027-07-30 a 2027-08-02. & 49 & Davor Serey\\\hline
\texttt{1.\allowbreak{}4.\allowbreak{}1.\allowbreak{}1.\allowbreak{}7.\allowbreak{}1} & Coordinación y documentación del núcleo: Movimientos de salida y regreso identificados — lote 1/2\newline Aceptación: lote versionado y trazable al flujo; revisión independiente sin defectos impeditivos.\newline Perfil J; 2027-08-02 a 2027-08-04. & 53 & Ignacio Silva\\\hline
\texttt{1.\allowbreak{}4.\allowbreak{}1.\allowbreak{}1.\allowbreak{}7.\allowbreak{}2} & Coordinación y documentación del núcleo: Movimientos de salida y regreso identificados — lote 2/2\newline Aceptación: lote versionado y trazable al flujo; revisión independiente sin defectos impeditivos.\newline Perfil J; 2027-08-04 a 2027-08-06. & 53 & Ignacio Silva\\\hline
\texttt{1.\allowbreak{}4.\allowbreak{}1.\allowbreak{}2.\allowbreak{}4.\allowbreak{}1} & Análisis de reglas, excepciones y protocolo: Planes voluntarios y discrepancias operacionales — lote 1/1\newline Aceptación: lote versionado y trazable al flujo; revisión independiente sin defectos impeditivos.\newline Perfil A; 2027-01-28 a 2027-02-02. & 16 & Ulysses Barreto\\\hline
\texttt{1.\allowbreak{}4.\allowbreak{}1.\allowbreak{}2.\allowbreak{}4.\allowbreak{}2} & Diseño de datos y restricciones del núcleo: Planes voluntarios y discrepancias operacionales — lote 1/1\newline Aceptación: lote versionado y trazable al flujo; revisión independiente sin defectos impeditivos.\newline Perfil D; 2027-01-29 a 2027-02-02. & 16 & Nicolás Torfan\\\hline
\texttt{1.\allowbreak{}4.\allowbreak{}1.\allowbreak{}2.\allowbreak{}4.\allowbreak{}3} & Diseño de contratos internos y composición: Planes voluntarios y discrepancias operacionales — lote 1/1\newline Aceptación: lote versionado y trazable al flujo; revisión independiente sin defectos impeditivos.\newline Perfil I; 2027-02-01 a 2027-02-03. & 23 & Fernando Guerrero\\\hline
\texttt{1.\allowbreak{}4.\allowbreak{}1.\allowbreak{}2.\allowbreak{}5.\allowbreak{}1} & Construcción funcional y pruebas unitarias: Planes voluntarios y discrepancias operacionales — lote 1/1\newline Aceptación: lote versionado y trazable al flujo; revisión independiente sin defectos impeditivos.\newline Perfil V; 2027-07-21 a 2027-07-23. & 62 & Nikolai Navea\\\hline
\texttt{1.\allowbreak{}4.\allowbreak{}1.\allowbreak{}2.\allowbreak{}6.\allowbreak{}1} & Verificación del núcleo, regresión y evidencia: Planes voluntarios y discrepancias operacionales — lote 1/1\newline Aceptación: lote versionado y trazable al flujo; revisión independiente sin defectos impeditivos.\newline Perfil Q; 2027-07-28 a 2027-07-30. & 45 & Davor Serey\\\hline
\texttt{1.\allowbreak{}4.\allowbreak{}1.\allowbreak{}2.\allowbreak{}7.\allowbreak{}1} & Coordinación y documentación del núcleo: Planes voluntarios y discrepancias operacionales — lote 1/1\newline Aceptación: lote versionado y trazable al flujo; revisión independiente sin defectos impeditivos.\newline Perfil J; 2027-07-30 a 2027-08-03. & 53 & Ignacio Silva\\\hline
\texttt{1.\allowbreak{}4.\allowbreak{}1.\allowbreak{}3.\allowbreak{}4.\allowbreak{}1} & Análisis de reglas, excepciones y protocolo: Atención y escalamiento de planes vencidos — lote 1/1\newline Aceptación: lote versionado y trazable al flujo; revisión independiente sin defectos impeditivos.\newline Perfil A; 2027-01-28 a 2027-02-02. & 16 & Ulysses Barreto\\\hline
\texttt{1.\allowbreak{}4.\allowbreak{}1.\allowbreak{}3.\allowbreak{}4.\allowbreak{}2} & Diseño de datos y restricciones del núcleo: Atención y escalamiento de planes vencidos — lote 1/1\newline Aceptación: lote versionado y trazable al flujo; revisión independiente sin defectos impeditivos.\newline Perfil D; 2027-01-29 a 2027-02-02. & 16 & Nicolás Torfan\\\hline
\texttt{1.\allowbreak{}4.\allowbreak{}1.\allowbreak{}3.\allowbreak{}4.\allowbreak{}3} & Diseño de contratos internos y composición: Atención y escalamiento de planes vencidos — lote 1/1\newline Aceptación: lote versionado y trazable al flujo; revisión independiente sin defectos impeditivos.\newline Perfil I; 2027-02-03 a 2027-02-05. & 23 & Fernando Guerrero\\\hline
\texttt{1.\allowbreak{}4.\allowbreak{}1.\allowbreak{}3.\allowbreak{}5.\allowbreak{}1} & Construcción funcional y pruebas unitarias: Atención y escalamiento de planes vencidos — lote 1/1\newline Aceptación: lote versionado y trazable al flujo; revisión independiente sin defectos impeditivos.\newline Perfil V; 2027-07-23 a 2027-07-28. & 61 & Nikolai Navea\\\hline
\texttt{1.\allowbreak{}4.\allowbreak{}1.\allowbreak{}3.\allowbreak{}6.\allowbreak{}1} & Verificación del núcleo, regresión y evidencia: Atención y escalamiento de planes vencidos — lote 1/1\newline Aceptación: lote versionado y trazable al flujo; revisión independiente sin defectos impeditivos.\newline Perfil Q; 2027-08-03 a 2027-08-05. & 45 & Davor Serey\\\hline
\texttt{1.\allowbreak{}4.\allowbreak{}1.\allowbreak{}3.\allowbreak{}7.\allowbreak{}1} & Coordinación y documentación del núcleo: Atención y escalamiento de planes vencidos — lote 1/1\newline Aceptación: lote versionado y trazable al flujo; revisión independiente sin defectos impeditivos.\newline Perfil J; 2027-08-05 a 2027-08-09. & 53 & Ignacio Silva\\\hline
\texttt{1.\allowbreak{}4.\allowbreak{}1.\allowbreak{}4.\allowbreak{}4.\allowbreak{}1} & Análisis de reglas, excepciones y protocolo: Autorización operacional de amarra y compatibilidad — lote 1/1\newline Aceptación: lote versionado y trazable al flujo; revisión independiente sin defectos impeditivos.\newline Perfil A; 2027-01-27 a 2027-01-29. & 16 & Ulysses Barreto\\\hline
\texttt{1.\allowbreak{}4.\allowbreak{}1.\allowbreak{}4.\allowbreak{}4.\allowbreak{}2} & Diseño de datos y restricciones del núcleo: Autorización operacional de amarra y compatibilidad — lote 1/1\newline Aceptación: lote versionado y trazable al flujo; revisión independiente sin defectos impeditivos.\newline Perfil D; 2027-01-29 a 2027-02-02. & 16 & Nicolás Torfan\\\hline
\texttt{1.\allowbreak{}4.\allowbreak{}1.\allowbreak{}4.\allowbreak{}4.\allowbreak{}3} & Diseño de contratos internos y composición: Autorización operacional de amarra y compatibilidad — lote 1/1\newline Aceptación: lote versionado y trazable al flujo; revisión independiente sin defectos impeditivos.\newline Perfil I; 2027-02-05 a 2027-02-08. & 23 & Fernando Guerrero\\\hline
\texttt{1.\allowbreak{}4.\allowbreak{}1.\allowbreak{}4.\allowbreak{}5.\allowbreak{}1} & Construcción funcional y pruebas unitarias: Autorización operacional de amarra y compatibilidad — lote 1/1\newline Aceptación: lote versionado y trazable al flujo; revisión independiente sin defectos impeditivos.\newline Perfil V; 2027-07-21 a 2027-07-23. & 61 & Nikolai Navea\\\hline
\texttt{1.\allowbreak{}4.\allowbreak{}1.\allowbreak{}4.\allowbreak{}6.\allowbreak{}1} & Verificación del núcleo, regresión y evidencia: Autorización operacional de amarra y compatibilidad — lote 1/1\newline Aceptación: lote versionado y trazable al flujo; revisión independiente sin defectos impeditivos.\newline Perfil Q; 2027-08-03 a 2027-08-05. & 45 & Davor Serey\\\hline
\texttt{1.\allowbreak{}4.\allowbreak{}1.\allowbreak{}4.\allowbreak{}7.\allowbreak{}1} & Coordinación y documentación del núcleo: Autorización operacional de amarra y compatibilidad — lote 1/1\newline Aceptación: lote versionado y trazable al flujo; revisión independiente sin defectos impeditivos.\newline Perfil J; 2027-08-06 a 2027-08-10. & 53 & Ignacio Silva\\\hline
\texttt{1.\allowbreak{}4.\allowbreak{}2.\allowbreak{}1.\allowbreak{}4.\allowbreak{}1} & Análisis de reglas, excepciones y protocolo: Clases, nóminas y participación por intervalo — lote 1/1\newline Aceptación: lote versionado y trazable al flujo; revisión independiente sin defectos impeditivos.\newline Perfil A; 2027-01-29 a 2027-02-02. & 51 & Ulysses Barreto\\\hline
\texttt{1.\allowbreak{}4.\allowbreak{}2.\allowbreak{}1.\allowbreak{}4.\allowbreak{}2} & Diseño de datos y restricciones del núcleo: Clases, nóminas y participación por intervalo — lote 1/1\newline Aceptación: lote versionado y trazable al flujo; revisión independiente sin defectos impeditivos.\newline Perfil D; 2027-02-02 a 2027-02-04. & 44 & Nicolás Torfan\\\hline
\texttt{1.\allowbreak{}4.\allowbreak{}2.\allowbreak{}1.\allowbreak{}4.\allowbreak{}3} & Diseño de contratos internos y composición: Clases, nóminas y participación por intervalo — lote 1/2\newline Aceptación: lote versionado y trazable al flujo; revisión independiente sin defectos impeditivos.\newline Perfil I; 2027-02-01 a 2027-02-03. & 33 & Fernando Guerrero\\\hline
\texttt{1.\allowbreak{}4.\allowbreak{}2.\allowbreak{}1.\allowbreak{}4.\allowbreak{}4} & Diseño de contratos internos y composición: Clases, nóminas y participación por intervalo — lote 2/2\newline Aceptación: lote versionado y trazable al flujo; revisión independiente sin defectos impeditivos.\newline Perfil I; 2027-02-03 a 2027-02-05. & 32 & Fernando Guerrero\\\hline
\texttt{1.\allowbreak{}4.\allowbreak{}2.\allowbreak{}1.\allowbreak{}5.\allowbreak{}1} & Construcción funcional y pruebas unitarias: Clases, nóminas y participación por intervalo — lote 1/4\newline Aceptación: lote versionado y trazable al flujo; revisión independiente sin defectos impeditivos.\newline Perfil V; 2027-07-23 a 2027-07-27. & 51 & Nikolai Navea\\\hline
\texttt{1.\allowbreak{}4.\allowbreak{}2.\allowbreak{}1.\allowbreak{}5.\allowbreak{}2} & Construcción funcional y pruebas unitarias: Clases, nóminas y participación por intervalo — lote 2/4\newline Aceptación: lote versionado y trazable al flujo; revisión independiente sin defectos impeditivos.\newline Perfil V; 2027-07-21 a 2027-07-23. & 51 & Nikolai Navea\\\hline
\texttt{1.\allowbreak{}4.\allowbreak{}2.\allowbreak{}1.\allowbreak{}5.\allowbreak{}3} & Construcción funcional y pruebas unitarias: Clases, nóminas y participación por intervalo — lote 3/4\newline Aceptación: lote versionado y trazable al flujo; revisión independiente sin defectos impeditivos.\newline Perfil V; 2027-07-23 a 2027-07-27. & 51 & Nikolai Navea\\\hline
\texttt{1.\allowbreak{}4.\allowbreak{}2.\allowbreak{}1.\allowbreak{}5.\allowbreak{}4} & Construcción funcional y pruebas unitarias: Clases, nóminas y participación por intervalo — lote 4/4\newline Aceptación: lote versionado y trazable al flujo; revisión independiente sin defectos impeditivos.\newline Perfil V; 2027-07-27 a 2027-07-29. & 50 & Nikolai Navea\\\hline
\texttt{1.\allowbreak{}4.\allowbreak{}2.\allowbreak{}1.\allowbreak{}6.\allowbreak{}1} & Verificación del núcleo, regresión y evidencia: Clases, nóminas y participación por intervalo — lote 1/2\newline Aceptación: lote versionado y trazable al flujo; revisión independiente sin defectos impeditivos.\newline Perfil Q; 2027-08-05 a 2027-08-09. & 49 & Davor Serey\\\hline
\texttt{1.\allowbreak{}4.\allowbreak{}2.\allowbreak{}1.\allowbreak{}6.\allowbreak{}2} & Verificación del núcleo, regresión y evidencia: Clases, nóminas y participación por intervalo — lote 2/2\newline Aceptación: lote versionado y trazable al flujo; revisión independiente sin defectos impeditivos.\newline Perfil Q; 2027-08-05 a 2027-08-09. & 49 & Davor Serey\\\hline
\texttt{1.\allowbreak{}4.\allowbreak{}2.\allowbreak{}1.\allowbreak{}7.\allowbreak{}1} & Coordinación y documentación del núcleo: Clases, nóminas y participación por intervalo — lote 1/2\newline Aceptación: lote versionado y trazable al flujo; revisión independiente sin defectos impeditivos.\newline Perfil J; 2027-08-09 a 2027-08-11. & 53 & Ignacio Silva\\\hline
\texttt{1.\allowbreak{}4.\allowbreak{}2.\allowbreak{}1.\allowbreak{}7.\allowbreak{}2} & Coordinación y documentación del núcleo: Clases, nóminas y participación por intervalo — lote 2/2\newline Aceptación: lote versionado y trazable al flujo; revisión independiente sin defectos impeditivos.\newline Perfil J; 2027-08-09 a 2027-08-11. & 53 & Ignacio Silva\\\hline
\texttt{1.\allowbreak{}4.\allowbreak{}2.\allowbreak{}2.\allowbreak{}4.\allowbreak{}1} & Análisis de reglas, excepciones y protocolo: Retiros coordinados y entrega única de alumno — lote 1/1\newline Aceptación: lote versionado y trazable al flujo; revisión independiente sin defectos impeditivos.\newline Perfil A; 2027-01-27 a 2027-01-29. & 15 & Ulysses Barreto\\\hline
\texttt{1.\allowbreak{}4.\allowbreak{}2.\allowbreak{}2.\allowbreak{}4.\allowbreak{}2} & Diseño de datos y restricciones del núcleo: Retiros coordinados y entrega única de alumno — lote 1/1\newline Aceptación: lote versionado y trazable al flujo; revisión independiente sin defectos impeditivos.\newline Perfil D; 2027-02-04 a 2027-02-08. & 16 & Nicolás Torfan\\\hline
\texttt{1.\allowbreak{}4.\allowbreak{}2.\allowbreak{}2.\allowbreak{}4.\allowbreak{}3} & Diseño de contratos internos y composición: Retiros coordinados y entrega única de alumno — lote 1/1\newline Aceptación: lote versionado y trazable al flujo; revisión independiente sin defectos impeditivos.\newline Perfil I; 2027-02-05 a 2027-02-09. & 23 & Fernando Guerrero\\\hline
\texttt{1.\allowbreak{}4.\allowbreak{}2.\allowbreak{}2.\allowbreak{}5.\allowbreak{}1} & Construcción funcional y pruebas unitarias: Retiros coordinados y entrega única de alumno — lote 1/1\newline Aceptación: lote versionado y trazable al flujo; revisión independiente sin defectos impeditivos.\newline Perfil V; 2027-07-22 a 2027-07-27. & 61 & Nikolai Navea\\\hline
\texttt{1.\allowbreak{}4.\allowbreak{}2.\allowbreak{}2.\allowbreak{}6.\allowbreak{}1} & Verificación del núcleo, regresión y evidencia: Retiros coordinados y entrega única de alumno — lote 1/1\newline Aceptación: lote versionado y trazable al flujo; revisión independiente sin defectos impeditivos.\newline Perfil Q; 2027-08-09 a 2027-08-11. & 45 & Davor Serey\\\hline
\texttt{1.\allowbreak{}4.\allowbreak{}2.\allowbreak{}2.\allowbreak{}7.\allowbreak{}1} & Coordinación y documentación del núcleo: Retiros coordinados y entrega única de alumno — lote 1/1\newline Aceptación: lote versionado y trazable al flujo; revisión independiente sin defectos impeditivos.\newline Perfil J; 2027-08-11 a 2027-08-13. & 53 & Ignacio Silva\\\hline
\texttt{1.\allowbreak{}4.\allowbreak{}2.\allowbreak{}3.\allowbreak{}4.\allowbreak{}1} & Análisis de reglas, excepciones y protocolo: Retorno de participantes y vista mínima de salud — lote 1/1\newline Aceptación: lote versionado y trazable al flujo; revisión independiente sin defectos impeditivos.\newline Perfil A; 2027-02-02 a 2027-02-03. & 51 & Ulysses Barreto\\\hline
\texttt{1.\allowbreak{}4.\allowbreak{}2.\allowbreak{}3.\allowbreak{}4.\allowbreak{}2} & Diseño de datos y restricciones del núcleo: Retorno de participantes y vista mínima de salud — lote 1/1\newline Aceptación: lote versionado y trazable al flujo; revisión independiente sin defectos impeditivos.\newline Perfil D; 2027-02-04 a 2027-02-08. & 44 & Nicolás Torfan\\\hline
\texttt{1.\allowbreak{}4.\allowbreak{}2.\allowbreak{}3.\allowbreak{}4.\allowbreak{}3} & Diseño de contratos internos y composición: Retorno de participantes y vista mínima de salud — lote 1/2\newline Aceptación: lote versionado y trazable al flujo; revisión independiente sin defectos impeditivos.\newline Perfil I; 2027-02-09 a 2027-02-10. & 33 & Fernando Guerrero\\\hline
\texttt{1.\allowbreak{}4.\allowbreak{}2.\allowbreak{}3.\allowbreak{}4.\allowbreak{}4} & Diseño de contratos internos y composición: Retorno de participantes y vista mínima de salud — lote 2/2\newline Aceptación: lote versionado y trazable al flujo; revisión independiente sin defectos impeditivos.\newline Perfil I; 2027-02-03 a 2027-02-05. & 32 & Fernando Guerrero\\\hline
\texttt{1.\allowbreak{}4.\allowbreak{}2.\allowbreak{}3.\allowbreak{}5.\allowbreak{}1} & Construcción funcional y pruebas unitarias: Retorno de participantes y vista mínima de salud — lote 1/4\newline Aceptación: lote versionado y trazable al flujo; revisión independiente sin defectos impeditivos.\newline Perfil V; 2027-07-27 a 2027-07-28. & 51 & Nikolai Navea\\\hline
\texttt{1.\allowbreak{}4.\allowbreak{}2.\allowbreak{}3.\allowbreak{}5.\allowbreak{}2} & Construcción funcional y pruebas unitarias: Retorno de participantes y vista mínima de salud — lote 2/4\newline Aceptación: lote versionado y trazable al flujo; revisión independiente sin defectos impeditivos.\newline Perfil V; 2027-07-28 a 2027-07-30. & 51 & Nikolai Navea\\\hline
\texttt{1.\allowbreak{}4.\allowbreak{}2.\allowbreak{}3.\allowbreak{}5.\allowbreak{}3} & Construcción funcional y pruebas unitarias: Retorno de participantes y vista mínima de salud — lote 3/4\newline Aceptación: lote versionado y trazable al flujo; revisión independiente sin defectos impeditivos.\newline Perfil V; 2027-07-26 a 2027-07-27. & 51 & Nikolai Navea\\\hline
\texttt{1.\allowbreak{}4.\allowbreak{}2.\allowbreak{}3.\allowbreak{}5.\allowbreak{}4} & Construcción funcional y pruebas unitarias: Retorno de participantes y vista mínima de salud — lote 4/4\newline Aceptación: lote versionado y trazable al flujo; revisión independiente sin defectos impeditivos.\newline Perfil V; 2027-07-27 a 2027-07-29. & 50 & Nikolai Navea\\\hline
\texttt{1.\allowbreak{}4.\allowbreak{}2.\allowbreak{}3.\allowbreak{}6.\allowbreak{}1} & Verificación del núcleo, regresión y evidencia: Retorno de participantes y vista mínima de salud — lote 1/2\newline Aceptación: lote versionado y trazable al flujo; revisión independiente sin defectos impeditivos.\newline Perfil Q; 2027-08-09 a 2027-08-11. & 49 & Davor Serey\\\hline
\texttt{1.\allowbreak{}4.\allowbreak{}2.\allowbreak{}3.\allowbreak{}6.\allowbreak{}2} & Verificación del núcleo, regresión y evidencia: Retorno de participantes y vista mínima de salud — lote 2/2\newline Aceptación: lote versionado y trazable al flujo; revisión independiente sin defectos impeditivos.\newline Perfil Q; 2027-08-09 a 2027-08-11. & 49 & Davor Serey\\\hline
\texttt{1.\allowbreak{}4.\allowbreak{}2.\allowbreak{}3.\allowbreak{}7.\allowbreak{}1} & Coordinación y documentación del núcleo: Retorno de participantes y vista mínima de salud — lote 1/2\newline Aceptación: lote versionado y trazable al flujo; revisión independiente sin defectos impeditivos.\newline Perfil J; 2027-08-13 a 2027-08-17. & 53 & Ignacio Silva\\\hline
\texttt{1.\allowbreak{}4.\allowbreak{}2.\allowbreak{}3.\allowbreak{}7.\allowbreak{}2} & Coordinación y documentación del núcleo: Retorno de participantes y vista mínima de salud — lote 2/2\newline Aceptación: lote versionado y trazable al flujo; revisión independiente sin defectos impeditivos.\newline Perfil J; 2027-08-13 a 2027-08-17. & 52 & Ignacio Silva\\\hline
\texttt{1.\allowbreak{}4.\allowbreak{}3.\allowbreak{}1.\allowbreak{}4.\allowbreak{}1} & Análisis de reglas, excepciones y protocolo: Inventario y condiciones de puestos y fondeos — lote 1/1\newline Aceptación: lote versionado y trazable al flujo; revisión independiente sin defectos impeditivos.\newline Perfil A; 2027-01-27 a 2027-01-29. & 15 & Ulysses Barreto\\\hline
\texttt{1.\allowbreak{}4.\allowbreak{}3.\allowbreak{}1.\allowbreak{}4.\allowbreak{}2} & Diseño de datos y restricciones del núcleo: Inventario y condiciones de puestos y fondeos — lote 1/1\newline Aceptación: lote versionado y trazable al flujo; revisión independiente sin defectos impeditivos.\newline Perfil D; 2027-01-29 a 2027-02-02. & 16 & Nicolás Torfan\\\hline
\texttt{1.\allowbreak{}4.\allowbreak{}3.\allowbreak{}1.\allowbreak{}4.\allowbreak{}3} & Diseño de contratos internos y composición: Inventario y condiciones de puestos y fondeos — lote 1/1\newline Aceptación: lote versionado y trazable al flujo; revisión independiente sin defectos impeditivos.\newline Perfil I; 2027-02-05 a 2027-02-09. & 23 & Fernando Guerrero\\\hline
\texttt{1.\allowbreak{}4.\allowbreak{}3.\allowbreak{}1.\allowbreak{}5.\allowbreak{}1} & Construcción funcional y pruebas unitarias: Inventario y condiciones de puestos y fondeos — lote 1/1\newline Aceptación: lote versionado y trazable al flujo; revisión independiente sin defectos impeditivos.\newline Perfil V; 2027-07-29 a 2027-08-02. & 61 & Nikolai Navea\\\hline
\texttt{1.\allowbreak{}4.\allowbreak{}3.\allowbreak{}1.\allowbreak{}6.\allowbreak{}1} & Verificación del núcleo, regresión y evidencia: Inventario y condiciones de puestos y fondeos — lote 1/1\newline Aceptación: lote versionado y trazable al flujo; revisión independiente sin defectos impeditivos.\newline Perfil Q; 2027-08-11 a 2027-08-13. & 45 & Davor Serey\\\hline
\texttt{1.\allowbreak{}4.\allowbreak{}3.\allowbreak{}1.\allowbreak{}7.\allowbreak{}1} & Coordinación y documentación del núcleo: Inventario y condiciones de puestos y fondeos — lote 1/1\newline Aceptación: lote versionado y trazable al flujo; revisión independiente sin defectos impeditivos.\newline Perfil J; 2027-08-13 a 2027-08-17. & 53 & Ignacio Silva\\\hline
\texttt{1.\allowbreak{}4.\allowbreak{}3.\allowbreak{}2.\allowbreak{}4.\allowbreak{}1} & Análisis de reglas, excepciones y protocolo: Inspecciones de boyas y evidencia desconectada — lote 1/1\newline Aceptación: lote versionado y trazable al flujo; revisión independiente sin defectos impeditivos.\newline Perfil A; 2027-02-03 a 2027-02-05. & 51 & Ulysses Barreto\\\hline
\texttt{1.\allowbreak{}4.\allowbreak{}3.\allowbreak{}2.\allowbreak{}4.\allowbreak{}2} & Diseño de datos y restricciones del núcleo: Inspecciones de boyas y evidencia desconectada — lote 1/1\newline Aceptación: lote versionado y trazable al flujo; revisión independiente sin defectos impeditivos.\newline Perfil D; 2027-02-08 a 2027-02-10. & 44 & Nicolás Torfan\\\hline
\texttt{1.\allowbreak{}4.\allowbreak{}3.\allowbreak{}2.\allowbreak{}4.\allowbreak{}3} & Diseño de contratos internos y composición: Inspecciones de boyas y evidencia desconectada — lote 1/2\newline Aceptación: lote versionado y trazable al flujo; revisión independiente sin defectos impeditivos.\newline Perfil I; 2027-02-01 a 2027-02-03. & 33 & Fernando Guerrero\\\hline
\texttt{1.\allowbreak{}4.\allowbreak{}3.\allowbreak{}2.\allowbreak{}4.\allowbreak{}4} & Diseño de contratos internos y composición: Inspecciones de boyas y evidencia desconectada — lote 2/2\newline Aceptación: lote versionado y trazable al flujo; revisión independiente sin defectos impeditivos.\newline Perfil I; 2027-02-03 a 2027-02-05. & 32 & Fernando Guerrero\\\hline
\texttt{1.\allowbreak{}4.\allowbreak{}3.\allowbreak{}2.\allowbreak{}5.\allowbreak{}1} & Construcción funcional y pruebas unitarias: Inspecciones de boyas y evidencia desconectada — lote 1/4\newline Aceptación: lote versionado y trazable al flujo; revisión independiente sin defectos impeditivos.\newline Perfil V; 2027-07-26 a 2027-07-27. & 51 & Nikolai Navea\\\hline
\texttt{1.\allowbreak{}4.\allowbreak{}3.\allowbreak{}2.\allowbreak{}5.\allowbreak{}2} & Construcción funcional y pruebas unitarias: Inspecciones de boyas y evidencia desconectada — lote 2/4\newline Aceptación: lote versionado y trazable al flujo; revisión independiente sin defectos impeditivos.\newline Perfil V; 2027-07-27 a 2027-07-29. & 51 & Nikolai Navea\\\hline
\texttt{1.\allowbreak{}4.\allowbreak{}3.\allowbreak{}2.\allowbreak{}5.\allowbreak{}3} & Construcción funcional y pruebas unitarias: Inspecciones de boyas y evidencia desconectada — lote 3/4\newline Aceptación: lote versionado y trazable al flujo; revisión independiente sin defectos impeditivos.\newline Perfil V; 2027-07-29 a 2027-08-02. & 51 & Nikolai Navea\\\hline
\texttt{1.\allowbreak{}4.\allowbreak{}3.\allowbreak{}2.\allowbreak{}5.\allowbreak{}4} & Construcción funcional y pruebas unitarias: Inspecciones de boyas y evidencia desconectada — lote 4/4\newline Aceptación: lote versionado y trazable al flujo; revisión independiente sin defectos impeditivos.\newline Perfil V; 2027-07-26 a 2027-07-27. & 50 & Nikolai Navea\\\hline
\texttt{1.\allowbreak{}4.\allowbreak{}3.\allowbreak{}2.\allowbreak{}6.\allowbreak{}1} & Verificación del núcleo, regresión y evidencia: Inspecciones de boyas y evidencia desconectada — lote 1/2\newline Aceptación: lote versionado y trazable al flujo; revisión independiente sin defectos impeditivos.\newline Perfil Q; 2027-08-11 a 2027-08-13. & 49 & Davor Serey\\\hline
\texttt{1.\allowbreak{}4.\allowbreak{}3.\allowbreak{}2.\allowbreak{}6.\allowbreak{}2} & Verificación del núcleo, regresión y evidencia: Inspecciones de boyas y evidencia desconectada — lote 2/2\newline Aceptación: lote versionado y trazable al flujo; revisión independiente sin defectos impeditivos.\newline Perfil Q; 2027-08-05 a 2027-08-09. & 48 & Davor Serey\\\hline
\texttt{1.\allowbreak{}4.\allowbreak{}3.\allowbreak{}2.\allowbreak{}7.\allowbreak{}1} & Coordinación y documentación del núcleo: Inspecciones de boyas y evidencia desconectada — lote 1/2\newline Aceptación: lote versionado y trazable al flujo; revisión independiente sin defectos impeditivos.\newline Perfil J; 2027-08-17 a 2027-08-18. & 53 & Ignacio Silva\\\hline
\texttt{1.\allowbreak{}4.\allowbreak{}3.\allowbreak{}2.\allowbreak{}7.\allowbreak{}2} & Coordinación y documentación del núcleo: Inspecciones de boyas y evidencia desconectada — lote 2/2\newline Aceptación: lote versionado y trazable al flujo; revisión independiente sin defectos impeditivos.\newline Perfil J; 2027-08-17 a 2027-08-18. & 52 & Ignacio Silva\\\hline
\texttt{1.\allowbreak{}4.\allowbreak{}3.\allowbreak{}3.\allowbreak{}4.\allowbreak{}1} & Análisis de reglas, excepciones y protocolo: Asociaciones temporales de medidores e instalaciones — lote 1/1\newline Aceptación: lote versionado y trazable al flujo; revisión independiente sin defectos impeditivos.\newline Perfil A; 2027-01-29 a 2027-02-02. & 15 & Ulysses Barreto\\\hline
\texttt{1.\allowbreak{}4.\allowbreak{}3.\allowbreak{}3.\allowbreak{}4.\allowbreak{}2} & Diseño de datos y restricciones del núcleo: Asociaciones temporales de medidores e instalaciones — lote 1/1\newline Aceptación: lote versionado y trazable al flujo; revisión independiente sin defectos impeditivos.\newline Perfil D; 2027-02-01 a 2027-02-03. & 16 & Nicolás Torfan\\\hline
\texttt{1.\allowbreak{}4.\allowbreak{}3.\allowbreak{}3.\allowbreak{}4.\allowbreak{}3} & Diseño de contratos internos y composición: Asociaciones temporales de medidores e instalaciones — lote 1/1\newline Aceptación: lote versionado y trazable al flujo; revisión independiente sin defectos impeditivos.\newline Perfil I; 2027-02-05 a 2027-02-08. & 23 & Fernando Guerrero\\\hline
\texttt{1.\allowbreak{}4.\allowbreak{}3.\allowbreak{}3.\allowbreak{}5.\allowbreak{}1} & Construcción funcional y pruebas unitarias: Asociaciones temporales de medidores e instalaciones — lote 1/1\newline Aceptación: lote versionado y trazable al flujo; revisión independiente sin defectos impeditivos.\newline Perfil V; 2027-07-27 a 2027-07-29. & 61 & Nikolai Navea\\\hline
\texttt{1.\allowbreak{}4.\allowbreak{}3.\allowbreak{}3.\allowbreak{}6.\allowbreak{}1} & Verificación del núcleo, regresión y evidencia: Asociaciones temporales de medidores e instalaciones — lote 1/1\newline Aceptación: lote versionado y trazable al flujo; revisión independiente sin defectos impeditivos.\newline Perfil Q; 2027-08-03 a 2027-08-05. & 45 & Davor Serey\\\hline
\texttt{1.\allowbreak{}4.\allowbreak{}3.\allowbreak{}3.\allowbreak{}7.\allowbreak{}1} & Coordinación y documentación del núcleo: Asociaciones temporales de medidores e instalaciones — lote 1/1\newline Aceptación: lote versionado y trazable al flujo; revisión independiente sin defectos impeditivos.\newline Perfil J; 2027-08-05 a 2027-08-09. & 53 & Ignacio Silva\\\hline
\texttt{1.\allowbreak{}5.\allowbreak{}1.\allowbreak{}1.\allowbreak{}4.\allowbreak{}1} & Análisis de reglas, excepciones y protocolo: Maestros de faena y habilitaciones de consulta inicial — lote 1/1\newline Aceptación: lote versionado y trazable al flujo; revisión independiente sin defectos impeditivos.\newline Perfil A; 2027-01-29 a 2027-02-02. & 15 & Ulysses Barreto\\\hline
\texttt{1.\allowbreak{}5.\allowbreak{}1.\allowbreak{}1.\allowbreak{}4.\allowbreak{}2} & Diseño de datos y restricciones del núcleo: Maestros de faena y habilitaciones de consulta inicial — lote 1/1\newline Aceptación: lote versionado y trazable al flujo; revisión independiente sin defectos impeditivos.\newline Perfil D; 2027-02-01 a 2027-02-03. & 16 & Nicolás Torfan\\\hline
\texttt{1.\allowbreak{}5.\allowbreak{}1.\allowbreak{}1.\allowbreak{}4.\allowbreak{}3} & Diseño de contratos internos y composición: Maestros de faena y habilitaciones de consulta inicial — lote 1/1\newline Aceptación: lote versionado y trazable al flujo; revisión independiente sin defectos impeditivos.\newline Perfil I; 2027-02-08 a 2027-02-10. & 23 & Fernando Guerrero\\\hline
\texttt{1.\allowbreak{}5.\allowbreak{}1.\allowbreak{}1.\allowbreak{}5.\allowbreak{}1} & Construcción funcional y pruebas unitarias: Maestros de faena y habilitaciones de consulta inicial — lote 1/1\newline Aceptación: lote versionado y trazable al flujo; revisión independiente sin defectos impeditivos.\newline Perfil V; 2027-07-29 a 2027-08-02. & 61 & Nikolai Navea\\\hline
\texttt{1.\allowbreak{}5.\allowbreak{}1.\allowbreak{}1.\allowbreak{}6.\allowbreak{}1} & Verificación del núcleo, regresión y evidencia: Maestros de faena y habilitaciones de consulta inicial — lote 1/1\newline Aceptación: lote versionado y trazable al flujo; revisión independiente sin defectos impeditivos.\newline Perfil Q; 2027-08-09 a 2027-08-12. & 45 & Davor Serey\\\hline
\texttt{1.\allowbreak{}5.\allowbreak{}1.\allowbreak{}1.\allowbreak{}7.\allowbreak{}1} & Coordinación y documentación del núcleo: Maestros de faena y habilitaciones de consulta inicial — lote 1/1\newline Aceptación: lote versionado y trazable al flujo; revisión independiente sin defectos impeditivos.\newline Perfil J; 2027-08-12 a 2027-08-13. & 53 & Ignacio Silva\\\hline
\texttt{1.\allowbreak{}5.\allowbreak{}2.\allowbreak{}1.\allowbreak{}4.\allowbreak{}1} & Análisis de reglas, excepciones y protocolo: Despacho identificado y captura segura de combustible — lote 1/1\newline Aceptación: lote versionado y trazable al flujo; revisión independiente sin defectos impeditivos.\newline Perfil A; 2027-01-29 a 2027-02-02. & 15 & Ulysses Barreto\\\hline
\texttt{1.\allowbreak{}5.\allowbreak{}2.\allowbreak{}1.\allowbreak{}4.\allowbreak{}2} & Diseño de datos y restricciones del núcleo: Despacho identificado y captura segura de combustible — lote 1/1\newline Aceptación: lote versionado y trazable al flujo; revisión independiente sin defectos impeditivos.\newline Perfil D; 2027-02-03 a 2027-02-05. & 16 & Nicolás Torfan\\\hline
\texttt{1.\allowbreak{}5.\allowbreak{}2.\allowbreak{}1.\allowbreak{}4.\allowbreak{}3} & Diseño de contratos internos y composición: Despacho identificado y captura segura de combustible — lote 1/1\newline Aceptación: lote versionado y trazable al flujo; revisión independiente sin defectos impeditivos.\newline Perfil I; 2027-01-27 a 2027-01-28. & 23 & Fernando Guerrero\\\hline
\texttt{1.\allowbreak{}5.\allowbreak{}2.\allowbreak{}1.\allowbreak{}5.\allowbreak{}1} & Construcción funcional y pruebas unitarias: Despacho identificado y captura segura de combustible — lote 1/1\newline Aceptación: lote versionado y trazable al flujo; revisión independiente sin defectos impeditivos.\newline Perfil V; 2027-07-21 a 2027-07-23. & 61 & Nikolai Navea\\\hline
\texttt{1.\allowbreak{}5.\allowbreak{}2.\allowbreak{}1.\allowbreak{}6.\allowbreak{}1} & Verificación del núcleo, regresión y evidencia: Despacho identificado y captura segura de combustible — lote 1/1\newline Aceptación: lote versionado y trazable al flujo; revisión independiente sin defectos impeditivos.\newline Perfil Q; 2027-08-02 a 2027-08-04. & 45 & Davor Serey\\\hline
\texttt{1.\allowbreak{}5.\allowbreak{}2.\allowbreak{}1.\allowbreak{}7.\allowbreak{}1} & Coordinación y documentación del núcleo: Despacho identificado y captura segura de combustible — lote 1/1\newline Aceptación: lote versionado y trazable al flujo; revisión independiente sin defectos impeditivos.\newline Perfil J; 2027-08-04 a 2027-08-06. & 53 & Ignacio Silva\\\hline
\texttt{1.\allowbreak{}5.\allowbreak{}2.\allowbreak{}2.\allowbreak{}4.\allowbreak{}1} & Análisis de reglas, excepciones y protocolo: Conciliación de tramos y rectificación de despachos — lote 1/1\newline Aceptación: lote versionado y trazable al flujo; revisión independiente sin defectos impeditivos.\newline Perfil A; 2027-01-27 a 2027-01-29. & 15 & Ulysses Barreto\\\hline
\texttt{1.\allowbreak{}5.\allowbreak{}2.\allowbreak{}2.\allowbreak{}4.\allowbreak{}2} & Diseño de datos y restricciones del núcleo: Conciliación de tramos y rectificación de despachos — lote 1/1\newline Aceptación: lote versionado y trazable al flujo; revisión independiente sin defectos impeditivos.\newline Perfil D; 2027-02-03 a 2027-02-05. & 16 & Nicolás Torfan\\\hline
\texttt{1.\allowbreak{}5.\allowbreak{}2.\allowbreak{}2.\allowbreak{}4.\allowbreak{}3} & Diseño de contratos internos y composición: Conciliación de tramos y rectificación de despachos — lote 1/1\newline Aceptación: lote versionado y trazable al flujo; revisión independiente sin defectos impeditivos.\newline Perfil I; 2027-01-27 a 2027-01-28. & 23 & Fernando Guerrero\\\hline
\texttt{1.\allowbreak{}5.\allowbreak{}2.\allowbreak{}2.\allowbreak{}5.\allowbreak{}1} & Construcción funcional y pruebas unitarias: Conciliación de tramos y rectificación de despachos — lote 1/1\newline Aceptación: lote versionado y trazable al flujo; revisión independiente sin defectos impeditivos.\newline Perfil V; 2027-07-23 a 2027-07-28. & 61 & Nikolai Navea\\\hline
\texttt{1.\allowbreak{}5.\allowbreak{}2.\allowbreak{}2.\allowbreak{}6.\allowbreak{}1} & Verificación del núcleo, regresión y evidencia: Conciliación de tramos y rectificación de despachos — lote 1/1\newline Aceptación: lote versionado y trazable al flujo; revisión independiente sin defectos impeditivos.\newline Perfil Q; 2027-08-04 a 2027-08-06. & 45 & Davor Serey\\\hline
\texttt{1.\allowbreak{}5.\allowbreak{}2.\allowbreak{}2.\allowbreak{}7.\allowbreak{}1} & Coordinación y documentación del núcleo: Conciliación de tramos y rectificación de despachos — lote 1/1\newline Aceptación: lote versionado y trazable al flujo; revisión independiente sin defectos impeditivos.\newline Perfil J; 2027-08-09 a 2027-08-11. & 53 & Ignacio Silva\\\hline
\texttt{1.\allowbreak{}5.\allowbreak{}3.\allowbreak{}1.\allowbreak{}4.\allowbreak{}1} & Análisis de reglas, excepciones y protocolo: Intervalos y series conciliadas de agua y energía — lote 1/1\newline Aceptación: lote versionado y trazable al flujo; revisión independiente sin defectos impeditivos.\newline Perfil A; 2027-01-29 a 2027-02-02. & 15 & Ulysses Barreto\\\hline
\texttt{1.\allowbreak{}5.\allowbreak{}3.\allowbreak{}1.\allowbreak{}4.\allowbreak{}2} & Diseño de datos y restricciones del núcleo: Intervalos y series conciliadas de agua y energía — lote 1/1\newline Aceptación: lote versionado y trazable al flujo; revisión independiente sin defectos impeditivos.\newline Perfil D; 2027-01-29 a 2027-02-02. & 15 & Nicolás Torfan\\\hline
\texttt{1.\allowbreak{}5.\allowbreak{}3.\allowbreak{}1.\allowbreak{}4.\allowbreak{}3} & Diseño de contratos internos y composición: Intervalos y series conciliadas de agua y energía — lote 1/1\newline Aceptación: lote versionado y trazable al flujo; revisión independiente sin defectos impeditivos.\newline Perfil I; 2027-01-29 a 2027-02-02. & 23 & Fernando Guerrero\\\hline
\texttt{1.\allowbreak{}5.\allowbreak{}3.\allowbreak{}1.\allowbreak{}5.\allowbreak{}1} & Construcción funcional y pruebas unitarias: Intervalos y series conciliadas de agua y energía — lote 1/1\newline Aceptación: lote versionado y trazable al flujo; revisión independiente sin defectos impeditivos.\newline Perfil V; 2027-07-28 a 2027-07-30. & 61 & Nikolai Navea\\\hline
\texttt{1.\allowbreak{}5.\allowbreak{}3.\allowbreak{}1.\allowbreak{}6.\allowbreak{}1} & Verificación del núcleo, regresión y evidencia: Intervalos y series conciliadas de agua y energía — lote 1/1\newline Aceptación: lote versionado y trazable al flujo; revisión independiente sin defectos impeditivos.\newline Perfil Q; 2027-08-06 a 2027-08-10. & 45 & Davor Serey\\\hline
\texttt{1.\allowbreak{}5.\allowbreak{}3.\allowbreak{}1.\allowbreak{}7.\allowbreak{}1} & Coordinación y documentación del núcleo: Intervalos y series conciliadas de agua y energía — lote 1/1\newline Aceptación: lote versionado y trazable al flujo; revisión independiente sin defectos impeditivos.\newline Perfil J; 2027-08-13 a 2027-08-17. & 53 & Ignacio Silva\\\hline
\texttt{1.\allowbreak{}5.\allowbreak{}3.\allowbreak{}2.\allowbreak{}4.\allowbreak{}1} & Análisis de reglas, excepciones y protocolo: Consultas de consumo y evidencia ambiental inicial — lote 1/1\newline Aceptación: lote versionado y trazable al flujo; revisión independiente sin defectos impeditivos.\newline Perfil A; 2027-01-29 a 2027-02-02. & 15 & Ulysses Barreto\\\hline
\texttt{1.\allowbreak{}5.\allowbreak{}3.\allowbreak{}2.\allowbreak{}4.\allowbreak{}2} & Diseño de datos y restricciones del núcleo: Consultas de consumo y evidencia ambiental inicial — lote 1/1\newline Aceptación: lote versionado y trazable al flujo; revisión independiente sin defectos impeditivos.\newline Perfil D; 2027-01-27 a 2027-01-29. & 15 & Nicolás Torfan\\\hline
\texttt{1.\allowbreak{}5.\allowbreak{}3.\allowbreak{}2.\allowbreak{}4.\allowbreak{}3} & Diseño de contratos internos y composición: Consultas de consumo y evidencia ambiental inicial — lote 1/1\newline Aceptación: lote versionado y trazable al flujo; revisión independiente sin defectos impeditivos.\newline Perfil I; 2027-01-27 a 2027-01-28. & 23 & Fernando Guerrero\\\hline
\texttt{1.\allowbreak{}5.\allowbreak{}3.\allowbreak{}2.\allowbreak{}5.\allowbreak{}1} & Construcción funcional y pruebas unitarias: Consultas de consumo y evidencia ambiental inicial — lote 1/1\newline Aceptación: lote versionado y trazable al flujo; revisión independiente sin defectos impeditivos.\newline Perfil V; 2027-07-21 a 2027-07-23. & 61 & Nikolai Navea\\\hline
\texttt{1.\allowbreak{}5.\allowbreak{}3.\allowbreak{}2.\allowbreak{}6.\allowbreak{}1} & Verificación del núcleo, regresión y evidencia: Consultas de consumo y evidencia ambiental inicial — lote 1/1\newline Aceptación: lote versionado y trazable al flujo; revisión independiente sin defectos impeditivos.\newline Perfil Q; 2027-08-02 a 2027-08-04. & 45 & Davor Serey\\\hline
\texttt{1.\allowbreak{}5.\allowbreak{}3.\allowbreak{}2.\allowbreak{}7.\allowbreak{}1} & Coordinación y documentación del núcleo: Consultas de consumo y evidencia ambiental inicial — lote 1/1\newline Aceptación: lote versionado y trazable al flujo; revisión independiente sin defectos impeditivos.\newline Perfil J; 2027-08-04 a 2027-08-06. & 53 & Ignacio Silva\\\hline
\texttt{1.\allowbreak{}6.\allowbreak{}1.\allowbreak{}1.\allowbreak{}4.\allowbreak{}1} & Análisis de reglas, excepciones y protocolo: Maestros de personas, organizaciones y embarcaciones — lote 1/1\newline Aceptación: lote versionado y trazable al flujo; revisión independiente sin defectos impeditivos.\newline Perfil A; 2027-01-27 a 2027-01-29. & 15 & Ulysses Barreto\\\hline
\texttt{1.\allowbreak{}6.\allowbreak{}1.\allowbreak{}1.\allowbreak{}4.\allowbreak{}2} & Diseño de datos y restricciones del núcleo: Maestros de personas, organizaciones y embarcaciones — lote 1/1\newline Aceptación: lote versionado y trazable al flujo; revisión independiente sin defectos impeditivos.\newline Perfil D; 2027-01-27 a 2027-01-29. & 15 & Nicolás Torfan\\\hline
\texttt{1.\allowbreak{}6.\allowbreak{}1.\allowbreak{}1.\allowbreak{}4.\allowbreak{}3} & Diseño de contratos internos y composición: Maestros de personas, organizaciones y embarcaciones — lote 1/1\newline Aceptación: lote versionado y trazable al flujo; revisión independiente sin defectos impeditivos.\newline Perfil I; 2027-02-02 a 2027-02-03. & 23 & Fernando Guerrero\\\hline
\texttt{1.\allowbreak{}6.\allowbreak{}1.\allowbreak{}1.\allowbreak{}5.\allowbreak{}1} & Construcción funcional y pruebas unitarias: Maestros de personas, organizaciones y embarcaciones — lote 1/1\newline Aceptación: lote versionado y trazable al flujo; revisión independiente sin defectos impeditivos.\newline Perfil V; 2027-07-23 a 2027-07-28. & 61 & Nikolai Navea\\\hline
\texttt{1.\allowbreak{}6.\allowbreak{}1.\allowbreak{}1.\allowbreak{}6.\allowbreak{}1} & Verificación del núcleo, regresión y evidencia: Maestros de personas, organizaciones y embarcaciones — lote 1/1\newline Aceptación: lote versionado y trazable al flujo; revisión independiente sin defectos impeditivos.\newline Perfil Q; 2027-08-10 a 2027-08-12. & 45 & Davor Serey\\\hline
\texttt{1.\allowbreak{}6.\allowbreak{}1.\allowbreak{}1.\allowbreak{}7.\allowbreak{}1} & Coordinación y documentación del núcleo: Maestros de personas, organizaciones y embarcaciones — lote 1/1\newline Aceptación: lote versionado y trazable al flujo; revisión independiente sin defectos impeditivos.\newline Perfil J; 2027-08-12 a 2027-08-16. & 53 & Ignacio Silva\\\hline
\texttt{1.\allowbreak{}6.\allowbreak{}1.\allowbreak{}2.\allowbreak{}4.\allowbreak{}1} & Análisis de reglas, excepciones y protocolo: Contratos y documentos vigentes de consulta E1 — lote 1/1\newline Aceptación: lote versionado y trazable al flujo; revisión independiente sin defectos impeditivos.\newline Perfil A; 2027-01-27 a 2027-01-29. & 15 & Ulysses Barreto\\\hline
\texttt{1.\allowbreak{}6.\allowbreak{}1.\allowbreak{}2.\allowbreak{}4.\allowbreak{}2} & Diseño de datos y restricciones del núcleo: Contratos y documentos vigentes de consulta E1 — lote 1/1\newline Aceptación: lote versionado y trazable al flujo; revisión independiente sin defectos impeditivos.\newline Perfil D; 2027-02-02 a 2027-02-04. & 15 & Nicolás Torfan\\\hline
\texttt{1.\allowbreak{}6.\allowbreak{}1.\allowbreak{}2.\allowbreak{}4.\allowbreak{}3} & Diseño de contratos internos y composición: Contratos y documentos vigentes de consulta E1 — lote 1/1\newline Aceptación: lote versionado y trazable al flujo; revisión independiente sin defectos impeditivos.\newline Perfil I; 2027-02-03 a 2027-02-05. & 23 & Fernando Guerrero\\\hline
\texttt{1.\allowbreak{}6.\allowbreak{}1.\allowbreak{}2.\allowbreak{}5.\allowbreak{}1} & Construcción funcional y pruebas unitarias: Contratos y documentos vigentes de consulta E1 — lote 1/1\newline Aceptación: lote versionado y trazable al flujo; revisión independiente sin defectos impeditivos.\newline Perfil V; 2027-07-28 a 2027-07-30. & 61 & Nikolai Navea\\\hline
\texttt{1.\allowbreak{}6.\allowbreak{}1.\allowbreak{}2.\allowbreak{}6.\allowbreak{}1} & Verificación del núcleo, regresión y evidencia: Contratos y documentos vigentes de consulta E1 — lote 1/1\newline Aceptación: lote versionado y trazable al flujo; revisión independiente sin defectos impeditivos.\newline Perfil Q; 2027-08-03 a 2027-08-05. & 45 & Davor Serey\\\hline
\texttt{1.\allowbreak{}6.\allowbreak{}1.\allowbreak{}2.\allowbreak{}7.\allowbreak{}1} & Coordinación y documentación del núcleo: Contratos y documentos vigentes de consulta E1 — lote 1/1\newline Aceptación: lote versionado y trazable al flujo; revisión independiente sin defectos impeditivos.\newline Perfil J; 2027-08-06 a 2027-08-10. & 53 & Ignacio Silva\\\hline
\texttt{1.\allowbreak{}6.\allowbreak{}1.\allowbreak{}3.\allowbreak{}4.\allowbreak{}1} & Análisis de reglas, excepciones y protocolo: Atención a visitas y asignación transitoria inicial — lote 1/1\newline Aceptación: lote versionado y trazable al flujo; revisión independiente sin defectos impeditivos.\newline Perfil A; 2027-01-27 a 2027-01-29. & 15 & Ulysses Barreto\\\hline
\texttt{1.\allowbreak{}6.\allowbreak{}1.\allowbreak{}3.\allowbreak{}4.\allowbreak{}2} & Diseño de datos y restricciones del núcleo: Atención a visitas y asignación transitoria inicial — lote 1/1\newline Aceptación: lote versionado y trazable al flujo; revisión independiente sin defectos impeditivos.\newline Perfil D; 2027-02-04 a 2027-02-08. & 15 & Nicolás Torfan\\\hline
\texttt{1.\allowbreak{}6.\allowbreak{}1.\allowbreak{}3.\allowbreak{}4.\allowbreak{}3} & Diseño de contratos internos y composición: Atención a visitas y asignación transitoria inicial — lote 1/1\newline Aceptación: lote versionado y trazable al flujo; revisión independiente sin defectos impeditivos.\newline Perfil I; 2027-02-05 a 2027-02-08. & 23 & Fernando Guerrero\\\hline
\texttt{1.\allowbreak{}6.\allowbreak{}1.\allowbreak{}3.\allowbreak{}5.\allowbreak{}1} & Construcción funcional y pruebas unitarias: Atención a visitas y asignación transitoria inicial — lote 1/1\newline Aceptación: lote versionado y trazable al flujo; revisión independiente sin defectos impeditivos.\newline Perfil V; 2027-07-21 a 2027-07-23. & 61 & Nikolai Navea\\\hline
\texttt{1.\allowbreak{}6.\allowbreak{}1.\allowbreak{}3.\allowbreak{}6.\allowbreak{}1} & Verificación del núcleo, regresión y evidencia: Atención a visitas y asignación transitoria inicial — lote 1/1\newline Aceptación: lote versionado y trazable al flujo; revisión independiente sin defectos impeditivos.\newline Perfil Q; 2027-07-28 a 2027-07-30. & 45 & Davor Serey\\\hline
\texttt{1.\allowbreak{}6.\allowbreak{}1.\allowbreak{}3.\allowbreak{}7.\allowbreak{}1} & Coordinación y documentación del núcleo: Atención a visitas y asignación transitoria inicial — lote 1/1\newline Aceptación: lote versionado y trazable al flujo; revisión independiente sin defectos impeditivos.\newline Perfil J; 2027-08-02 a 2027-08-04. & 53 & Ignacio Silva\\\hline
\texttt{1.\allowbreak{}5.\allowbreak{}1.\allowbreak{}2.\allowbreak{}4.\allowbreak{}1} & Análisis de reglas, excepciones y protocolo: Agenda de varadero y posiciones compatibles — lote 1/1\newline Aceptación: lote versionado y trazable al flujo; revisión independiente sin defectos impeditivos.\newline Perfil A; 2027-12-03 a 2027-12-07. & 16 & Ulysses Barreto\\\hline
\texttt{1.\allowbreak{}5.\allowbreak{}1.\allowbreak{}2.\allowbreak{}4.\allowbreak{}2} & Diseño de datos y restricciones del núcleo: Agenda de varadero y posiciones compatibles — lote 1/1\newline Aceptación: lote versionado y trazable al flujo; revisión independiente sin defectos impeditivos.\newline Perfil D; 2027-12-03 a 2027-12-07. & 16 & Nicolás Torfan\\\hline
\texttt{1.\allowbreak{}5.\allowbreak{}1.\allowbreak{}2.\allowbreak{}4.\allowbreak{}3} & Diseño de contratos internos y composición: Agenda de varadero y posiciones compatibles — lote 1/1\newline Aceptación: lote versionado y trazable al flujo; revisión independiente sin defectos impeditivos.\newline Perfil I; 2027-12-03 a 2027-12-06. & 23 & Fernando Guerrero\\\hline
\texttt{1.\allowbreak{}5.\allowbreak{}1.\allowbreak{}2.\allowbreak{}5.\allowbreak{}1} & Construcción funcional y pruebas unitarias: Agenda de varadero y posiciones compatibles — lote 1/1\newline Aceptación: lote versionado y trazable al flujo; revisión independiente sin defectos impeditivos.\newline Perfil V; 2028-01-20 a 2028-01-24. & 62 & Nikolai Navea\\\hline
\texttt{1.\allowbreak{}5.\allowbreak{}1.\allowbreak{}2.\allowbreak{}6.\allowbreak{}1} & Verificación del núcleo, regresión y evidencia: Agenda de varadero y posiciones compatibles — lote 1/1\newline Aceptación: lote versionado y trazable al flujo; revisión independiente sin defectos impeditivos.\newline Perfil Q; 2028-02-08 a 2028-02-10. & 45 & Davor Serey\\\hline
\texttt{1.\allowbreak{}5.\allowbreak{}1.\allowbreak{}2.\allowbreak{}7.\allowbreak{}1} & Coordinación y documentación del núcleo: Agenda de varadero y posiciones compatibles — lote 1/1\newline Aceptación: lote versionado y trazable al flujo; revisión independiente sin defectos impeditivos.\newline Perfil J; 2028-02-10 a 2028-02-11. & 53 & Ignacio Silva\\\hline
\texttt{1.\allowbreak{}5.\allowbreak{}1.\allowbreak{}3.\allowbreak{}4.\allowbreak{}1} & Análisis de reglas, excepciones y protocolo: Plan de izaje, validación y suspensión de faena — lote 1/1\newline Aceptación: lote versionado y trazable al flujo; revisión independiente sin defectos impeditivos.\newline Perfil A; 2027-12-20 a 2027-12-22. & 16 & Ulysses Barreto\\\hline
\texttt{1.\allowbreak{}5.\allowbreak{}1.\allowbreak{}3.\allowbreak{}4.\allowbreak{}2} & Diseño de datos y restricciones del núcleo: Plan de izaje, validación y suspensión de faena — lote 1/1\newline Aceptación: lote versionado y trazable al flujo; revisión independiente sin defectos impeditivos.\newline Perfil D; 2027-12-20 a 2027-12-22. & 16 & Nicolás Torfan\\\hline
\texttt{1.\allowbreak{}5.\allowbreak{}1.\allowbreak{}3.\allowbreak{}4.\allowbreak{}3} & Diseño de contratos internos y composición: Plan de izaje, validación y suspensión de faena — lote 1/1\newline Aceptación: lote versionado y trazable al flujo; revisión independiente sin defectos impeditivos.\newline Perfil I; 2027-12-20 a 2027-12-21. & 23 & Fernando Guerrero\\\hline
\texttt{1.\allowbreak{}5.\allowbreak{}1.\allowbreak{}3.\allowbreak{}5.\allowbreak{}1} & Construcción funcional y pruebas unitarias: Plan de izaje, validación y suspensión de faena — lote 1/1\newline Aceptación: lote versionado y trazable al flujo; revisión independiente sin defectos impeditivos.\newline Perfil V; 2028-01-26 a 2028-01-28. & 61 & Nikolai Navea\\\hline
\texttt{1.\allowbreak{}5.\allowbreak{}1.\allowbreak{}3.\allowbreak{}6.\allowbreak{}1} & Verificación del núcleo, regresión y evidencia: Plan de izaje, validación y suspensión de faena — lote 1/1\newline Aceptación: lote versionado y trazable al flujo; revisión independiente sin defectos impeditivos.\newline Perfil Q; 2028-02-18 a 2028-02-22. & 45 & Davor Serey\\\hline
\texttt{1.\allowbreak{}5.\allowbreak{}1.\allowbreak{}3.\allowbreak{}7.\allowbreak{}1} & Coordinación y documentación del núcleo: Plan de izaje, validación y suspensión de faena — lote 1/1\newline Aceptación: lote versionado y trazable al flujo; revisión independiente sin defectos impeditivos.\newline Perfil J; 2028-02-22 a 2028-02-24. & 53 & Ignacio Silva\\\hline
\texttt{1.\allowbreak{}5.\allowbreak{}1.\allowbreak{}4.\allowbreak{}4.\allowbreak{}1} & Análisis de reglas, excepciones y protocolo: Órdenes, ejecución y reprogramación de trabajos — lote 1/1\newline Aceptación: lote versionado y trazable al flujo; revisión independiente sin defectos impeditivos.\newline Perfil A; 2027-12-03 a 2027-12-07. & 15 & Ulysses Barreto\\\hline
\texttt{1.\allowbreak{}5.\allowbreak{}1.\allowbreak{}4.\allowbreak{}4.\allowbreak{}2} & Diseño de datos y restricciones del núcleo: Órdenes, ejecución y reprogramación de trabajos — lote 1/1\newline Aceptación: lote versionado y trazable al flujo; revisión independiente sin defectos impeditivos.\newline Perfil D; 2027-12-07 a 2027-12-10. & 16 & Nicolás Torfan\\\hline
\texttt{1.\allowbreak{}5.\allowbreak{}1.\allowbreak{}4.\allowbreak{}4.\allowbreak{}3} & Diseño de contratos internos y composición: Órdenes, ejecución y reprogramación de trabajos — lote 1/1\newline Aceptación: lote versionado y trazable al flujo; revisión independiente sin defectos impeditivos.\newline Perfil I; 2027-12-06 a 2027-12-09. & 23 & Fernando Guerrero\\\hline
\texttt{1.\allowbreak{}5.\allowbreak{}1.\allowbreak{}4.\allowbreak{}5.\allowbreak{}1} & Construcción funcional y pruebas unitarias: Órdenes, ejecución y reprogramación de trabajos — lote 1/1\newline Aceptación: lote versionado y trazable al flujo; revisión independiente sin defectos impeditivos.\newline Perfil V; 2028-01-24 a 2028-01-26. & 61 & Nikolai Navea\\\hline
\texttt{1.\allowbreak{}5.\allowbreak{}1.\allowbreak{}4.\allowbreak{}6.\allowbreak{}1} & Verificación del núcleo, regresión y evidencia: Órdenes, ejecución y reprogramación de trabajos — lote 1/1\newline Aceptación: lote versionado y trazable al flujo; revisión independiente sin defectos impeditivos.\newline Perfil Q; 2028-02-03 a 2028-02-07. & 45 & Davor Serey\\\hline
\texttt{1.\allowbreak{}5.\allowbreak{}1.\allowbreak{}4.\allowbreak{}7.\allowbreak{}1} & Coordinación y documentación del núcleo: Órdenes, ejecución y reprogramación de trabajos — lote 1/1\newline Aceptación: lote versionado y trazable al flujo; revisión independiente sin defectos impeditivos.\newline Perfil J; 2028-02-07 a 2028-02-09. & 53 & Ignacio Silva\\\hline
\texttt{1.\allowbreak{}5.\allowbreak{}1.\allowbreak{}5.\allowbreak{}4.\allowbreak{}1} & Análisis de reglas, excepciones y protocolo: Acreditación, inducción y acceso a intervención — lote 1/1\newline Aceptación: lote versionado y trazable al flujo; revisión independiente sin defectos impeditivos.\newline Perfil A; 2027-12-20 a 2027-12-22. & 15 & Ulysses Barreto\\\hline
\texttt{1.\allowbreak{}5.\allowbreak{}1.\allowbreak{}5.\allowbreak{}4.\allowbreak{}2} & Diseño de datos y restricciones del núcleo: Acreditación, inducción y acceso a intervención — lote 1/1\newline Aceptación: lote versionado y trazable al flujo; revisión independiente sin defectos impeditivos.\newline Perfil D; 2027-12-20 a 2027-12-22. & 16 & Nicolás Torfan\\\hline
\texttt{1.\allowbreak{}5.\allowbreak{}1.\allowbreak{}5.\allowbreak{}4.\allowbreak{}3} & Diseño de contratos internos y composición: Acreditación, inducción y acceso a intervención — lote 1/1\newline Aceptación: lote versionado y trazable al flujo; revisión independiente sin defectos impeditivos.\newline Perfil I; 2027-12-20 a 2027-12-21. & 23 & Fernando Guerrero\\\hline
\texttt{1.\allowbreak{}5.\allowbreak{}1.\allowbreak{}5.\allowbreak{}5.\allowbreak{}1} & Construcción funcional y pruebas unitarias: Acreditación, inducción y acceso a intervención — lote 1/1\newline Aceptación: lote versionado y trazable al flujo; revisión independiente sin defectos impeditivos.\newline Perfil V; 2028-01-25 a 2028-01-27. & 61 & Nikolai Navea\\\hline
\texttt{1.\allowbreak{}5.\allowbreak{}1.\allowbreak{}5.\allowbreak{}6.\allowbreak{}1} & Verificación del núcleo, regresión y evidencia: Acreditación, inducción y acceso a intervención — lote 1/1\newline Aceptación: lote versionado y trazable al flujo; revisión independiente sin defectos impeditivos.\newline Perfil Q; 2028-02-18 a 2028-02-22. & 45 & Davor Serey\\\hline
\texttt{1.\allowbreak{}5.\allowbreak{}1.\allowbreak{}5.\allowbreak{}7.\allowbreak{}1} & Coordinación y documentación del núcleo: Acreditación, inducción y acceso a intervención — lote 1/1\newline Aceptación: lote versionado y trazable al flujo; revisión independiente sin defectos impeditivos.\newline Perfil J; 2028-02-22 a 2028-02-24. & 53 & Ignacio Silva\\\hline
\texttt{1.\allowbreak{}5.\allowbreak{}3.\allowbreak{}3.\allowbreak{}4.\allowbreak{}1} & Análisis de reglas, excepciones y protocolo: Recepción de residuos y discrepancias de origen — lote 1/1\newline Aceptación: lote versionado y trazable al flujo; revisión independiente sin defectos impeditivos.\newline Perfil A; 2027-12-03 a 2027-12-07. & 15 & Ulysses Barreto\\\hline
\texttt{1.\allowbreak{}5.\allowbreak{}3.\allowbreak{}3.\allowbreak{}4.\allowbreak{}2} & Diseño de datos y restricciones del núcleo: Recepción de residuos y discrepancias de origen — lote 1/1\newline Aceptación: lote versionado y trazable al flujo; revisión independiente sin defectos impeditivos.\newline Perfil D; 2027-12-03 a 2027-12-07. & 16 & Nicolás Torfan\\\hline
\texttt{1.\allowbreak{}5.\allowbreak{}3.\allowbreak{}3.\allowbreak{}4.\allowbreak{}3} & Diseño de contratos internos y composición: Recepción de residuos y discrepancias de origen — lote 1/1\newline Aceptación: lote versionado y trazable al flujo; revisión independiente sin defectos impeditivos.\newline Perfil I; 2027-12-09 a 2027-12-13. & 23 & Fernando Guerrero\\\hline
\texttt{1.\allowbreak{}5.\allowbreak{}3.\allowbreak{}3.\allowbreak{}5.\allowbreak{}1} & Construcción funcional y pruebas unitarias: Recepción de residuos y discrepancias de origen — lote 1/1\newline Aceptación: lote versionado y trazable al flujo; revisión independiente sin defectos impeditivos.\newline Perfil V; 2028-01-24 a 2028-01-26. & 61 & Nikolai Navea\\\hline
\texttt{1.\allowbreak{}5.\allowbreak{}3.\allowbreak{}3.\allowbreak{}6.\allowbreak{}1} & Verificación del núcleo, regresión y evidencia: Recepción de residuos y discrepancias de origen — lote 1/1\newline Aceptación: lote versionado y trazable al flujo; revisión independiente sin defectos impeditivos.\newline Perfil Q; 2028-02-07 a 2028-02-09. & 45 & Davor Serey\\\hline
\texttt{1.\allowbreak{}5.\allowbreak{}3.\allowbreak{}3.\allowbreak{}7.\allowbreak{}1} & Coordinación y documentación del núcleo: Recepción de residuos y discrepancias de origen — lote 1/1\newline Aceptación: lote versionado y trazable al flujo; revisión independiente sin defectos impeditivos.\newline Perfil J; 2028-02-11 a 2028-02-15. & 53 & Ignacio Silva\\\hline
\texttt{1.\allowbreak{}5.\allowbreak{}3.\allowbreak{}4.\allowbreak{}4.\allowbreak{}1} & Análisis de reglas, excepciones y protocolo: Entregas RESPEL parciales y certificados vinculados — lote 1/1\newline Aceptación: lote versionado y trazable al flujo; revisión independiente sin defectos impeditivos.\newline Perfil A; 2027-12-20 a 2027-12-21. & 51 & Ulysses Barreto\\\hline
\texttt{1.\allowbreak{}5.\allowbreak{}3.\allowbreak{}4.\allowbreak{}4.\allowbreak{}2} & Diseño de datos y restricciones del núcleo: Entregas RESPEL parciales y certificados vinculados — lote 1/1\newline Aceptación: lote versionado y trazable al flujo; revisión independiente sin defectos impeditivos.\newline Perfil D; 2027-12-20 a 2027-12-22. & 44 & Nicolás Torfan\\\hline
\texttt{1.\allowbreak{}5.\allowbreak{}3.\allowbreak{}4.\allowbreak{}4.\allowbreak{}3} & Diseño de contratos internos y composición: Entregas RESPEL parciales y certificados vinculados — lote 1/2\newline Aceptación: lote versionado y trazable al flujo; revisión independiente sin defectos impeditivos.\newline Perfil I; 2027-12-21 a 2027-12-23. & 33 & Fernando Guerrero\\\hline
\texttt{1.\allowbreak{}5.\allowbreak{}3.\allowbreak{}4.\allowbreak{}4.\allowbreak{}4} & Diseño de contratos internos y composición: Entregas RESPEL parciales y certificados vinculados — lote 2/2\newline Aceptación: lote versionado y trazable al flujo; revisión independiente sin defectos impeditivos.\newline Perfil I; 2027-12-21 a 2027-12-24. & 32 & Fernando Guerrero\\\hline
\texttt{1.\allowbreak{}5.\allowbreak{}3.\allowbreak{}4.\allowbreak{}5.\allowbreak{}1} & Construcción funcional y pruebas unitarias: Entregas RESPEL parciales y certificados vinculados — lote 1/4\newline Aceptación: lote versionado y trazable al flujo; revisión independiente sin defectos impeditivos.\newline Perfil V; 2028-01-20 a 2028-01-24. & 51 & Nikolai Navea\\\hline
\texttt{1.\allowbreak{}5.\allowbreak{}3.\allowbreak{}4.\allowbreak{}5.\allowbreak{}2} & Construcción funcional y pruebas unitarias: Entregas RESPEL parciales y certificados vinculados — lote 2/4\newline Aceptación: lote versionado y trazable al flujo; revisión independiente sin defectos impeditivos.\newline Perfil V; 2028-01-24 a 2028-01-26. & 51 & Nikolai Navea\\\hline
\texttt{1.\allowbreak{}5.\allowbreak{}3.\allowbreak{}4.\allowbreak{}5.\allowbreak{}3} & Construcción funcional y pruebas unitarias: Entregas RESPEL parciales y certificados vinculados — lote 3/4\newline Aceptación: lote versionado y trazable al flujo; revisión independiente sin defectos impeditivos.\newline Perfil V; 2028-01-26 a 2028-01-28. & 51 & Nikolai Navea\\\hline
\texttt{1.\allowbreak{}5.\allowbreak{}3.\allowbreak{}4.\allowbreak{}5.\allowbreak{}4} & Construcción funcional y pruebas unitarias: Entregas RESPEL parciales y certificados vinculados — lote 4/4\newline Aceptación: lote versionado y trazable al flujo; revisión independiente sin defectos impeditivos.\newline Perfil V; 2028-01-20 a 2028-01-24. & 50 & Nikolai Navea\\\hline
\texttt{1.\allowbreak{}5.\allowbreak{}3.\allowbreak{}4.\allowbreak{}6.\allowbreak{}1} & Verificación del núcleo, regresión y evidencia: Entregas RESPEL parciales y certificados vinculados — lote 1/2\newline Aceptación: lote versionado y trazable al flujo; revisión independiente sin defectos impeditivos.\newline Perfil Q; 2028-02-18 a 2028-02-22. & 49 & Davor Serey\\\hline
\texttt{1.\allowbreak{}5.\allowbreak{}3.\allowbreak{}4.\allowbreak{}6.\allowbreak{}2} & Verificación del núcleo, regresión y evidencia: Entregas RESPEL parciales y certificados vinculados — lote 2/2\newline Aceptación: lote versionado y trazable al flujo; revisión independiente sin defectos impeditivos.\newline Perfil Q; 2028-02-22 a 2028-02-24. & 49 & Davor Serey\\\hline
\texttt{1.\allowbreak{}5.\allowbreak{}3.\allowbreak{}4.\allowbreak{}7.\allowbreak{}1} & Coordinación y documentación del núcleo: Entregas RESPEL parciales y certificados vinculados — lote 1/2\newline Aceptación: lote versionado y trazable al flujo; revisión independiente sin defectos impeditivos.\newline Perfil J; 2028-02-24 a 2028-02-28. & 53 & Ignacio Silva\\\hline
\texttt{1.\allowbreak{}5.\allowbreak{}3.\allowbreak{}4.\allowbreak{}7.\allowbreak{}2} & Coordinación y documentación del núcleo: Entregas RESPEL parciales y certificados vinculados — lote 2/2\newline Aceptación: lote versionado y trazable al flujo; revisión independiente sin defectos impeditivos.\newline Perfil J; 2028-02-24 a 2028-02-28. & 53 & Ignacio Silva\\\hline
\texttt{1.\allowbreak{}5.\allowbreak{}3.\allowbreak{}5.\allowbreak{}4.\allowbreak{}1} & Análisis de reglas, excepciones y protocolo: Indicadores ambientales e imputación comercial de consumo — lote 1/1\newline Aceptación: lote versionado y trazable al flujo; revisión independiente sin defectos impeditivos.\newline Perfil A; 2027-12-03 a 2027-12-07. & 15 & Ulysses Barreto\\\hline
\texttt{1.\allowbreak{}5.\allowbreak{}3.\allowbreak{}5.\allowbreak{}4.\allowbreak{}2} & Diseño de datos y restricciones del núcleo: Indicadores ambientales e imputación comercial de consumo — lote 1/1\newline Aceptación: lote versionado y trazable al flujo; revisión independiente sin defectos impeditivos.\newline Perfil D; 2027-12-07 a 2027-12-10. & 16 & Nicolás Torfan\\\hline
\texttt{1.\allowbreak{}5.\allowbreak{}3.\allowbreak{}5.\allowbreak{}4.\allowbreak{}3} & Diseño de contratos internos y composición: Indicadores ambientales e imputación comercial de consumo — lote 1/1\newline Aceptación: lote versionado y trazable al flujo; revisión independiente sin defectos impeditivos.\newline Perfil I; 2027-12-13 a 2027-12-14. & 23 & Fernando Guerrero\\\hline
\texttt{1.\allowbreak{}5.\allowbreak{}3.\allowbreak{}5.\allowbreak{}5.\allowbreak{}1} & Construcción funcional y pruebas unitarias: Indicadores ambientales e imputación comercial de consumo — lote 1/1\newline Aceptación: lote versionado y trazable al flujo; revisión independiente sin defectos impeditivos.\newline Perfil V; 2028-01-20 a 2028-01-24. & 61 & Nikolai Navea\\\hline
\texttt{1.\allowbreak{}5.\allowbreak{}3.\allowbreak{}5.\allowbreak{}6.\allowbreak{}1} & Verificación del núcleo, regresión y evidencia: Indicadores ambientales e imputación comercial de consumo — lote 1/1\newline Aceptación: lote versionado y trazable al flujo; revisión independiente sin defectos impeditivos.\newline Perfil Q; 2028-02-09 a 2028-02-14. & 45 & Davor Serey\\\hline
\texttt{1.\allowbreak{}5.\allowbreak{}3.\allowbreak{}5.\allowbreak{}7.\allowbreak{}1} & Coordinación y documentación del núcleo: Indicadores ambientales e imputación comercial de consumo — lote 1/1\newline Aceptación: lote versionado y trazable al flujo; revisión independiente sin defectos impeditivos.\newline Perfil J; 2028-02-15 a 2028-02-17. & 53 & Ignacio Silva\\\hline
\texttt{1.\allowbreak{}6.\allowbreak{}1.\allowbreak{}4.\allowbreak{}4.\allowbreak{}1} & Análisis de reglas, excepciones y protocolo: Espera y gestión contractual completa — lote 1/1\newline Aceptación: lote versionado y trazable al flujo; revisión independiente sin defectos impeditivos.\newline Perfil A; 2027-12-21 a 2027-12-23. & 51 & Ulysses Barreto\\\hline
\texttt{1.\allowbreak{}6.\allowbreak{}1.\allowbreak{}4.\allowbreak{}4.\allowbreak{}2} & Diseño de datos y restricciones del núcleo: Espera y gestión contractual completa — lote 1/1\newline Aceptación: lote versionado y trazable al flujo; revisión independiente sin defectos impeditivos.\newline Perfil D; 2027-12-22 a 2027-12-24. & 44 & Nicolás Torfan\\\hline
\texttt{1.\allowbreak{}6.\allowbreak{}1.\allowbreak{}4.\allowbreak{}4.\allowbreak{}3} & Diseño de contratos internos y composición: Espera y gestión contractual completa — lote 1/2\newline Aceptación: lote versionado y trazable al flujo; revisión independiente sin defectos impeditivos.\newline Perfil I; 2027-12-23 a 2027-12-24. & 33 & Fernando Guerrero\\\hline
\texttt{1.\allowbreak{}6.\allowbreak{}1.\allowbreak{}4.\allowbreak{}4.\allowbreak{}4} & Diseño de contratos internos y composición: Espera y gestión contractual completa — lote 2/2\newline Aceptación: lote versionado y trazable al flujo; revisión independiente sin defectos impeditivos.\newline Perfil I; 2027-12-24 a 2027-12-28. & 32 & Fernando Guerrero\\\hline
\texttt{1.\allowbreak{}6.\allowbreak{}1.\allowbreak{}4.\allowbreak{}5.\allowbreak{}1} & Construcción funcional y pruebas unitarias: Espera y gestión contractual completa — lote 1/4\newline Aceptación: lote versionado y trazable al flujo; revisión independiente sin defectos impeditivos.\newline Perfil V; 2028-01-24 a 2028-01-26. & 51 & Nikolai Navea\\\hline
\texttt{1.\allowbreak{}6.\allowbreak{}1.\allowbreak{}4.\allowbreak{}5.\allowbreak{}2} & Construcción funcional y pruebas unitarias: Espera y gestión contractual completa — lote 2/4\newline Aceptación: lote versionado y trazable al flujo; revisión independiente sin defectos impeditivos.\newline Perfil V; 2028-01-26 a 2028-01-28. & 51 & Nikolai Navea\\\hline
\texttt{1.\allowbreak{}6.\allowbreak{}1.\allowbreak{}4.\allowbreak{}5.\allowbreak{}3} & Construcción funcional y pruebas unitarias: Espera y gestión contractual completa — lote 3/4\newline Aceptación: lote versionado y trazable al flujo; revisión independiente sin defectos impeditivos.\newline Perfil V; 2028-01-20 a 2028-01-24. & 51 & Nikolai Navea\\\hline
\texttt{1.\allowbreak{}6.\allowbreak{}1.\allowbreak{}4.\allowbreak{}5.\allowbreak{}4} & Construcción funcional y pruebas unitarias: Espera y gestión contractual completa — lote 4/4\newline Aceptación: lote versionado y trazable al flujo; revisión independiente sin defectos impeditivos.\newline Perfil V; 2028-01-24 a 2028-01-26. & 50 & Nikolai Navea\\\hline
\texttt{1.\allowbreak{}6.\allowbreak{}1.\allowbreak{}4.\allowbreak{}6.\allowbreak{}1} & Verificación del núcleo, regresión y evidencia: Espera y gestión contractual completa — lote 1/2\newline Aceptación: lote versionado y trazable al flujo; revisión independiente sin defectos impeditivos.\newline Perfil Q; 2028-02-22 a 2028-02-24. & 49 & Davor Serey\\\hline
\texttt{1.\allowbreak{}6.\allowbreak{}1.\allowbreak{}4.\allowbreak{}6.\allowbreak{}2} & Verificación del núcleo, regresión y evidencia: Espera y gestión contractual completa — lote 2/2\newline Aceptación: lote versionado y trazable al flujo; revisión independiente sin defectos impeditivos.\newline Perfil Q; 2028-02-24 a 2028-02-28. & 48 & Davor Serey\\\hline
\texttt{1.\allowbreak{}6.\allowbreak{}1.\allowbreak{}4.\allowbreak{}7.\allowbreak{}1} & Coordinación y documentación del núcleo: Espera y gestión contractual completa — lote 1/2\newline Aceptación: lote versionado y trazable al flujo; revisión independiente sin defectos impeditivos.\newline Perfil J; 2028-02-28 a 2028-03-01. & 53 & Ignacio Silva\\\hline
\texttt{1.\allowbreak{}6.\allowbreak{}1.\allowbreak{}4.\allowbreak{}7.\allowbreak{}2} & Coordinación y documentación del núcleo: Espera y gestión contractual completa — lote 2/2\newline Aceptación: lote versionado y trazable al flujo; revisión independiente sin defectos impeditivos.\newline Perfil J; 2028-02-29 a 2028-03-02. & 53 & Ignacio Silva\\\hline
\texttt{1.\allowbreak{}6.\allowbreak{}2.\allowbreak{}1.\allowbreak{}4.\allowbreak{}1} & Análisis de reglas, excepciones y protocolo: Hechos facturables, tarifas y lotes de envío — lote 1/1\newline Aceptación: lote versionado y trazable al flujo; revisión independiente sin defectos impeditivos.\newline Perfil A; 2027-12-03 a 2027-12-07. & 15 & Ulysses Barreto\\\hline
\texttt{1.\allowbreak{}6.\allowbreak{}2.\allowbreak{}1.\allowbreak{}4.\allowbreak{}2} & Diseño de datos y restricciones del núcleo: Hechos facturables, tarifas y lotes de envío — lote 1/1\newline Aceptación: lote versionado y trazable al flujo; revisión independiente sin defectos impeditivos.\newline Perfil D; 2027-12-03 a 2027-12-07. & 15 & Nicolás Torfan\\\hline
\texttt{1.\allowbreak{}6.\allowbreak{}2.\allowbreak{}1.\allowbreak{}4.\allowbreak{}3} & Diseño de contratos internos y composición: Hechos facturables, tarifas y lotes de envío — lote 1/1\newline Aceptación: lote versionado y trazable al flujo; revisión independiente sin defectos impeditivos.\newline Perfil I; 2027-12-03 a 2027-12-06. & 23 & Fernando Guerrero\\\hline
\texttt{1.\allowbreak{}6.\allowbreak{}2.\allowbreak{}1.\allowbreak{}5.\allowbreak{}1} & Construcción funcional y pruebas unitarias: Hechos facturables, tarifas y lotes de envío — lote 1/1\newline Aceptación: lote versionado y trazable al flujo; revisión independiente sin defectos impeditivos.\newline Perfil V; 2028-01-24 a 2028-01-26. & 61 & Nikolai Navea\\\hline
\texttt{1.\allowbreak{}6.\allowbreak{}2.\allowbreak{}1.\allowbreak{}6.\allowbreak{}1} & Verificación del núcleo, regresión y evidencia: Hechos facturables, tarifas y lotes de envío — lote 1/1\newline Aceptación: lote versionado y trazable al flujo; revisión independiente sin defectos impeditivos.\newline Perfil Q; 2028-02-03 a 2028-02-07. & 45 & Davor Serey\\\hline
\texttt{1.\allowbreak{}6.\allowbreak{}2.\allowbreak{}1.\allowbreak{}7.\allowbreak{}1} & Coordinación y documentación del núcleo: Hechos facturables, tarifas y lotes de envío — lote 1/1\newline Aceptación: lote versionado y trazable al flujo; revisión independiente sin defectos impeditivos.\newline Perfil J; 2028-02-07 a 2028-02-09. & 53 & Ignacio Silva\\\hline
\texttt{1.\allowbreak{}6.\allowbreak{}2.\allowbreak{}2.\allowbreak{}4.\allowbreak{}1} & Análisis de reglas, excepciones y protocolo: Respuesta contable, pagos y saldos informados — lote 1/1\newline Aceptación: lote versionado y trazable al flujo; revisión independiente sin defectos impeditivos.\newline Perfil A; 2027-12-20 a 2027-12-22. & 15 & Ulysses Barreto\\\hline
\texttt{1.\allowbreak{}6.\allowbreak{}2.\allowbreak{}2.\allowbreak{}4.\allowbreak{}2} & Diseño de datos y restricciones del núcleo: Respuesta contable, pagos y saldos informados — lote 1/1\newline Aceptación: lote versionado y trazable al flujo; revisión independiente sin defectos impeditivos.\newline Perfil D; 2027-12-22 a 2027-12-24. & 15 & Nicolás Torfan\\\hline
\texttt{1.\allowbreak{}6.\allowbreak{}2.\allowbreak{}2.\allowbreak{}4.\allowbreak{}3} & Diseño de contratos internos y composición: Respuesta contable, pagos y saldos informados — lote 1/1\newline Aceptación: lote versionado y trazable al flujo; revisión independiente sin defectos impeditivos.\newline Perfil I; 2027-12-24 a 2027-12-28. & 23 & Fernando Guerrero\\\hline
\texttt{1.\allowbreak{}6.\allowbreak{}2.\allowbreak{}2.\allowbreak{}5.\allowbreak{}1} & Construcción funcional y pruebas unitarias: Respuesta contable, pagos y saldos informados — lote 1/1\newline Aceptación: lote versionado y trazable al flujo; revisión independiente sin defectos impeditivos.\newline Perfil V; 2028-01-26 a 2028-01-28. & 61 & Nikolai Navea\\\hline
\texttt{1.\allowbreak{}6.\allowbreak{}2.\allowbreak{}2.\allowbreak{}6.\allowbreak{}1} & Verificación del núcleo, regresión y evidencia: Respuesta contable, pagos y saldos informados — lote 1/1\newline Aceptación: lote versionado y trazable al flujo; revisión independiente sin defectos impeditivos.\newline Perfil Q; 2028-02-24 a 2028-02-28. & 45 & Davor Serey\\\hline
\texttt{1.\allowbreak{}6.\allowbreak{}2.\allowbreak{}2.\allowbreak{}7.\allowbreak{}1} & Coordinación y documentación del núcleo: Respuesta contable, pagos y saldos informados — lote 1/1\newline Aceptación: lote versionado y trazable al flujo; revisión independiente sin defectos impeditivos.\newline Perfil J; 2028-03-01 a 2028-03-03. & 53 & Ignacio Silva\\\hline
\texttt{1.\allowbreak{}6.\allowbreak{}2.\allowbreak{}3.\allowbreak{}4.\allowbreak{}1} & Análisis de reglas, excepciones y protocolo: Cobranza y medidas con decisión competente — lote 1/1\newline Aceptación: lote versionado y trazable al flujo; revisión independiente sin defectos impeditivos.\newline Perfil A; 2027-12-03 a 2027-12-07. & 15 & Ulysses Barreto\\\hline
\texttt{1.\allowbreak{}6.\allowbreak{}2.\allowbreak{}3.\allowbreak{}4.\allowbreak{}2} & Diseño de datos y restricciones del núcleo: Cobranza y medidas con decisión competente — lote 1/1\newline Aceptación: lote versionado y trazable al flujo; revisión independiente sin defectos impeditivos.\newline Perfil D; 2027-12-07 a 2027-12-10. & 15 & Nicolás Torfan\\\hline
\texttt{1.\allowbreak{}6.\allowbreak{}2.\allowbreak{}3.\allowbreak{}4.\allowbreak{}3} & Diseño de contratos internos y composición: Cobranza y medidas con decisión competente — lote 1/1\newline Aceptación: lote versionado y trazable al flujo; revisión independiente sin defectos impeditivos.\newline Perfil I; 2027-12-06 a 2027-12-09. & 23 & Fernando Guerrero\\\hline
\texttt{1.\allowbreak{}6.\allowbreak{}2.\allowbreak{}3.\allowbreak{}5.\allowbreak{}1} & Construcción funcional y pruebas unitarias: Cobranza y medidas con decisión competente — lote 1/1\newline Aceptación: lote versionado y trazable al flujo; revisión independiente sin defectos impeditivos.\newline Perfil V; 2028-01-24 a 2028-01-26. & 61 & Nikolai Navea\\\hline
\texttt{1.\allowbreak{}6.\allowbreak{}2.\allowbreak{}3.\allowbreak{}6.\allowbreak{}1} & Verificación del núcleo, regresión y evidencia: Cobranza y medidas con decisión competente — lote 1/1\newline Aceptación: lote versionado y trazable al flujo; revisión independiente sin defectos impeditivos.\newline Perfil Q; 2028-02-08 a 2028-02-10. & 45 & Davor Serey\\\hline
\texttt{1.\allowbreak{}6.\allowbreak{}2.\allowbreak{}3.\allowbreak{}7.\allowbreak{}1} & Coordinación y documentación del núcleo: Cobranza y medidas con decisión competente — lote 1/1\newline Aceptación: lote versionado y trazable al flujo; revisión independiente sin defectos impeditivos.\newline Perfil J; 2028-02-15 a 2028-02-17. & 53 & Ignacio Silva\\\hline
\texttt{1.\allowbreak{}6.\allowbreak{}2.\allowbreak{}4.\allowbreak{}4.\allowbreak{}1} & Análisis de reglas, excepciones y protocolo: Expedientes de custodia y exportación íntegra — lote 1/1\newline Aceptación: lote versionado y trazable al flujo; revisión independiente sin defectos impeditivos.\newline Perfil A; 2027-12-20 a 2027-12-22. & 15 & Ulysses Barreto\\\hline
\texttt{1.\allowbreak{}6.\allowbreak{}2.\allowbreak{}4.\allowbreak{}4.\allowbreak{}2} & Diseño de datos y restricciones del núcleo: Expedientes de custodia y exportación íntegra — lote 1/1\newline Aceptación: lote versionado y trazable al flujo; revisión independiente sin defectos impeditivos.\newline Perfil D; 2027-12-22 a 2027-12-24. & 15 & Nicolás Torfan\\\hline
\texttt{1.\allowbreak{}6.\allowbreak{}2.\allowbreak{}4.\allowbreak{}4.\allowbreak{}3} & Diseño de contratos internos y composición: Expedientes de custodia y exportación íntegra — lote 1/1\newline Aceptación: lote versionado y trazable al flujo; revisión independiente sin defectos impeditivos.\newline Perfil I; 2027-12-20 a 2027-12-21. & 23 & Fernando Guerrero\\\hline
\texttt{1.\allowbreak{}6.\allowbreak{}2.\allowbreak{}4.\allowbreak{}5.\allowbreak{}1} & Construcción funcional y pruebas unitarias: Expedientes de custodia y exportación íntegra — lote 1/1\newline Aceptación: lote versionado y trazable al flujo; revisión independiente sin defectos impeditivos.\newline Perfil V; 2028-01-28 a 2028-02-01. & 61 & Nikolai Navea\\\hline
\texttt{1.\allowbreak{}6.\allowbreak{}2.\allowbreak{}4.\allowbreak{}6.\allowbreak{}1} & Verificación del núcleo, regresión y evidencia: Expedientes de custodia y exportación íntegra — lote 1/1\newline Aceptación: lote versionado y trazable al flujo; revisión independiente sin defectos impeditivos.\newline Perfil Q; 2028-02-24 a 2028-02-28. & 45 & Davor Serey\\\hline
\texttt{1.\allowbreak{}6.\allowbreak{}2.\allowbreak{}4.\allowbreak{}7.\allowbreak{}1} & Coordinación y documentación del núcleo: Expedientes de custodia y exportación íntegra — lote 1/1\newline Aceptación: lote versionado y trazable al flujo; revisión independiente sin defectos impeditivos.\newline Perfil J; 2028-03-02 a 2028-03-06. & 53 & Ignacio Silva\\\hline
\end{longtable}\endgroup
Los recursos son DEV/QA, contratos internos, modelo de datos, escenarios del flujo y ejecutores propios asignados. Los hitos son H4/H9 y H5/H10 según etapa. Las cantidades derivan de UCP y de los créditos por actividad; las exclusiones y condiciones de provisión del modelo son sus supuestos. Los costos se imputarán por perfil y paquete exclusivamente en la Oferta Económica.


\section{Diccionario y ensayos de migración}
La cuenta 1.8 comprende diagnóstico, transformación, saneamiento, carga y archivo histórico. Los ensayos integrales se registran en 1.10.1 para comprobar la migración sin duplicar su esfuerzo de construcción.
\subsection{Componentes de migración y frontera con pruebas e implantación}
La rama 1.8 entregará datos interpretables, conciliados y consultables, conforme al alcance de SD5. El diagnóstico 1.8.1 será la entrada para especificar los componentes de la Tabla~\ref{tab:sd7-componentes-migracion}. Estos códigos identifican resultados superiores; T-14 desarrolla sus paquetes de ingeniería por conjunto y sus lotes de revisión documental. Las 120 horas de diagnóstico no incluyen la ejecución de los componentes posteriores.

\begin{table}[htbp]
\centering
\begin{tabularx}{\linewidth}{L{13mm}L{35mm}Y}
\hline
\EncabezadoTabla{Código} & \EncabezadoTabla{Resultado} & \EncabezadoTabla{Condición de recepción}\\\hline
1.8.2 & Correspondencias y transformaciones especificadas & Cada conjunto tiene origen, destino, claves, reglas, tratamiento de ausencias y validador; no se convierten antecedentes antiguos en autorizaciones vigentes.\\\hline
1.8.3 & Datos saneados con decisiones trazables & Cada defecto conserva origen, decisión, responsable y resultado; los cambios se aplican a copias controladas sin fabricar pagos ni evidencia histórica.\\\hline
1.8.4 & Herramientas de carga y conciliación reproducibles & Cargas por lote con manifiestos, correspondencias, recuentos y controles monetarios; la reejecución no duplica efectos y conserva evidencia de cada ejecución.\\\hline
1.8.5 & Archivo histórico independiente consultable & Índice, relaciones, permisos, integridad, retención y exportación verificables; la consulta no depende del software de 2014 como única vía.\\\hline
1.8.6 & Herramientas de delta y retorno disponibles & Detección de altas, cambios y bajas; preservación del diario y tratamiento de respuestas externas. Incluye herramientas para retorno previo y posterior a la nueva escritura.\\\hline
\end{tabularx}
\caption{Componentes posteriores al diagnóstico de migración}
\label{tab:sd7-componentes-migracion}
\FuenteTabla{SD5, secciones 5.3.1--5.3.3; SD2, RNF-MIG-01 y DEC-02.}
\end{table}

El Líder de Datos responderá por los componentes 1.8.1--1.8.6. Los dos ensayos completos se imputarán a \texttt{1.10.1}, como componente de verificación de migración a cargo del Líder de Calidad y Pruebas, con ejecución de Datos y participación de Seguridad y los validadores funcionales. La producción de herramientas se contabilizará en 1.8 y su empleo en ensayos integrales en 1.10.1; una misma actividad no se cargará a ambas ramas.

El componente \texttt{1.11.1} reunirá la implantación de los grupos de datos, el corte, la convivencia y la estabilización bajo el Líder de Implantación y Gestión del Cambio. Los grupos de SD5 ---maestros e inventario; antecedentes operacionales por zona o proceso; y contratos, saldos y sustitución administrativa--- ordenan esas activaciones. La consulta de los contratos necesarios para E1 no transfiere anticipadamente la autoridad de escritura administrativa del legado. Cada activación conservará su recepción, su alcance y su responsable de negocio.

Antes de programar las cargas se medirán las fuentes y se revisará la hipótesis de volumen de EST-10 y el vacío VAC-07 del SD2. El saneamiento documental, la digitalización que resulte necesaria y la comprobación semántica completa de los 14 expedientes y de las autorizaciones escolares vigentes tendrán paquetes y esfuerzos propios. El perfilado automatizado del diagnóstico no absorbe esas revisiones.

\subsection{Trazabilidad de requisitos y condiciones de recepción de migración}
La Tabla~\ref{tab:sd7-trazabilidad-migracion} vincula la rama de datos con las obligaciones de migración y los controles que habilitan el cambio de servicio. Un requisito puede necesitar varios componentes; cada uno aportará una evidencia distinta. La recepción se decidirá sobre el conjunto de evidencias y no por el solo término de un paquete de diagnóstico.

\begingroup
\setlength{\tabcolsep}{3pt}
\begin{longtable}{L{\dimexpr.20\linewidth-2\tabcolsep\relax}L{\dimexpr.23\linewidth-2\tabcolsep\relax}L{\dimexpr.22\linewidth-2\tabcolsep\relax}L{\dimexpr.35\linewidth-2\tabcolsep\relax}}
\caption{Trazabilidad de obligaciones de migración}\label{tab:sd7-trazabilidad-migracion}\\
\hline
\EncabezadoTabla{Origen} & \EncabezadoTabla{Resultado exigido} & \EncabezadoTabla{Componentes} & \EncabezadoTabla{Evidencia de recepción}\\\hline
\endfirsthead
\hline
\EncabezadoTabla{Origen} & \EncabezadoTabla{Resultado exigido} & \EncabezadoTabla{Componentes} & \EncabezadoTabla{Evidencia de recepción}\\\hline
\endhead
C07 RT-05.15; SD2 RNF-MIG-01 & Seis conjuntos obligatorios y sus períodos & 1.8.1--1.8.5; 1.10.1 & Catálogo, matriz de cobertura, manifiestos y conciliación por conjunto.\\\hline
BT RT-05.11; SD2 DEC-02 & Plan de migración y retorno & 1.8.1--1.8.6; 1.11.1 & Plan versionado con origen, volumen medido, reglas, calidad, ejecución y reversión.\\\hline
BT RT-05.12 & Perfilado y saneamiento previo & 1.8.1.3; 1.8.2--1.8.3 & Informe de defectos y registro de decisiones y correcciones por lote.\\\hline
BT RT-05.13 & Dos ensayos completos en PREPROD & 1.8.4--1.8.6; 1.10.1 & Dos expedientes de ejecución con todos los conjuntos, duración total, conciliación y retorno.\\\hline
BT RT-05.14; SD5 5.3.2 & Conciliación cuantitativa & 1.8.4; 1.10.1; 1.11.1 & Recuentos, claves, sumas, revisión dirigida y diferencias explicadas; cero diferencias críticas de carga.\\\hline
BT RT-05.15; SD2 DEC-02 & Consulta histórica protegida & 1.8.5; 1.10.1 & Consulta y exportación con índice, relaciones, permisos e integridad independientes del legado.\\\hline
SD3 CA-22; SD5 5.3 & 14 expedientes completos & 1.8.2--1.8.5; 1.10.1 & Índice de cada pieza, revisión completa, legibilidad, folio, integridad y exportación PDF/A.\\\hline
SD5 5.3.2 & Autorizaciones escolares vigentes & 1.8.2--1.8.4; 1.10.1 & Revisión completa de vínculos, vigencia y destino de las autorizaciones que se utilizarán al corte.\\\hline
BT RT-20.02--04; SD5 5.3.3 & Cambio gradual, reversible y conciliado & 1.8.6; 1.10.1; 1.11.1 & Ensayos, decisión de avance, actas por grupo y conciliación diaria de marcha blanca.\\\hline
\end{longtable}
\FuenteTabla{Bases Técnicas Transversales, RT-05.11--RT-05.15 y RT-20.02--RT-20.04; Caso 07, RT-05.15; SD2, SD3 y SD5.}
\endgroup

Los ensayos comprobarán los umbrales del SD5: cero omisiones inexplicadas del manifiesto, cero diferencias monetarias de carga, cero referencias críticas inválidas y cero efectos duplicados. Para atributos no críticos se conservará la revisión de todos los registros anómalos y la muestra reproducible del mayor entre 30 registros y el 5\,\% de cada fuente y categoría, limitada a su población. Los expedientes y las autorizaciones vigentes tendrán revisión completa, sin sustituirla por esa muestra \parencite[cap.~5, RT-05.12--RT-05.14]{bases-transversales}.

La implantación conservará la ventana administrativa propuesta de 22:00 a 02:00, hora de America/Santiago. La decisión de continuar se registrará antes de las 00:30 y mantendrá una reserva exclusiva de 90 minutos para retorno previo a nuevas escrituras. El escenario posterior a nuevas operaciones conserva la estimación inicial de 120 minutos de SD5 y su obligación de ensayo; no amplía esa reserva ni autoriza a descartar operaciones. Los resultados de PREPROD determinarán si es necesario ajustar precarga y herramientas antes de fijar una fecha habilitada de corte.

Esta trazabilidad corresponde a migración. Las demás ramas se vincularán con sus propios requisitos, pruebas y criterios de recepción; la matriz general del contrato no se dará por completa con estas filas.

\subsection{Síntesis de estimación de ingeniería de migración}
El T-14 descompone los componentes anteriores en 28 paquetes de ingeniería de 24--40 horas-persona, incluidos los cuatro paquetes de diagnóstico. Su subtotal propuesto es 936 horas: Análisis Funcional 96, Datos 552, Seguridad 32, Calidad 112 y Desarrollo 144. La construcción de las herramientas reutiliza los componentes base de SD4; las cifras son estimaciones ascendentes bajo supuestos de formatos y clases de defecto, no productividad medida.

La revisión completa de los expedientes y autorizaciones escolares se estima por lotes separados, según el inventario real, y no se da por cubierta por las 936 horas. Tampoco se incluyen los dos ensayos completos, la ejecución de los cortes, la marcha blanca ni los 36 meses de operación. El archivo independiente \path{SYNAPTIX-Formulario-T-15.pdf} registra las cargas y su distribución mensual inicial para comprobarlas junto con las demás ramas.

\section{Diccionario de los paquetes de diagnóstico}
Las fichas siguientes fijan el contenido, la frontera y la recepción de cada paquete. El Líder de Datos responde por su resultado; Análisis Funcional, Seguridad y Calidad aportan las tareas indicadas, con los titulares y funciones del SD1. El Líder de Datos ejerce la responsabilidad de migración prevista en SD5; su suplencia técnica se asignará entre los perfiles habilitados del equipo y se contabilizará en T-15. La contraparte funcional del cliente valida el significado de los antecedentes y las decisiones de negocio, sin asumir la ejecución técnica de la extracción. Administración designará su suplente funcional para reproducir el procedimiento durante el segundo ensayo completo, según SD5, sección 5.3.2.

\paragraph*{1.8.1.1 --- Catálogo de fuentes y accesos}
El entregable será un catálogo versionado de sistemas, archivos y expedientes con su custodio, ubicación, período, formato, cantidad de registros, tamaño, sensibilidad y mecanismo de acceso. Las cantidades que todavía no puedan comprobarse se identificarán como pendientes de extracción diagnóstica. El trabajo comprende preparar entrevistas, registrar fuentes y contrastar su cobertura con los seis conjuntos obligatorios; excluye digitalizar todos los documentos y sanear registros.

\textbf{Aceptación y responsable.} El Líder de Datos presentará una relación para cada conjunto, con sus fuentes identificadas o su ausencia documentada. Cada fuente tendrá custodio, acceso solicitado y condición de disponibilidad; las autorizaciones se conservarán como evidencia y las pendientes impedirán habilitar la extracción correspondiente. Administración, Escuela y Operaciones validarán los antecedentes de su ámbito.

\textbf{Esfuerzo y recursos.} Análisis Funcional: 16 horas para preparación y entrevistas; Datos: 8 horas para consolidación y contraste. Se requieren acceso autorizado al sistema de 2014, padrones, carpetas y registros escolares, y un repositorio protegido. 

\textbf{Supuestos, hitos y referencias.} Se estiman cuatro sesiones coordinadas: las dos iniciales de 90 minutos con quien conoce el legado, ya previstas en SD5, sección 5.3.2, y dos adicionales de hasta dos horas para contrastar las otras fuentes con sus custodios. Las 16 horas de Análisis Funcional incluyen hasta siete horas de sesiones y nueve de preparación y registro; no se supone dedicación exclusiva del cliente. El cierre mensual asistido previo al cambio administrativo se programará en los paquetes de conciliación, fuera de este diagnóstico. El catálogo identifica lotes y custodios, no revisa página por página los expedientes. Su recepción habilitará el diseño de extracción; depende de SUP-01 y SUP-02 del SD2 y de SD5, sección 5.3.1.

\paragraph*{1.8.1.2 --- Procedimiento de extracción diagnóstica}
El entregable será un procedimiento versionado con las herramientas necesarias para obtener una copia de diagnóstico y su manifiesto. Datos asignará 8 horas a examinar el mecanismo disponible, 16 a preparar la extracción y 8 a repetirla y documentarla; Seguridad asignará 8 horas a permisos, protección de copias y revisión del registro de acceso. El trabajo no modificará la fuente ni comprenderá cargas productivas.

\textbf{Aceptación y responsable.} El Líder de Datos demostrará dos ejecuciones repetibles sobre un origen o instantánea estable y autorizada. Los manifiestos permitirán comparar archivos o tablas, cantidades, períodos y sumas de comprobación; el procedimiento registrará versión de herramienta, parámetros, errores y ubicación protegida. Seguridad verificará que los permisos y las copias correspondan al alcance autorizado. Estas ejecuciones diagnósticas no sustituyen los dos ensayos completos de migración.

\textbf{Recursos.} Se requieren una copia protegida del origen, almacenamiento cifrado, herramientas de extracción y un técnico suplente que pueda reproducir el procedimiento. La repetición estará incluida en las 32 horas de Datos. Las licencias o insumos adicionales se identificarán por separado cuando el mecanismo elegido los requiera.

\textbf{Supuestos, hitos y referencias.} Las 40 horas suponen acceso a un mecanismo utilizable sin ingeniería inversa del ejecutable ni recuperación de soportes dañados. Esas condiciones se comprobarán al recibir el catálogo; su incumplimiento exigirá estimar los paquetes adicionales y revisar la programación. La disponibilidad del único operador conocedor del sistema se coordinará con un suplente de Synaptix. El procedimiento habilitará el perfilado y atenderá SD5, sección 5.3, y SUP-02 del SD2.

\paragraph*{1.8.1.3 --- Informe reproducible de calidad del origen}
El entregable será un informe acompañado de consultas o reglas versionadas para identificar campos obligatorios ausentes, duplicados, relaciones rotas, formatos incompatibles y diferencias entre registros y manifiestos. Datos dedicará 8 horas a reglas, 12 a ejecución y análisis y 4 a documentación; Calidad dedicará 8 horas a repetir comprobaciones y revisar la evidencia. El trabajo cuantificará incidencias sin corregir los registros de origen.

\textbf{Aceptación y responsable.} El Líder de Datos entregará resultados por conjunto con denominadores, cantidad de incidencias, referencia a registros afectados y limitaciones identificadas. Calidad reproducirá las comprobaciones sobre la misma copia y obtendrá los mismos resultados. Las diferencias monetarias se analizarán por moneda y período, sin compensar errores positivos y negativos. Los adjuntos se comprobarán por existencia y vínculo; su lectura semántica exhaustiva no pertenece a este paquete.

\textbf{Recursos.} Se requieren el catálogo, la copia de diagnóstico, sus manifiestos y un ambiente aislado de análisis. Las licencias o insumos adicionales se identificarán por separado cuando el mecanismo elegido los requiera.

\textbf{Supuestos, hitos y referencias.} Se supone que las fuentes tabulares pueden procesarse con las herramientas habilitadas en 1.8.1.2. La transcripción manual, OCR y reconstrucción documental se estimarán fuera de estas 32 horas. El informe habilitará las decisiones de saneamiento y conciliación; su referencia es SD5, secciones 5.2.4 y 5.3.2.

\paragraph*{1.8.1.4 --- Matriz de cobertura y decisiones de migración}
El entregable será una matriz que relacione cada conjunto obligatorio con sus fuentes, cobertura temporal, destino operativo o histórico, incidencias, reglas propuestas y validadores. Análisis Funcional dedicará 8 horas a decisiones y validación funcional; Datos, 8 a la matriz y sus vínculos; y Calidad, 8 a comprobar cobertura y frontera con los siguientes componentes.

\textbf{Aceptación y responsable.} El Líder de Datos entregará una fila por cada conjunto obligatorio y sus fuentes vinculadas, incluidos los 14 expedientes de abandono. Cada incidencia tendrá tratamiento propuesto y responsable de decisión. Los valores ausentes permanecerán identificados como tales; ninguna regla inventará pagos, autorizaciones ni movimientos históricos. Las decisiones aprobadas y las pendientes se distinguirán con su evidencia. La matriz se recibirá como diagnóstico completo cuando cubra todos los conjuntos; no autorizará una carga definitiva mientras existan decisiones bloqueantes sin resolver.

\textbf{Recursos.} Se requieren los tres entregables anteriores y la participación programada de los validadores del cliente. 

\textbf{Supuestos, hitos y referencias.} Las 24 horas comprenden una ronda consolidada de decisiones; los tiempos de espera del cliente se incorporarán al calendario. La matriz habilitará la especificación de saneamiento, carga y consulta histórica, con trazabilidad a SD3, sección 3.2, SD5, sección 5.3 y C07 RT-05.15.

Los hitos de recepción de los cuatro paquetes se enlazarán en el cronograma, manteniendo separados la entrega técnica de Synaptix y el tiempo de validación del cliente.

\section{Diccionario de ingeniería y revisión documental de migración}
La estimación distingue las herramientas reproducibles de la revisión humana de antecedentes. La Tabla~\ref{tab:mig-cargas} desarrolla 28 paquetes de ingeniería, incluidos los cuatro de diagnóstico ya descritos. La numeración final $c$ de los paquetes 1.8.2.$c$, 1.8.3.$c$ y 1.8.4.$c$ identifica, respectivamente, socios y armadores; contratos vigentes y documentos; seis años de pagos y saldos; embarcaciones e invernadas; los 14 expedientes de abandono; y tres años de matrícula escolar. Cada paquete produce un resultado propio de su conjunto.

\begingroup
\setlength{\tabcolsep}{2pt}
\begin{longtable}{L{\dimexpr.15\linewidth-2\tabcolsep\relax}L{\dimexpr.43\linewidth-2\tabcolsep\relax}R{\dimexpr.06\linewidth-2\tabcolsep\relax}R{\dimexpr.06\linewidth-2\tabcolsep\relax}R{\dimexpr.06\linewidth-2\tabcolsep\relax}R{\dimexpr.06\linewidth-2\tabcolsep\relax}R{\dimexpr.06\linewidth-2\tabcolsep\relax}R{\dimexpr.07\linewidth-2\tabcolsep\relax}R{\dimexpr.05\linewidth-2\tabcolsep\relax}}
\caption{Cargas iniciales de ingeniería de migración por paquete}\label{tab:mig-cargas}\\
\hline
 \EncabezadoTabla{Código} & \EncabezadoTabla{Resultado} & \EncabezadoTabla{A} & \EncabezadoTabla{D} & \EncabezadoTabla{S} & \EncabezadoTabla{Q} & \EncabezadoTabla{V} & \EncabezadoTabla{HH} & \EncabezadoTabla{Mes}\\\hline
\endfirsthead
\hline
 \EncabezadoTabla{Código} & \EncabezadoTabla{Resultado} & \EncabezadoTabla{A} & \EncabezadoTabla{D} & \EncabezadoTabla{S} & \EncabezadoTabla{Q} & \EncabezadoTabla{V} & \EncabezadoTabla{HH} & \EncabezadoTabla{Mes}\\\hline
\endhead
1.8.1.1 & Catálogo de fuentes y accesos & 16 & 8 & 0 & 0 & 0 & 24 & 1 \\\hline
1.8.1.2 & Extracción diagnóstica reproducible & 0 & 32 & 8 & 0 & 0 & 40 & 2\\\hline
1.8.1.3 & Informe de calidad del origen & 0 & 24 & 0 & 8 & 0 & 32 & 2 \\\hline
1.8.1.4 & Matriz de cobertura y decisiones & 8 & 8 & 0 & 8 & 0 & 24 & 3\\\hline
1.8.2.1 & Reglas: Socios y armadores & 8 & 16 & 0 & 0 & 0 & 24 & 3 \\\hline
1.8.2.2 & Reglas: Contratos vigentes y documentos & 8 & 16 & 0 & 0 & 0 & 24 & 3\\\hline
1.8.2.3 & Reglas: Seis años de pagos y saldos & 8 & 16 & 0 & 0 & 0 & 24 & 3 \\\hline
1.8.2.4 & Reglas: Embarcaciones e invernadas & 8 & 16 & 0 & 0 & 0 & 24 & 3\\\hline
1.8.2.5 & Reglas: 14 expedientes de abandono & 8 & 16 & 0 & 0 & 0 & 24 & 3 \\\hline
1.8.2.6 & Reglas: Tres años de matrícula escolar & 8 & 16 & 0 & 0 & 0 & 24 & 3\\\hline
1.8.3.1 & Saneamiento automatizado: Socios y armadores & 4 & 24 & 0 & 4 & 0 & 32 & 7 \\\hline
1.8.3.2 & Saneamiento automatizado: Contratos vigentes y documentos & 4 & 24 & 0 & 4 & 0 & 32 & 7\\\hline
1.8.3.3 & Saneamiento automatizado: Seis años de pagos y saldos & 4 & 24 & 0 & 4 & 0 & 32 & 7 \\\hline
1.8.3.4 & Saneamiento automatizado: Embarcaciones e invernadas & 4 & 24 & 0 & 4 & 0 & 32 & 7\\\hline
1.8.3.5 & Saneamiento automatizado: 14 expedientes de abandono & 4 & 24 & 0 & 4 & 0 & 32 & 7 \\\hline
1.8.3.6 & Saneamiento automatizado: Tres años de matrícula escolar & 4 & 24 & 0 & 4 & 0 & 32 & 7\\\hline
1.8.4.1 & Carga y conciliación: Socios y armadores & 0 & 24 & 0 & 8 & 8 & 40 & 8 \\\hline
1.8.4.2 & Carga y conciliación: Contratos vigentes y documentos & 0 & 24 & 0 & 8 & 8 & 40 & 8\\\hline
1.8.4.3 & Carga y conciliación: Seis años de pagos y saldos & 0 & 24 & 0 & 8 & 8 & 40 & 8 \\\hline
1.8.4.4 & Carga y conciliación: Embarcaciones e invernadas & 0 & 24 & 0 & 8 & 8 & 40 & 8\\\hline
1.8.4.5 & Carga y conciliación: 14 expedientes de abandono & 0 & 24 & 0 & 8 & 8 & 40 & 8 \\\hline
1.8.4.6 & Carga y conciliación: Tres años de matrícula escolar & 0 & 24 & 0 & 8 & 8 & 40 & 8\\\hline
1.8.5.1 & Repositorio e índice histórico & 0 & 8 & 4 & 4 & 24 & 40 & 9 \\\hline
1.8.5.2 & Consulta histórica con permisos & 0 & 8 & 4 & 4 & 24 & 40 & 9\\\hline
1.8.5.3 & Exportación y comprobación de integridad & 0 & 8 & 4 & 4 & 24 & 40 & 9 \\\hline
1.8.6.1 & Detección y carga idempotente de delta & 0 & 24 & 4 & 4 & 8 & 40 & 9\\\hline
1.8.6.2 & Herramientas de retorno previo & 0 & 24 & 4 & 4 & 8 & 40 & 9 \\\hline
1.8.6.3 & Herramientas de retorno posterior y diario & 0 & 24 & 4 & 4 & 8 & 40 & 9\\\hline
\textbf{Subtotal} & \textbf{Ingeniería de migración} & 96 & 552 & 32 & 112 & 144 & \textbf{936} & \\\hline
\end{longtable}
\FuenteTabla{Estimación ascendente de Synaptix. A: Análisis Funcional; D: Datos; S: Seguridad; Q: Calidad; V: Desarrollo. Horas-persona; meses relativos propuestos.}
\endgroup


\subsection{Correspondencias y transformaciones: 1.8.2.1--1.8.2.6}
Cada paquete entregará una especificación versionada de origen, destino, claves, relaciones y reglas del conjunto. Comprende 8 horas de Análisis Funcional para significado y validación, y 16 de Datos: 8 para correspondencias y 8 para reglas y ejemplos. La estimación de 24 horas por conjunto supone una fuente principal con hasta cinco formatos asociados y una ronda consolidada de decisión; formatos adicionales o reglas controvertidas se medirán al cerrar el diagnóstico.

El Líder de Datos responderá por el resultado. La recepción exigirá cobertura de los campos y vínculos obligatorios, casos de valores ausentes, destinos operativos o históricos y conformidad funcional de las decisiones que condicionen la carga. Las decisiones pendientes quedarán identificadas y bloquearán la transformación afectada. Se requieren el catálogo, el informe de calidad, el modelo de SD5 y los validadores del cliente. El hito será la recepción de la especificación del conjunto, anterior a su saneamiento y construcción de carga. Referencias: SD5, secciones 5.1 y 5.3; RNF-MIG-01; BT RT-05.11.

\subsection{Saneamiento automatizado: 1.8.3.1--1.8.3.6}
Cada paquete entregará reglas ejecutables y una copia saneada con registro de cambios, conservando el origen. Datos dedicará 16 horas a implementar reglas y 8 a ejecutar y documentar su resultado; Análisis Funcional dedicará 4 a decisiones y Calidad 4 a comprobar reproducibilidad. Las 32 horas por conjunto suponen hasta cinco clases de defecto tratables mediante reglas; no incluyen lectura manual exhaustiva, OCR ni recuperación de archivos dañados.

El Líder de Datos responderá por el resultado. La aceptación comprobará que cada corrección identifica regla, versión y registros afectados; los defectos no corregibles conservarán decisión y responsable. No se inventarán pagos ni autorizaciones. Se requieren copia protegida y reglas recibidas en 1.8.2. Su hito habilitará la prueba del cargador correspondiente. Referencias: SD5, secciones 5.2.4 y 5.3.2; BT RT-05.12.

\subsection{Herramientas de carga y conciliación: 1.8.4.1--1.8.4.6}
Cada paquete entregará un cargador y sus controles de lote para el conjunto correspondiente. Datos dedicará 16 horas a transformación y carga y 8 a manifiestos y conciliación; Desarrollo dedicará 8 a automatización y manejo de errores; Calidad dedicará 8 a pruebas locales de reejecución y controles. El esfuerzo será 40 horas por conjunto. Las herramientas utilizarán el diagnóstico y las interfaces especificadas; la estimación supone una ruta de carga principal por conjunto, sin incluir extracción por ingeniería inversa.

El Líder de Datos responderá por el resultado. La aceptación verificará correspondencias, recuentos, sumas monetarias separadas por moneda, integridad y reejecución sin efectos duplicados. La prueba local de un cargador no será un ensayo completo de migración. Se requieren reglas recibidas, copia saneada y esquema destino. El hito será la recepción del cargador para su integración en 1.10.1. Referencias: SD5, secciones 5.3.1--5.3.2; BT RT-05.13--RT-05.14.

\subsection{Archivo histórico: 1.8.5.1--1.8.5.3}
Los paquetes entregarán, respectivamente, repositorio e índice; consulta con permisos; y exportación con comprobación de integridad. Cada uno se estima en 40 horas: Desarrollo 24, Datos 8, Seguridad 4 y Calidad 4. Desarrollo distribuirá su trabajo en 16 horas de implementación y 8 de pruebas y documentación. La reutilización de almacenamiento, identidad y componentes de consulta de SD4 es un supuesto de estimación; su construcción base pertenece a otras ramas y no se suma nuevamente aquí.

El Líder de Datos responderá por los tres resultados, con ejecución de Desarrollo. La recepción del repositorio comprobará índice, relaciones y categorías de retención; la consulta comprobará perfiles, búsqueda y acceso a antecedentes; la exportación comprobará formatos abiertos, manifiestos e integridad. La consulta será independiente del ejecutable de 2014. La exportación de los 14 expedientes incluirá los controles PDF/A de CA-22; no se supondrá que la descarga de un archivo equivale a esa validación. Se requieren contratos de almacenamiento e identidad, y la clasificación de 1.8.2. Los hitos serán recepción del repositorio, de la consulta y de la exportación, anteriores a su ensayo integral. Referencias: SD5, secciones 5.2.5 y 5.3; BT RT-05.15; SD3 CA-22.

\subsection{Delta y retorno: 1.8.6.1--1.8.6.3}
Los paquetes entregarán detección y carga de altas, cambios y bajas; herramientas de retorno previo a nueva escritura; y herramientas de retorno posterior con diario conciliado. Cada paquete comprende 24 horas de Datos ---16 de implementación y 8 de pruebas locales---, 8 de Desarrollo para integración, 4 de Seguridad para permisos y protección y 4 de Calidad para revisión del resultado: 40 horas en total.

El Líder de Datos responderá por las herramientas. La aceptación del delta comprobará cambios sobre copias sucesivas sin presumir captura automática disponible. El retorno previo preservará la autoridad anterior y el diario; el posterior identificará hechos nuevos, respuestas externas pendientes y efectos que el legado no puede representar. Las herramientas no descartarán hechos para recuperar una copia anterior. Se requieren copia recuperable, correspondencias y contratos administrativos; la estimación supone que los cargadores ya están recibidos. La recepción habilitará los ensayos cronometrados de 1.10.1; las 40 horas no demuestran los umbrales de 90 o 120 minutos de SD5. Referencias: SD5, sección 5.3.3; BT RT-05.11 y RT-20.02.

\subsection{Revisión completa de expedientes y autorizaciones}
La revisión documental tendrá paquetes adicionales de tamaño controlable. Para cada expediente de abandono $e$, se define \texttt{1.8.3.7.e.b}, donde $b$ identifica un lote de hasta 24 piezas inventariadas. Cada lote se estima en 8 horas: Datos 2 para índice y vínculos, Análisis Funcional 4 para legibilidad y significado, y Calidad 2 para comprobación y evidencia. Estas tasas suponen piezas legibles y disponibles; digitalización, recuperación y aportación de antecedentes ausentes se tratarán por separado sin omitirlos del universo de recepción.

Sea $N_e$ la cantidad inventariada de piezas del expediente $e$. El número de lotes será $K=\sum_{e=1}^{14}\max(1,\lceil N_e/24\rceil)$. El lote de un expediente sin piezas entregará la constatación de ausencia y sus bloqueos; no acreditará un expediente completo. El esfuerzo propuesto será $8K$ horas. El Líder de Datos responderá por el registro; la aprobación del contenido probatorio corresponderá a Administración y la instancia competente del cliente. La recepción completa exigirá las piezas de los 14 expedientes con destino explicado y los controles de CA-22.

Para las autorizaciones escolares se define \texttt{1.8.3.8.b}, por lotes de 32--256 autorizaciones; un remanente inferior a 32 se incorporará a otro lote compatible. Se propone 0,25 horas por autorización: 0,15 de Análisis Funcional, 0,05 de Datos y 0,05 de Calidad. Con $N_A$ autorizaciones exigibles al corte, el esfuerzo será $0{,}25N_A$ horas. Si el universo completo es menor que 32, se formulará un paquete conjunto de revisión de al menos 8 horas que incluya sus decisiones y documentación, sin declarar una población ficticia.

El Líder de Datos responderá por los vínculos y el registro; Escuela validará vigencia y facultades de quienes autorizan. La aceptación verificará el 100\,\% de las autorizaciones que se utilizarán al corte y distinguirá antecedentes históricos de permisos vigentes. El inventario de tres años de matrícula no equivale a $N_A$. La revisión se reservará antes de cada activación escolar y se actualizará si cambian participantes o autorizaciones durante la marcha blanca; esa actualización pertenecerá a los paquetes de implantación. Referencia: SD5, secciones 5.3.1--5.3.3.

\subsection{Perímetro de la estimación}
Los 28 paquetes de ingeniería suman 936 horas-persona: Análisis Funcional 96, Datos 552, Seguridad 32, Calidad 112 y Desarrollo 144. El subtotal incluye las 120 horas de diagnóstico y sus pruebas locales. Con autorizaciones suficientes para aplicar los lotes indicados, el esfuerzo de esta parte de 1.8 será $936+8K+0{,}25N_A$; el caso de un universo escolar menor que 32 conservará el mínimo de paquete descrito. No se asigna todavía un valor observado a $K$ ni a $N_A$.

El subtotal excluye insumos, licencias no compartidas, respuesta a defectos fuera de los supuestos, los dos ensayos completos, los cortes de producción, marcha blanca, estabilización y operación. Los resultados del diagnóstico determinarán el número real de formatos, reglas y lotes y permitirán actualizar la estimación antes de cerrar la línea base. La oficina de proyecto registrará su impacto en costo, recursos y calendario sin modificar unilateralmente los hitos contractuales.

\subsection{Ensayos completos y recepción previa a los cortes}
Synaptix realizará dos ensayos completos de migración en PREPROD antes del corte definitivo de cada etapa. E1 comprobará los seis conjuntos de SD5 distinguiendo datos activos y preparados para E2. Antes del corte E2 se repetirán los dos ensayos con versiones y delta efectivos; la evidencia E1 no sustituirá los cambios posteriores \parencite[cap.~5, RT-05.13--RT-05.14; cap.~20, secc.~20.1]{bases-transversales}.

Se define \texttt{1.10.1.e.j.k}: $e$ es la etapa, $j$ el ensayo 1 o 2 y $k$ uno de los cinco paquetes de la Tabla~\ref{tab:sd7-ensayo-paquetes}. El Líder de Calidad y Pruebas responde por evidencias y Datos ejecuta procedimientos. Todos los ensayos recorren extracción, transformación, carga, conciliación y retorno completos.
\begin{table}[htbp]
\centering
\begin{tabularx}{\linewidth}{L{9mm}YL{48mm}R{12mm}}
\hline
\EncabezadoTabla{$k$} & \EncabezadoTabla{Entregable} & \EncabezadoTabla{Esfuerzo por perfil} & \EncabezadoTabla{HH}\\\hline
1 & PREPROD, origen y protocolo preparados & Datos 4; Calidad 8; Seguridad 4 & 16\\\hline
2 & Ejecución, manifiestos y tiempos & Datos 16; Calidad 16; Integración 8; Operación 8 & 48\\\hline
3 & Conciliación y revisión funcional & Funcional 12; Datos 12; Calidad 16; Seguridad 8 & 48\\\hline
4 & Retorno previo y posterior comprobados & Datos 12; Calidad 12; Operación 8; Seguridad 8 & 40\\\hline
5 & Expediente presentado a recepción & Funcional 4; Calidad 8; Proyecto 4 & 16\\\hline
\textbf{Total} & \textbf{Un ensayo} & & \textbf{168}\\\hline
\end{tabularx}
\caption{Paquetes de cada ensayo completo}\label{tab:sd7-ensayo-paquetes}
\FuenteTabla{Estimación de Synaptix; SD5, sección 5.3, y BT RT-05.13--RT-05.14.}
\end{table}

Dos ensayos por etapa suman 336 horas; cuatro, 672. Por ensayo: Funcional 16, Datos 44, Calidad 60, Seguridad 20, Integración 8, Operación 16 y Proyecto 4. Se supone una ejecución y revisión con herramientas recibidas y fuentes accesibles. Correcciones de herramientas se imputarán a construcción y reensayos adicionales se estimarán aparte, sin reducir los obligatorios. 

La línea base utiliza inicio el 1 de diciembre de 2026. Los ensayos E1 se ejecutarán del 12 al 26 de agosto y del 24 de septiembre al 8 de octubre de 2027; revisión y subsanación cierran el 9 de noviembre. Los ensayos E2 se ejecutarán del 24 de enero al 7 de febrero y del 6 al 20 de marzo de 2028; revisión y subsanación cierran el 18 de abril. Cada ensayo reserva diez días hábiles de ejecución, diez de revisión y diez de subsanación, calculados con el calendario de feriados del escenario, sin convertir todo mes en veinte días hábiles. VM-E1 se libera el 12 de agosto de 2027 y VM-E2 el 24 de enero de 2028 para migración; VS-E1 se libera el 18 de agosto y VS-E2 el 6 de marzo para pruebas integrales. Son puertas internas de distinta cobertura, no recepciones contractuales anticipadas. H4/H9 conservan entrega para pruebas y H5/H10 certificación en sus meses E-25. Los ensayos de PREPROD durante congelamiento no intervendrán producción, maniobras ni usuarios del sitio. Toda modificación que invalide migración obliga a repetir integralmente la evidencia afectada; los cambios restantes requieren regresión. La reserva temporal no acredita HH ejecutadas ni permite aceptar observaciones pendientes. Las fechas y predecesoras se conservan en \path{componentes/red_global.csv} y se contrastan con los recursos de T-15.\parencite[arts.~17--18; Formulario~E-25]{bases-admin}

Cada paquete se recibirá con: protocolo y origen identificados; manifiestos y tiempos completos; diferencias explicadas y controles aceptados; retorno sin pérdida demostrado; y expediente trazable con observaciones y decisión. Los umbrales serán los de SD5, incluyendo expedientes y autorizaciones íntegros. El segundo ensayo incluirá suplentes técnico y funcional. Recursos: PREPROD aislado, copia protegida, herramientas, perfiles de prueba y validadores. El cliente autoriza el corte; las 168 horas no demuestran previamente los tiempos de retorno.

\subsection{Activación de grupos, convivencia y estabilización}
Se define \texttt{1.11.1.g.k}, con grupos $g=1,2,3$: maestros e inventario; antecedentes operacionales; contratos, saldos y sustitución administrativa. La Tabla~\ref{tab:sd7-implantacion-paquetes} estima una activación consolidada por grupo, bajo el Líder de Implantación y Gestión del Cambio. Activaciones adicionales por zona serán paquetes hijos con esfuerzo propio.
\begin{table}[htbp]
\centering
\begin{tabularx}{\linewidth}{L{9mm}YL{48mm}R{12mm}}
\hline
\EncabezadoTabla{$k$} & \EncabezadoTabla{Entregable} & \EncabezadoTabla{Esfuerzo por perfil} & \EncabezadoTabla{HH}\\\hline
1 & Lista de habilitación y formación & Funcional 8; Datos 8; Calidad 8; Implantación 8 & 32\\\hline
2 & Activación y decisión documentadas & Datos 12; Calidad 8; Seguridad 4; Implantación 8; Proyecto 4 & 36\\\hline
3 & Conciliación inicial del grupo & Funcional 8; Datos 8; Calidad 8; Implantación 8 & 32\\\hline
\textbf{Total} & \textbf{Una activación} & & \textbf{100}\\\hline
\end{tabularx}
\caption{Paquetes por grupo de implantación}\label{tab:sd7-implantacion-paquetes}
\FuenteTabla{SD5, sección 5.3; BT RT-20.01--RT-20.05; estimación de Synaptix.}
\end{table}

Las tres activaciones suman 300 horas. La habilitación se recibirá con datos conciliados, permisos, usuarios preparados y ventana autorizada; la activación, con versión, tiempos y decisión; la conciliación inicial, con autoridad y diferencias comprobadas. Recursos: operadores, copia recuperable y contingencia. La formación corresponde al cambio de datos y excluye la capacitación de todos los procesos. Se proponen preparación de grupos 1 y 2 en el mes 12 y activación en 13; preparación del grupo 3 en 18 y activación en 19, sujetas a aceptación y ventanas. Producción oficial conserva meses 16 y 21. La precarga y las escrituras supervisadas de marcha blanca no anticipan la autoridad administrativa oficial. El legado seguirá siendo el registro oficial del alcance E2 hasta el acta de cierre y habilitación del mes 21; los efectos externos se ejecutarán una sola vez por la autoridad vigente y se contrastarán en la solución supervisada. El corte de activación del mes 19 y el paso oficial del mes 21 son actuaciones diferentes.

Se define \texttt{1.11.1.MB.e.m.q} por quincena de marcha blanca: seis paquetes E1 y cuatro E2. Cada uno comprende Funcional 8, Datos 8, Calidad 8 y Operación 24 horas, total 48; diez suman 480. Se recibirá con registros diarios, incidencias, denominadores, indicadores y conciliación. Las comprobaciones particulares del corte pertenecen a 1.11.1.g.3; MB incluye las revisiones diarias posteriores, sin doble conteo. Las horas son esfuerzo de tareas, no cobertura presencial continua ni guardia técnica.

La estabilización de datos será \texttt{1.11.1.EST.e.q}, dos paquetes quincenales por etapa durante las cuatro semanas siguientes al paso oficial: inicio E1 al cierre de marzo de 2028 y continuidad en abril; E2 desde el 1 de agosto de 2028. La fecha de cada contribución será la de la asignación integrada, aunque su código conserve la quincena de origen. Cada uno comprende Datos 8, Calidad 8 y Operación 32, total 48; cuatro suman 192. Se recibirá con incidentes, correcciones, estabilidad diaria y traspaso al servicio ordinario. Recursos: telemetría, soporte y responsables de cada conjunto. Esta dotación de datos es parcial: pantalanes, varadero y fines de semana requieren turnos adicionales del plan general.

El responsable de MB y estabilización será Operación/SRE. Operación e Implantación comparten titular de SD1 y se sumarán por persona. El costo de cada paquete será horas por tarifas de sus perfiles, sin repetir gestión. Referencias: BA artículo 17.3, BT RT-20.03--RT-20.06 y SD5 5.3. El acta firmada no se sustituye por un informe quincenal.

Ensayos, activaciones, convivencia y estabilización suman 1.644 horas: 672, 300, 480 y 192. Con las 936 de ingeniería, el subtotal es 2.580, más revisión documental y correcciones adicionales. Las 96 de estabilización E2 del mes 21 pertenecen a Operación y no se sumarán nuevamente al servicio ordinario.


\section{Descomposición y control de las innovaciones}
Las cinco innovaciones conservan su correspondencia con SD13. Sus fichas distinguen resultados incrementales, pilotos y seguimiento, con responsabilidades y aceptación específicas.
\begin{figure}[htbp]\centering
\begin{tikzpicture}[font=\fontsize{9}{11}\selectfont,
 raiz/.style={draw=SynNavy,fill=SynSurface,rounded corners=2pt,text width=54mm,align=center,minimum height=10mm},
 hijo/.style={draw=SynBlue,fill=SynSurface,rounded corners=2pt,text width=27mm,align=center,minimum height=17mm}]
\node[raiz] (root) at (0,0) {1.9 Innovaciones incorporadas y evaluadas};
\node[hijo] (a) at (-6.4,-2) {1.9.1\newline Solicitudes de mantenimiento};
\node[hijo] (b) at (-3.2,-2) {1.9.2\newline Mantenimiento preventivo en invernada};
\node[hijo] (c) at (0,-2) {1.9.3\newline Lectura asistida};
\node[hijo] (d) at (3.2,-2) {1.9.4\newline Servicio gestionado por resultados};
\node[hijo] (e) at (6.4,-2) {1.9.5\newline Pasaporte Ambiental Náutico};
\foreach \n in {a,b,c,d,e}{\draw[SynNavy] (root.south) -- (0,-0.8) -| (\n.north);}
\end{tikzpicture}
\caption{Descomposición por entregables de la rama de innovación}
\label{fig:sd7-edt-innovaciones}
\FuenteFigura{Synaptix; correspondencia con las cinco innovaciones de SD13. Las líneas expresan pertenencia; las precedencias se registran aparte.}
\end{figure}

\subsection{Componentes incrementales de las cinco innovaciones}
La rama 1.9 conserva las cinco innovaciones de SD13 y sus entregables de diseño, construcción, pruebas específicas y evaluación. Cada código terminal de origen mantiene su correspondencia; cuando describe un período superior a una quincena se generan hijos por resultado y período. El esfuerzo se suma desde los hijos, sin agregar otra estimación al padre. Los componentes de IN4 e IN5 reciben paquetes terminales propios sin alterar los nombres de sus innovaciones.

La descomposición cuantifica 84 paquetes y 2.680 HH: IN1, 624; IN2, 712; IN3, 568; IN4, 352; IN5, 424. De ellas, 2.352 corresponden al alcance inicial con mes fijo y 328 al piloto estacional de IN2, ahora situado en 30--37 con invernada en 31. Los complementos de seguimiento y botadura se suman en el catálogo integrado. Estas horas no incluyen la revisión previa de cada embarcación ni el servicio recurrente descritos a continuación. Las horas de aceptación final de innovaciones del mes 21 son 120 y pertenecen al cierre de implementación; su soporte posterior se registra en operación. Los pilotos de IN2 desde $M_{\mathrm{inv}}$ pertenecen a la fase de operación y no se anticipan al precio de implementación.

Los paquetes específicos incluyen la evidencia de privacidad, accesibilidad, seguridad, continuidad y uso de IA que corresponde a la innovación. La construcción obligatoria de esos servicios comunes permanece en 1.2, 1.3 y 1.7; las pruebas integradas y la certificación contractual permanecen en 1.10. Una evidencia específica podrá alimentar el expediente general, sin imputar dos veces la misma ejecución \parencite[caps.~19, 20 y 26]{bases-transversales}.

\subsection{Revisión previa a botaduras y operación incremental}
\paragraph{1.9.2.5.4.k: revisión de pendientes previa a botadura}
Este paquete conserva el origen \texttt{IN2.5.4}. Se seleccionan cinco lotes $k=1,\ldots,5$ de cuatro embarcaciones: dos en 33, dos en 34 y uno en 35, para veinte embarcaciones. La distribución se ajustará al calendario confirmado de botaduras, conservando tres semanas de anticipación y sin duplicar el lote. La fecha de cada revisión deberá ser tres semanas anterior a su botadura prevista, conforme a SD13. El entregable es una lista por embarcación con plan aceptado, trabajos, pendientes, responsable de decisión y evidencia disponible. Excluye ejecutar mantenimiento físico, autorizar una navegación o asumir una certificación técnica ajena a Synaptix.

José Lara Arce responde por el paquete; Funcional revisa el plan y Calidad su evidencia. Esfuerzo: Funcional 4, Operación 8 y Calidad 4 HH; total 16. Se requieren planes, registros del piloto, calendario de botaduras confirmado, contratistas y validadores del cliente. Se recibirá con una revisión por cada embarcación del lote, diferencias identificadas y comunicación de decisiones pendientes; un pendiente sin resolver no equivale a autorización de botadura. Se estima una revisión documental consolidada por embarcación con antecedentes accesibles, no una inspección física ni una investigación de incidentes.

Para $N_{o,q}$ embarcaciones que requieren revisión, se define $B_{o,q}=\lceil N_{o,q}/4\rceil$ si $N_{o,q}>0$, y cero cuando el calendario confirmado no contiene revisiones. El esfuerzo es $16\sum B_{o,q}$ HH. La falta de calendario o antecedentes no justifica declarar cero: constituye una habilitación pendiente. Cambios de fecha obligan a actualizar la revisión y dimensionar el trabajo adicional. Predecesoras: planes aceptados e informes de seguimiento disponibles; sucesoras: resultados de botaduras y evaluación IN2.5.6. Hito: evaluación de la campaña, sin nuevo porcentaje de pago de implementación. Referencias: SD13, IN2.5.4, y ventanas del caso 07.

\paragraph{1.12.2.i.m.q: servicio quincenal de IN1, IN2 e IN3}
Se define un informe por innovación $i=1,2,3$, mes $m$ y quincena $q$. Los orígenes son \texttt{IN1.7.m.q}, \texttt{IN2.7} y \texttt{IN3.7}. José Lara Arce responde por el resultado. IN1 registra solicitudes, avisos fallidos e incidencias: Operación 4, Funcional 2 y Calidad 2 HH. IN2 utiliza un expediente mensual de 16 HH: Operación 8, Funcional 4 y Calidad 4, código 1.12.2.2.m. Conserva el mismo esfuerzo equivalente a dos informes de 8 HH y sus evidencias quincenales, sin sumar ambos formatos. IN3 registra procesamiento, consumo, formatos y fallos: Operación 4, Datos 2 y Calidad 2 HH. Los paquetes quincenales IN1/IN3 tienen 8 HH.

Se recibirá con versión del servicio, universo observado, incidencias, consumo o cambios aplicables, acciones y responsables trazados. Requiere telemetría, registros protegidos y servicio habilitado. Excluye guardias continuas, NOC/SOC, soporte base, nuevos desarrollos y los informes de piloto: estos últimos se imputan sólo a 1.9.2.5. IN1 e IN3 generan 72 paquetes cada una entre 21 y 56: 576 HH por innovación. IN2 genera $(57-M_{\mathrm{inv}})$ expedientes mensuales desde su habilitación estacional hasta 56: $16(57-M_{\mathrm{inv}})$ HH. Referencias: SD13, operación de cada innovación; contrato de operación y SLA. El volumen real de incidencias se contrastará para dimensionar soporte adicional.

\paragraph{1.12.2.4.m: expediente mensual de conciliación de resultados}
Conserva el trabajo recurrente de \texttt{IN4.4}, separado de la habilitación \texttt{1.9.4.4.1}. Ulysses Barreto responde por el expediente: Funcional 8, Calidad 4 y Proyecto 4 HH, total 16. El expediente conserva muestra, intervalos aceptados, cálculo versionado, condiciones de calidad, factor resultante, observaciones y decisión contractual. La aprobación del incentivo corresponde al cliente y no se sustituye por la preparación del informe.

Se recibirá con cálculo reproducible y vínculo a todas sus entradas; sin evidencia suficiente, muestra admisible o mejora atribuible el factor será cero. Se requieren registros validados de IN1, regla aceptada y responsables de preparación y aprobación distintos. No incluye de nuevo la captura del piloto, cobro automático ni un ahorro ya atribuido a IN1. Se genera un paquete mensual en 21--56: 36 paquetes, 576 HH. La espera de aprobación se programa aparte. Referencias: SD13, IN4.4, y E-25 para la separación de pagos de operación.

\paragraph{1.12.2.5.m: informe mensual de seguimiento ambiental}
Conserva el trabajo recurrente de \texttt{IN5.4}, separado de la habilitación \texttt{1.9.5.4.1}. Nicolás Torfan responde por el resultado: Datos 8, Calidad 4 y Operación 4 HH, total 16. Entregará cobertura de lecturas, permanencias conciliadas, consentimiento, acciones y comparación ambiental válida. Se recibirá con denominadores, períodos y exclusiones trazados; datos incompletos o incomparables impedirán declarar una reducción.

Requiere medición individual validada, historial, participantes autorizados y telemetría. Excluye compra y obras de medidores a cargo del cliente, captura base de consumo y mantenimiento físico. Se genera un paquete mensual en 21--56: 36 paquetes, 576 HH. El informe previo a la temporada 2029 se obtiene de esta serie y sus evidencias; no añade otra vez su esfuerzo de seguimiento. Referencias: SD13, IN5.4; SD2, EXC-06/EXC-10 y SUP-04; caso 07, capítulos 10 y 13.

Los cuatro servicios con calendario fijo suman 2.304 HH de operación; IN2 añade $16(57-M_{\mathrm{inv}})$ y la revisión de botaduras añade $16\sum B_{o,q}$. Los expedientes mensuales IN4/IN5 se preparan y controlan en la segunda quincena a partir de registros continuos del servicio; la espera de cierre, validación o aprobación del mes se enlaza aparte. Su frecuencia mensual no convierte el paquete de 16 HH en una tarea activa durante todo el mes. Son cargas incrementales por resultado, sin equivalencia a cobertura presencial ni atención continua. Se agregarán por persona a los paquetes generales de operación y a la estabilización de datos, conservando una imputación por tarea.

\section{Diccionario de los paquetes incrementales de innovación}
Cada ficha identifica un único resultado, su responsable último, aceptación, esfuerzo por perfil, recursos, supuestos y referencia. Los códigos IN1--IN3 conservan los terminales de SD13; los períodos amplios se subdividen en informes quincenales. IN4 e IN5 conservan sus componentes de origen y reciben terminales para planificación y control. La aceptación de cada ficha describe una obligación futura; la Contraparte Técnica mantiene la recepción contractual.

Se supone reutilización de identidad, portales, repositorio, contratos, telemetría y sincronización de la solución base. Sus costos de construcción no se repiten aquí. Cada paquete de software incluye sus pruebas unitarias y documentación; las pruebas específicas de innovación se registran en sus fichas y alimentan 1.10 sin repetir su ejecución. Las cifras son estimaciones iniciales, sin mediciones de productividad ejecutadas. Los pilotos incluyen registro y evaluación por parte de Synaptix; la participación y decisiones del cliente son habilitaciones externas, sin horas de empleados del cliente imputadas como personal de Synaptix.

Los perfiles I (Arquitectura/Integración) y O/L (Operación/Formación) comparten la disponibilidad de sus respectivos titulares; una ficha no reserva otra persona por utilizar dos perfiles.

La recepción contractual corresponde a la contraparte autorizada. Todas las fichas requieren insumos y versiones recibidos, registros y entornos autorizados, herramientas y contraparte programada. Las HH excluyen dedicación del cliente y cobertura externa continua. La falta de insumos o evidencia bloquea la recepción y exige reprogramar dentro de las condiciones contractuales. Estas condiciones comunes se aplican además de los criterios específicos de cada ficha.

Las fichas de la Tabla~\ref{tab:dic-compacto-innovaciones} presentan resultado, responsable, esfuerzo y aceptación. Las condiciones particulares, recursos y ventanas por código se conservan en \path{componentes/condiciones_fichas_t14.csv}, junto con la asignación y precedencias del diccionario técnico integrado. Las actividades de T-15 desarrollan los aportes por perfil y el T-18 desarrolla la ejecución de implantación.


\begingroup\small\setlength{\tabcolsep}{3pt}
\begin{longtable}{L{52mm}L{91mm}}
\caption{Diccionario de innovaciones}\label{tab:dic-compacto-innovaciones}\\\hline
 \EncabezadoTabla{Paquete, resultado y responsable} & \EncabezadoTabla{Contenido y criterio de aceptación}\\\hline\endfirsthead
\hline \EncabezadoTabla{Paquete, resultado y responsable} & \EncabezadoTabla{Contenido y criterio de aceptación}\\\hline\endhead
\texttt{1.\allowbreak{}9.\allowbreak{}1.\allowbreak{}1.\allowbreak{}1}\newline Flujo, estados y matriz de roles revisados\newline Fernando Guerrero; 32 HH. & Se entregará flujo, estados y matriz de roles revisados, correspondiente a \texttt{IN1.1.1}. Excluye los resultados de las fichas predecesoras, la construcción de capacidades base y el servicio recurrente.  Cada transición tiene actor, precondición, rechazo y autoridad identificados; se valida el aislamiento entre armadores y contratistas. Fernando Guerrero, Arquitectura/Integración, responde por el resultado; los demás perfiles ejecutan sus aportes y el cliente valida las decisiones de su ámbito.\\\hline
\texttt{1.\allowbreak{}9.\allowbreak{}1.\allowbreak{}1.\allowbreak{}2}\newline Especificación de las interfaces\newline Fernando Guerrero; 40 HH. & Se entregará especificación de las interfaces, correspondiente a \texttt{IN1.1.2}. Excluye los resultados de las fichas predecesoras, la construcción de capacidades base y el servicio recurrente.  Contratos OpenAPI y eventos versionados identifican errores, idempotencia, correlación y pruebas de compatibilidad. Fernando Guerrero, Arquitectura/Integración, responde por el resultado; los demás perfiles ejecutan sus aportes y el cliente valida las decisiones de su ámbito.\\\hline
\texttt{1.\allowbreak{}9.\allowbreak{}1.\allowbreak{}1.\allowbreak{}3}\newline Protocolo de medición y aceptación\newline Davor Serey; 24 HH. & Se entregará protocolo de medición y aceptación, correspondiente a \texttt{IN1.1.3}. Excluye los resultados de las fichas predecesoras, la construcción de capacidades base y el servicio recurrente.  Quedan fijados cohorte, registro de primera respuesta, tratamiento de casos abiertos y criterio de aceptación antes de medir. Davor Serey, Calidad, responde por el resultado; los demás perfiles ejecutan sus aportes y el cliente valida las decisiones de su ámbito.\\\hline
\texttt{1.\allowbreak{}9.\allowbreak{}1.\allowbreak{}2.\allowbreak{}1}\newline Servicio de solicitudes y transiciones\newline Nikolai Navea; 56 HH. & Se entregará servicio de solicitudes y transiciones, correspondiente a \texttt{IN1.2.1}. Comprende implementación, pruebas unitarias y documentación del resultado. Excluye los resultados de las fichas predecesoras, la construcción de capacidades base y el servicio recurrente.  Casos válidos e inválidos y escrituras repetidas conservan una sola transición autorizada y su evidencia. Nikolai Navea, Desarrollo, responde por el resultado; los demás perfiles ejecutan sus aportes y el cliente valida las decisiones de su ámbito.\\\hline
\texttt{1.\allowbreak{}9.\allowbreak{}1.\allowbreak{}2.\allowbreak{}2}\newline Formularios y vistas en AUT\newline Nikolai Navea; 48 HH. & Se entregará formularios y vistas en aut, correspondiente a \texttt{IN1.2.2}. Comprende implementación, pruebas unitarias y documentación del resultado. Excluye los resultados de las fichas predecesoras, la construcción de capacidades base y el servicio recurrente.  Armador y contratista sólo ven y modifican su ámbito; las vistas muestran estados y errores y admiten navegación accesible. Nikolai Navea, Desarrollo, responde por el resultado; los demás perfiles ejecutan sus aportes y el cliente valida las decisiones de su ámbito.\\\hline
\texttt{1.\allowbreak{}9.\allowbreak{}1.\allowbreak{}2.\allowbreak{}3}\newline Historial y registro asistido auditable\newline Nikolai Navea; 48 HH. & Se entregará historial y registro asistido auditable, correspondiente a \texttt{IN1.2.3}. Comprende implementación, pruebas unitarias y documentación del resultado. Excluye los resultados de las fichas predecesoras, la construcción de capacidades base y el servicio recurrente.  Cada cambio registra actor, fecha, estado anterior y nuevo; el ingreso asistido identifica quién lo realizó. Nikolai Navea, Desarrollo, responde por el resultado; los demás perfiles ejecutan sus aportes y el cliente valida las decisiones de su ámbito.\\\hline
\texttt{1.\allowbreak{}9.\allowbreak{}1.\allowbreak{}3.\allowbreak{}1}\newline Conexiones con CTR y VAR verificadas\newline Fernando Guerrero; 48 HH. & Se entregará conexiones con ctr y var verificadas, correspondiente a \texttt{IN1.3.1}. Comprende implementación, pruebas unitarias y documentación del resultado. Excluye los resultados de las fichas predecesoras, la construcción de capacidades base y el servicio recurrente.  Solicitudes y faenas se vinculan por contrato sin escribir tablas de otro dominio y toleran reintentos. Fernando Guerrero, Arquitectura/Integración, responde por el resultado; los demás perfiles ejecutan sus aportes y el cliente valida las decisiones de su ámbito.\\\hline
\texttt{1.\allowbreak{}9.\allowbreak{}1.\allowbreak{}3.\allowbreak{}2}\newline Identidad y vínculo de armador integrados\newline Bruno Toro; 40 HH. & Se entregará identidad y vínculo de armador integrados, correspondiente a \texttt{IN1.3.2}. Comprende implementación, pruebas unitarias y documentación del resultado. Excluye los resultados de las fichas predecesoras, la construcción de capacidades base y el servicio recurrente.  Los casos de suplantación y acceso cruzado se rechazan y los vínculos autorizados se comprueban. Bruno Toro, Seguridad, responde por el resultado; los demás perfiles ejecutan sus aportes y el cliente valida las decisiones de su ámbito.\\\hline
\texttt{1.\allowbreak{}9.\allowbreak{}1.\allowbreak{}3.\allowbreak{}3}\newline Notificaciones y reintentos integrados\newline Nikolai Navea; 40 HH. & Se entregará notificaciones y reintentos integrados, correspondiente a \texttt{IN1.3.3}. Comprende implementación, pruebas unitarias y documentación del resultado. Excluye los resultados de las fichas predecesoras, la construcción de capacidades base y el servicio recurrente.  Avisos fallidos quedan visibles y los reintentos no duplican efectos ni conceden una habilitación. Nikolai Navea, Desarrollo, responde por el resultado; los demás perfiles ejecutan sus aportes y el cliente valida las decisiones de su ámbito.\\\hline
\texttt{1.\allowbreak{}9.\allowbreak{}1.\allowbreak{}4.\allowbreak{}1}\newline Informe de pruebas funcionales\newline Davor Serey; 32 HH. & Se entregará informe de pruebas funcionales, correspondiente a \texttt{IN1.4.1}. Comprende preparación, ejecución de comprobaciones y expediente. Excluye los resultados de las fichas predecesoras, la construcción de capacidades base y el servicio recurrente.  Se comprueban estados, duplicados, concurrencia, habilitación y flujos completos; cada defecto queda trazado. Davor Serey, Calidad, responde por el resultado; los demás perfiles ejecutan sus aportes y el cliente valida las decisiones de su ámbito.\\\hline
\texttt{1.\allowbreak{}9.\allowbreak{}1.\allowbreak{}4.\allowbreak{}2}\newline Informe de seguridad y recuperación\newline Bruno Toro; 40 HH. & Se entregará informe de seguridad y recuperación, correspondiente a \texttt{IN1.4.2}. Comprende preparación, ejecución de comprobaciones y expediente. Excluye los resultados de las fichas predecesoras, la construcción de capacidades base y el servicio recurrente.  No quedan vulnerabilidades críticas o altas abiertas; aislamiento y recuperación se prueban sobre versión identificada. Bruno Toro, Seguridad, responde por el resultado; los demás perfiles ejecutan sus aportes y el cliente valida las decisiones de su ámbito.\\\hline
\texttt{1.\allowbreak{}9.\allowbreak{}1.\allowbreak{}5.\allowbreak{}1}\newline Guía operativa y capacitación registrada\newline José Lara Arce; 32 HH. & Se entregará guía operativa y capacitación registrada, correspondiente a \texttt{IN1.5.1}. Comprende registro, validación y emisión del resultado. Excluye los resultados de las fichas predecesoras, la construcción de capacidades base y el servicio recurrente.  Los participantes reproducen los flujos; quedan material, asistencia, evaluación y soporte identificados. José Lara Arce, Implantación, responde por el resultado; los demás perfiles ejecutan sus aportes y el cliente valida las decisiones de su ámbito.\\\hline
\texttt{1.\allowbreak{}9.\allowbreak{}1.\allowbreak{}5.\allowbreak{}2}\newline Registro base de la primera quincena\newline Davor Serey; 16 HH. & Se entregará registro base de la primera quincena, correspondiente a \texttt{IN1.5.2}. Comprende registro, validación y emisión del resultado. Excluye los resultados de las fichas predecesoras, la construcción de capacidades base y el servicio recurrente.  Las solicitudes de la base conservan apertura, respuesta y casos sin cerrar; el flujo adicional sigue desactivado para la cohorte. Davor Serey, Calidad, responde por el resultado; los demás perfiles ejecutan sus aportes y el cliente valida las decisiones de su ámbito.\\\hline
\texttt{1.\allowbreak{}9.\allowbreak{}1.\allowbreak{}5.\allowbreak{}3}\newline Registro base de la segunda quincena\newline Davor Serey; 16 HH. & Se entregará registro base de la segunda quincena, correspondiente a \texttt{IN1.5.3}. Comprende registro, validación y emisión del resultado. Excluye los resultados de las fichas predecesoras, la construcción de capacidades base y el servicio recurrente.  El segundo registro usa el mismo método y cohorte, identifica exclusiones y cierra el período base. Davor Serey, Calidad, responde por el resultado; los demás perfiles ejecutan sus aportes y el cliente valida las decisiones de su ámbito.\\\hline
\texttt{1.\allowbreak{}9.\allowbreak{}1.\allowbreak{}5.\allowbreak{}4}\newline Registro de piloto de la primera quincena\newline José Lara Arce; 24 HH. & Se entregará registro de piloto de la primera quincena, correspondiente a \texttt{IN1.5.4}. Comprende registro, validación y emisión del resultado. Excluye los resultados de las fichas predecesoras, la construcción de capacidades base y el servicio recurrente.  Se conserva el universo de solicitudes y la primera respuesta de cada una, incluidos rechazos y casos abiertos. José Lara Arce, Implantación, responde por el resultado; los demás perfiles ejecutan sus aportes y el cliente valida las decisiones de su ámbito.\\\hline
\texttt{1.\allowbreak{}9.\allowbreak{}1.\allowbreak{}5.\allowbreak{}5}\newline Registro de piloto de la segunda quincena\newline José Lara Arce; 24 HH. & Se entregará registro de piloto de la segunda quincena, correspondiente a \texttt{IN1.5.5}. Comprende registro, validación y emisión del resultado. Excluye los resultados de las fichas predecesoras, la construcción de capacidades base y el servicio recurrente.  El registro conserva método comparable, incidencias, cobertura y pérdidas de observación explicadas. José Lara Arce, Implantación, responde por el resultado; los demás perfiles ejecutan sus aportes y el cliente valida las decisiones de su ámbito.\\\hline
\texttt{1.\allowbreak{}9.\allowbreak{}1.\allowbreak{}5.\allowbreak{}6}\newline Informe comparativo y acta de evaluación\newline Davor Serey; 32 HH. & Se entregará informe comparativo y acta de evaluación, correspondiente a \texttt{IN1.5.6}. Comprende registro, validación y emisión del resultado. Excluye los resultados de las fichas predecesoras, la construcción de capacidades base y el servicio recurrente.  La mediana y la trazabilidad se comparan con la base sin inventar beneficio; la decisión y observaciones quedan registradas. Davor Serey, Calidad, responde por el resultado; los demás perfiles ejecutan sus aportes y el cliente valida las decisiones de su ámbito.\\\hline
\texttt{1.\allowbreak{}9.\allowbreak{}1.\allowbreak{}6.\allowbreak{}1}\newline Servicio desplegado y transferido a operación\newline José Lara Arce; 32 HH. & Se entregará servicio desplegado y transferido a operación, correspondiente a \texttt{IN1.6.1}. Comprende registro, validación y emisión del resultado. Excluye los resultados de las fichas predecesoras, la construcción de capacidades base y el servicio recurrente.  Versión, permisos, procedimiento, soporte y decisión de habilitación quedan recibidos; no se presenta un piloto como aceptación. José Lara Arce, Operación/SRE, responde por el resultado; los demás perfiles ejecutan sus aportes y el cliente valida las decisiones de su ámbito.\\\hline
\texttt{1.\allowbreak{}9.\allowbreak{}2.\allowbreak{}1.\allowbreak{}1}\newline Catálogo de mantenimiento y reglas de planificación definidos\newline Ulysses Barreto; 32 HH. & Se entregará catálogo de mantenimiento y reglas de planificación definidos, correspondiente a \texttt{IN2.1.1}. Excluye los resultados de las fichas predecesoras, la construcción de capacidades base y el servicio recurrente.  Cada tarea tiene condición, frecuencia, antecedente, responsable y límite de recomendación; no se diagnostica una embarcación sin evidencia. Ulysses Barreto, Funcional, responde por el resultado; los demás perfiles ejecutan sus aportes y el cliente valida las decisiones de su ámbito.\\\hline
\texttt{1.\allowbreak{}9.\allowbreak{}2.\allowbreak{}1.\allowbreak{}2}\newline Ficha, responsabilidades e interfaces especificadas\newline Fernando Guerrero; 32 HH. & Se entregará ficha, responsabilidades e interfaces especificadas, correspondiente a \texttt{IN2.1.2}. Excluye los resultados de las fichas predecesoras, la construcción de capacidades base y el servicio recurrente.  Se distinguen recomendación, aceptación del armador y ejecución por contratista; cada interfaz y permiso tiene autoridad. Fernando Guerrero, Arquitectura/Integración, responde por el resultado; los demás perfiles ejecutan sus aportes y el cliente valida las decisiones de su ámbito.\\\hline
\texttt{1.\allowbreak{}9.\allowbreak{}2.\allowbreak{}2.\allowbreak{}1}\newline Ficha de ingreso de antecedentes disponible en AUT\newline Nikolai Navea; 48 HH. & Se entregará ficha de ingreso de antecedentes disponible en aut, correspondiente a \texttt{IN2.2.1}. Comprende implementación, pruebas unitarias y documentación del resultado. Excluye los resultados de las fichas predecesoras, la construcción de capacidades base y el servicio recurrente.  Campos ausentes se identifican y no se convierten en datos técnicos confirmados; acceso restringido al armador autorizado. Nikolai Navea, Desarrollo, responde por el resultado; los demás perfiles ejecutan sus aportes y el cliente valida las decisiones de su ámbito.\\\hline
\texttt{1.\allowbreak{}9.\allowbreak{}2.\allowbreak{}2.\allowbreak{}2}\newline Catálogo configurable y aceptación de propuestas implementados\newline Nikolai Navea; 56 HH. & Se entregará catálogo configurable y aceptación de propuestas implementados, correspondiente a \texttt{IN2.2.2}. Comprende implementación, pruebas unitarias y documentación del resultado. Excluye los resultados de las fichas predecesoras, la construcción de capacidades base y el servicio recurrente.  Cambios de catálogo se versionan y ninguna propuesta genera una orden sin aceptación explícita. Nikolai Navea, Desarrollo, responde por el resultado; los demás perfiles ejecutan sus aportes y el cliente valida las decisiones de su ámbito.\\\hline
\texttt{1.\allowbreak{}9.\allowbreak{}2.\allowbreak{}3.\allowbreak{}1}\newline Calendario integrado con posiciones, botaduras y habilitación\newline Fernando Guerrero; 56 HH. & Se entregará calendario integrado con posiciones, botaduras y habilitación, correspondiente a \texttt{IN2.3.1}. Comprende implementación, pruebas unitarias y documentación del resultado. Excluye los resultados de las fichas predecesoras, la construcción de capacidades base y el servicio recurrente.  Se detectan incompatibilidades de posición, fecha y habilitación y las decisiones conservan su origen. Fernando Guerrero, Arquitectura/Integración, responde por el resultado; los demás perfiles ejecutan sus aportes y el cliente valida las decisiones de su ámbito.\\\hline
\texttt{1.\allowbreak{}9.\allowbreak{}2.\allowbreak{}3.\allowbreak{}2}\newline Trabajos aceptados vinculados a solicitudes IN1\newline Nikolai Navea; 48 HH. & Se entregará trabajos aceptados vinculados a solicitudes in1, correspondiente a \texttt{IN2.3.2}. Comprende implementación, pruebas unitarias y documentación del resultado. Excluye los resultados de las fichas predecesoras, la construcción de capacidades base y el servicio recurrente.  Cada trabajo aceptado crea o vincula una solicitud una sola vez y conserva actor, consentimiento y correlación. Nikolai Navea, Desarrollo, responde por el resultado; los demás perfiles ejecutan sus aportes y el cliente valida las decisiones de su ámbito.\\\hline
\texttt{1.\allowbreak{}9.\allowbreak{}2.\allowbreak{}3.\allowbreak{}3}\newline Permisos y notificaciones configurados\newline Bruno Toro; 40 HH. & Se entregará permisos y notificaciones configurados, correspondiente a \texttt{IN2.3.3}. Comprende implementación, pruebas unitarias y documentación del resultado. Excluye los resultados de las fichas predecesoras, la construcción de capacidades base y el servicio recurrente.  El contratista ve sólo su encargo y los avisos no equivalen a autorización ni duplican solicitudes. Bruno Toro, Seguridad, responde por el resultado; los demás perfiles ejecutan sus aportes y el cliente valida las decisiones de su ámbito.\\\hline
\texttt{1.\allowbreak{}9.\allowbreak{}2.\allowbreak{}4.\allowbreak{}1}\newline Evidencias de pruebas funcionales y de integración\newline Davor Serey; 32 HH. & Se entregará evidencias de pruebas funcionales y de integración, correspondiente a \texttt{IN2.4.1}. Comprende preparación, ejecución de comprobaciones y expediente. Excluye los resultados de las fichas predecesoras, la construcción de capacidades base y el servicio recurrente.  Se prueban propuestas, rechazo, aceptación, cambio de fecha, contratista no habilitado y reintentos. Davor Serey, Calidad, responde por el resultado; los demás perfiles ejecutan sus aportes y el cliente valida las decisiones de su ámbito.\\\hline
\texttt{1.\allowbreak{}9.\allowbreak{}2.\allowbreak{}4.\allowbreak{}2}\newline Seguridad, recuperación y aceptación técnica\newline Bruno Toro; 40 HH. & Se entregará seguridad, recuperación y aceptación técnica, correspondiente a \texttt{IN2.4.2}. Comprende preparación, ejecución de comprobaciones y expediente. Excluye los resultados de las fichas predecesoras, la construcción de capacidades base y el servicio recurrente.  Se comprueban aislamiento, continuidad y retorno; no quedan defectos críticos o altos abiertos y se registra decisión técnica. Bruno Toro, Seguridad, responde por el resultado; los demás perfiles ejecutan sus aportes y el cliente valida las decisiones de su ámbito.\\\hline
\texttt{1.\allowbreak{}9.\allowbreak{}2.\allowbreak{}5.\allowbreak{}1}\newline Participantes inscritos, fichas y capacitación registrada\newline José Lara Arce; 32 HH. & Se entregará participantes inscritos, fichas y capacitación registrada, correspondiente a \texttt{IN2.5.1}. Comprende registro, validación y emisión del resultado. Excluye los resultados de las fichas predecesoras, la construcción de capacidades base y el servicio recurrente.  Participación autorizada, fichas completas o ausencias identificadas y evaluación de capacitación por participante. José Lara Arce, Implantación, responde por el resultado; los demás perfiles ejecutan sus aportes y el cliente valida las decisiones de su ámbito.\\\hline
\texttt{1.\allowbreak{}9.\allowbreak{}2.\allowbreak{}5.\allowbreak{}2}\newline Planes aceptados y calendario confirmado\newline Ulysses Barreto; 40 HH. & Se entregará planes aceptados y calendario confirmado, correspondiente a \texttt{IN2.5.2}. Comprende registro, validación y emisión del resultado. Excluye los resultados de las fichas predecesoras, la construcción de capacidades base y el servicio recurrente.  Cada trabajo tiene aceptación, contratista habilitado, fecha y vínculo de solicitud; el cliente confirma ventanas. Ulysses Barreto, Funcional, responde por el resultado; los demás perfiles ejecutan sus aportes y el cliente valida las decisiones de su ámbito.\\\hline
\texttt{1.\allowbreak{}9.\allowbreak{}2.\allowbreak{}5.\allowbreak{}3.\allowbreak{}0.\allowbreak{}1}\newline Informe quincenal del piloto invernal\newline José Lara Arce; 16 HH. & Se entregará informe quincenal del piloto invernal, correspondiente a \texttt{IN2.5.3.0.1}. Comprende registro, validación y emisión del resultado. Excluye los resultados de las fichas predecesoras, la construcción de capacidades base y el servicio recurrente.  El informe conserva planes, solicitudes, trabajos aceptados y pendientes, incidencias y cambios de calendario sin atribuir trabajos no realizados. José Lara Arce, Implantación, responde por el resultado; los demás perfiles ejecutan sus aportes y el cliente valida las decisiones de su ámbito.\\\hline
\texttt{1.\allowbreak{}9.\allowbreak{}2.\allowbreak{}5.\allowbreak{}3.\allowbreak{}0.\allowbreak{}2}\newline Informe quincenal del piloto invernal\newline José Lara Arce; 16 HH. & Se entregará informe quincenal del piloto invernal, correspondiente a \texttt{IN2.5.3.0.2}. Comprende registro, validación y emisión del resultado. Excluye los resultados de las fichas predecesoras, la construcción de capacidades base y el servicio recurrente.  El informe conserva planes, solicitudes, trabajos aceptados y pendientes, incidencias y cambios de calendario sin atribuir trabajos no realizados. José Lara Arce, Implantación, responde por el resultado; los demás perfiles ejecutan sus aportes y el cliente valida las decisiones de su ámbito.\\\hline
\texttt{1.\allowbreak{}9.\allowbreak{}2.\allowbreak{}5.\allowbreak{}3.\allowbreak{}1.\allowbreak{}1}\newline Informe quincenal del piloto invernal\newline José Lara Arce; 16 HH. & Se entregará informe quincenal del piloto invernal, correspondiente a \texttt{IN2.5.3.1.1}. Excluye los resultados de las fichas predecesoras, la construcción de capacidades base y el servicio recurrente.  El informe conserva planes, solicitudes, trabajos aceptados y pendientes, incidencias y cambios de calendario sin atribuir trabajos no realizados. José Lara Arce, Implantación, responde por el resultado; los demás perfiles ejecutan sus aportes y el cliente valida las decisiones de su ámbito.\\\hline
\texttt{1.\allowbreak{}9.\allowbreak{}2.\allowbreak{}5.\allowbreak{}3.\allowbreak{}1.\allowbreak{}2}\newline Informe quincenal del piloto invernal\newline José Lara Arce; 16 HH. & Se entregará informe quincenal del piloto invernal, correspondiente a \texttt{IN2.5.3.1.2}. Excluye los resultados de las fichas predecesoras, la construcción de capacidades base y el servicio recurrente.  El informe conserva planes, solicitudes, trabajos aceptados y pendientes, incidencias y cambios de calendario sin atribuir trabajos no realizados. José Lara Arce, Implantación, responde por el resultado; los demás perfiles ejecutan sus aportes y el cliente valida las decisiones de su ámbito.\\\hline
\texttt{1.\allowbreak{}9.\allowbreak{}2.\allowbreak{}5.\allowbreak{}3.\allowbreak{}2.\allowbreak{}1}\newline Informe quincenal del piloto invernal\newline José Lara Arce; 16 HH. & Se entregará informe quincenal del piloto invernal, correspondiente a \texttt{IN2.5.3.2.1}. Comprende registro, validación y emisión del resultado. Excluye los resultados de las fichas predecesoras, la construcción de capacidades base y el servicio recurrente.  El informe conserva planes, solicitudes, trabajos aceptados y pendientes, incidencias y cambios de calendario sin atribuir trabajos no realizados. José Lara Arce, Implantación, responde por el resultado; los demás perfiles ejecutan sus aportes y el cliente valida las decisiones de su ámbito.\\\hline
\texttt{1.\allowbreak{}9.\allowbreak{}2.\allowbreak{}5.\allowbreak{}3.\allowbreak{}2.\allowbreak{}2}\newline Informe quincenal del piloto invernal\newline José Lara Arce; 16 HH. & Se entregará informe quincenal del piloto invernal, correspondiente a \texttt{IN2.5.3.2.2}. Comprende registro, validación y emisión del resultado. Excluye los resultados de las fichas predecesoras, la construcción de capacidades base y el servicio recurrente.  El informe conserva planes, solicitudes, trabajos aceptados y pendientes, incidencias y cambios de calendario sin atribuir trabajos no realizados. José Lara Arce, Implantación, responde por el resultado; los demás perfiles ejecutan sus aportes y el cliente valida las decisiones de su ámbito.\\\hline
\texttt{1.\allowbreak{}9.\allowbreak{}2.\allowbreak{}5.\allowbreak{}5.\allowbreak{}3.\allowbreak{}1}\newline Registro quincenal de resultados de botaduras\newline Davor Serey; 16 HH. & Se entregará registro quincenal de resultados de botaduras, correspondiente a \texttt{IN2.5.5.3.1}. Comprende registro, validación y emisión del resultado. Excluye los resultados de las fichas predecesoras, la construcción de capacidades base y el servicio recurrente.  Cada botadura observada se relaciona con el plan y su revisión previa; se registran resultados, ausencias y reprogramaciones. Davor Serey, Calidad, responde por el resultado; los demás perfiles ejecutan sus aportes y el cliente valida las decisiones de su ámbito.\\\hline
\texttt{1.\allowbreak{}9.\allowbreak{}2.\allowbreak{}5.\allowbreak{}5.\allowbreak{}3.\allowbreak{}2}\newline Registro quincenal de resultados de botaduras\newline Davor Serey; 16 HH. & Se entregará registro quincenal de resultados de botaduras, correspondiente a \texttt{IN2.5.5.3.2}. Comprende registro, validación y emisión del resultado. Excluye los resultados de las fichas predecesoras, la construcción de capacidades base y el servicio recurrente.  Cada botadura observada se relaciona con el plan y su revisión previa; se registran resultados, ausencias y reprogramaciones. Davor Serey, Calidad, responde por el resultado; los demás perfiles ejecutan sus aportes y el cliente valida las decisiones de su ámbito.\\\hline
\texttt{1.\allowbreak{}9.\allowbreak{}2.\allowbreak{}5.\allowbreak{}5.\allowbreak{}4.\allowbreak{}1}\newline Registro quincenal de resultados de botaduras\newline Davor Serey; 16 HH. & Se entregará registro quincenal de resultados de botaduras, correspondiente a \texttt{IN2.5.5.4.1}. Comprende preparación, ejecución de comprobaciones y expediente. Excluye los resultados de las fichas predecesoras, la construcción de capacidades base y el servicio recurrente.  Cada botadura observada se relaciona con el plan y su revisión previa; se registran resultados, ausencias y reprogramaciones. Davor Serey, Calidad, responde por el resultado; los demás perfiles ejecutan sus aportes y el cliente valida las decisiones de su ámbito.\\\hline
\texttt{1.\allowbreak{}9.\allowbreak{}2.\allowbreak{}5.\allowbreak{}5.\allowbreak{}4.\allowbreak{}2}\newline Registro quincenal de resultados de botaduras\newline Davor Serey; 16 HH. & Se entregará registro quincenal de resultados de botaduras, correspondiente a \texttt{IN2.5.5.4.2}. Comprende preparación, ejecución de comprobaciones y expediente. Excluye los resultados de las fichas predecesoras, la construcción de capacidades base y el servicio recurrente.  Cada botadura observada se relaciona con el plan y su revisión previa; se registran resultados, ausencias y reprogramaciones. Davor Serey, Calidad, responde por el resultado; los demás perfiles ejecutan sus aportes y el cliente valida las decisiones de su ámbito.\\\hline
\texttt{1.\allowbreak{}9.\allowbreak{}2.\allowbreak{}5.\allowbreak{}5.\allowbreak{}5.\allowbreak{}1}\newline Registro quincenal de resultados de botaduras\newline Davor Serey; 16 HH. & Se entregará registro quincenal de resultados de botaduras, correspondiente a \texttt{IN2.5.5.5.1}. Comprende registro, validación y emisión del resultado. Excluye los resultados de las fichas predecesoras, la construcción de capacidades base y el servicio recurrente.  Cada botadura observada se relaciona con el plan y su revisión previa; se registran resultados, ausencias y reprogramaciones. Davor Serey, Calidad, responde por el resultado; los demás perfiles ejecutan sus aportes y el cliente valida las decisiones de su ámbito.\\\hline
\texttt{1.\allowbreak{}9.\allowbreak{}2.\allowbreak{}5.\allowbreak{}5.\allowbreak{}5.\allowbreak{}2}\newline Registro quincenal de resultados de botaduras\newline Davor Serey; 16 HH. & Se entregará registro quincenal de resultados de botaduras, correspondiente a \texttt{IN2.5.5.5.2}. Comprende registro, validación y emisión del resultado. Excluye los resultados de las fichas predecesoras, la construcción de capacidades base y el servicio recurrente.  Cada botadura observada se relaciona con el plan y su revisión previa; se registran resultados, ausencias y reprogramaciones. Davor Serey, Calidad, responde por el resultado; los demás perfiles ejecutan sus aportes y el cliente valida las decisiones de su ámbito.\\\hline
\texttt{1.\allowbreak{}9.\allowbreak{}2.\allowbreak{}5.\allowbreak{}6}\newline Informe comparativo de resultados y acta de evaluación\newline Davor Serey; 32 HH. & Se entregará informe comparativo de resultados y acta de evaluación, correspondiente a \texttt{IN2.5.6}. Comprende registro, validación y emisión del resultado. Excluye los resultados de las fichas predecesoras, la construcción de capacidades base y el servicio recurrente.  Se comparan planes, cumplimiento y botaduras con casos pendientes incluidos; acta identifica límites y decisión. Davor Serey, Calidad, responde por el resultado; los demás perfiles ejecutan sus aportes y el cliente valida las decisiones de su ámbito.\\\hline
\texttt{1.\allowbreak{}9.\allowbreak{}2.\allowbreak{}6.\allowbreak{}1}\newline Procedimiento, documentación y transferencia a operación\newline José Lara Arce; 32 HH. & Se entregará procedimiento, documentación y transferencia a operación, correspondiente a \texttt{IN2.6.1}. Comprende registro, validación y emisión del resultado. Excluye los resultados de las fichas predecesoras, la construcción de capacidades base y el servicio recurrente.  Quedan versión del procedimiento, catálogo, responsables de actualización y decisión de transferencia. José Lara Arce, Operación/SRE, responde por el resultado; los demás perfiles ejecutan sus aportes y el cliente valida las decisiones de su ámbito.\\\hline
\texttt{1.\allowbreak{}9.\allowbreak{}3.\allowbreak{}1.\allowbreak{}1}\newline Inventario de formatos, volumen y campos requeridos\newline Nicolás Torfan; 32 HH. & Se entregará inventario de formatos, volumen y campos requeridos, correspondiente a \texttt{IN3.1.1}. Excluye los resultados de las fichas predecesoras, la construcción de capacidades base y el servicio recurrente.  Se identifican formatos, campos obligatorios, volumen observado y casos ilegibles sin asumir una muestra representativa. Nicolás Torfan, Datos, responde por el resultado; los demás perfiles ejecutan sus aportes y el cliente valida las decisiones de su ámbito.\\\hline
\texttt{1.\allowbreak{}9.\allowbreak{}3.\allowbreak{}1.\allowbreak{}2}\newline Interfaces y controles de IA especificados\newline Fernando Guerrero; 40 HH. & Se entregará interfaces y controles de ia especificados, correspondiente a \texttt{IN3.1.2}. Excluye los resultados de las fichas predecesoras, la construcción de capacidades base y el servicio recurrente.  Se define procesamiento, revisión humana, finalidades, errores, retención y frontera con SD4 y requisitos de IA. Fernando Guerrero, Arquitectura/Integración, responde por el resultado; los demás perfiles ejecutan sus aportes y el cliente valida las decisiones de su ámbito.\\\hline
\texttt{1.\allowbreak{}9.\allowbreak{}3.\allowbreak{}2.\allowbreak{}1}\newline Adaptador del extractor y procesamiento asíncrono\newline Nikolai Navea; 64 HH. & Se entregará adaptador del extractor y procesamiento asíncrono, correspondiente a \texttt{IN3.2.1}. Comprende implementación, pruebas unitarias y documentación del resultado. Excluye los resultados de las fichas predecesoras, la construcción de capacidades base y el servicio recurrente.  Cada trabajo identifica versión, documento y respuesta; resultados tardíos y reintentos no alteran datos aceptados. Nikolai Navea, Desarrollo, responde por el resultado; los demás perfiles ejecutan sus aportes y el cliente valida las decisiones de su ámbito.\\\hline
\texttt{1.\allowbreak{}9.\allowbreak{}3.\allowbreak{}2.\allowbreak{}2}\newline Reglas de campos, discrepancias y datos pendientes\newline Nicolás Torfan; 48 HH. & Se entregará reglas de campos, discrepancias y datos pendientes, correspondiente a \texttt{IN3.2.2}. Comprende implementación, pruebas unitarias y documentación del resultado. Excluye los resultados de las fichas predecesoras, la construcción de capacidades base y el servicio recurrente.  Campos faltantes o contradictorios quedan pendientes de revisión; no se completan como hechos del documento. Nicolás Torfan, Datos, responde por el resultado; los demás perfiles ejecutan sus aportes y el cliente valida las decisiones de su ámbito.\\\hline
\texttt{1.\allowbreak{}9.\allowbreak{}3.\allowbreak{}3.\allowbreak{}1}\newline Formulario de revisión y aprobación humana\newline Nikolai Navea; 56 HH. & Se entregará formulario de revisión y aprobación humana, correspondiente a \texttt{IN3.3.1}. Comprende implementación, pruebas unitarias y documentación del resultado. Excluye los resultados de las fichas predecesoras, la construcción de capacidades base y el servicio recurrente.  El revisor compara fuente y campos y confirma o corrige antes de registrar; cada decisión queda atribuida. Nikolai Navea, Desarrollo, responde por el resultado; los demás perfiles ejecutan sus aportes y el cliente valida las decisiones de su ámbito.\\\hline
\texttt{1.\allowbreak{}9.\allowbreak{}3.\allowbreak{}3.\allowbreak{}2}\newline Integración con expediente, NAV y auditoría\newline Fernando Guerrero; 56 HH. & Se entregará integración con expediente, nav y auditoría, correspondiente a \texttt{IN3.3.2}. Comprende implementación, pruebas unitarias y documentación del resultado. Excluye los resultados de las fichas predecesoras, la construcción de capacidades base y el servicio recurrente.  El registro aprobado conserva documento, actor y versión, y accede a NAV por su contrato autorizado. Fernando Guerrero, Arquitectura/Integración, responde por el resultado; los demás perfiles ejecutan sus aportes y el cliente valida las decisiones de su ámbito.\\\hline
\texttt{1.\allowbreak{}9.\allowbreak{}3.\allowbreak{}3.\allowbreak{}3}\newline Ingreso manual, reintentos y resultados tardíos\newline Nikolai Navea; 40 HH. & Se entregará ingreso manual, reintentos y resultados tardíos, correspondiente a \texttt{IN3.3.3}. Comprende implementación, pruebas unitarias y documentación del resultado. Excluye los resultados de las fichas predecesoras, la construcción de capacidades base y el servicio recurrente.  El ingreso manual mantiene los mismos controles y un resultado tardío no sobrescribe una decisión humana. Nikolai Navea, Desarrollo, responde por el resultado; los demás perfiles ejecutan sus aportes y el cliente valida las decisiones de su ámbito.\\\hline
\texttt{1.\allowbreak{}9.\allowbreak{}3.\allowbreak{}4.\allowbreak{}1}\newline Pruebas de extracción, integración y continuidad\newline Davor Serey; 40 HH. & Se entregará pruebas de extracción, integración y continuidad, correspondiente a \texttt{IN3.4.1}. Comprende preparación, ejecución de comprobaciones y expediente. Excluye los resultados de las fichas predecesoras, la construcción de capacidades base y el servicio recurrente.  Se comprueban legibilidad, discrepancias, duplicados, fallos del servicio, ingreso manual y recuperación. Davor Serey, Calidad, responde por el resultado; los demás perfiles ejecutan sus aportes y el cliente valida las decisiones de su ámbito.\\\hline
\texttt{1.\allowbreak{}9.\allowbreak{}3.\allowbreak{}4.\allowbreak{}2}\newline Matriz y evidencias de controles de IA\newline Bruno Toro; 32 HH. & Se entregará matriz y evidencias de controles de ia, correspondiente a \texttt{IN3.4.2}. Comprende preparación, ejecución de comprobaciones y expediente. Excluye los resultados de las fichas predecesoras, la construcción de capacidades base y el servicio recurrente.  Cada control aplicable se vincula con diseño y evidencia; se verifican aislamiento, revisión y trazabilidad sin declarar cumplimiento por usar IA. Bruno Toro, Seguridad, responde por el resultado; los demás perfiles ejecutan sus aportes y el cliente valida las decisiones de su ámbito.\\\hline
\texttt{1.\allowbreak{}9.\allowbreak{}3.\allowbreak{}5.\allowbreak{}1}\newline Muestra autorizada, protocolo y capacitación\newline Nicolás Torfan; 32 HH. & Se entregará muestra autorizada, protocolo y capacitación, correspondiente a \texttt{IN3.5.1}. Comprende registro, validación y emisión del resultado. Excluye los resultados de las fichas predecesoras, la construcción de capacidades base y el servicio recurrente.  Muestra, autorizaciones, método manual/asistido y evaluación de usuarios están documentados antes de medir. Nicolás Torfan, Datos, responde por el resultado; los demás perfiles ejecutan sus aportes y el cliente valida las decisiones de su ámbito.\\\hline
\texttt{1.\allowbreak{}9.\allowbreak{}3.\allowbreak{}5.\allowbreak{}2.\allowbreak{}19.\allowbreak{}1}\newline Registro quincenal de ejecución manual y asistida\newline José Lara Arce; 16 HH. & Se entregará registro quincenal de ejecución manual y asistida, correspondiente a \texttt{IN3.5.2.19.1}. Comprende registro, validación y emisión del resultado. Excluye los resultados de las fichas predecesoras, la construcción de capacidades base y el servicio recurrente.  El registro conserva documentos autorizados, tiempos, errores, correcciones y costo de procesamiento; incluye casos no completados. José Lara Arce, Implantación, responde por el resultado; los demás perfiles ejecutan sus aportes y el cliente valida las decisiones de su ámbito.\\\hline
\texttt{1.\allowbreak{}9.\allowbreak{}3.\allowbreak{}5.\allowbreak{}2.\allowbreak{}19.\allowbreak{}2}\newline Registro quincenal de ejecución manual y asistida\newline José Lara Arce; 16 HH. & Se entregará registro quincenal de ejecución manual y asistida, correspondiente a \texttt{IN3.5.2.19.2}. Comprende registro, validación y emisión del resultado. Excluye los resultados de las fichas predecesoras, la construcción de capacidades base y el servicio recurrente.  El registro conserva documentos autorizados, tiempos, errores, correcciones y costo de procesamiento; incluye casos no completados. José Lara Arce, Implantación, responde por el resultado; los demás perfiles ejecutan sus aportes y el cliente valida las decisiones de su ámbito.\\\hline
\texttt{1.\allowbreak{}9.\allowbreak{}3.\allowbreak{}5.\allowbreak{}2.\allowbreak{}20.\allowbreak{}1}\newline Registro quincenal de ejecución manual y asistida\newline José Lara Arce; 16 HH. & Se entregará registro quincenal de ejecución manual y asistida, correspondiente a \texttt{IN3.5.2.20.1}. Comprende registro, validación y emisión del resultado. Excluye los resultados de las fichas predecesoras, la construcción de capacidades base y el servicio recurrente.  El registro conserva documentos autorizados, tiempos, errores, correcciones y costo de procesamiento; incluye casos no completados. José Lara Arce, Implantación, responde por el resultado; los demás perfiles ejecutan sus aportes y el cliente valida las decisiones de su ámbito.\\\hline
\texttt{1.\allowbreak{}9.\allowbreak{}3.\allowbreak{}5.\allowbreak{}2.\allowbreak{}20.\allowbreak{}2}\newline Registro quincenal de ejecución manual y asistida\newline José Lara Arce; 16 HH. & Se entregará registro quincenal de ejecución manual y asistida, correspondiente a \texttt{IN3.5.2.20.2}. Comprende registro, validación y emisión del resultado. Excluye los resultados de las fichas predecesoras, la construcción de capacidades base y el servicio recurrente.  El registro conserva documentos autorizados, tiempos, errores, correcciones y costo de procesamiento; incluye casos no completados. José Lara Arce, Implantación, responde por el resultado; los demás perfiles ejecutan sus aportes y el cliente valida las decisiones de su ámbito.\\\hline
\texttt{1.\allowbreak{}9.\allowbreak{}3.\allowbreak{}5.\allowbreak{}3}\newline Informe de tiempos, errores, costos y beneficio\newline Davor Serey; 32 HH. & Se entregará informe de tiempos, errores, costos y beneficio, correspondiente a \texttt{IN3.5.3}. Comprende registro, validación y emisión del resultado. Excluye los resultados de las fichas predecesoras, la construcción de capacidades base y el servicio recurrente.  La comparación conserva errores, rechazos, revisión humana y costos; evidencia insuficiente no se transforma en mejora demostrada. Davor Serey, Calidad, responde por el resultado; los demás perfiles ejecutan sus aportes y el cliente valida las decisiones de su ámbito.\\\hline
\texttt{1.\allowbreak{}9.\allowbreak{}3.\allowbreak{}6.\allowbreak{}1}\newline Configuración, manuales y transferencia\newline José Lara Arce; 32 HH. & Se entregará configuración, manuales y transferencia, correspondiente a \texttt{IN3.6.1}. Comprende registro, validación y emisión del resultado. Excluye los resultados de las fichas predecesoras, la construcción de capacidades base y el servicio recurrente.  Se reciben configuración, límites, ingreso manual, monitoreo de consumo y procedimiento de incidentes. José Lara Arce, Operación/SRE, responde por el resultado; los demás perfiles ejecutan sus aportes y el cliente valida las decisiones de su ámbito.\\\hline
\texttt{1.\allowbreak{}9.\allowbreak{}4.\allowbreak{}1.\allowbreak{}1}\newline Ficha de medición y atribución del resultado\newline Ulysses Barreto; 24 HH. & Se entregará ficha de medición y atribución del resultado, correspondiente a \texttt{IN4.1.1}. Excluye los resultados de las fichas predecesoras, la construcción de capacidades base y el servicio recurrente.  Se define tiempo activo, tareas comparables, calidad y exclusiones; las esperas y ahorros de otras innovaciones no se duplican. Ulysses Barreto, Funcional, responde por el resultado; los demás perfiles ejecutan sus aportes y el cliente valida las decisiones de su ámbito.\\\hline
\texttt{1.\allowbreak{}9.\allowbreak{}4.\allowbreak{}1.\allowbreak{}2}\newline Regla comercial y escenarios de incentivo\newline Ignacio Silva; 32 HH. & Se entregará regla comercial y escenarios de incentivo, correspondiente a \texttt{IN4.1.2}. Excluye los resultados de las fichas predecesoras, la construcción de capacidades base y el servicio recurrente.  La fórmula, tope, rechazo y factor cero están documentados; no modifica las obligaciones ni los hitos E-25. Ignacio Silva, Proyecto, responde por el resultado; los demás perfiles ejecutan sus aportes y el cliente valida las decisiones de su ámbito.\\\hline
\texttt{1.\allowbreak{}9.\allowbreak{}4.\allowbreak{}2.\allowbreak{}1}\newline Interfaces y expediente de medición especificados\newline Fernando Guerrero; 32 HH. & Se entregará interfaces y expediente de medición especificados, correspondiente a \texttt{IN4.2.1}. Comprende registro, validación y emisión del resultado. Excluye los resultados de las fichas predecesoras, la construcción de capacidades base y el servicio recurrente.  Se especifican captura de intervalos, solicitudes, autoridad, versiones y separación entre preparación y aprobación. Fernando Guerrero, Arquitectura/Integración, responde por el resultado; los demás perfiles ejecutan sus aportes y el cliente valida las decisiones de su ámbito.\\\hline
\texttt{1.\allowbreak{}9.\allowbreak{}4.\allowbreak{}2.\allowbreak{}2}\newline Captura de intervalos y validaciones implementadas\newline Nikolai Navea; 48 HH. & Se entregará captura de intervalos y validaciones implementadas, correspondiente a \texttt{IN4.2.2}. Comprende implementación, pruebas unitarias y documentación del resultado. Excluye los resultados de las fichas predecesoras, la construcción de capacidades base y el servicio recurrente.  Cada intervalo tiene solicitud, persona, inicio, término y origen; se detectan superposiciones y datos incompletos. Nikolai Navea, Desarrollo, responde por el resultado; los demás perfiles ejecutan sus aportes y el cliente valida las decisiones de su ámbito.\\\hline
\texttt{1.\allowbreak{}9.\allowbreak{}4.\allowbreak{}2.\allowbreak{}3}\newline Cálculo versionado y expediente de conciliación\newline Nikolai Navea; 56 HH. & Se entregará cálculo versionado y expediente de conciliación, correspondiente a \texttt{IN4.2.3}. Comprende implementación, pruebas unitarias y documentación del resultado. Excluye los resultados de las fichas predecesoras, la construcción de capacidades base y el servicio recurrente.  El cálculo es reproducible, conserva sus entradas y versiones y no emite cobro automático. Nikolai Navea, Desarrollo, responde por el resultado; los demás perfiles ejecutan sus aportes y el cliente valida las decisiones de su ámbito.\\\hline
\texttt{1.\allowbreak{}9.\allowbreak{}4.\allowbreak{}2.\allowbreak{}4}\newline Evidencia de seguridad y cálculo del incentivo\newline Davor Serey; 40 HH. & Se entregará evidencia de seguridad y cálculo del incentivo, correspondiente a \texttt{IN4.2.4}. Comprende registro, validación y emisión del resultado. Excluye los resultados de las fichas predecesoras, la construcción de capacidades base y el servicio recurrente.  Los casos sin evidencia, base cero, calidad fallida y muestra insuficiente producen factor cero; no hay cambios silenciosos. Davor Serey, Calidad, responde por el resultado; los demás perfiles ejecutan sus aportes y el cliente valida las decisiones de su ámbito.\\\hline
\texttt{1.\allowbreak{}9.\allowbreak{}4.\allowbreak{}3.\allowbreak{}1.\allowbreak{}1}\newline Registro base del tiempo activo\newline Davor Serey; 16 HH. & Se entregará registro base del tiempo activo, correspondiente a \texttt{IN4.3.1.1}. Excluye los resultados de las fichas predecesoras, la construcción de capacidades base y el servicio recurrente.  Cada intervalo identifica persona, solicitud y tarea; esperas y tiempos superpuestos quedan excluidos sin duplicar la medición de IN1. Davor Serey, Calidad, responde por el resultado; los demás perfiles ejecutan sus aportes y el cliente valida las decisiones de su ámbito.\\\hline
\texttt{1.\allowbreak{}9.\allowbreak{}4.\allowbreak{}3.\allowbreak{}1.\allowbreak{}2}\newline Registro base del tiempo activo\newline Davor Serey; 16 HH. & Se entregará registro base del tiempo activo, correspondiente a \texttt{IN4.3.1.2}. Excluye los resultados de las fichas predecesoras, la construcción de capacidades base y el servicio recurrente.  Cada intervalo identifica persona, solicitud y tarea; esperas y tiempos superpuestos quedan excluidos sin duplicar la medición de IN1. Davor Serey, Calidad, responde por el resultado; los demás perfiles ejecutan sus aportes y el cliente valida las decisiones de su ámbito.\\\hline
\texttt{1.\allowbreak{}9.\allowbreak{}4.\allowbreak{}3.\allowbreak{}2.\allowbreak{}1}\newline Registro de evaluación del tiempo activo\newline Davor Serey; 16 HH. & Se entregará registro de evaluación del tiempo activo, correspondiente a \texttt{IN4.3.2.1}. Comprende registro, validación y emisión del resultado. Excluye los resultados de las fichas predecesoras, la construcción de capacidades base y el servicio recurrente.  Cada intervalo identifica persona, solicitud y tarea; esperas y tiempos superpuestos quedan excluidos sin duplicar la medición de IN1. Davor Serey, Calidad, responde por el resultado; los demás perfiles ejecutan sus aportes y el cliente valida las decisiones de su ámbito.\\\hline
\texttt{1.\allowbreak{}9.\allowbreak{}4.\allowbreak{}3.\allowbreak{}2.\allowbreak{}2}\newline Registro de evaluación del tiempo activo\newline Davor Serey; 16 HH. & Se entregará registro de evaluación del tiempo activo, correspondiente a \texttt{IN4.3.2.2}. Comprende registro, validación y emisión del resultado. Excluye los resultados de las fichas predecesoras, la construcción de capacidades base y el servicio recurrente.  Cada intervalo identifica persona, solicitud y tarea; esperas y tiempos superpuestos quedan excluidos sin duplicar la medición de IN1. Davor Serey, Calidad, responde por el resultado; los demás perfiles ejecutan sus aportes y el cliente valida las decisiones de su ámbito.\\\hline
\texttt{1.\allowbreak{}9.\allowbreak{}4.\allowbreak{}3.\allowbreak{}3}\newline Evaluación del piloto sin efectos de pago\newline Davor Serey; 32 HH. & Se entregará evaluación del piloto sin efectos de pago, correspondiente a \texttt{IN4.3.3}. Comprende registro, validación y emisión del resultado. Excluye los resultados de las fichas predecesoras, la construcción de capacidades base y el servicio recurrente.  Se compara tiempo activo y calidad con la regla aprobada y se conserva decisión; el piloto no genera pagos variables. Davor Serey, Calidad, responde por el resultado; los demás perfiles ejecutan sus aportes y el cliente valida las decisiones de su ámbito.\\\hline
\texttt{1.\allowbreak{}9.\allowbreak{}4.\allowbreak{}4.\allowbreak{}1}\newline Procedimiento de conciliación contractual habilitado\newline Ignacio Silva; 24 HH. & Se entregará procedimiento de conciliación contractual habilitado, correspondiente a \texttt{IN4.4.1}. Comprende registro, validación y emisión del resultado. Excluye los resultados de las fichas predecesoras, la construcción de capacidades base y el servicio recurrente.  Cliente y Synaptix disponen de procedimiento, evidencia, reclamación y decisión registrada antes de habilitar el incentivo. Ignacio Silva, Proyecto, responde por el resultado; los demás perfiles ejecutan sus aportes y el cliente valida las decisiones de su ámbito.\\\hline
\texttt{1.\allowbreak{}9.\allowbreak{}5.\allowbreak{}1.\allowbreak{}1}\newline Plan ambiental y protocolo de comparación\newline Ulysses Barreto; 24 HH. & Se entregará plan ambiental y protocolo de comparación, correspondiente a \texttt{IN5.1.1}. Excluye los resultados de las fichas predecesoras, la construcción de capacidades base y el servicio recurrente.  Se fijan indicador de consumo por permanencia, ventanas de 28 días, criterios de comparabilidad y evidencia de acciones. Ulysses Barreto, Funcional, responde por el resultado; los demás perfiles ejecutan sus aportes y el cliente valida las decisiones de su ámbito.\\\hline
\texttt{1.\allowbreak{}9.\allowbreak{}5.\allowbreak{}1.\allowbreak{}2}\newline Consentimiento, modelo y contratos definidos\newline Fernando Guerrero; 40 HH. & Se entregará consentimiento, modelo y contratos definidos, correspondiente a \texttt{IN5.1.2}. Excluye los resultados de las fichas predecesoras, la construcción de capacidades base y el servicio recurrente.  Adhesión, plan, acción y evaluación tienen estados, conservación y permisos; se especifica revocación y autoridad de lecturas. Fernando Guerrero, Arquitectura/Integración, responde por el resultado; los demás perfiles ejecutan sus aportes y el cliente valida las decisiones de su ámbito.\\\hline
\texttt{1.\allowbreak{}9.\allowbreak{}5.\allowbreak{}2.\allowbreak{}1}\newline Historial validado y cálculo por permanencia\newline Nikolai Navea; 56 HH. & Se entregará historial validado y cálculo por permanencia, correspondiente a \texttt{IN5.2.1}. Comprende implementación, pruebas unitarias y documentación del resultado. Excluye los resultados de las fichas predecesoras, la construcción de capacidades base y el servicio recurrente.  El historial usa intervalos conciliados, conserva originales y no suma lecturas acumulativas ni períodos repetidos. Nikolai Navea, Desarrollo, responde por el resultado; los demás perfiles ejecutan sus aportes y el cliente valida las decisiones de su ámbito.\\\hline
\texttt{1.\allowbreak{}9.\allowbreak{}5.\allowbreak{}2.\allowbreak{}2}\newline Adhesión, acciones y portal ambiental\newline Nikolai Navea; 56 HH. & Se entregará adhesión, acciones y portal ambiental, correspondiente a \texttt{IN5.2.2}. Comprende implementación, pruebas unitarias y documentación del resultado. Excluye los resultados de las fichas predecesoras, la construcción de capacidades base y el servicio recurrente.  El armador acepta el plan y consulta sólo sus datos; cada acción tiene responsable, fechas y evidencia. Nikolai Navea, Desarrollo, responde por el resultado; los demás perfiles ejecutan sus aportes y el cliente valida las decisiones de su ámbito.\\\hline
\texttt{1.\allowbreak{}9.\allowbreak{}5.\allowbreak{}2.\allowbreak{}3}\newline Integración, consentimiento y actualización de cortes\newline Fernando Guerrero; 48 HH. & Se entregará integración, consentimiento y actualización de cortes, correspondiente a \texttt{IN5.2.3}. Comprende implementación, pruebas unitarias y documentación del resultado. Excluye los resultados de las fichas predecesoras, la construcción de capacidades base y el servicio recurrente.  Revocación y correcciones propagan límites de uso; la vista identifica cobertura y último corte validado. Fernando Guerrero, Arquitectura/Integración, responde por el resultado; los demás perfiles ejecutan sus aportes y el cliente valida las decisiones de su ámbito.\\\hline
\texttt{1.\allowbreak{}9.\allowbreak{}5.\allowbreak{}2.\allowbreak{}4}\newline Pruebas de medición, privacidad y accesibilidad\newline Davor Serey; 40 HH. & Se entregará pruebas de medición, privacidad y accesibilidad, correspondiente a \texttt{IN5.2.4}. Comprende registro, validación y emisión del resultado. Excluye los resultados de las fichas predecesoras, la construcción de capacidades base y el servicio recurrente.  Se prueban acceso entre flotas, exportaciones, revocación, diferencias de lectura, permanencia y vistas accesibles. Davor Serey, Calidad, responde por el resultado; los demás perfiles ejecutan sus aportes y el cliente valida las decisiones de su ámbito.\\\hline
\texttt{1.\allowbreak{}9.\allowbreak{}5.\allowbreak{}3.\allowbreak{}1}\newline Participantes y acciones preparados para el piloto\newline José Lara Arce; 32 HH. & Se entregará participantes y acciones preparados para el piloto, correspondiente a \texttt{IN5.3.1}. Comprende registro, validación y emisión del resultado. Excluye los resultados de las fichas predecesoras, la construcción de capacidades base y el servicio recurrente.  Hay adhesión voluntaria, medición validada, permanencias conciliadas y usuarios preparados; no se presume participación efectiva. José Lara Arce, Implantación, responde por el resultado; los demás perfiles ejecutan sus aportes y el cliente valida las decisiones de su ámbito.\\\hline
\texttt{1.\allowbreak{}9.\allowbreak{}5.\allowbreak{}3.\allowbreak{}2.\allowbreak{}1}\newline Registro base del consumo normalizado\newline Davor Serey; 16 HH. & Se entregará registro base del consumo normalizado, correspondiente a \texttt{IN5.3.2.1}. Comprende registro, validación y emisión del resultado. Excluye los resultados de las fichas predecesoras, la construcción de capacidades base y el servicio recurrente.  Se conservan intervalos de consumo conciliados y permanencia de la cohorte; no se afirma reducción con datos faltantes o acumulativos duplicados. Davor Serey, Calidad, responde por el resultado; los demás perfiles ejecutan sus aportes y el cliente valida las decisiones de su ámbito.\\\hline
\texttt{1.\allowbreak{}9.\allowbreak{}5.\allowbreak{}3.\allowbreak{}2.\allowbreak{}2}\newline Registro base del consumo normalizado\newline Davor Serey; 16 HH. & Se entregará registro base del consumo normalizado, correspondiente a \texttt{IN5.3.2.2}. Comprende registro, validación y emisión del resultado. Excluye los resultados de las fichas predecesoras, la construcción de capacidades base y el servicio recurrente.  Se conservan intervalos de consumo conciliados y permanencia de la cohorte; no se afirma reducción con datos faltantes o acumulativos duplicados. Davor Serey, Calidad, responde por el resultado; los demás perfiles ejecutan sus aportes y el cliente valida las decisiones de su ámbito.\\\hline
\texttt{1.\allowbreak{}9.\allowbreak{}5.\allowbreak{}3.\allowbreak{}3.\allowbreak{}1}\newline Registro de evaluación del consumo normalizado\newline Davor Serey; 16 HH. & Se entregará registro de evaluación del consumo normalizado, correspondiente a \texttt{IN5.3.3.1}. Comprende registro, validación y emisión del resultado. Excluye los resultados de las fichas predecesoras, la construcción de capacidades base y el servicio recurrente.  Se conservan intervalos de consumo conciliados y permanencia de la cohorte; no se afirma reducción con datos faltantes o acumulativos duplicados. Davor Serey, Calidad, responde por el resultado; los demás perfiles ejecutan sus aportes y el cliente valida las decisiones de su ámbito.\\\hline
\texttt{1.\allowbreak{}9.\allowbreak{}5.\allowbreak{}3.\allowbreak{}3.\allowbreak{}2}\newline Registro de evaluación del consumo normalizado\newline Davor Serey; 16 HH. & Se entregará registro de evaluación del consumo normalizado, correspondiente a \texttt{IN5.3.3.2}. Comprende registro, validación y emisión del resultado. Excluye los resultados de las fichas predecesoras, la construcción de capacidades base y el servicio recurrente.  Se conservan intervalos de consumo conciliados y permanencia de la cohorte; no se afirma reducción con datos faltantes o acumulativos duplicados. Davor Serey, Calidad, responde por el resultado; los demás perfiles ejecutan sus aportes y el cliente valida las decisiones de su ámbito.\\\hline
\texttt{1.\allowbreak{}9.\allowbreak{}5.\allowbreak{}3.\allowbreak{}4}\newline Evaluación ambiental y acta del piloto\newline Davor Serey; 32 HH. & Se entregará evaluación ambiental y acta del piloto, correspondiente a \texttt{IN5.3.4}. Comprende registro, validación y emisión del resultado. Excluye los resultados de las fichas predecesoras, la construcción de capacidades base y el servicio recurrente.  Se comparan ventanas equivalentes y acciones evidenciadas; faltantes o muestra insuficiente impiden afirmar reducción. Davor Serey, Calidad, responde por el resultado; los demás perfiles ejecutan sus aportes y el cliente valida las decisiones de su ámbito.\\\hline
\texttt{1.\allowbreak{}9.\allowbreak{}5.\allowbreak{}4.\allowbreak{}1}\newline Pasaporte y seguimiento transferidos a operación\newline José Lara Arce; 32 HH. & Se entregará pasaporte y seguimiento transferidos a operación, correspondiente a \texttt{IN5.4.1}. Comprende registro, validación y emisión del resultado. Excluye los resultados de las fichas predecesoras, la construcción de capacidades base y el servicio recurrente.  Se reciben configuración, consentimiento, monitoreo de calidad de lecturas y responsables del informe previo a temporada 2029. José Lara Arce, Operación/SRE, responde por el resultado; los demás perfiles ejecutan sus aportes y el cliente valida las decisiones de su ámbito.\\\hline
\end{longtable}
\endgroup
\FuenteTabla{Elaboración propia. Las horas y los códigos corresponden al catálogo técnico integrado.}


\section{Gobierno, formación, operación y salida}
Estas cuentas reúnen dirección, verificación integral, formación, servicio y transferencia. Las fichas comunes definen las obligaciones y las instancias periódicas conservan su programación individual.
\subsection{Gobierno, verificación, formación y servicio}
Las cuentas restantes distinguen el control del proyecto 1.1, las pruebas integrales 1.10.2, la formación 1.11.2/3, el cierre de marcha blanca 1.11.5 y el servicio/salida 1.12.3--11. Se instancian 317 paquetes fijos, con 7.504 HH: 2.528 de implementación y 4.976 de operación. Las fichas comunes de trabajos periódicos se aplican a cada código/ventana del registro \path{componentes/paquetes_gobierno_servicio.csv}; no se suma otra vez el componente padre.

El expediente quincenal de gobierno incluye actas, tablero, acuerdos, riesgos, cambios y valor ganado con SPI/CPI. La reversibilidad se presenta dentro de los primeros 90 días, se actualiza y se prepara para una transferencia acompañada al menos 90 días reales. La salida se programa en 53--56 para reservar más de tres meses y absorber la duración real; el calendario verificará fechas efectivas. La transferencia contractual no se confunde con reconstruir por completo una plataforma alternativa estimada en SD4.

Las pruebas integrales incluyen QA/cobertura, regresión, carga/estrés, operación desconectada, resiliencia, conmutación real de DR, seguridad ofensiva, accesibilidad y UAT/T-17. Antes de cada paso a producción, incluidas liberaciones de Operación, se comprobarán carga/estrés a 1,5 veces el peak declarado, WCAG 2.2 AA documentada, ofensiva independiente sin críticos/altos abiertos, resiliencia y DR real conforme a T-10/PRUEBAS-6. Los dos ensayos por etapa permanecen en 1.10.1 y no se repiten en sus costos. En operación, DR/resiliencia se repiten además al menos semestralmente y ofensiva al menos anualmente, con fechas relativas sujetas a ventanas permitidas y sin exceder la frecuencia comprometida. Proyecto verifica esfuerzo y revisión disponibles antes de seleccionar una liberación; si las bolsas previstas no bastan, Synaptix aporta refuerzo a su costo sin reducir las pruebas \parencite[cap.~20, secc.~20.1, pp.~34--35]{bases-transversales}.

\subsection{Sesiones de capacitación por perfil, sitio y cohorte}
Se distinguen cinco perfiles: usuarios finales, avanzados, administradores, responsables técnicos designados por el cliente y soporte niveles 1/2. Synaptix no presume un área TI interna. Los materiales incluyen fuentes editables en español, guías, FAQ, guion/video y base de conocimiento; su disponibilidad en autoformación registra avance. Se combinarán sesiones presenciales por sitio, sincrónicas y acompañamiento de tarea. Los perfiles administradores/técnicos deberán evaluarse y certificarse antes de cada cierre \parencite[RT-22.01--07]{bases-transversales}.

\paragraph{1.11.7.e.p.s.b: sesión y evaluación de cohorte}
Por etapa $e$, perfil $p$, sitio/turno $s$ y lote $b$ de hasta doce asistentes, el entregable incluye preparación, dos sesiones de dos horas, ejercicio y registro individual. José Lara Arce responde por el paquete: Implantación 8, Funcional 2 y Calidad 2 HH, total 12. La aceptación requiere ejercicio reproducido, asistencia, resultado y recuperación asignada; asistencia sin evaluación no acredita certificación. Se requieren material recibido, entorno seguro, usuarios autorizados y ventana acordada. Excluye participación del cliente, traslado, formación incremental de innovación y sesiones de recuperación adicionales no realizadas.

El número de paquetes será $B_e=\sum_{p,s}\lceil N_{e,p,s}/12\rceil$, sobre personas y turnos confirmados, sin asumir que el mismo material reduce automáticamente el número de sesiones. Se agregan $12(B_1+B_2)$ HH de implementación. En operación se usará \texttt{1.12.10.m.p.s.b}, con mes $m$, para incorporación de personal nuevo o recuperación requerida, con la misma ficha y costo incluidos en servicio. Las dos jornadas anuales fijas no sustituyen esa incorporación. La coordinación de personal estacional y monitores nuevos debe completar formación durante noviembre y antes del 15 de diciembre, sin capacitación dentro del peak protegido. Cero participantes sólo se registrará para una cohorte confirmada sin personas.

\subsection{Presencia por ola y estabilización}
\paragraph{1.11.4.e.w.d.z: acompañamiento diario en terreno}
Cada etapa $e$ y ola $w$ genera paquetes por día $d$ y zona $z$. Se propone un puesto de acompañamiento de ocho horas por zona durante los primeros catorce días y cuatro horas por zona en los siguientes catorce, en ventanas seguras, con disminución condicionada a adopción. El paquete diario tendrá 8 HH: las primeras dos semanas, ocho de presencia; las siguientes, cuatro de presencia y cuatro de preparación, seguimiento y evidencia. Excluye reuniones de control, capacitación formal y conciliación de datos ya estimadas. El referente de coordinación será José Lara Arce; la ejecución requerirá personal de implantación de refuerzo con competencias y turnos acreditados, sin atribuir esas jornadas al titular además de sus otros paquetes.

Se recibirá registro de zona, horario, tareas acompañadas, incidencias, pendientes y entrega al relevo. Debe existir dotación por turno, transporte y suplencia; se cubrirán fines de semana aplicables. El traslado se planifica por separado de las horas de acompañamiento. Para dos zonas por ola y 28 días son 56 paquetes y 448 HH, con 336 horas de presencia y 112 de preparación/seguimiento. Si una ola requiere otras zonas o más cobertura, se instanciarán sus paquetes adicionales.

La estabilización de cuatro semanas después de cada paso a producción empleará \texttt{1.11.6.e.d.z}, dos puestos de ocho horas diarios, uno en pantalanes y otro en varadero, durante 28 días consecutivos: 56 paquetes/448 HH por etapa, incluidos fines de semana. Se recibirá evidencia de atención, relevo y resolución. E1 pertenece a implementación; E2 al servicio de operación del mes 21. La dotación en esos puestos requiere reemplazos y horarios autorizados; dos puestos no significan que dos personas puedan trabajar 28 días sin descanso. La presencia queda separada de las 96 HH de estabilización de datos por etapa, que realizan tareas técnicas distintas.

\subsection{Cobertura externa y especialistas}
La solución reservará cobertura NOC y SOC continua y centro de atención bilingüe 06:00--24:00 en temporada alta y 08:00--20:00 el resto, con emergencia y planes vencidos 24$\times$7. Se propone un puesto activo de NOC, uno de SOC y al menos uno de atención por turno; capacidad de cola, competencias, supervisión y relevo deberán demostrar los SLA. No se presentan proveedores, ubicaciones ni contratos como seleccionados o firmados \parencite[RT-21.01/06]{bases-transversales}\parencite[RT-21.06]{caso07}.

Para un mes con $D_m$ días, cada puesto continuo requiere $24D_m$ horas de cobertura; atención requiere $18D_{H,m}+12D_{B,m}$, donde temporada alta es 15 de noviembre--31 de marzo. Los puestos continuo de NOC/SOC deben recibir alertas y escalar críticos/altos fuera del horario de atención, sin espera hasta apertura. Con 144 horas de cobertura efectiva por equivalente de jornada como hipótesis de dimensionamiento, las cotas son $\lceil24D_m/144\rceil$ por puesto continuo y $\lceil(18D_{H,m}+12D_{B,m})/144\rceil$ para atención. Ese supuesto se contrastará con turnos legales, descansos, ausencias y reemplazos del prestador; no acredita dotación contratada.

Las horas de cobertura se presupuestan como capacidad externa reservada y no se atribuyen a los ocho titulares de SD1. Su precio incluye recepción y tareas de primer nivel dentro de la cobertura contratada; esas horas no se vuelven a sumar como atención individual. El gerente dedicado y especialistas atienden coordinación, diagnóstico y resolución especializada dentro de sus paquetes y disponibilidad. Volumen de incidencias, tiempos de servicio, concurrencia y reemplazo definirán puestos adicionales para sostener 80\,\% antes de veinte segundos, resolución al primer contacto de 70\,\% y abandono máximo 5\,\%; un solo puesto no demuestra esos resultados.

\section{Diccionario de gobierno, pruebas, formación y servicio}
Los paquetes siguientes identifican resultados terminales entre 8 y 80 HH y ventanas quincenales de ejecución. Sus estimaciones iniciales suponen procedimientos, herramientas y versiones disponibles. Los períodos de revisión/subsanación del artículo 18 se programan fuera de las horas de trabajo; las fechas son entradas para integrar la red, no pruebas de cabida. Licencias, servicios externos, traslados y reservas tienen una imputación única en la planificación.

Las pruebas específicas de construcción e innovación y los ensayos 1.10.1 no se duplican. 1.10.2 verifica conjuntamente la solución y consolida recepción. Las correcciones de código y reejecuciones necesarias se imputarán a sus paquetes, con recursos dimensionados; una ventana de corrección no aporta horas de Desarrollo implícitas. La ejecución de una prueba de larga duración puede exigir turnos adicionales de observación: sus horas de cobertura no se presumen dentro de un informe.

La recepción contractual corresponde a la contraparte autorizada. Todas las fichas requieren insumos y versiones recibidos, registros y entornos autorizados, herramientas y contraparte programada. Las HH excluyen dedicación del cliente y cobertura externa continua. La falta de insumos o evidencia bloquea la recepción y exige reprogramar dentro de las condiciones contractuales. Estas condiciones comunes se aplican además de los criterios específicos de cada ficha.

Las fichas de la Tabla~\ref{tab:dic-compacto-gobierno_servicio} presentan resultado, responsable, esfuerzo y aceptación. Las condiciones particulares, recursos y ventanas por código se conservan en \path{componentes/condiciones_fichas_t14.csv}, junto con la asignación y precedencias del diccionario técnico integrado. Las actividades de T-15 desarrollan los aportes por perfil y el T-18 desarrolla la ejecución de implantación.


\begingroup\small\setlength{\tabcolsep}{3pt}
\begin{longtable}{L{52mm}L{91mm}}
\caption{Diccionario de gobierno, formación y servicio}\label{tab:dic-compacto-gobierno_servicio}\\\hline
 \EncabezadoTabla{Paquete, resultado y responsable} & \EncabezadoTabla{Contenido y criterio de aceptación}\\\hline\endfirsthead
\hline \EncabezadoTabla{Paquete, resultado y responsable} & \EncabezadoTabla{Contenido y criterio de aceptación}\\\hline\endhead
\texttt{1.\allowbreak{}1.\allowbreak{}1.\allowbreak{}1}\newline Acta de arranque, interlocutores y decisiones\newline Ignacio Silva; 32 HH. & Responsabilidades, reuniones y habilitaciones; excluye ejecución de obra y construcción. Interlocutores, autoridad, calendario de decisión y pendientes quedan registrados. Ignacio Silva responde por el paquete.\\\hline
\texttt{1.\allowbreak{}1.\allowbreak{}1.\allowbreak{}2}\newline Línea base de alcance y trazabilidad\newline Ignacio Silva; 40 HH. & Vinculación alcance-requisitos-entregables y control de cambios; excluye pruebas ejecutadas. Se conservan IDs/enunciados y decisiones de cobertura; aprobación H1 se documenta sin convertir pendientes en conformidad. Ignacio Silva responde por el paquete.\\\hline
\texttt{1.\allowbreak{}1.\allowbreak{}1.\allowbreak{}3}\newline Definición de terminado y recepción\newline Davor Serey; 32 HH. & Criterios código, pruebas, documentación, seguridad y despliegue por entregable. Criterios medibles y protocolo T-17 tienen responsable y tratamiento de observaciones. Davor Serey responde por el paquete.\\\hline
\texttt{1.\allowbreak{}1.\allowbreak{}2.\allowbreak{}1}\newline Registro de riesgos y respuestas presupuestadas\newline Ignacio Silva; 32 HH. & Riesgos, disparadores y respuestas con cuenta de costo; excluye sumar otra vez acciones ya estimadas. Cada respuesta tiene responsable, recurso, disparador, cuenta y efecto de calendario. Ignacio Silva responde por el paquete.\\\hline
\texttt{1.\allowbreak{}1.\allowbreak{}2.\allowbreak{}2}\newline Línea base integrada de plazo, costo y recursos\newline Ignacio Silva; 40 HH. & Consolidación EDT, red, cargas, calendario y presupuesto; excluye aceptar sobrecargas como disponibilidad. Cobertura, recursos, holguras y reservas están demostrados; si fallan, se registra rechazo de línea base y plan de corrección. Ignacio Silva responde por el paquete.\\\hline
\texttt{1.\allowbreak{}1.\allowbreak{}3.\allowbreak{}1}\newline Plan inicial de reversibilidad\newline Ignacio Silva; 48 HH. & Inventario, formatos, custodia, acompañamiento y sustitución; excluye ejecutar traslado a proveedor alternativo. Plan se presenta dentro de 90 días de inicio e identifica trabajo portable/no portable y transferencia sin cobro adicional. Ignacio Silva responde por el paquete.\\\hline
\texttt{1.\allowbreak{}12.\allowbreak{}1.\allowbreak{}1}\newline Modelo operativo y contrato de cobertura preparados\newline José Lara Arce; 48 HH. & Cobertura, proveedores, escalamiento y suplencia; excluye afirmar contrato externo firmado. NOC, SOC, mesa y especialistas tienen ubicación, dotación por turno, capacidad, condiciones y responsables verificables antes de servicio. José Lara Arce responde por el paquete.\\\hline
\texttt{1.\allowbreak{}12.\allowbreak{}1.\allowbreak{}2}\newline Catálogo, runbooks y traspaso al servicio\newline José Lara Arce; 56 HH. & Procedimientos de continuidad, atención, mantenimiento y acceso; excluye nuevas capacidades de software. Soporte reproduce procedimiento, identifica autoridad y evidencia, y recibe fuentes, accesos y pendientes. José Lara Arce responde por el paquete.\\\hline
\texttt{1.\allowbreak{}1.\allowbreak{}3.\allowbreak{}2}\newline Expediente de custodia independiente de fuentes\newline Ignacio Silva; 24 HH. & Preparación de contrato, fuentes y condiciones de liberación con tercero independiente; excluye atribuir custodia ya contratada. Custodio y mecanismo tienen contrato, ubicación, acceso, actualización semestral y liberación por insolvencia/incumplimiento verificables. Ignacio Silva responde por el paquete.\\\hline
\texttt{1.\allowbreak{}1.\allowbreak{}3.\allowbreak{}3.\allowbreak{}12}\newline Actualización semestral de custodia de fuentes\newline Ignacio Silva; 16 HH. & Depósito de fuentes y manifiesto ante custodio independiente; excluye reconstruir fuentes no versionadas. Custodio confirma versión, integridad y condiciones de acceso/liberación y se conserva comprobante. Ignacio Silva responde por el paquete.\\\hline
\texttt{1.\allowbreak{}1.\allowbreak{}3.\allowbreak{}4}\newline Plan de reversibilidad actualizado durante implementación\newline Ignacio Silva; 32 HH. & Actualiza dependencias, inventario y estrategia de transferencia; excluye ejecutar traslado completo. La actualización anual conserva versión, custodios, formatos y recursos requeridos. Ignacio Silva responde por el paquete.\\\hline
\texttt{1.\allowbreak{}1.\allowbreak{}4.\allowbreak{}1.\allowbreak{}1}\newline Expediente quincenal de control\newline Ignacio Silva; 16 HH. & Comités, acuerdos, riesgos, cambios, avance físico y costo, tablero y seguimiento; excluye ejecución de paquetes técnicos. Se reportan valor ganado, SPI/CPI, acuerdos y desviaciones; actas en dos días hábiles y desvío mayor a 10\% comunicado en cinco días con recuperación. Ignacio Silva responde por el paquete.\\\hline
\texttt{1.\allowbreak{}10.\allowbreak{}2.\allowbreak{}1.\allowbreak{}1}\newline Protocolo integral y datos de prueba E1\newline Davor Serey; 32 HH. & Casos, requisitos, datos, versiones y criterios; excluye unidades ya incluidas en construcción. Cada caso y criterio CA tiene escenario, umbral, autoridad y evidencia prevista. Davor Serey responde por el paquete.\\\hline
\texttt{1.\allowbreak{}10.\allowbreak{}2.\allowbreak{}1.\allowbreak{}2}\newline Puerta de QA, cobertura y contratos E1\newline Davor Serey; 40 HH. & Consolidación de cobertura y flujos QA; excluye construir unidades y contratos otra vez. Lógica de negocio alcanza 70\% con puerta bloqueante y sin pruebas fallidas; integraciones de la etapa ejecutadas sin error. Davor Serey responde por el paquete.\\\hline
\texttt{1.\allowbreak{}10.\allowbreak{}2.\allowbreak{}1.\allowbreak{}3}\newline Regresión de sistema y convivencia E1\newline Davor Serey; 40 HH. & Batería completa por versión y protección de E1; excluye regresión específica ya ejecutada de innovación. No hay regresión abierta y los flujos/canales conservan permisos, hechos y autoridad. Davor Serey responde por el paquete.\\\hline
\texttt{1.\allowbreak{}10.\allowbreak{}2.\allowbreak{}1.\allowbreak{}4}\newline Carga y estrés con mezcla operacional E1\newline Davor Serey; 48 HH. & Preparación y ejecución de carga/estrés, curvas y recuperación; excluye fijar límites por promedios simples. Umbrales se cumplen a 1.5 veces peak declarado y se registra saturación y degradación sin pérdida silenciosa. Davor Serey responde por el paquete.\\\hline
\texttt{1.\allowbreak{}10.\allowbreak{}2.\allowbreak{}1.\allowbreak{}5}\newline Continuidad local y terminales desconectados E1\newline Davor Serey; 56 HH. & Prueba conjunta de autoridad y reconciliación; excluye medir sólo la cola vacía. Recinto opera 24 horas, terminales 12 e inspección 8 con mezcla y fallos previstos, sin pérdida ni duplicados. Davor Serey responde por el paquete.\\\hline
\texttt{1.\allowbreak{}10.\allowbreak{}2.\allowbreak{}1.\allowbreak{}6}\newline Resiliencia de componentes y recuperación E1\newline José Lara Arce; 48 HH. & Falla controlada de enlace, nodo y dependencia, con registro y recuperación. Se conserva operación/degradación controlada y recuperación, con tiempos, transacciones y brechas documentados. José Lara Arce responde por el paquete.\\\hline
\texttt{1.\allowbreak{}10.\allowbreak{}2.\allowbreak{}1.\allowbreak{}7}\newline Conmutación real de desastre y retorno E1\newline José Lara Arce; 64 HH. & Conmutación regional real, verificación de datos/identidad y retorno; excluye usar existencia de réplica como evidencia. RTO/RPO comprometidos se alcanzan y evidencia identifica último dato, tiempos, permisos y conciliación. José Lara Arce responde por el paquete.\\\hline
\texttt{1.\allowbreak{}10.\allowbreak{}2.\allowbreak{}1.\allowbreak{}8}\newline Prueba ofensiva y cierre de hallazgos E1\newline Bruno Toro; 56 HH. & Ejecución ofensiva de superficie completa y comprobación de resolución; correcciones de código se imputan a construcción. No quedan hallazgos críticos ni altos abiertos y cada resultado corresponde a versión recibida. Bruno Toro responde por el paquete.\\\hline
\texttt{1.\allowbreak{}10.\allowbreak{}2.\allowbreak{}1.\allowbreak{}9}\newline Verificación de accesibilidad WCAG 2.2 AA E1\newline Davor Serey; 40 HH. & Auditoría automática y manual de tareas/canales; excluye acreditar AA sólo por biblioteca usada. Criterios aplicables AA están verificados, barreras corregidas y evidencia de teclado, foco, mensajes y ayudas registrada. Davor Serey responde por el paquete.\\\hline
\texttt{1.\allowbreak{}10.\allowbreak{}2.\allowbreak{}1.\allowbreak{}10}\newline Aceptación de usuario y expediente T-17 E1\newline Davor Serey; 48 HH. & Ejecución conjunta de escenarios CA y expediente, sin repetir ensayos de migración 1.10.1. Contraparte decide y firma aceptación; observaciones resueltas, versiones y pruebas trazadas antes de certificación. Davor Serey responde por el paquete.\\\hline
\texttt{1.\allowbreak{}11.\allowbreak{}2.\allowbreak{}1.\allowbreak{}1}\newline Materiales y ejercicios perfil 1 E1\newline José Lara Arce; 32 HH. & Manual, guía, FAQ, guion/video y ejercicios editables de su perfil; excluye impartir todas las sesiones y materiales incrementales de innovación. Cliente recibe fuentes en español, ejemplos, contingencia y ejercicios; autoformación registra avance y material se vincula a versión. José Lara Arce responde por el paquete.\\\hline
\texttt{1.\allowbreak{}11.\allowbreak{}2.\allowbreak{}1.\allowbreak{}2}\newline Materiales y ejercicios perfil 2 E1\newline José Lara Arce; 32 HH. & Manual, guía, FAQ, guion/video y ejercicios editables de su perfil; excluye impartir todas las sesiones y materiales incrementales de innovación. Cliente recibe fuentes en español, ejemplos, contingencia y ejercicios; autoformación registra avance y material se vincula a versión. José Lara Arce responde por el paquete.\\\hline
\texttt{1.\allowbreak{}11.\allowbreak{}2.\allowbreak{}1.\allowbreak{}3}\newline Materiales y ejercicios perfil 3 E1\newline José Lara Arce; 32 HH. & Manual, guía, FAQ, guion/video y ejercicios editables de su perfil; excluye impartir todas las sesiones y materiales incrementales de innovación. Cliente recibe fuentes en español, ejemplos, contingencia y ejercicios; autoformación registra avance y material se vincula a versión. José Lara Arce responde por el paquete.\\\hline
\texttt{1.\allowbreak{}11.\allowbreak{}2.\allowbreak{}1.\allowbreak{}4}\newline Materiales y ejercicios perfil 4 E1\newline José Lara Arce; 32 HH. & Manual, guía, FAQ, guion/video y ejercicios editables de su perfil; excluye impartir todas las sesiones y materiales incrementales de innovación. Cliente recibe fuentes en español, ejemplos, contingencia y ejercicios; autoformación registra avance y material se vincula a versión. José Lara Arce responde por el paquete.\\\hline
\texttt{1.\allowbreak{}11.\allowbreak{}2.\allowbreak{}1.\allowbreak{}5}\newline Materiales y ejercicios perfil 5 E1\newline José Lara Arce; 32 HH. & Manual, guía, FAQ, guion/video y ejercicios editables de su perfil; excluye impartir todas las sesiones y materiales incrementales de innovación. Cliente recibe fuentes en español, ejemplos, contingencia y ejercicios; autoformación registra avance y material se vincula a versión. José Lara Arce responde por el paquete.\\\hline
\texttt{1.\allowbreak{}11.\allowbreak{}3.\allowbreak{}1.\allowbreak{}1}\newline Plan de formación y turnos E1\newline José Lara Arce; 32 HH. & Cohortes, sitio, turnos, validadores, calendario y evaluaciones; excluye horas de participación de empleados del cliente. Cada perfil/canal tiene cohorte y sitio; calendario evita congelamiento/peak y prepara estacionales en noviembre antes del 15 de diciembre. José Lara Arce responde por el paquete.\\\hline
\texttt{1.\allowbreak{}11.\allowbreak{}3.\allowbreak{}1.\allowbreak{}2}\newline Expediente de evaluación y certificación E1\newline José Lara Arce; 32 HH. & Resultados, brechas y decisión para administradores/equipo técnico y soporte; excluye impartir sesiones de recuperación no dimensionadas. Roles administradores/técnicos designados son evaluados y certificados antes del cierre; ausencia de área TI no elimina los responsables que recibirán transferencia. José Lara Arce responde por el paquete.\\\hline
\texttt{1.\allowbreak{}11.\allowbreak{}5.\allowbreak{}1.\allowbreak{}1}\newline Expediente de cierre de marcha blanca E1\newline Ignacio Silva; 40 HH. & Consolidación de indicadores y acta sin volver a estimar informes de datos o pruebas. Últimas cuatro semanas tienen volumen real comprometido, disponibilidad/tiempos, conciliación, formación y ningún crítico/alto abierto; contraparte firma. Ignacio Silva responde por el paquete.\\\hline
\texttt{1.\allowbreak{}10.\allowbreak{}2.\allowbreak{}2.\allowbreak{}1}\newline Protocolo integral y datos de prueba E2\newline Davor Serey; 32 HH. & Casos, requisitos, datos, versiones y criterios; excluye unidades ya incluidas en construcción. Cada caso y criterio CA tiene escenario, umbral, autoridad y evidencia prevista. Davor Serey responde por el paquete.\\\hline
\texttt{1.\allowbreak{}10.\allowbreak{}2.\allowbreak{}2.\allowbreak{}2}\newline Puerta de QA, cobertura y contratos E2\newline Davor Serey; 40 HH. & Consolidación de cobertura y flujos QA; excluye construir unidades y contratos otra vez. Lógica de negocio alcanza 70\% con puerta bloqueante y sin pruebas fallidas; integraciones de la etapa ejecutadas sin error. Davor Serey responde por el paquete.\\\hline
\texttt{1.\allowbreak{}10.\allowbreak{}2.\allowbreak{}2.\allowbreak{}3}\newline Regresión de sistema y convivencia E2\newline Davor Serey; 40 HH. & Batería completa por versión y protección de E1; excluye regresión específica ya ejecutada de innovación. No hay regresión abierta y los flujos/canales conservan permisos, hechos y autoridad. Davor Serey responde por el paquete.\\\hline
\texttt{1.\allowbreak{}10.\allowbreak{}2.\allowbreak{}2.\allowbreak{}4}\newline Carga y estrés con mezcla operacional E2\newline Davor Serey; 48 HH. & Preparación y ejecución de carga/estrés, curvas y recuperación; excluye fijar límites por promedios simples. Umbrales se cumplen a 1.5 veces peak declarado y se registra saturación y degradación sin pérdida silenciosa. Davor Serey responde por el paquete.\\\hline
\texttt{1.\allowbreak{}10.\allowbreak{}2.\allowbreak{}2.\allowbreak{}5}\newline Continuidad local y terminales desconectados E2\newline Davor Serey; 56 HH. & Prueba conjunta de autoridad y reconciliación; excluye medir sólo la cola vacía. Recinto opera 24 horas, terminales 12 e inspección 8 con mezcla y fallos previstos, sin pérdida ni duplicados. Davor Serey responde por el paquete.\\\hline
\texttt{1.\allowbreak{}10.\allowbreak{}2.\allowbreak{}2.\allowbreak{}6}\newline Resiliencia de componentes y recuperación E2\newline José Lara Arce; 48 HH. & Falla controlada de enlace, nodo y dependencia, con registro y recuperación. Se conserva operación/degradación controlada y recuperación, con tiempos, transacciones y brechas documentados. José Lara Arce responde por el paquete.\\\hline
\texttt{1.\allowbreak{}10.\allowbreak{}2.\allowbreak{}2.\allowbreak{}7}\newline Conmutación real de desastre y retorno E2\newline José Lara Arce; 64 HH. & Conmutación regional real, verificación de datos/identidad y retorno; excluye usar existencia de réplica como evidencia. RTO/RPO comprometidos se alcanzan y evidencia identifica último dato, tiempos, permisos y conciliación. José Lara Arce responde por el paquete.\\\hline
\texttt{1.\allowbreak{}10.\allowbreak{}2.\allowbreak{}2.\allowbreak{}8}\newline Prueba ofensiva y cierre de hallazgos E2\newline Bruno Toro; 56 HH. & Ejecución ofensiva de superficie completa y comprobación de resolución; correcciones de código se imputan a construcción. No quedan hallazgos críticos ni altos abiertos y cada resultado corresponde a versión recibida. Bruno Toro responde por el paquete.\\\hline
\texttt{1.\allowbreak{}10.\allowbreak{}2.\allowbreak{}2.\allowbreak{}9}\newline Verificación de accesibilidad WCAG 2.2 AA E2\newline Davor Serey; 40 HH. & Auditoría automática y manual de tareas/canales; excluye acreditar AA sólo por biblioteca usada. Criterios aplicables AA están verificados, barreras corregidas y evidencia de teclado, foco, mensajes y ayudas registrada. Davor Serey responde por el paquete.\\\hline
\texttt{1.\allowbreak{}10.\allowbreak{}2.\allowbreak{}2.\allowbreak{}10}\newline Aceptación de usuario y expediente T-17 E2\newline Davor Serey; 48 HH. & Ejecución conjunta de escenarios CA y expediente, sin repetir ensayos de migración 1.10.1. Contraparte decide y firma aceptación; observaciones resueltas, versiones y pruebas trazadas antes de certificación. Davor Serey responde por el paquete.\\\hline
\texttt{1.\allowbreak{}11.\allowbreak{}2.\allowbreak{}2.\allowbreak{}1}\newline Materiales y ejercicios perfil 1 E2\newline José Lara Arce; 32 HH. & Manual, guía, FAQ, guion/video y ejercicios editables de su perfil; excluye impartir todas las sesiones y materiales incrementales de innovación. Cliente recibe fuentes en español, ejemplos, contingencia y ejercicios; autoformación registra avance y material se vincula a versión. José Lara Arce responde por el paquete.\\\hline
\texttt{1.\allowbreak{}11.\allowbreak{}2.\allowbreak{}2.\allowbreak{}2}\newline Materiales y ejercicios perfil 2 E2\newline José Lara Arce; 32 HH. & Manual, guía, FAQ, guion/video y ejercicios editables de su perfil; excluye impartir todas las sesiones y materiales incrementales de innovación. Cliente recibe fuentes en español, ejemplos, contingencia y ejercicios; autoformación registra avance y material se vincula a versión. José Lara Arce responde por el paquete.\\\hline
\texttt{1.\allowbreak{}11.\allowbreak{}2.\allowbreak{}2.\allowbreak{}3}\newline Materiales y ejercicios perfil 3 E2\newline José Lara Arce; 32 HH. & Manual, guía, FAQ, guion/video y ejercicios editables de su perfil; excluye impartir todas las sesiones y materiales incrementales de innovación. Cliente recibe fuentes en español, ejemplos, contingencia y ejercicios; autoformación registra avance y material se vincula a versión. José Lara Arce responde por el paquete.\\\hline
\texttt{1.\allowbreak{}11.\allowbreak{}2.\allowbreak{}2.\allowbreak{}4}\newline Materiales y ejercicios perfil 4 E2\newline José Lara Arce; 32 HH. & Manual, guía, FAQ, guion/video y ejercicios editables de su perfil; excluye impartir todas las sesiones y materiales incrementales de innovación. Cliente recibe fuentes en español, ejemplos, contingencia y ejercicios; autoformación registra avance y material se vincula a versión. José Lara Arce responde por el paquete.\\\hline
\texttt{1.\allowbreak{}11.\allowbreak{}2.\allowbreak{}2.\allowbreak{}5}\newline Materiales y ejercicios perfil 5 E2\newline José Lara Arce; 32 HH. & Manual, guía, FAQ, guion/video y ejercicios editables de su perfil; excluye impartir todas las sesiones y materiales incrementales de innovación. Cliente recibe fuentes en español, ejemplos, contingencia y ejercicios; autoformación registra avance y material se vincula a versión. José Lara Arce responde por el paquete.\\\hline
\texttt{1.\allowbreak{}11.\allowbreak{}3.\allowbreak{}2.\allowbreak{}1}\newline Plan de formación y turnos E2\newline José Lara Arce; 32 HH. & Cohortes, sitio, turnos, validadores, calendario y evaluaciones; excluye horas de participación de empleados del cliente. Cada perfil/canal tiene cohorte y sitio; calendario evita congelamiento/peak y prepara estacionales en noviembre antes del 15 de diciembre. José Lara Arce responde por el paquete.\\\hline
\texttt{1.\allowbreak{}11.\allowbreak{}3.\allowbreak{}2.\allowbreak{}2}\newline Expediente de evaluación y certificación E2\newline José Lara Arce; 32 HH. & Resultados, brechas y decisión para administradores/equipo técnico y soporte; excluye impartir sesiones de recuperación no dimensionadas. Roles administradores/técnicos designados son evaluados y certificados antes del cierre; ausencia de área TI no elimina los responsables que recibirán transferencia. José Lara Arce responde por el paquete.\\\hline
\texttt{1.\allowbreak{}11.\allowbreak{}5.\allowbreak{}2.\allowbreak{}1}\newline Expediente de cierre de marcha blanca E2\newline Ignacio Silva; 40 HH. & Consolidación de indicadores y acta sin volver a estimar informes de datos o pruebas. Últimas cuatro semanas tienen volumen real comprometido, disponibilidad/tiempos, conciliación, formación y ningún crítico/alto abierto; contraparte firma. Ignacio Silva responde por el paquete.\\\hline
\texttt{1.\allowbreak{}12.\allowbreak{}3.\allowbreak{}1.\allowbreak{}21}\newline Control mensual de servicio y SLA\newline José Lara Arce; 16 HH. & SLA, incidentes, costos, compromisos y causas; excluye turnos de recepción continua. Denominadores, tiempos, multas y acciones son trazables; una ausencia de registros no acredita cumplimiento. José Lara Arce responde por el paquete.\\\hline
\texttt{1.\allowbreak{}12.\allowbreak{}3.\allowbreak{}2.\allowbreak{}21}\newline Comprobación de respaldos y restauración\newline José Lara Arce; 16 HH. & Revisión de copias, prueba de muestra y evidencia; excluye conmutación semestral. Restauración tiene integridad, accesos y tiempo registrados y brechas con responsable. José Lara Arce responde por el paquete.\\\hline
\texttt{1.\allowbreak{}12.\allowbreak{}3.\allowbreak{}3.\allowbreak{}21}\newline Revisión de seguridad y acceso\newline Bruno Toro; 16 HH. & Accesos, alertas, agentes y cambios de permisos; excluye guardia SOC y ofensiva anual. Cobertura por activo y cuentas efectivas quedan verificadas y desvíos tratados. Bruno Toro responde por el paquete.\\\hline
\texttt{1.\allowbreak{}12.\allowbreak{}3.\allowbreak{}4.\allowbreak{}21}\newline Mantenimiento, deuda y liberación controlada\newline Nikolai Navea; 24 HH. & Mantenimiento preventivo, revisión y liberación acotada; excluye asumir ilimitadas correcciones/evoluciones en 24 HH. Cambio tiene pruebas de cada producción según PRUEBAS-6, preaviso mínimo de diez días hábiles al CLIENTE si es mantenimiento programado, ventana, versión, depósito actualizado de fuentes/manifiesto en repositorio del CLIENTE y retorno. 20\% de sus horas se reserva a deuda técnica y se cuantifica avance. Nikolai Navea responde por el paquete; Proyecto dimensiona refuerzo propio si los controles exceden su capacidad.\\\hline
\texttt{1.\allowbreak{}12.\allowbreak{}3.\allowbreak{}5.\allowbreak{}21}\newline Proyección de capacidad e infraestructura\newline Fernando Guerrero; 16 HH. & Seguimiento mensual y consolidación trimestral sin segunda imputación. Crecimiento, cuello de botella, alertas y propuesta de costo se vinculan a datos observados. Fernando Guerrero responde por el paquete.\\\hline
\texttt{1.\allowbreak{}12.\allowbreak{}4.\allowbreak{}1.\allowbreak{}6}\newline Resiliencia de componentes y recuperación: control semestral\newline José Lara Arce; 48 HH. & Falla controlada de enlace, nodo y dependencia, con registro y recuperación. Se conserva operación/degradación controlada y recuperación, con tiempos, transacciones y brechas documentados. José Lara Arce responde por el paquete.\\\hline
\texttt{1.\allowbreak{}12.\allowbreak{}4.\allowbreak{}1.\allowbreak{}7}\newline Conmutación real de desastre y retorno: control semestral\newline José Lara Arce; 64 HH. & Conmutación regional real, verificación de datos/identidad y retorno; excluye usar existencia de réplica como evidencia. RTO/RPO comprometidos se alcanzan y evidencia identifica último dato, tiempos, permisos y conciliación. José Lara Arce responde por el paquete.\\\hline
\texttt{1.\allowbreak{}12.\allowbreak{}11.\allowbreak{}1}\newline Actualización semestral de custodia de fuentes\newline Ignacio Silva; 16 HH. & Depósito de fuentes y manifiesto ante custodio independiente; excluye reconstruir fuentes no versionadas. Custodio confirma versión, integridad y condiciones de acceso/liberación y se conserva comprobante. Ignacio Silva responde por el paquete.\\\hline
\texttt{1.\allowbreak{}12.\allowbreak{}5.\allowbreak{}1}\newline Prueba ofensiva y cierre de hallazgos: control anual\newline Bruno Toro; 56 HH. & Ejecución ofensiva de superficie completa y comprobación de resolución; correcciones de código se imputan a construcción. No quedan hallazgos críticos ni altos abiertos y cada resultado corresponde a versión recibida. Bruno Toro responde por el paquete.\\\hline
\texttt{1.\allowbreak{}12.\allowbreak{}6.\allowbreak{}1}\newline Jornada de actualización y transferencia\newline José Lara Arce; 32 HH. & Actualización, ejercicios y evaluación en ventana acordada; excluye incorporación adicional y capacitación en peak protegido. Se conserva asistencia, evaluación, materiales actualizados y brechas; se cumple mínimo de dos jornadas por año de operación. José Lara Arce responde por el paquete.\\\hline
\texttt{1.\allowbreak{}12.\allowbreak{}7.\allowbreak{}1}\newline Plan de reversibilidad actualizado\newline Ignacio Silva; 32 HH. & Actualiza activos, dependencias, formatos, esfuerzo y custodia; excluye ejecutar sustitución total de proveedor. Inventario/credenciales/exportaciones y estrategia de traslado reflejan versión efectiva sin copiar claves no exportables. Ignacio Silva responde por el paquete.\\\hline
\texttt{1.\allowbreak{}12.\allowbreak{}8.\allowbreak{}1}\newline Inventario de salida y responsables\newline Ignacio Silva; 40 HH. & Transferencia de salida incluida; excluye confundir entrega contractual con reconstrucción completa de otra plataforma. Cliente valida activos, interfaces, custodios y destino de transferencia. Ignacio Silva responde por el paquete.\\\hline
\texttt{1.\allowbreak{}12.\allowbreak{}8.\allowbreak{}2}\newline Exportación abierta y manifiestos de datos\newline Nicolás Torfan; 56 HH. & Transferencia de salida incluida; excluye confundir entrega contractual con reconstrucción completa de otra plataforma. Datos, objetos, relaciones, versiones y huellas se exportan y concilian sin depender del software legado. Nicolás Torfan responde por el paquete.\\\hline
\texttt{1.\allowbreak{}12.\allowbreak{}8.\allowbreak{}3}\newline Fuentes, IaC, artefactos y documentación transferidos\newline Fernando Guerrero; 48 HH. & Transferencia de salida incluida; excluye confundir entrega contractual con reconstrucción completa de otra plataforma. Repositorio del cliente contiene fuentes, builds, configuración y procedimientos reproducibles; escrow se verifica. Fernando Guerrero responde por el paquete.\\\hline
\texttt{1.\allowbreak{}12.\allowbreak{}8.\allowbreak{}4}\newline Custodia, credenciales y recifrado autorizado\newline Bruno Toro; 56 HH. & Transferencia de salida incluida; excluye confundir entrega contractual con reconstrucción completa de otra plataforma. Claves no exportables se sustituyen por procedimiento autorizado y acceso queda transferido sin divulgar secretos. Bruno Toro responde por el paquete.\\\hline
\texttt{1.\allowbreak{}12.\allowbreak{}8.\allowbreak{}5}\newline Aceptación de transferencia y eliminación autorizada\newline Ignacio Silva; 48 HH. & Transferencia de salida incluida; excluye confundir entrega contractual con reconstrucción completa de otra plataforma. Receptor reproduce operación, pendientes tienen dueño y sólo se elimina origen tras autorización y evidencia. Ignacio Silva responde por el paquete.\\\hline
\texttt{1.\allowbreak{}12.\allowbreak{}9.\allowbreak{}53.\allowbreak{}1}\newline Acompañamiento quincenal de transición\newline José Lara Arce; 24 HH. & Soporte de transición, reproducción y atención al receptor; excluye guardias continuas y doble imputación de exportación. Consultas, ejecución observada, brechas y responsables se registran; acompañamiento abarca al menos 90 días reales sin cobro adicional. José Lara Arce responde por el paquete.\\\hline
\end{longtable}
\endgroup
\FuenteTabla{Elaboración propia. Las horas y los códigos corresponden al catálogo técnico integrado.}

\subsection{Ficha de mentoría técnica mensual 1.12.13}
Las seis instancias 1.12.13.21--26 entregan mentoría a referentes técnicos y suplentes, con agenda/casos, documentación, acta y verificación independiente. Responsable: José Lara Arce. Cada instancia: L16, I4 y Q4, 24 HH; total 144 HH. Recursos MEN-IL-01 y MEN-Q-01 propios, independientes de turnos; referentes del CLIENTE autorizados. Excluye mesa, estabilización, salida y formación nueva/intervención productiva durante congelamiento. La ficha y calendario del componente de mentoría desarrollan recepción, continuidad y suplencia. Referencias: RT-22.07--08 y BA89--90.


\section{Marco contractual de programación}
La Tabla~\ref{tab:t14-hitos} fija los intervalos obligatorios para desarrollar la red de actividades y la carta Gantt. Estos intervalos son restricciones contractuales; sus duraciones no provienen de una estimación de paquetes ni constituyen por sí solas un cálculo de ruta crítica.

\begin{table}[htbp]
\centering
\begin{tabularx}{\linewidth}{L{23mm}Y}
\hline
\EncabezadoTabla{Meses} & \EncabezadoTabla{Fase o hito contractual}\\\hline
1--12 & Desarrollo de Etapa 1 y certificación de su alcance.\\\hline
13--15 & Marcha blanca de Etapa 1 en convivencia con la operación vigente.\\\hline
16 & Paso a producción de Etapa 1.\\\hline
13--18 & Desarrollo de Etapa 2, sin degradar Etapa 1.\\\hline
19--20 & Marcha blanca de Etapa 2 con Etapa 1 en producción.\\\hline
21 & Paso a producción de Etapa 2 y aceptación final de implementación.\\\hline
21--56 & Operación de ambos alcances durante 36 meses.\\\hline
\end{tabularx}
\caption{Restricciones contractuales de programación}
\label{tab:t14-hitos}
\FuenteTabla{Bases Administrativas, artículo 17.1.}
\end{table}

La ejecución de las actividades se contrastará con la asignación por recurso y quincena del T-15. Sus puestos de refuerzo permiten dimensionar capacidad; el reparto de HH entre ejecutores exige comprobar precedencias, duración, coordinación y calendario antes de incorporarlo a la línea base.

\subsection{Método de estimación, secuenciamiento y control}
Synaptix utiliza estimación ascendente por resultado: el diccionario fija contenido, exclusiones, recepción, responsable último, HH por perfil, recursos y supuestos; sus actividades asignan ejecutor, dedicación, lugar, calendario y predecesoras. La descomposición se detiene en resultados estimables, asignables y verificables. Las familias repetitivas comparten su ficha, mientras cada período o lote conserva código, responsable, esfuerzo y programación; el rango observado de 8--80 HH no es una obligación normativa. Un código identifica pertenencia a la EDT y no una fecha: se conservan los códigos históricos aunque cambie la ventana de ejecución. La etapa y el mes son atributos separados.

La ingeniería separa trabajo técnico y consolidación de Calidad; una aportación grande puede repartirse en subtareas entre personas competentes, sin sumar otra vez las HH. Lógica, persistencia e interfaz tienen integración asignada en la ficha. La consolidación de Calidad recibe los aportes previos; la validación final de marcha blanca permite en paralelo la verificación de datos por Calidad y la lista contractual por Proyecto, porque son resultados independientes. Las actividades y contribuciones individuales se identifican en \path{actividades_integradas.csv}.

La programación pasa por cinco controles: cobertura requisito--entregable--paquete; precedencias completas sin ciclos y puertas con responsable; CPM de avance/retroceso con holgura total y libre; escenario PERT de tres duraciones; y nivelación con calendario y capacidad. No se suma un porcentaje general de gestión o pruebas a trabajos ya estimados. Los servicios por turnos, las reservas y las esperas se muestran por separado. PERT no equivale a una probabilidad de entrega mientras no se sustenten distribuciones e independencia.

El calendario usa un escenario de inicio el 1 de diciembre de 2026 y feriados chilenos. La fecha contractual, los decretos de feriados, los catorce eventos náuticos y las campañas efectivas se incorporarán al control de cambios antes de aprobar la línea base. El método de este capítulo constituye la especificación que deberá recoger el capítulo metodológico de la oferta; no se atribuye aprobación o concordancia documental aún no realizada.


El esfuerzo de los núcleos se selecciona mediante UCP y se distribuye entre paquetes identificados, conservando una sola imputación del trabajo original. Las asignaciones y curvas corresponden a este alcance conciliado. La fecha de inicio sigue siendo una hipótesis explícita de calendario; la aprobación contractual de la línea base y la revisión humana técnica son controles futuros separados de estos cálculos.

\subsection{Contraste funcional independiente de los núcleos}
El modelo de estimación delimita los treinta componentes de los núcleos Operacional, Servicios y Administrativo. Un caso representa un objetivo observable del actor; las transacciones son intercambios completos de solicitud y respuesta. Las validaciones internas, clics y sentencias de persistencia no reciben puntos propios. Salida y regreso, preparación y participación de clase, consulta de salud y retorno, espera y cambio contractual se cuentan por objetivos distintos. Las alternativas de rechazo, duplicación o desconexión pertenecen al mismo objetivo. El modelo conserva actor, flujo, fuente y correspondencia con los noventa paquetes iniciales en \path{componentes/ucp_casos_nucleos.csv}.

La frontera incluye reglas y vistas propias de 1.4--1.6. Excluye construcción de identidad, canales, adaptadores y sincronización comunes, infraestructura, migración histórica, pilotos, formación y operación; estas partidas conservan su imputación ascendente. Los sistemas contables y dispositivos quedan fuera de la frontera de este conteo: sus adaptadores se estiman en 1.7. Los roles humanos se cuentan una sola vez aunque utilicen varios canales. El modelo es una descomposición de diseño de la oferta, no un levantamiento validado ni productividad observada.

La clasificación de casos y actores y los ajustes técnicos y ambientales siguen el método UCP descrito por Cohn; sus factores se aplican al modelo particular de la marina.~\parencite{cohn-ucp} Se aplica el método de puntos de casos de uso: actores simples/medios/complejos con pesos 1/2/3; casos con hasta tres, cuatro a siete y más de siete transacciones con pesos 5/10/15. $UUCP=UAW+UUCW$, $TCF=0,6+0,01\sum w_iT_i$ y $EF=1,4-0,03\sum w_iE_i$. Los trece factores técnicos y ocho de ambiente tienen valoración y justificación individual en \path{componentes/ucp_factores_nucleos.csv}. Las capacidades de ambiente son condiciones de provisión propuestas; no acreditan experiencia de personas ni contratos ya ejecutados.

El conteo contiene 34 casos y 12 roles complejos: $UAW=36$, $UUCW=180$ y $UUCP=216$. Se obtiene $TCF=1,165$, $EF=1,01$ y $UCP=254,16$. Con 1 factor de ambiente desfavorable se usa $CF=20$ HH/punto como referencia no calibrada.

Se declara la interpretación de esfuerzo de programación: $E_P=UCP\,CF$, y el esfuerzo total de referencia es $E_T=E_P/0,40$. Como supuesto de planificación de Synaptix, la distribución de referencia es análisis 10\,\%, diseño 20\,\%, programación 40\,\%, pruebas 15\,\% y sobrecarga 15\,\%. El reparto 10/20/40/15/15 y la lectura de programación son decisiones explícitas de esta estimación; no se atribuyen a una productividad histórica de Synaptix ni a una distribución universal del método UCP. Es una comparación independiente; no se agrega ese total al de la EDT ni se vuelven a sumar gestión y pruebas ya imputadas. Para evitar duplicar actores entre etapas, su peso se distribuye proporcionalmente al peso funcional de E1 y E2; los dos subtotales conservan exactamente el conteo global.

La Tabla~\ref{tab:sd7-ucp-contraste} presenta contraste de esfuerzo directo y referencia UCP.
\begin{table}[H]
\centering
\begin{tabularx}{\linewidth}{Y R{29mm} R{29mm}}\hline
 \EncabezadoTabla{Magnitud} & \EncabezadoTabla{EDT directa} & \EncabezadoTabla{Referencia UCP}\\\hline
Programación/Desarrollo de los núcleos & 2.400 HH & 5.083,13 HH\\\hline
Esfuerzo de los noventa paquetes / total de referencia & 4.200 HH & 12.707,82 HH\\\hline
\end{tabularx}
\caption{Contraste de esfuerzo directo y referencia UCP}
\label{tab:sd7-ucp-contraste}
\FuenteTabla{Elaboración propia: conteo de actores y casos, factores técnicos y ambientales y EDT inicial de los núcleos.}
\end{table}\par\smallskip
El contraste de programación alcanza 2,12 veces las HH de Desarrollo iniciales. La comparación de totales requiere además imputar análisis, diseño, pruebas y gobierno compartidos del mismo alcance; no se atribuye toda esa diferencia a subestimación porque la EDT directa excluye trabajos comunes. Aun así, la diferencia de programación no queda explicada por la mera existencia de bibliotecas reutilizables: no se dispone de un factor de productividad medido que permita reducir $CF$ hasta coincidir con la EDT.

La sensibilidad deteriora E1 y E3 de 3 a 2 y E7 de 3 a 4, manteniendo las restantes valoraciones. Con cuatro factores desfavorables, $EF=1,115$ y $CF=28$, la programación asciende a 7.856,20 HH y el total de referencia a 19.640,50 HH. Los registros de contraste por etapa y sensibilidad permiten reproducir estos resultados.

\paragraph{Modelo de interacción y regla de conteo.} El conteo se basa en el modelo de interacción propuesto para la oferta. Cada transacción identifica una entrada del actor y la respuesta que devuelve el sistema antes de esperar otro estímulo. Validar permisos, comprobar restricciones, calcular resultados y persistir datos son procesamiento de esa respuesta, salvo que se identifique otra entrada independiente. Una confirmación sólo se cuenta cuando el flujo requiere una decisión posterior del actor; no se añade una petición de validación para elevar el peso. El detalle de los 34 objetivos se conserva en \path{componentes/ucp_transacciones_nucleos.csv} y en su desarrollo detallado del Formulario T-15. Son decisiones de diseño propuestas, sujetas a validación con los usuarios, no interacciones ya observadas. Los rechazos dentro de una respuesta no reciben puntos adicionales; las resoluciones que exigen otra entrada están identificadas en el flujo.
\ifdefined\EnCuerpoSD
El detalle de las 34 transacciones por objetivo se presenta en el Formulario T-15, Tabla~\textit{Transacciones del modelo funcional propuesto}; T-14 conserva su correspondencia con los paquetes. En el capítulo, el conteo y las tablas de esfuerzo permiten seguir el cálculo sin repetir el listado completo.
\else
La Tabla~\ref{tab:t15-ucp-transacciones} documenta las entradas y respuestas que sustentan el peso de cada objetivo.
\begingroup\small\begin{longtable}{L{14mm}L{27mm}L{97mm}R{9mm}}
\caption{Transacciones del modelo funcional propuesto}\label{tab:t15-ucp-transacciones}\\\hline
 \EncabezadoTabla{Caso} & \EncabezadoTabla{Objetivo} & \EncabezadoTabla{Entradas y respuestas independientes} & \EncabezadoTabla{Peso}\\\hline\endfirsthead
\hline \EncabezadoTabla{Caso} & \EncabezadoTabla{Objetivo} & \EncabezadoTabla{Entradas y respuestas independientes} & \EncabezadoTabla{Peso}\\\hline\endhead
CU-01 & Registrar salida & 1. Actor: Identifica embarcación y patrón.\newline Sistema: Presenta contexto de salida y formulario.\newline 2. Actor: Presenta nómina y consentimiento y solicita registrar salida.\newline Sistema: Valida condición de registro y devuelve salida y hora registradas. & 5\\\hline
CU-02 & Confirmar regreso seguro & 1. Actor: Identifica nave e informante.\newline Sistema: Presenta salida pendiente.\newline 2. Actor: Presenta evidencia de regreso.\newline Sistema: Contrasta reporte y presenta discrepancias.\newline 3. Actor: Resuelve discrepancia y confirma regreso.\newline Sistema: Registra estado y responsable. & 5\\\hline
CU-03 & Declarar plan voluntario & 1. Actor: Identifica nave y responsable.\newline Sistema: Presenta formulario del plan.\newline 2. Actor: Presenta destino, duración y consentimiento.\newline Sistema: Valida y devuelve plan registrado. & 5\\\hline
CU-04 & Atender plan vencido & 1. Actor: Consulta vencimiento.\newline Sistema: Presenta plan vencido y evidencia.\newline 2. Actor: Registra contacto o intento fallido.\newline Sistema: Devuelve contacto registrado.\newline 3. Actor: Decide escalamiento según regla.\newline Sistema: Registra actuación y estado. & 5\\\hline
CU-05 & Autorizar amarra & 1. Actor: Consulta nave y puesto.\newline Sistema: Presenta compatibilidad y ocupación.\newline 2. Actor: Solicita asignación.\newline Sistema: Comprueba exclusividad local y confirma o rechaza.\newline 3. Actor: Presenta evidencia de maniobra segura.\newline Sistema: Devuelve evidencia registrada. & 5\\\hline
CU-06 & Preparar clase y nómina & 1. Actor: Selecciona clase y monitor.\newline Sistema: Presenta ficha de clase.\newline 2. Actor: Incorpora alumnos y habilitaciones.\newline Sistema: Comprueba consentimientos y restricciones y devuelve nómina revisable.\newline 3. Actor: Publica nómina.\newline Sistema: Devuelve nómina publicada. & 5\\\hline
CU-07 & Registrar participación por intervalo & 1. Actor: Consulta clase y nómina.\newline Sistema: Presenta participantes.\newline 2. Actor: Registra alumno, bote, monitor e intervalo.\newline Sistema: Valida solapamientos y devuelve cambios registrados. & 5\\\hline
CU-08 & Confirmar retiro de alumno & 1. Actor: Presenta solicitud y autorizado.\newline Sistema: Comprueba consentimiento e identidad y devuelve resultado.\newline 2. Actor: Coordina punto de entrega.\newline Sistema: Presenta coordinación registrada.\newline 3. Actor: Confirma entrega.\newline Sistema: Devuelve entrega única registrada. & 5\\\hline
CU-09 & Consultar salud mínima de participantes & 1. Actor: Identifica monitor y clase y consulta salud.\newline Sistema: Comprueba participación efectiva y devuelve datos mínimos autorizados. & 5\\\hline
CU-10 & Conciliar retorno de clase & 1. Actor: Presenta retorno por bote.\newline Sistema: Contrasta participantes y presenta diferencias.\newline 2. Actor: Resuelve diferencia o registra alerta.\newline Sistema: Devuelve conciliación o alerta registrada.\newline 3. Actor: Confirma retorno completo.\newline Sistema: Devuelve cierre de retorno. & 5\\\hline
CU-11 & Actualizar condición de infraestructura & 1. Actor: Identifica puesto, boya o fondeo.\newline Sistema: Presenta estado vigente.\newline 2. Actor: Presenta condición y evidencia.\newline Sistema: Valida relaciones y vigencia y devuelve inventario actualizado. & 5\\\hline
CU-12 & Presentar inspección desconectada & 1. Actor: Identifica boya y orden.\newline Sistema: Presenta formulario de inspección.\newline 2. Actor: Registra hallazgos y evidencia sin enlace.\newline Sistema: Devuelve resultado guardado localmente.\newline 3. Actor: Solicita sincronización.\newline Sistema: Devuelve recepción o conflicto.\newline 4. Actor: Resuelve duplicado o conflicto.\newline Sistema: Devuelve conciliación registrada. & 10\\\hline
CU-13 & Asociar medidor temporalmente & 1. Actor: Identifica medidor e instalación.\newline Sistema: Presenta asociaciones vigentes.\newline 2. Actor: Presenta intervalo de asociación.\newline Sistema: Valida exclusividad temporal y devuelve asociación trazable. & 5\\\hline
CU-14 & Consultar habilitación inicial de faena & 1. Actor: Identifica faena o contratista y consulta habilitación.\newline Sistema: Comprueba vigencia y autorización y devuelve estado y bloqueos. & 5\\\hline
CU-15 & Programar faena de varadero & 1. Actor: Presenta nave, trabajo y fechas.\newline Sistema: Presenta puestos compatibles y restricciones de campaña y maniobra.\newline 2. Actor: Selecciona puesto y solicita reserva.\newline Sistema: Devuelve reserva o alternativa. & 5\\\hline
CU-16 & Autorizar izaje & 1. Actor: Identifica orden y nave.\newline Sistema: Presenta orden y requisitos de izaje.\newline 2. Actor: Presenta plan, medios y responsables.\newline Sistema: Comprueba restricciones y devuelve estado de revisión.\newline 3. Actor: Autoriza o suspende con evidencia.\newline Sistema: Devuelve decisión registrada. & 5\\\hline
CU-17 & Cerrar orden de trabajo & 1. Actor: Consulta orden.\newline Sistema: Presenta orden y habilitación.\newline 2. Actor: Presenta avance, incidencias y evidencia y solicita cierre.\newline Sistema: Valida resultado y devuelve cierre o necesidad de reprogramación. & 5\\\hline
CU-18 & Habilitar acceso a intervención & 1. Actor: Presenta identidad y orden.\newline Sistema: Comprueba acreditación, inducción, vigencia y zona y devuelve autorización o denegación.\newline 2. Actor: Registra entrada.\newline Sistema: Devuelve entrada registrada.\newline 3. Actor: Registra salida.\newline Sistema: Devuelve salida registrada. & 5\\\hline
CU-19 & Registrar despacho de combustible & 1. Actor: Identifica nave y operador.\newline Sistema: Presenta contexto del despacho.\newline 2. Actor: Presenta producto, volumen y evidencia en momento seguro.\newline Sistema: Valida y devuelve hecho local único. & 5\\\hline
CU-20 & Rectificar despacho con conciliación & 1. Actor: Consulta despacho.\newline Sistema: Presenta registro y evidencia.\newline 2. Actor: Presenta rectificación motivada.\newline Sistema: Comprueba autoridad y presenta diferencias.\newline 3. Actor: Confirma rectificación.\newline Sistema: Devuelve nueva versión conservando origen. & 5\\\hline
CU-21 & Conciliar serie de consumo & 1. Actor: Selecciona medidor e intervalo.\newline Sistema: Presenta lecturas, calidad y asociación temporal.\newline 2. Actor: Resuelve faltantes y anomalías.\newline Sistema: Devuelve serie conciliada.\newline 3. Actor: Publica serie.\newline Sistema: Devuelve serie validada publicada. & 5\\\hline
CU-22 & Consultar consumo e historia ambiental & 1. Actor: Solicita período y contrato.\newline Sistema: Valida identidad y pertenencia y devuelve serie común y cobertura. & 5\\\hline
CU-23 & Recibir residuo con trazabilidad & 1. Actor: Identifica origen y categoría.\newline Sistema: Presenta datos de recepción.\newline 2. Actor: Presenta cantidad, fecha y evidencia.\newline Sistema: Presenta discrepancias.\newline 3. Actor: Concilia discrepancias y confirma recepción.\newline Sistema: Devuelve recepción y retención registradas. & 5\\\hline
CU-24 & Cerrar entrega a gestor autorizado & 1. Actor: Selecciona residuos y gestor.\newline Sistema: Presenta saldo y autorización del gestor.\newline 2. Actor: Presenta entrega parcial.\newline Sistema: Valida cantidades y devuelve entrega registrada.\newline 3. Actor: Vincula certificado sanitario.\newline Sistema: Devuelve certificado vinculado.\newline 4. Actor: Confirma cierre.\newline Sistema: Devuelve saldo y expediente cerrado. & 10\\\hline
CU-25 & Publicar indicador e imputación de consumo & 1. Actor: Selecciona serie validada.\newline Sistema: Comprueba contrato y permanencia y devuelve indicador y cargo propuesto.\newline 2. Actor: Revisa exclusiones y resuelve diferencias.\newline Sistema: Devuelve propuesta ajustada.\newline 3. Actor: Publica resultado.\newline Sistema: Devuelve resultado trazable publicado. & 5\\\hline
CU-26 & Mantener maestro contractual & 1. Actor: Identifica persona, organización o nave.\newline Sistema: Presenta maestro vigente.\newline 2. Actor: Presenta datos y documentos.\newline Sistema: Valida duplicados y relaciones y devuelve versión y autoridad. & 5\\\hline
CU-27 & Consultar contrato vigente & 1. Actor: Identifica contrato y período y solicita consulta.\newline Sistema: Comprueba permisos y vigencia y devuelve documentos y límites. & 5\\\hline
CU-28 & Atender visita náutica & 1. Actor: Presenta nave visitante y responsable.\newline Sistema: Presenta disponibilidad comercial.\newline 2. Actor: Solicita autorización operacional.\newline Sistema: Devuelve autorización o rechazo de la autoridad local.\newline 3. Actor: Confirma asignación transitoria.\newline Sistema: Devuelve asignación y condiciones. & 5\\\hline
CU-29 & Gestionar lista de espera & 1. Actor: Presenta solicitud y prioridades.\newline Sistema: Valida elegibilidad y presenta posición o alternativa.\newline 2. Actor: Acepta alternativa o mantiene espera.\newline Sistema: Devuelve estado registrado. & 5\\\hline
CU-30 & Formalizar cambio contractual & 1. Actor: Identifica contrato y cambio.\newline Sistema: Presenta estado contractual.\newline 2. Actor: Presenta documentos y fechas y confirma cambio.\newline Sistema: Comprueba exclusividad operacional y devuelve versión contractual. & 5\\\hline
CU-31 & Presentar lote de hechos facturables & 1. Actor: Selecciona hechos conciliados.\newline Sistema: Aplica tarifa y comprueba duplicados y período y presenta lote.\newline 2. Actor: Presenta lote contable.\newline Sistema: Devuelve aceptación o rechazo registrado. & 5\\\hline
CU-32 & Conciliar respuesta contable & 1. Actor: Presenta emisión, pagos y saldo recibidos.\newline Sistema: Valida referencia y origen y presenta diferencias por moneda.\newline 2. Actor: Resuelve diferencia o solicita reintento.\newline Sistema: Devuelve conciliación o resultado del reintento.\newline 3. Actor: Confirma conciliación.\newline Sistema: Devuelve estado importado actualizado. & 5\\\hline
CU-33 & Resolver medida de cobranza & 1. Actor: Consulta deuda.\newline Sistema: Presenta deuda y evidencia.\newline 2. Actor: Presenta fundamento y propuesta.\newline Sistema: Devuelve propuesta pendiente de decisión.\newline 3. Actor: Registra decisión competente.\newline Sistema: Devuelve actuación sin sanción automática. & 5\\\hline
CU-34 & Exportar expediente de custodia & 1. Actor: Selecciona expediente y alcance y solicita exportación.\newline Sistema: Comprueba permisos y piezas, genera formatos y manifiesto, verifica integridad y PDF/A y devuelve exportación trazable. & 5\\\hline
\end{longtable}\endgroup

\FuenteTabla{Elaboración propia de Synaptix: modelo de interacción y reglas de conteo UCP.}
El peso se asigna por intercambios independientes, sin sumar validaciones internas; los registros preservan la correspondencia con componentes y paquetes del mismo alcance.
\fi

\paragraph{Esfuerzo seleccionado y conciliación.} Synaptix adopta la referencia UCP sin reducir el factor de conversión por reutilización no medida: 7.766 HH en E1 y 4.942 HH en E2, 12.708 HH en total. El redondeo por etapa deja la programación en 3.107 y 1.977 HH, respectivamente, por encima de los 3.106,36 y 1.976,77 HH de referencia. Las 4.200 HH originales se acreditan una sola vez y se añaden 8.508 HH en 216 paquetes de 8--80 HH, identificados en \path{componentes/paquetes_conciliacion_ucp.csv}. Cada lote tiene resultado, responsable y aceptación; sus precedencias completan especificación/diseño, construcción, verificación y documentación antes del manifiesto de versión.

La imputación por componente es proporcional a su peso funcional; el peso de los actores conserva el reparto E1/E2 del modelo. Análisis acredita las HH Funcional existentes; diseño acredita Datos, Integración y Seguridad propios del núcleo; programación acredita Desarrollo; pruebas acredita Calidad; sobrecarga cubre coordinación y documentación propias del componente. La matriz \path{componentes/conciliacion_ucp_actividades.csv} concilia exactamente el esfuerzo seleccionado. La seguridad compartida, los canales, adaptadores, migración, UAT integral, formación y operación permanecen fuera de esta frontera; las comprobaciones del núcleo no se presupuestan otra vez como UAT integral.

La Tabla~\ref{tab:sd7-ucp-seleccion} presenta distribución del esfuerzo seleccionado por actividad y etapa.
\begin{table}[H]
\centering
\begin{tabularx}{\linewidth}{Y R{21mm}R{21mm}R{22mm}}\hline
 \EncabezadoTabla{Actividad} & \EncabezadoTabla{E1 HH} & \EncabezadoTabla{E2 HH} & \EncabezadoTabla{Total HH}\\\hline
Análisis & 777 & 494 & 1.271\\\hline
Diseño & 1.553 & 988 & 2.541\\\hline
Programación & 3.107 & 1.977 & 5.084\\\hline
Pruebas & 1.165 & 741 & 1.906\\\hline
Sobrecarga & 1.164 & 742 & 1.906\\\hline
Total & 7.766 & 4.942 & 12.708\\\hline
\end{tabularx}
\caption{Distribución del esfuerzo seleccionado por actividad y etapa}
\label{tab:sd7-ucp-seleccion}
\FuenteTabla{Elaboración propia: conciliación UCP y redondeo por etapa, sin duplicar esfuerzo compartido.}
\end{table}\par\smallskip
La decisión fija el esfuerzo de diseño de la oferta y sus recursos recalculados; no acredita productividad histórica ni aprobación humana o contractual. La sensibilidad de 19.640,50 HH es un escenario de deterioro, separado del esfuerzo seleccionado: si cambian los factores de provisión, Proyecto recalculará esfuerzo, perfiles y calendario antes de comprometer los puestos afectados. El modelo delimita los núcleos, y no se presenta como conteo UCP de todo el software de las cuentas 1.3, 1.7 y 1.9.

\subsection{Resumen integrado de esfuerzo y capacidad}
Synaptix selecciona 49.396 HH en 2181 paquetes y 122.070 HH de cobertura propia: 171.466 HH profesionales. Los 1.959 paquetes de origen y sus 40.744 HH se conservan; 216 paquetes de conciliación funcional añaden 8.508 HH. Se añaden seis paquetes mensuales de mentoría, 144 HH, con dos puestos parciales propios e independientes de cobertura. Los padres no reciben horas. Los códigos enlazan resultado, responsable, perfiles, actividades, red y recepción; las reservas representan capacidad futura y no trabajos ejecutados.
La Tabla~\ref{tab:sd7-esfuerzo-integral} presenta esfuerzo de entregables y cobertura por fase.
\begin{table}[H]
\centering
\begin{tabularx}{\linewidth}{Y R{25mm}}\hline \EncabezadoTabla{Concepto} & \EncabezadoTabla{HH}\\\hline
Implementación: entregables e ingeniería & 30.872\\\hline
Cobertura inicial y marchas blancas (13--20) & 22.692\\\hline
Operación: entregables e ingeniería & 18.524\\\hline
Operación: NOC/SOC y mesa (21--56) & 99.378\\\hline
Total profesional programado & 171.466\\\hline
\end{tabularx}
\caption{Esfuerzo de entregables y cobertura por fase}
\label{tab:sd7-esfuerzo-integral}
\FuenteTabla{Elaboración propia: diccionario y asignaciones integradas de T-14/T-15.}
\end{table}\par\smallskip
La suma distingue 49.396 HH de entregables e ingeniería de 122.070 HH de cobertura; ambas partidas concilian las 171.466 HH sin duplicar turnos. Ingeniería dispone de 7,2 HH productivas por día y un máximo de 50 HH de paquetes por persona y quincena. Los turnos y el control natural final disponen de ejecutores propios distintos, con máximo de 144 HH al mes, 40 HH por siete días y doce horas de descanso. El máximo mensual de dotación comprometida es 106 puestos, incluidos los ocho titulares; el máximo de refuerzos de ingeniería es 98. Es el resultado de esta asignación, sin atribuir contratación ni optimización global. Synaptix proveerá estaciones, accesos, entornos, incorporación y suplencias para cada puesto antes de su intervención; los ocho equipos de los titulares no cubren la dotación adicional.
La habilitación común inicial se imputa a E1 y el gobierno posterior al mes 12 al desarrollo/cierre E2. La estabilización, aceptación final y cobertura de E1 en producción se mantienen como partidas adicionales identificadas. Los máximos por fase no se suman como personas diferentes. La valoración profesional se desarrolla exclusivamente en la Oferta Económica.

\subsection{Mentoría técnica posterior a producción E2}
Synaptix acompañará a los referentes técnicos y sus suplentes designados por la marina durante los seis meses 21--26, agosto de 2028 a enero de 2029. Este servicio es diferente de la estabilización de 28 días, la mesa ordinaria y el acompañamiento de salida. José Lara Arce responde por coordinación y transferencia; el refuerzo MEN-IL-01 ejecutará integración e implantación y MEN-Q-01 verificará evidencia, con capacidades independientes de turnos y de ingeniería ya asignada. Son puestos previstos, sin atribuir contratación ni aceptación ejecutadas.
Cada paquete \texttt{1.12.13.m} reserva 24 HH mensuales: Implantación 16, Integración 4 y Calidad 4. Comprende cuatro contactos de acompañamiento de hasta tres horas, cuatro horas de preparación/documentación, cuatro de resolución técnica y cuatro de revisión independiente. Son HH de proveedor; no incluyen como esfuerzo propio las horas del CLIENTE. El total es 144 HH adicionales de Operación, incluidas sin cargo adicional y valorables por los perfiles en la Oferta Económica. La cobertura NOC/SOC/mesa y las bolsas de ingeniería no se reducen ni vuelven a contarse.
Durante agosto--noviembre y hasta el 14 de diciembre los casos podrán incluir ejercicio supervisado en entorno seguro. Durante el congelamiento, en particular enero, la mentoría se limitará a acompañamiento remoto de procedimientos ya operativos, revisión de casos y documentación, sin capacitación formal nueva, cambios de producción ni interferencia con maniobras. Si hace falta formación nueva, se preparará antes del congelamiento o en la siguiente ventana permitida, conservando la atención a los referentes durante los seis meses.
\begin{tabularx}{\linewidth}{L{15mm}Y R{15mm}}\hline
Mes & Resultado mensual & HH\\\hline
21 & Referentes/suplentes habilitados, accesos y protocolo de escalamiento & 24\\\hline
22 & Casos de incidentes, recuperación y conciliación acompañados & 24\\\hline
23 & Solicitudes, cambios y lectura de indicadores practicados & 24\\\hline
24 & Preparación estacional y transferencia de procedimientos vigentes & 24\\\hline
25 & Revisión de casos y continuidad de responsables durante congelamiento & 24\\\hline
26 & Cierre de transferencia, suplencia y plan de continuidad del conocimiento & 24\\\hline
\end{tabularx}
Cada recepción exige agenda, participantes y suplentes, casos, acta, documentación actualizada, resultado de verificación y pendientes con responsable/plazo. El CLIENTE valida las decisiones de negocio; Synaptix conserva la administración técnica del servicio. Ausencia del referente activa suplente y reprogramación dentro del mes, sin interrupción de atención. Capacidad adicional necesaria para cumplir el servicio será provista por Synaptix, sin usar esas HH como límite excluyente.\parencite[RT-22.07--08]{bases-transversales}\parencite[arts.~89--90]{bases-admin}

Las cantidades, cohortes, turnos, capacidades y asignaciones se desarrollan en el Formulario T-15. Las familias de este diccionario conservan sus cantidades y esfuerzo; la estrategia de despliegue y reversión se desarrolla en el Formulario T-18.
\section{Diccionario de familias repetitivas y capacidades}
Los lotes, cohortes, turnos y períodos que comparten resultado, aceptación, responsable y esfuerzo se presentan mediante una ficha de familia. Cada instancia conserva su código, período y predecesoras en el registro técnico \path{componentes/diccionario_integrado.csv}. Cada recepción permanece individual y el esfuerzo de T-15 conserva su imputación. Las reservas sólo acreditan resultados cuando se ejecutan y verifican.
La Tabla~\ref{tab:sd7-dic-complementario} identifica las familias. Los códigos de instancia se relacionan de forma exhaustiva con cada ficha en \path{componentes/familias_t14.csv}. Las fechas y cantidades son las del escenario de programación; su ajuste conserva el alcance obligatorio y las condiciones de recepción.
\begingroup\small\setlength{\tabcolsep}{3pt}
\begin{longtable}{L{40mm}L{103mm}}
\caption{Diccionario de familias: resultado, aceptación y responsabilidad}\label{tab:sd7-dic-complementario}\\\hline
Familia y alcance & Entregable y criterio de aceptación\\\hline\endfirsthead
\hline Familia y alcance & Entregable y criterio de aceptación\\\hline\endhead
\texttt{F01}\newline \texttt{1.\allowbreak{}8.\allowbreak{}7.\allowbreak{}1}\newline 14 instancias; 112 HH\newline Mes 7 & Expediente de abandono revisado y vinculado. Folio, origen, integridad y vínculos recibidos; faltantes y decisiones jurídicas identificados.\newline Responsable: Nicolás Torfan. 8 HH por instancia.\\\hline
\texttt{F02}\newline \texttt{1.\allowbreak{}8.\allowbreak{}7.\allowbreak{}2}\newline 56 instancias; 448 HH\newline Mes 7 & Lote de hasta 32 autorizaciones escolares. Lectura humana, vigencia, menor/apoderado y evidencia; autorización ausente bloquea participación.\newline Responsable: Nicolás Torfan. 8 HH por instancia.\\\hline
\texttt{F03}\newline \texttt{1.\allowbreak{}2.\allowbreak{}3.\allowbreak{}7}\newline 10 instancias; 160 HH\newline Mes 6 & Lote de enrolamiento de hasta cuatro terminales. Equipo nominal inscrito, cifrado, pérdida, protección, licencia y revisión de política comprobados.\newline Responsable: José Lara Arce. 16 HH por instancia.\\\hline
\texttt{F04}\newline \texttt{1.\allowbreak{}7.\allowbreak{}2.\allowbreak{}5}\newline 16 instancias; 384 HH\newline Mes 5 & Recepción funcional de 40 puntos de medición. Se exigirá la recepción de los 640 puntos: identificación por servicio/amarra, calibración, error, flexión, tiempo, corte y lectura trazados; instalación eléctrica a cargo del cliente.\newline Responsable: Fernando Guerrero. 24 HH por instancia.\\\hline
\texttt{F05}\newline \texttt{1.\allowbreak{}8.\allowbreak{}8}\newline 4 instancias; 224 HH\newline Mes 14 & Adaptación de ensayo al universo E2. Manifiesto E2 identifica esquemas/contratos, transformaciones y universos incluidos, sin usar vistas como modelo.\newline Responsable: Nicolás Torfan. 56 HH por instancia.\\\hline
\texttt{F06}\newline \texttt{1.\allowbreak{}11.\allowbreak{}7.\allowbreak{}1.\allowbreak{}1.\allowbreak{}1}\newline 100 instancias; 1200 HH\newline Mes 12 & Sesión y certificación E1, perfil 1, cohorte. Dos sesiones de dos horas, ejercicio, registro individual y recuperación; máximo doce asistentes y sin interrumpir turnos.\newline Responsable: José Lara Arce. 12 HH por instancia.\\\hline
\texttt{F07}\newline \texttt{1.\allowbreak{}11.\allowbreak{}7.\allowbreak{}1.\allowbreak{}2.\allowbreak{}1}\newline 2 instancias; 24 HH\newline Mes 12 & Sesión y certificación E1, perfil 2, cohorte. Dos sesiones de dos horas, ejercicio, registro individual y recuperación; máximo doce asistentes y sin interrumpir turnos.\newline Responsable: José Lara Arce. 12 HH por instancia.\\\hline
\texttt{F08}\newline \texttt{1.\allowbreak{}11.\allowbreak{}7.\allowbreak{}1.\allowbreak{}3.\allowbreak{}1.\allowbreak{}1}\newline 1 instancia; 12 HH\newline Mes 12 & Sesión y certificación E1, perfil 3, cohorte. Dos sesiones de dos horas, ejercicio, registro individual y recuperación; máximo doce asistentes y sin interrumpir turnos.\newline Responsable: José Lara Arce. 12 HH por instancia.\\\hline
\texttt{F09}\newline \texttt{1.\allowbreak{}11.\allowbreak{}7.\allowbreak{}1.\allowbreak{}4.\allowbreak{}1.\allowbreak{}1}\newline 1 instancia; 12 HH\newline Mes 12 & Sesión y certificación E1, perfil 4, cohorte. Dos sesiones de dos horas, ejercicio, registro individual y recuperación; máximo doce asistentes y sin interrumpir turnos.\newline Responsable: José Lara Arce. 12 HH por instancia.\\\hline
\texttt{F10}\newline \texttt{1.\allowbreak{}11.\allowbreak{}7.\allowbreak{}1.\allowbreak{}5.\allowbreak{}1}\newline 3 instancias; 36 HH\newline Mes 12 & Sesión y certificación E1, perfil 5, cohorte. Dos sesiones de dos horas, ejercicio, registro individual y recuperación; máximo doce asistentes y sin interrumpir turnos.\newline Responsable: José Lara Arce. 12 HH por instancia.\\\hline
\texttt{F11}\newline \texttt{1.\allowbreak{}11.\allowbreak{}4.\allowbreak{}1}\newline 224 instancias; 1792 HH\newline Mes 13 & Acompañamiento E1, por ola, por día, por zona. Registro diario, turno, zona, tarea, incidencia y relevo; 8 h presencia días 1--14, luego 4 h presencia +4 h preparación.\newline Responsable: José Lara Arce. 8 HH por instancia.\\\hline
\texttt{F12}\newline \texttt{1.\allowbreak{}11.\allowbreak{}6.\allowbreak{}1}\newline 56 instancias; 448 HH\newline Mes 16 & Estabilización presencial E1, por día, por zona. Dos puestos diarios, pantalanes y varadero, 28 días consecutivos, incluidos fines de semana; relevo e incidencias registrados.\newline Responsable: José Lara Arce. 8 HH por instancia.\\\hline
\texttt{F13}\newline \texttt{1.\allowbreak{}11.\allowbreak{}7.\allowbreak{}2.\allowbreak{}1.\allowbreak{}1}\newline 100 instancias; 1200 HH\newline Mes 18 & Sesión y certificación E2, perfil 1, cohorte. Dos sesiones de dos horas, ejercicio, registro individual y recuperación; máximo doce asistentes y sin interrumpir turnos.\newline Responsable: José Lara Arce. 12 HH por instancia.\\\hline
\texttt{F14}\newline \texttt{1.\allowbreak{}11.\allowbreak{}7.\allowbreak{}2.\allowbreak{}2.\allowbreak{}1}\newline 2 instancias; 24 HH\newline Mes 18 & Sesión y certificación E2, perfil 2, cohorte. Dos sesiones de dos horas, ejercicio, registro individual y recuperación; máximo doce asistentes y sin interrumpir turnos.\newline Responsable: José Lara Arce. 12 HH por instancia.\\\hline
\texttt{F15}\newline \texttt{1.\allowbreak{}11.\allowbreak{}7.\allowbreak{}2.\allowbreak{}3.\allowbreak{}1.\allowbreak{}1}\newline 1 instancia; 12 HH\newline Mes 18 & Sesión y certificación E2, perfil 3, cohorte. Dos sesiones de dos horas, ejercicio, registro individual y recuperación; máximo doce asistentes y sin interrumpir turnos.\newline Responsable: José Lara Arce. 12 HH por instancia.\\\hline
\texttt{F16}\newline \texttt{1.\allowbreak{}11.\allowbreak{}7.\allowbreak{}2.\allowbreak{}4.\allowbreak{}1.\allowbreak{}1}\newline 1 instancia; 12 HH\newline Mes 18 & Sesión y certificación E2, perfil 4, cohorte. Dos sesiones de dos horas, ejercicio, registro individual y recuperación; máximo doce asistentes y sin interrumpir turnos.\newline Responsable: José Lara Arce. 12 HH por instancia.\\\hline
\texttt{F17}\newline \texttt{1.\allowbreak{}11.\allowbreak{}7.\allowbreak{}2.\allowbreak{}5.\allowbreak{}1}\newline 3 instancias; 36 HH\newline Mes 18 & Sesión y certificación E2, perfil 5, cohorte. Dos sesiones de dos horas, ejercicio, registro individual y recuperación; máximo doce asistentes y sin interrumpir turnos.\newline Responsable: José Lara Arce. 12 HH por instancia.\\\hline
\texttt{F18}\newline \texttt{1.\allowbreak{}11.\allowbreak{}4.\allowbreak{}2}\newline 224 instancias; 1792 HH\newline Mes 19 & Acompañamiento E2, por ola, por día, por zona. Registro diario, turno, zona, tarea, incidencia y relevo; 8 h presencia días 1--14, luego 4 h presencia +4 h preparación.\newline Responsable: José Lara Arce. 8 HH por instancia.\\\hline
\texttt{F19}\newline \texttt{1.\allowbreak{}11.\allowbreak{}6.\allowbreak{}2}\newline 56 instancias; 448 HH\newline Mes 21 & Estabilización presencial E2, por día, por zona. Dos puestos diarios, pantalanes y varadero, 28 días consecutivos, incluidos fines de semana; relevo e incidencias registrados.\newline Responsable: José Lara Arce. 8 HH por instancia.\\\hline
\texttt{F20}\newline \texttt{1.\allowbreak{}12.\allowbreak{}10}\newline 75 instancias; 900 HH\newline Mes 24--48 & Incorporación estacional mensual, perfil p, cohorte. Incorporación y certificación antes del 15 de diciembre, incluidos monitores; 300 plazas anuales reservadas, ampliación sin costo si se requiere.\newline Responsable: José Lara Arce. 12 HH por instancia.\\\hline
\texttt{F21}\newline \texttt{1.\allowbreak{}12.\allowbreak{}12.\allowbreak{}correctiva}\newline 36 instancias; 2880 HH\newline Mes 21--56 & Capacidad reservada de ingeniería correctiva, mensual. Backlog, estimación, autorización, versión y evidencia; se liquida sólo el trabajo realizado, sin cobrar dos veces actividades ya presupuestadas.\newline Responsable: Nikolai Navea. 80 HH por instancia.\\\hline
\texttt{F22}\newline \texttt{1.\allowbreak{}12.\allowbreak{}12.\allowbreak{}normativa}\newline 36 instancias; 1728 HH\newline Mes 21--56 & Capacidad reservada de ingeniería normativa, mensual. Backlog, estimación, autorización, versión y evidencia; se liquida sólo el trabajo realizado, sin cobrar dos veces actividades ya presupuestadas.\newline Responsable: Nikolai Navea. 48 HH por instancia.\\\hline
\texttt{F23}\newline \texttt{1.\allowbreak{}12.\allowbreak{}12.\allowbreak{}evolutiva}\newline 36 instancias; 2880 HH\newline Mes 21--56 & Capacidad reservada de ingeniería evolutiva, mensual. Backlog, estimación, autorización, versión y evidencia; se liquida sólo el trabajo realizado, sin cobrar dos veces actividades ya presupuestadas.\newline Responsable: Nikolai Navea. 80 HH por instancia.\\\hline
\texttt{F24}\newline \texttt{1.\allowbreak{}12.\allowbreak{}12.\allowbreak{}deuda}\newline 36 instancias; 1152 HH\newline Mes 21--56 & Capacidad reservada de ingeniería deuda, mensual. Backlog, estimación, autorización, versión y evidencia; se liquida sólo el trabajo realizado, sin cobrar dos veces actividades ya presupuestadas.\newline Responsable: Nikolai Navea. 32 HH por instancia.\\\hline
\texttt{F25}\newline \texttt{1.\allowbreak{}12.\allowbreak{}2.\allowbreak{}2}\newline 26 instancias; 416 HH\newline Mes 31--56 & Seguimiento incremental IN2 mensual. Estado de planes/trabajos y ventanas recibidos; seguimiento sin duplicar informes generales.\newline Responsable: José Lara Arce. 16 HH por instancia.\\\hline
\texttt{F26}\newline \texttt{1.\allowbreak{}9.\allowbreak{}2.\allowbreak{}5.\allowbreak{}4}\newline 5 instancias; 80 HH\newline Mes 33--35 & Revisión previa a botadura, por lote de cuatro embarcaciones. Plan, responsabilidades y trabajos revisados al menos tres semanas antes de botadura, sin interacción durante izaje.\newline Responsable: José Lara Arce. 16 HH por instancia.\\\hline
\texttt{F27}\newline \texttt{1.\allowbreak{}10.\allowbreak{}3.\allowbreak{}1}\newline 4 instancias; 320 HH\newline Mes 10 & Reserva de subsanación y reensayo E1, por bloque. Hallazgos identificados cerrados, regresión y evidencia sobre misma versión; horas no ejecutadas quedan como reserva.\newline Responsable: Nikolai Navea. 80 HH por instancia.\\\hline
\texttt{F28}\newline \texttt{1.\allowbreak{}10.\allowbreak{}3.\allowbreak{}2}\newline 4 instancias; 320 HH\newline Mes 15 & Reserva de subsanación y reensayo E2, por bloque. Hallazgos identificados cerrados, regresión y evidencia sobre misma versión; horas no ejecutadas quedan como reserva.\newline Responsable: Nikolai Navea. 80 HH por instancia.\\\hline
\texttt{F29}\newline \texttt{1.\allowbreak{}1.\allowbreak{}5.\allowbreak{}1}\newline 1 instancia; 32 HH\newline Mes 2 & Expediente transversal de alcance y contratos. Cláusulas, documentos, responsables y control de vigencias; se distingue evidencia de oferta de obligaciones futuras.\newline Responsable: Ignacio Silva. 32 HH por instancia.\\\hline
\texttt{F30}\newline \texttt{1.\allowbreak{}1.\allowbreak{}5.\allowbreak{}2}\newline 1 instancia; 32 HH\newline Mes 3 & Expediente transversal de seguridad y privacidad. Cláusulas, documentos, responsables y control de vigencias; se distingue evidencia de oferta de obligaciones futuras.\newline Responsable: Ignacio Silva. 32 HH por instancia.\\\hline
\texttt{F31}\newline \texttt{1.\allowbreak{}1.\allowbreak{}5.\allowbreak{}3}\newline 1 instancia; 32 HH\newline Mes 3 & Expediente transversal de IA y trazabilidad. Cláusulas, documentos, responsables y control de vigencias; se distingue evidencia de oferta de obligaciones futuras.\newline Responsable: Ignacio Silva. 32 HH por instancia.\\\hline
\texttt{F32}\newline \texttt{1.\allowbreak{}1.\allowbreak{}5.\allowbreak{}4}\newline 1 instancia; 32 HH\newline Mes 3 & Expediente transversal de licencias y soporte. Cláusulas, documentos, responsables y control de vigencias; se distingue evidencia de oferta de obligaciones futuras.\newline Responsable: Ignacio Silva. 32 HH por instancia.\\\hline
\texttt{F33}\newline \texttt{1.\allowbreak{}1.\allowbreak{}5.\allowbreak{}5}\newline 1 instancia; 32 HH\newline Mes 2 & Expediente transversal de custodia y propiedad. Cláusulas, documentos, responsables y control de vigencias; se distingue evidencia de oferta de obligaciones futuras.\newline Responsable: Ignacio Silva. 32 HH por instancia.\\\hline
\texttt{F34}\newline \texttt{1.\allowbreak{}1.\allowbreak{}5.\allowbreak{}6}\newline 1 instancia; 32 HH\newline Mes 3 & Expediente transversal de garantías y oferta económica. Cláusulas, documentos, responsables y control de vigencias; se distingue evidencia de oferta de obligaciones futuras.\newline Responsable: Ignacio Silva. 32 HH por instancia.\\\hline
\texttt{F35}\newline \texttt{1.\allowbreak{}11.\allowbreak{}5.\allowbreak{}1.\allowbreak{}2}\newline 1 instancia; 8 HH\newline Mes 16 & Validación final del expediente vivo de marcha blanca E1. Los datos de las cuatro últimas semanas y las condiciones copulativas se incorporan al expediente; no anticipa aceptación.\newline Responsable: Davor Serey. 8 HH por instancia.\\\hline
\texttt{F36}\newline \texttt{1.\allowbreak{}11.\allowbreak{}5.\allowbreak{}2.\allowbreak{}2}\newline 1 instancia; 8 HH\newline Mes 21 & Validación final del expediente vivo de marcha blanca E2. Los datos de las cuatro últimas semanas y las condiciones copulativas se incorporan al expediente; no anticipa aceptación.\newline Responsable: Davor Serey. 8 HH por instancia.\\\hline
\texttt{F37}\newline \texttt{1.\allowbreak{}12.\allowbreak{}4}\newline 3 instancias; 144 HH\newline Mes 48--56 & Control de continuidad adicional: resiliencia. Prueba sobre versión identificada y ventana autorizada, sin degradar E1; recuperación y retorno medidos.\newline Responsable: Davor Serey. 48 HH por instancia.\\\hline
\texttt{F38}\newline \texttt{1.\allowbreak{}12.\allowbreak{}4}\newline 3 instancias; 192 HH\newline Mes 48--56 & Control de continuidad adicional: DR real y retorno. Prueba sobre versión identificada y ventana autorizada, sin degradar E1; recuperación y retorno medidos.\newline Responsable: José Lara Arce. 64 HH por instancia.\\\hline
\end{longtable}
\endgroup
\FuenteTabla{Elaboración propia a partir de los registros de paquetes y de su programación integrada.}

\subsection{Cobertura integral de requisitos y frontera contractual}
El registro \texttt{trazabilidad\_integral.csv} conserva los 576 identificadores, enunciados y declaraciones del T-12 y relaciona cada alcance ofrecido con cuentas, paquetes y recepción prevista. Los 27 complementos C07 se conservan íntegros en \texttt{trazabilidad\_complementos\_c07.csv}; no se reemplaza una cláusula transversal por la específica. El cruce acredita asignación del trabajo, no ejecución ni cumplimiento automático: cada protocolo debe verificar todas las cláusulas del enunciado. Las 29 exclusiones Deseables permanecen excluidas, sin reinterpretarlas como obligaciones cumplidas.
Las obligaciones corporativas de sitio web y video pertenecen a la oferta anterior al contrato; su preparación no se cuenta como desarrollo contratado de la marina. La cuenta 1.1.5 gestiona durante el contrato la vigencia y coherencia de esos documentos, licencias, compromisos y contratos. Esta frontera permite revisar los 576 IDs sin inventar entregables ya producidos. El T-12 conserva sus declaraciones y enunciados originales.
\begin{longtable}{L{23mm}L{46mm}L{72mm}}
\caption{Índice de requisitos hacia cuentas de trabajo}\label{tab:sd7-trazabilidad-total}\\\hline
ID & Cuentas EDT & Recepción / tratamiento\\\hline\endfirsthead
\hline ID & Cuentas EDT & Recepción / tratamiento\\\hline\endhead
RF-NAV-01 & 1.4.1.1; 1.4.1.2; 1.4.1.3 & 1.10.2.1.10 / 1.10.2.2.10\\\hline
RF-NAV-02 & 1.4.1.1; 1.4.1.2; 1.4.1.3 & 1.10.2.1.10 / 1.10.2.2.10\\\hline
RF-NAV-03 & 1.4.1.1; 1.4.1.2; 1.4.1.3 & 1.10.2.1.10 / 1.10.2.2.10\\\hline
RF-NAV-04 & 1.4.1.1; 1.4.1.2; 1.4.1.3 & 1.10.2.1.10 / 1.10.2.2.10\\\hline
RF-NAV-05 & 1.4.1.1; 1.4.1.2; 1.4.1.3 & 1.10.2.1.10 / 1.10.2.2.10\\\hline
RF-NAV-06 & 1.4.1.1; 1.4.1.2; 1.4.1.3 & 1.10.2.1.10 / 1.10.2.2.10\\\hline
RF-NAV-07 & 1.4.1.1; 1.4.1.2; 1.4.1.3 & 1.10.2.1.10 / 1.10.2.2.10\\\hline
RF-NAV-08 & 1.4.1.1; 1.4.1.2; 1.4.1.3 & 1.10.2.1.10 / 1.10.2.2.10\\\hline
RF-NAV-09 & 1.4.1.1; 1.4.1.2; 1.4.1.3 & 1.10.2.1.10 / 1.10.2.2.10\\\hline
RF-NAV-10 & 1.4.1.1; 1.4.1.2; 1.4.1.3 & 1.10.2.1.10 / 1.10.2.2.10\\\hline
RF-NAV-11 & 1.4.1.1; 1.4.1.2; 1.4.1.3 & 1.10.2.1.10 / 1.10.2.2.10\\\hline
RF-NAV-12 & 1.4.1.1; 1.4.1.2; 1.4.1.3 & 1.10.2.1.10 / 1.10.2.2.10\\\hline
RF-NAV-13 & 1.4.1.1; 1.4.1.2; 1.4.1.3 & 1.10.2.1.10 / 1.10.2.2.10\\\hline
RF-ALE-01 & 1.7.4.1; 1.4.1.3 & 1.10.2.1.10 / 1.10.2.2.10\\\hline
RF-ALE-02 & 1.7.4.1; 1.4.1.3 & 1.10.2.1.10 / 1.10.2.2.10\\\hline
RF-ALE-03 & 1.7.4.1; 1.4.1.3 & 1.10.2.1.10 / 1.10.2.2.10\\\hline
RF-ALE-04 & 1.7.4.1; 1.4.1.3 & 1.10.2.1.10 / 1.10.2.2.10\\\hline
RF-ALE-05 & 1.7.4.1; 1.4.1.3 & 1.10.2.1.10 / 1.10.2.2.10\\\hline
RF-ALE-06 & 1.7.4.1; 1.4.1.3 & 1.10.2.1.10 / 1.10.2.2.10\\\hline
RF-ALE-07 & 1.7.4.1; 1.4.1.3 & 1.10.2.1.10 / 1.10.2.2.10\\\hline
RF-ALE-08 & 1.7.4.1; 1.4.1.3 & 1.10.2.1.10 / 1.10.2.2.10\\\hline
RF-ALE-09 & 1.7.4.1; 1.4.1.3 & 1.10.2.1.10 / 1.10.2.2.10\\\hline
RF-ALE-10 & 1.7.4.1; 1.4.1.3 & 1.10.2.1.10 / 1.10.2.2.10\\\hline
RF-ALE-11 & 1.7.4.1; 1.4.1.3 & 1.10.2.1.10 / 1.10.2.2.10\\\hline
RF-ALE-12 & 1.7.4.1; 1.4.1.3 & 1.10.2.1.10 / 1.10.2.2.10\\\hline
RF-ESC-01 & 1.4.2 & 1.10.2.1.10 / 1.10.2.2.10\\\hline
RF-ESC-02 & 1.4.2 & 1.10.2.1.10 / 1.10.2.2.10\\\hline
RF-ESC-03 & 1.4.2 & 1.10.2.1.10 / 1.10.2.2.10\\\hline
RF-ESC-04 & 1.4.2 & 1.10.2.1.10 / 1.10.2.2.10\\\hline
RF-ESC-05 & 1.4.2 & 1.10.2.1.10 / 1.10.2.2.10\\\hline
RF-ESC-06 & 1.4.2 & 1.10.2.1.10 / 1.10.2.2.10\\\hline
RF-ESC-07 & 1.4.2 & 1.10.2.1.10 / 1.10.2.2.10\\\hline
RF-ESC-08 & 1.4.2 & 1.10.2.1.10 / 1.10.2.2.10\\\hline
RF-ESC-09 & 1.4.2 & 1.10.2.1.10 / 1.10.2.2.10\\\hline
RF-ESC-10 & 1.4.2 & 1.10.2.1.10 / 1.10.2.2.10\\\hline
RF-ESC-11 & 1.4.2 & 1.10.2.1.10 / 1.10.2.2.10\\\hline
RF-ESC-12 & 1.4.2 & 1.10.2.1.10 / 1.10.2.2.10\\\hline
RF-ESC-13 & 1.4.2 & 1.10.2.1.10 / 1.10.2.2.10\\\hline
RF-ESC-14 & 1.4.2 & 1.10.2.1.10 / 1.10.2.2.10\\\hline
RF-AMA-01 & 1.4.1.4; 1.6.1 & 1.10.2.1.10 / 1.10.2.2.10\\\hline
RF-AMA-02 & 1.4.1.4; 1.6.1 & 1.10.2.1.10 / 1.10.2.2.10\\\hline
RF-AMA-03 & 1.4.1.4; 1.6.1 & 1.10.2.1.10 / 1.10.2.2.10\\\hline
RF-AMA-04 & 1.4.1.4; 1.6.1 & 1.10.2.1.10 / 1.10.2.2.10\\\hline
RF-AMA-05 & 1.4.1.4; 1.6.1 & 1.10.2.1.10 / 1.10.2.2.10\\\hline
RF-AMA-06 & 1.4.1.4; 1.6.1 & 1.10.2.1.10 / 1.10.2.2.10\\\hline
RF-AMA-07 & 1.4.1.4; 1.6.1 & 1.10.2.1.10 / 1.10.2.2.10\\\hline
RF-AMA-08 & 1.4.1.4; 1.6.1 & 1.10.2.1.10 / 1.10.2.2.10\\\hline
RF-AMA-09 & 1.4.1.4; 1.6.1 & 1.10.2.1.10 / 1.10.2.2.10\\\hline
RF-AMA-10 & 1.4.1.4; 1.6.1 & 1.10.2.1.10 / 1.10.2.2.10\\\hline
RF-AMA-11 & 1.4.1.4; 1.6.1 & 1.10.2.1.10 / 1.10.2.2.10\\\hline
RF-AMA-12 & 1.4.1.4; 1.6.1 & 1.10.2.1.10 / 1.10.2.2.10\\\hline
RF-VAR-01 & 1.5.1 & 1.10.2.1.10 / 1.10.2.2.10\\\hline
RF-VAR-02 & 1.5.1 & 1.10.2.1.10 / 1.10.2.2.10\\\hline
RF-VAR-03 & 1.5.1 & 1.10.2.1.10 / 1.10.2.2.10\\\hline
RF-VAR-04 & 1.5.1 & 1.10.2.1.10 / 1.10.2.2.10\\\hline
RF-VAR-05 & 1.5.1 & 1.10.2.1.10 / 1.10.2.2.10\\\hline
RF-VAR-06 & 1.5.1 & 1.10.2.1.10 / 1.10.2.2.10\\\hline
RF-VAR-07 & 1.5.1 & 1.10.2.1.10 / 1.10.2.2.10\\\hline
RF-VAR-08 & 1.5.1 & 1.10.2.1.10 / 1.10.2.2.10\\\hline
RF-CTR-01 & 1.5.1.5 & 1.10.2.1.10 / 1.10.2.2.10\\\hline
RF-CTR-02 & 1.5.1.5 & 1.10.2.1.10 / 1.10.2.2.10\\\hline
RF-CTR-03 & 1.5.1.5 & 1.10.2.1.10 / 1.10.2.2.10\\\hline
RF-CTR-04 & 1.5.1.5 & 1.10.2.1.10 / 1.10.2.2.10\\\hline
RF-CTR-05 & 1.5.1.5 & 1.10.2.1.10 / 1.10.2.2.10\\\hline
RF-CTR-06 & 1.5.1.5 & 1.10.2.1.10 / 1.10.2.2.10\\\hline
RF-CTR-07 & 1.5.1.5 & 1.10.2.1.10 / 1.10.2.2.10\\\hline
RF-AMB-01 & 1.5.3 & 1.10.2.1.10 / 1.10.2.2.10\\\hline
RF-AMB-02 & 1.5.3 & 1.10.2.1.10 / 1.10.2.2.10\\\hline
RF-AMB-03 & 1.5.3 & 1.10.2.1.10 / 1.10.2.2.10\\\hline
RF-AMB-04 & 1.5.3 & 1.10.2.1.10 / 1.10.2.2.10\\\hline
RF-AMB-05 & 1.5.3 & 1.10.2.1.10 / 1.10.2.2.10\\\hline
RF-AMB-06 & 1.5.3 & 1.10.2.1.10 / 1.10.2.2.10\\\hline
RF-AMB-07 & 1.5.3 & 1.10.2.1.10 / 1.10.2.2.10\\\hline
RF-CON-01 & 1.5.3; 1.4.3.3 & 1.10.2.1.10 / 1.10.2.2.10\\\hline
RF-CON-02 & 1.5.3; 1.4.3.3 & 1.10.2.1.10 / 1.10.2.2.10\\\hline
RF-CON-03 & 1.5.3; 1.4.3.3 & 1.10.2.1.10 / 1.10.2.2.10\\\hline
RF-CON-04 & 1.5.3; 1.4.3.3 & 1.10.2.1.10 / 1.10.2.2.10\\\hline
RF-CON-05 & 1.5.3; 1.4.3.3 & 1.10.2.1.10 / 1.10.2.2.10\\\hline
RF-CON-06 & 1.5.3; 1.4.3.3 & 1.10.2.1.10 / 1.10.2.2.10\\\hline
RF-CON-07 & 1.5.3; 1.4.3.3 & 1.10.2.1.10 / 1.10.2.2.10\\\hline
RF-CON-08 & 1.5.3; 1.4.3.3 & 1.10.2.1.10 / 1.10.2.2.10\\\hline
RF-CON-09 & 1.5.3; 1.4.3.3 & 1.10.2.1.10 / 1.10.2.2.10\\\hline
RF-BOY-01 & 1.4.3 & 1.10.2.1.10 / 1.10.2.2.10\\\hline
RF-BOY-02 & 1.4.3 & 1.10.2.1.10 / 1.10.2.2.10\\\hline
RF-BOY-03 & 1.4.3 & 1.10.2.1.10 / 1.10.2.2.10\\\hline
RF-BOY-04 & 1.4.3 & 1.10.2.1.10 / 1.10.2.2.10\\\hline
RF-BOY-05 & 1.4.3 & 1.10.2.1.10 / 1.10.2.2.10\\\hline
RF-BOY-06 & 1.4.3 & 1.10.2.1.10 / 1.10.2.2.10\\\hline
RF-BOY-07 & 1.4.3 & 1.10.2.1.10 / 1.10.2.2.10\\\hline
RF-FAC-01 & 1.6.2 & 1.10.2.1.10 / 1.10.2.2.10\\\hline
RF-FAC-02 & 1.6.2 & 1.10.2.1.10 / 1.10.2.2.10\\\hline
RF-FAC-03 & 1.6.2 & 1.10.2.1.10 / 1.10.2.2.10\\\hline
RF-FAC-04 & 1.6.2 & 1.10.2.1.10 / 1.10.2.2.10\\\hline
RF-FAC-05 & 1.6.2 & 1.10.2.1.10 / 1.10.2.2.10\\\hline
RF-FAC-06 & 1.6.2 & 1.10.2.1.10 / 1.10.2.2.10\\\hline
RF-FAC-07 & 1.6.2 & 1.10.2.1.10 / 1.10.2.2.10\\\hline
RF-FAC-08 & 1.6.2 & 1.10.2.1.10 / 1.10.2.2.10\\\hline
RF-FAC-09 & 1.6.2 & 1.10.2.1.10 / 1.10.2.2.10\\\hline
RF-ADM-01 & 1.6.1; 1.6.2; 1.3 & 1.10.2.1.10 / 1.10.2.2.10\\\hline
RF-ADM-02 & 1.6.1; 1.6.2; 1.3 & 1.10.2.1.10 / 1.10.2.2.10\\\hline
RF-ADM-03 & 1.6.1; 1.6.2; 1.3 & 1.10.2.1.10 / 1.10.2.2.10\\\hline
RF-ADM-04 & 1.6.1; 1.6.2; 1.3 & 1.10.2.1.10 / 1.10.2.2.10\\\hline
RF-ADM-05 & 1.6.1; 1.6.2; 1.3 & 1.10.2.1.10 / 1.10.2.2.10\\\hline
RF-ADM-06 & 1.6.1; 1.6.2; 1.3 & 1.10.2.1.10 / 1.10.2.2.10\\\hline
RF-ADM-07 & 1.6.1; 1.6.2; 1.3 & 1.10.2.1.10 / 1.10.2.2.10\\\hline
RF-ADM-08 & 1.6.1; 1.6.2; 1.3 & 1.10.2.1.10 / 1.10.2.2.10\\\hline
RF-ADM-09 & 1.6.1; 1.6.2; 1.3 & 1.10.2.1.10 / 1.10.2.2.10\\\hline
RF-ADM-10 & 1.6.1; 1.6.2; 1.3 & 1.10.2.1.10 / 1.10.2.2.10\\\hline
RF-ADM-11 & 1.6.1; 1.6.2; 1.3 & 1.10.2.1.10 / 1.10.2.2.10\\\hline
RF-ADM-12 & 1.6.1; 1.6.2; 1.3 & 1.10.2.1.10 / 1.10.2.2.10\\\hline
RF-ADM-13 & 1.6.1; 1.6.2; 1.3 & 1.10.2.1.10 / 1.10.2.2.10\\\hline
RF-ADM-14 & 1.6.1; 1.6.2; 1.3 & 1.10.2.1.10 / 1.10.2.2.10\\\hline
RF-ADM-15 & 1.6.1; 1.6.2; 1.3 & 1.10.2.1.10 / 1.10.2.2.10\\\hline
RF-ADM-16 & 1.6.1; 1.6.2; 1.3 & 1.10.2.1.10 / 1.10.2.2.10\\\hline
RF-ADM-17 & 1.6.1; 1.6.2; 1.3 & 1.10.2.1.10 / 1.10.2.2.10\\\hline
RF-ADM-18 & 1.6.1; 1.6.2; 1.3 & 1.10.2.1.10 / 1.10.2.2.10\\\hline
RF-ADM-19 & 1.6.1; 1.6.2; 1.3 & 1.10.2.1.10 / 1.10.2.2.10\\\hline
RF-ADM-20 & 1.6.1; 1.6.2; 1.3 & 1.10.2.1.10 / 1.10.2.2.10\\\hline
RF-ADM-21 & 1.6.1; 1.6.2; 1.3 & 1.10.2.1.10 / 1.10.2.2.10\\\hline
RF-ADM-22 & 1.6.1; 1.6.2; 1.3 & 1.10.2.1.10 / 1.10.2.2.10\\\hline
RF-ADM-23 & 1.6.1; 1.6.2; 1.3 & 1.10.2.1.10 / 1.10.2.2.10\\\hline
RF-ADM-24 & 1.6.1; 1.6.2; 1.3 & 1.10.2.1.10 / 1.10.2.2.10\\\hline
RF-ADM-25 & 1.6.1; 1.6.2; 1.3 & 1.10.2.1.10 / 1.10.2.2.10\\\hline
RF-ADM-26 & 1.6.1; 1.6.2; 1.3 & 1.10.2.1.10 / 1.10.2.2.10\\\hline
RF-ADM-27 & 1.6.1; 1.6.2; 1.3 & 1.10.2.1.10 / 1.10.2.2.10\\\hline
RF-ADM-28 & 1.6.1; 1.6.2; 1.3 & 1.10.2.1.10 / 1.10.2.2.10\\\hline
RF-ADM-29 & 1.6.1; 1.6.2; 1.3 & 1.10.2.1.10 / 1.10.2.2.10\\\hline
RF-ADM-30 & 1.6.1; 1.6.2; 1.3 & 1.10.2.1.10 / 1.10.2.2.10\\\hline
RF-ADM-31 & 1.6.1; 1.6.2; 1.3 & 1.10.2.1.10 / 1.10.2.2.10\\\hline
RF-ADM-32 & 1.6.1; 1.6.2; 1.3 & 1.10.2.1.10 / 1.10.2.2.10\\\hline
RF-ADM-33 & 1.6.1; 1.6.2; 1.3 & 1.10.2.1.10 / 1.10.2.2.10\\\hline
RF-ADM-34 & 1.6.1; 1.6.2; 1.3 & 1.10.2.1.10 / 1.10.2.2.10\\\hline
RF-ADM-35 & 1.6.1; 1.6.2; 1.3 & 1.10.2.1.10 / 1.10.2.2.10\\\hline
RF-ADM-36 & 1.6.1; 1.6.2; 1.3 & 1.10.2.1.10 / 1.10.2.2.10\\\hline
RF-ADM-37 & 1.6.1; 1.6.2; 1.3 & 1.10.2.1.10 / 1.10.2.2.10\\\hline
RF-AUT-01 & 1.7.3; 1.4.1.2 & 1.10.2.1.10 / 1.10.2.2.10\\\hline
RF-AUT-02 & 1.7.3; 1.4.1.2 & 1.10.2.1.10 / 1.10.2.2.10\\\hline
RF-AUT-03 & 1.7.3; 1.4.1.2 & 1.10.2.1.10 / 1.10.2.2.10\\\hline
RF-AUT-04 & 1.7.3; 1.4.1.2 & 1.10.2.1.10 / 1.10.2.2.10\\\hline
RF-AUT-05 & 1.7.3; 1.4.1.2 & 1.10.2.1.10 / 1.10.2.2.10\\\hline
RF-AUT-06 & 1.7.3; 1.4.1.2 & 1.10.2.1.10 / 1.10.2.2.10\\\hline
RF-AUT-07 & 1.7.3; 1.4.1.2 & 1.10.2.1.10 / 1.10.2.2.10\\\hline
RF-AUT-08 & 1.7.3; 1.4.1.2 & 1.10.2.1.10 / 1.10.2.2.10\\\hline
RF-AUT-09 & 1.7.3; 1.4.1.2 & 1.10.2.1.10 / 1.10.2.2.10\\\hline
RF-AUT-10 & 1.7.3; 1.4.1.2 & 1.10.2.1.10 / 1.10.2.2.10\\\hline
RF-AUT-11 & 1.7.3; 1.4.1.2 & 1.10.2.1.10 / 1.10.2.2.10\\\hline
RNF-DES-01 & 1.10.2 & 1.10.2.1.10 / 1.10.2.2.10\\\hline
RNF-DES-02 & 1.10.2 & 1.10.2.1.10 / 1.10.2.2.10\\\hline
RNF-DES-03 & 1.10.2 & 1.10.2.1.10 / 1.10.2.2.10\\\hline
RNF-DES-04 & 1.10.2 & 1.10.2.1.10 / 1.10.2.2.10\\\hline
RNF-DES-05 & 1.10.2 & 1.10.2.1.10 / 1.10.2.2.10\\\hline
RNF-DES-06 & 1.10.2 & 1.10.2.1.10 / 1.10.2.2.10\\\hline
RNF-DES-07 & 1.10.2 & 1.10.2.1.10 / 1.10.2.2.10\\\hline
RNF-DES-08 & 1.10.2 & 1.10.2.1.10 / 1.10.2.2.10\\\hline
RNF-DES-09 & 1.10.2 & 1.10.2.1.10 / 1.10.2.2.10\\\hline
RNF-DES-10 & 1.10.2 & 1.10.2.1.10 / 1.10.2.2.10\\\hline
RNF-DES-11 & 1.10.2 & 1.10.2.1.10 / 1.10.2.2.10\\\hline
RNF-DES-12 & 1.10.2 & 1.10.2.1.10 / 1.10.2.2.10\\\hline
RNF-CONT-01 & 1.2.2.5; 1.10.2; 1.12.4 & 1.10.2.1.10 / 1.10.2.2.10\\\hline
RNF-CONT-02 & 1.2.2.5; 1.10.2; 1.12.4 & 1.10.2.1.10 / 1.10.2.2.10\\\hline
RNF-CONT-03 & 1.2.2.5; 1.10.2; 1.12.4 & 1.10.2.1.10 / 1.10.2.2.10\\\hline
RNF-CONT-04 & 1.2.2.5; 1.10.2; 1.12.4 & 1.10.2.1.10 / 1.10.2.2.10\\\hline
RNF-CONT-05 & 1.2.2.5; 1.10.2; 1.12.4 & 1.10.2.1.10 / 1.10.2.2.10\\\hline
RNF-CONT-06 & 1.2.2.5; 1.10.2; 1.12.4 & 1.10.2.1.10 / 1.10.2.2.10\\\hline
RNF-SEG-01 & 1.3; 1.10.2 & 1.10.2.1.10 / 1.10.2.2.10\\\hline
RNF-SEG-02 & 1.3; 1.10.2 & 1.10.2.1.10 / 1.10.2.2.10\\\hline
RNF-SEG-03 & 1.3; 1.10.2 & 1.10.2.1.10 / 1.10.2.2.10\\\hline
RNF-SEG-04 & 1.3; 1.10.2 & 1.10.2.1.10 / 1.10.2.2.10\\\hline
RNF-SEG-05 & 1.3; 1.10.2 & 1.10.2.1.10 / 1.10.2.2.10\\\hline
RNF-SEG-06 & 1.3; 1.10.2 & 1.10.2.1.10 / 1.10.2.2.10\\\hline
RNF-SEG-07 & 1.3; 1.10.2 & 1.10.2.1.10 / 1.10.2.2.10\\\hline
RNF-USA-01 & 1.7.3; 1.10.2 & 1.10.2.1.10 / 1.10.2.2.10\\\hline
RNF-USA-02 & 1.7.3; 1.10.2 & 1.10.2.1.10 / 1.10.2.2.10\\\hline
RNF-USA-03 & 1.7.3; 1.10.2 & 1.10.2.1.10 / 1.10.2.2.10\\\hline
RNF-USA-04 & 1.7.3; 1.10.2 & 1.10.2.1.10 / 1.10.2.2.10\\\hline
RNF-IDM-01 & 1.3.1 & 1.10.2.1.10 / 1.10.2.2.10\\\hline
RNF-IDM-02 & 1.3.1 & 1.10.2.1.10 / 1.10.2.2.10\\\hline
RNF-ESCAL-01 & 1.2.1; 1.10.2 & 1.10.2.1.10 / 1.10.2.2.10\\\hline
RNF-ESCAL-02 & 1.2.1; 1.10.2 & 1.10.2.1.10 / 1.10.2.2.10\\\hline
RNF-INT-01 & 1.7.1; 1.7.4 & 1.10.2.1.10 / 1.10.2.2.10\\\hline
RNF-INT-02 & 1.7.1; 1.7.4 & 1.10.2.1.10 / 1.10.2.2.10\\\hline
RNF-OPE-01 & 1.12 & 1.10.2.1.10 / 1.10.2.2.10\\\hline
RNF-OPE-02 & 1.12 & 1.10.2.1.10 / 1.10.2.2.10\\\hline
RNF-OPE-03 & 1.12 & 1.10.2.1.10 / 1.10.2.2.10\\\hline
RNF-RET-01 & 1.8.5; 1.3.2 & 1.10.2.1.10 / 1.10.2.2.10\\\hline
RNF-RET-02 & 1.8.5; 1.3.2 & 1.10.2.1.10 / 1.10.2.2.10\\\hline
RNF-RET-03 & 1.8.5; 1.3.2 & 1.10.2.1.10 / 1.10.2.2.10\\\hline
RNF-RET-04 & 1.8.5; 1.3.2 & 1.10.2.1.10 / 1.10.2.2.10\\\hline
RNF-RET-05 & 1.8.5; 1.3.2 & 1.10.2.1.10 / 1.10.2.2.10\\\hline
RNF-RET-06 & 1.8.5; 1.3.2 & 1.10.2.1.10 / 1.10.2.2.10\\\hline
RNF-RET-07 & 1.8.5; 1.3.2 & 1.10.2.1.10 / 1.10.2.2.10\\\hline
RNF-RET-08 & 1.8.5; 1.3.2 & 1.10.2.1.10 / 1.10.2.2.10\\\hline
RNF-SOP-01 & 1.12 & 1.10.2.1.10 / 1.10.2.2.10\\\hline
RNF-SOP-02 & 1.12 & 1.10.2.1.10 / 1.10.2.2.10\\\hline
RNF-SOP-03 & 1.12 & 1.10.2.1.10 / 1.10.2.2.10\\\hline
RNF-DISP-01 & 1.2.4; 1.12.4 & 1.10.2.1.10 / 1.10.2.2.10\\\hline
RNF-MIG-01 & 1.8; 1.10.1 & 1.10.2.1.10 / 1.10.2.2.10\\\hline
RNF-RED-01 & 1.2.3 & 1.10.2.1.10 / 1.10.2.2.10\\\hline
RNF-TER-01 & 1.12.9 & 1.10.2.1.10 / 1.10.2.2.10\\\hline
RNF-VEN-01 & 1.11; 1.1.2 & 1.10.2.1.10 / 1.10.2.2.10\\\hline
RNF-CAP-01 & 1.11.7; 1.12.10 & 1.10.2.1.10 / 1.10.2.2.10\\\hline
RNF-CUMP-01 & 1.1.5; 1.3 & 1.10.2.1.10 / 1.10.2.2.10\\\hline
RT-02.01 & 1.1.5; 1.3 & Protocolo integral 1.10.2.1.1 / 1.10.2.2.1 y evidencia específica de la cláusula\\\hline
RT-02.02 & 1.1.5; 1.3 & Protocolo integral 1.10.2.1.1 / 1.10.2.2.1 y evidencia específica de la cláusula\\\hline
RT-02.03 & 1.1.5; 1.3 & Protocolo integral 1.10.2.1.1 / 1.10.2.2.1 y evidencia específica de la cláusula\\\hline
RT-02.04 & 1.1.5; 1.3 & Protocolo integral 1.10.2.1.1 / 1.10.2.2.1 y evidencia específica de la cláusula\\\hline
RT-02.05 & 1.1.5; 1.3 & Protocolo integral 1.10.2.1.1 / 1.10.2.2.1 y evidencia específica de la cláusula\\\hline
RT-02.06 & 1.1.5; 1.3 & Protocolo integral 1.10.2.1.1 / 1.10.2.2.1 y evidencia específica de la cláusula\\\hline
RT-02.07 & 1.1.5; 1.3 & Protocolo integral 1.10.2.1.1 / 1.10.2.2.1 y evidencia específica de la cláusula\\\hline
RT-02.08 & 1.1.5; 1.3 & Protocolo integral 1.10.2.1.1 / 1.10.2.2.1 y evidencia específica de la cláusula\\\hline
RT-02.09 & 1.1.5; 1.3 & Protocolo integral 1.10.2.1.1 / 1.10.2.2.1 y evidencia específica de la cláusula\\\hline
RT-02.10 & 1.1.5; 1.3 & Protocolo integral 1.10.2.1.1 / 1.10.2.2.1 y evidencia específica de la cláusula\\\hline
RT-02.11 & 1.1.5; 1.3 & Protocolo integral 1.10.2.1.1 / 1.10.2.2.1 y evidencia específica de la cláusula\\\hline
RT-02.12 & 1.1.5; 1.3 & Protocolo integral 1.10.2.1.1 / 1.10.2.2.1 y evidencia específica de la cláusula\\\hline
RT-02.13 & 1.1.5; 1.3 & Protocolo integral 1.10.2.1.1 / 1.10.2.2.1 y evidencia específica de la cláusula\\\hline
RT-02.14 & Excluido Deseable BT & No procede; conservar exclusión\\\hline
RT-03.01 & 1.2.1; 1.7.1 & Protocolo integral 1.10.2.1.1 / 1.10.2.2.1 y evidencia específica de la cláusula\\\hline
RT-03.02 & 1.2.1; 1.7.1 & Protocolo integral 1.10.2.1.1 / 1.10.2.2.1 y evidencia específica de la cláusula\\\hline
RT-03.03 & 1.2.1; 1.7.1 & Protocolo integral 1.10.2.1.1 / 1.10.2.2.1 y evidencia específica de la cláusula\\\hline
RT-03.04 & 1.2.1; 1.7.1 & Protocolo integral 1.10.2.1.1 / 1.10.2.2.1 y evidencia específica de la cláusula\\\hline
RT-03.05 & 1.2.1; 1.7.1 & Protocolo integral 1.10.2.1.1 / 1.10.2.2.1 y evidencia específica de la cláusula\\\hline
RT-03.06 & 1.2.1; 1.7.1 & Protocolo integral 1.10.2.1.1 / 1.10.2.2.1 y evidencia específica de la cláusula\\\hline
RT-03.07 & 1.2.1; 1.7.1 & Protocolo integral 1.10.2.1.1 / 1.10.2.2.1 y evidencia específica de la cláusula\\\hline
RT-03.08 & Excluido Deseable BT & No procede; conservar exclusión\\\hline
RT-03.09 & Excluido Deseable BT & No procede; conservar exclusión\\\hline
RT-03.10 & 1.2.1; 1.7.1 & Protocolo integral 1.10.2.1.1 / 1.10.2.2.1 y evidencia específica de la cláusula\\\hline
RT-03.11 & 1.2.1; 1.7.1 & Protocolo integral 1.10.2.1.1 / 1.10.2.2.1 y evidencia específica de la cláusula\\\hline
RT-03.12 & 1.2.1; 1.7.1 & Protocolo integral 1.10.2.1.1 / 1.10.2.2.1 y evidencia específica de la cláusula\\\hline
RT-03.13 & 1.2.1; 1.7.1 & Protocolo integral 1.10.2.1.1 / 1.10.2.2.1 y evidencia específica de la cláusula\\\hline
RT-03.14 & 1.2.1; 1.7.1 & Protocolo integral 1.10.2.1.1 / 1.10.2.2.1 y evidencia específica de la cláusula\\\hline
RT-03.15 & 1.2.1; 1.7.1 & Protocolo integral 1.10.2.1.1 / 1.10.2.2.1 y evidencia específica de la cláusula\\\hline
RT-03.16 & 1.2.1; 1.7.1 & Protocolo integral 1.10.2.1.1 / 1.10.2.2.1 y evidencia específica de la cláusula\\\hline
RT-03.17 & 1.2.1; 1.7.1 & Protocolo integral 1.10.2.1.1 / 1.10.2.2.1 y evidencia específica de la cláusula\\\hline
RT-03.18 & 1.2.1; 1.7.1 & Protocolo integral 1.10.2.1.1 / 1.10.2.2.1 y evidencia específica de la cláusula\\\hline
RT-03.19 & Excluido Deseable BT & No procede; conservar exclusión\\\hline
RT-03.20 & 1.2.1; 1.7.1 & Protocolo integral 1.10.2.1.1 / 1.10.2.2.1 y evidencia específica de la cláusula\\\hline
RT-03.21 & 1.2.1; 1.7.1 & Protocolo integral 1.10.2.1.1 / 1.10.2.2.1 y evidencia específica de la cláusula\\\hline
RT-03.22 & 1.2.1; 1.7.1 & Protocolo integral 1.10.2.1.1 / 1.10.2.2.1 y evidencia específica de la cláusula\\\hline
RT-03.23 & 1.2.1; 1.7.1 & Protocolo integral 1.10.2.1.1 / 1.10.2.2.1 y evidencia específica de la cláusula\\\hline
RT-03.24 & 1.2.1; 1.7.1 & Protocolo integral 1.10.2.1.1 / 1.10.2.2.1 y evidencia específica de la cláusula\\\hline
RT-04.01 & 1.2.2; 1.2.4 & Protocolo integral 1.10.2.1.1 / 1.10.2.2.1 y evidencia específica de la cláusula\\\hline
RT-04.02 & 1.2.2; 1.2.4 & Protocolo integral 1.10.2.1.1 / 1.10.2.2.1 y evidencia específica de la cláusula\\\hline
RT-04.03 & 1.2.2; 1.2.4 & Protocolo integral 1.10.2.1.1 / 1.10.2.2.1 y evidencia específica de la cláusula\\\hline
RT-04.04 & 1.2.2; 1.2.4 & Protocolo integral 1.10.2.1.1 / 1.10.2.2.1 y evidencia específica de la cláusula\\\hline
RT-04.05 & 1.2.2; 1.2.4 & Protocolo integral 1.10.2.1.1 / 1.10.2.2.1 y evidencia específica de la cláusula\\\hline
RT-04.06 & 1.2.2; 1.2.4 & Protocolo integral 1.10.2.1.1 / 1.10.2.2.1 y evidencia específica de la cláusula\\\hline
RT-04.07 & 1.2.2; 1.2.4 & Protocolo integral 1.10.2.1.1 / 1.10.2.2.1 y evidencia específica de la cláusula\\\hline
RT-04.08 & 1.2.2; 1.2.4 & Protocolo integral 1.10.2.1.1 / 1.10.2.2.1 y evidencia específica de la cláusula\\\hline
RT-04.09 & 1.2.2; 1.2.4 & Protocolo integral 1.10.2.1.1 / 1.10.2.2.1 y evidencia específica de la cláusula\\\hline
RT-04.10 & 1.2.2; 1.2.4 & Protocolo integral 1.10.2.1.1 / 1.10.2.2.1 y evidencia específica de la cláusula\\\hline
RT-04.11 & 1.2.2; 1.2.4 & Protocolo integral 1.10.2.1.1 / 1.10.2.2.1 y evidencia específica de la cláusula\\\hline
RT-04.12 & 1.2.2; 1.2.4 & Protocolo integral 1.10.2.1.1 / 1.10.2.2.1 y evidencia específica de la cláusula\\\hline
RT-04.13 & 1.2.2; 1.2.4 & Protocolo integral 1.10.2.1.1 / 1.10.2.2.1 y evidencia específica de la cláusula\\\hline
RT-04.14 & Excluido Deseable BT & No procede; conservar exclusión\\\hline
RT-05.01 & 1.8 & Protocolo integral 1.10.2.1.1 / 1.10.2.2.1 y evidencia específica de la cláusula\\\hline
RT-05.02 & 1.8 & Protocolo integral 1.10.2.1.1 / 1.10.2.2.1 y evidencia específica de la cláusula\\\hline
RT-05.03 & 1.8 & Protocolo integral 1.10.2.1.1 / 1.10.2.2.1 y evidencia específica de la cláusula\\\hline
RT-05.04 & 1.8 & Protocolo integral 1.10.2.1.1 / 1.10.2.2.1 y evidencia específica de la cláusula\\\hline
RT-05.05 & 1.8 & Protocolo integral 1.10.2.1.1 / 1.10.2.2.1 y evidencia específica de la cláusula\\\hline
RT-05.06 & 1.8 & Protocolo integral 1.10.2.1.1 / 1.10.2.2.1 y evidencia específica de la cláusula\\\hline
RT-05.07 & 1.8 & Protocolo integral 1.10.2.1.1 / 1.10.2.2.1 y evidencia específica de la cláusula\\\hline
RT-05.08 & 1.8 & Protocolo integral 1.10.2.1.1 / 1.10.2.2.1 y evidencia específica de la cláusula\\\hline
RT-05.09 & 1.8 & Protocolo integral 1.10.2.1.1 / 1.10.2.2.1 y evidencia específica de la cláusula\\\hline
RT-05.10 & 1.8 & Protocolo integral 1.10.2.1.1 / 1.10.2.2.1 y evidencia específica de la cláusula\\\hline
RT-05.11 & 1.8 & Protocolo integral 1.10.2.1.1 / 1.10.2.2.1 y evidencia específica de la cláusula\\\hline
RT-05.12 & 1.8 & Protocolo integral 1.10.2.1.1 / 1.10.2.2.1 y evidencia específica de la cláusula\\\hline
RT-05.13 & 1.8 & Protocolo integral 1.10.2.1.1 / 1.10.2.2.1 y evidencia específica de la cláusula\\\hline
RT-05.14 & 1.8 & Protocolo integral 1.10.2.1.1 / 1.10.2.2.1 y evidencia específica de la cláusula\\\hline
RT-05.15 & 1.8 & Protocolo integral 1.10.2.1.1 / 1.10.2.2.1 y evidencia específica de la cláusula\\\hline
RT-05.16 & 1.8 & Protocolo integral 1.10.2.1.1 / 1.10.2.2.1 y evidencia específica de la cláusula\\\hline
RT-05.17 & 1.8 & Protocolo integral 1.10.2.1.1 / 1.10.2.2.1 y evidencia específica de la cláusula\\\hline
RT-05.18 & 1.8 & Protocolo integral 1.10.2.1.1 / 1.10.2.2.1 y evidencia específica de la cláusula\\\hline
RT-05.19 & 1.8 & Protocolo integral 1.10.2.1.1 / 1.10.2.2.1 y evidencia específica de la cláusula\\\hline
RT-05.20 & 1.8 & Protocolo integral 1.10.2.1.1 / 1.10.2.2.1 y evidencia específica de la cláusula\\\hline
RT-05.21 & 1.8 & Protocolo integral 1.10.2.1.1 / 1.10.2.2.1 y evidencia específica de la cláusula\\\hline
RT-05.22 & 1.8 & Protocolo integral 1.10.2.1.1 / 1.10.2.2.1 y evidencia específica de la cláusula\\\hline
RT-05.23 & 1.8 & Protocolo integral 1.10.2.1.1 / 1.10.2.2.1 y evidencia específica de la cláusula\\\hline
RT-05.24 & Excluido Deseable BT & No procede; conservar exclusión\\\hline
RT-05.25 & 1.8 & Protocolo integral 1.10.2.1.1 / 1.10.2.2.1 y evidencia específica de la cláusula\\\hline
RT-05.26 & 1.8 & Protocolo integral 1.10.2.1.1 / 1.10.2.2.1 y evidencia específica de la cláusula\\\hline
RT-05.27 & 1.8 & Protocolo integral 1.10.2.1.1 / 1.10.2.2.1 y evidencia específica de la cláusula\\\hline
RT-05.28 & 1.8 & Protocolo integral 1.10.2.1.1 / 1.10.2.2.1 y evidencia específica de la cláusula\\\hline
RT-05.29 & 1.8 & Protocolo integral 1.10.2.1.1 / 1.10.2.2.1 y evidencia específica de la cláusula\\\hline
RT-05.30 & Excluido Deseable BT & No procede; conservar exclusión\\\hline
RT-06.01 & 1.7.1; 1.7.3; 1.4; 1.5; 1.6 & Protocolo integral 1.10.2.1.1 / 1.10.2.2.1 y evidencia específica de la cláusula\\\hline
RT-06.02 & 1.7.1; 1.7.3; 1.4; 1.5; 1.6 & Protocolo integral 1.10.2.1.1 / 1.10.2.2.1 y evidencia específica de la cláusula\\\hline
RT-06.03 & 1.7.1; 1.7.3; 1.4; 1.5; 1.6 & Protocolo integral 1.10.2.1.1 / 1.10.2.2.1 y evidencia específica de la cláusula\\\hline
RT-06.04 & 1.7.1; 1.7.3; 1.4; 1.5; 1.6 & Protocolo integral 1.10.2.1.1 / 1.10.2.2.1 y evidencia específica de la cláusula\\\hline
RT-06.05 & 1.7.1; 1.7.3; 1.4; 1.5; 1.6 & Protocolo integral 1.10.2.1.1 / 1.10.2.2.1 y evidencia específica de la cláusula\\\hline
RT-06.06 & 1.7.1; 1.7.3; 1.4; 1.5; 1.6 & Protocolo integral 1.10.2.1.1 / 1.10.2.2.1 y evidencia específica de la cláusula\\\hline
RT-06.07 & 1.7.1; 1.7.3; 1.4; 1.5; 1.6 & Protocolo integral 1.10.2.1.1 / 1.10.2.2.1 y evidencia específica de la cláusula\\\hline
RT-06.08 & 1.7.1; 1.7.3; 1.4; 1.5; 1.6 & Protocolo integral 1.10.2.1.1 / 1.10.2.2.1 y evidencia específica de la cláusula\\\hline
RT-06.09 & 1.7.1; 1.7.3; 1.4; 1.5; 1.6 & Protocolo integral 1.10.2.1.1 / 1.10.2.2.1 y evidencia específica de la cláusula\\\hline
RT-06.10 & 1.7.1; 1.7.3; 1.4; 1.5; 1.6 & Protocolo integral 1.10.2.1.1 / 1.10.2.2.1 y evidencia específica de la cláusula\\\hline
RT-06.11 & 1.7.1; 1.7.3; 1.4; 1.5; 1.6 & Protocolo integral 1.10.2.1.1 / 1.10.2.2.1 y evidencia específica de la cláusula\\\hline
RT-06.12 & Excluido Deseable BT & No procede; conservar exclusión\\\hline
RT-06.13 & 1.7.1; 1.7.3; 1.4; 1.5; 1.6 & Protocolo integral 1.10.2.1.1 / 1.10.2.2.1 y evidencia específica de la cláusula\\\hline
RT-06.14 & 1.7.1; 1.7.3; 1.4; 1.5; 1.6 & Protocolo integral 1.10.2.1.1 / 1.10.2.2.1 y evidencia específica de la cláusula\\\hline
RT-06.15 & Excluido Deseable BT & No procede; conservar exclusión\\\hline
RT-06.16 & 1.7.1; 1.7.3; 1.4; 1.5; 1.6 & Protocolo integral 1.10.2.1.1 / 1.10.2.2.1 y evidencia específica de la cláusula\\\hline
RT-06.17 & 1.7.1; 1.7.3; 1.4; 1.5; 1.6 & Protocolo integral 1.10.2.1.1 / 1.10.2.2.1 y evidencia específica de la cláusula\\\hline
RT-06.18 & 1.7.1; 1.7.3; 1.4; 1.5; 1.6 & Protocolo integral 1.10.2.1.1 / 1.10.2.2.1 y evidencia específica de la cláusula\\\hline
RT-06.19 & 1.7.1; 1.7.3; 1.4; 1.5; 1.6 & Protocolo integral 1.10.2.1.1 / 1.10.2.2.1 y evidencia específica de la cláusula\\\hline
RT-06.20 & 1.7.1; 1.7.3; 1.4; 1.5; 1.6 & Protocolo integral 1.10.2.1.1 / 1.10.2.2.1 y evidencia específica de la cláusula\\\hline
RT-06.21 & 1.7.1; 1.7.3; 1.4; 1.5; 1.6 & Protocolo integral 1.10.2.1.1 / 1.10.2.2.1 y evidencia específica de la cláusula\\\hline
RT-06.22 & 1.7.1; 1.7.3; 1.4; 1.5; 1.6 & Protocolo integral 1.10.2.1.1 / 1.10.2.2.1 y evidencia específica de la cláusula\\\hline
RT-06.23 & 1.7.1; 1.7.3; 1.4; 1.5; 1.6 & Protocolo integral 1.10.2.1.1 / 1.10.2.2.1 y evidencia específica de la cláusula\\\hline
RT-06.24 & 1.7.1; 1.7.3; 1.4; 1.5; 1.6 & Protocolo integral 1.10.2.1.1 / 1.10.2.2.1 y evidencia específica de la cláusula\\\hline
RT-06.25 & 1.7.1; 1.7.3; 1.4; 1.5; 1.6 & Protocolo integral 1.10.2.1.1 / 1.10.2.2.1 y evidencia específica de la cláusula\\\hline
RT-06.26 & 1.7.1; 1.7.3; 1.4; 1.5; 1.6 & Protocolo integral 1.10.2.1.1 / 1.10.2.2.1 y evidencia específica de la cláusula\\\hline
RT-06.27 & 1.7.1; 1.7.3; 1.4; 1.5; 1.6 & Protocolo integral 1.10.2.1.1 / 1.10.2.2.1 y evidencia específica de la cláusula\\\hline
RT-06.28 & 1.7.1; 1.7.3; 1.4; 1.5; 1.6 & Protocolo integral 1.10.2.1.1 / 1.10.2.2.1 y evidencia específica de la cláusula\\\hline
RT-06.29 & 1.7.1; 1.7.3; 1.4; 1.5; 1.6 & Protocolo integral 1.10.2.1.1 / 1.10.2.2.1 y evidencia específica de la cláusula\\\hline
RT-06.30 & 1.7.1; 1.7.3; 1.4; 1.5; 1.6 & Protocolo integral 1.10.2.1.1 / 1.10.2.2.1 y evidencia específica de la cláusula\\\hline
RT-06.31 & 1.7.1; 1.7.3; 1.4; 1.5; 1.6 & Protocolo integral 1.10.2.1.1 / 1.10.2.2.1 y evidencia específica de la cláusula\\\hline
RT-06.32 & 1.7.1; 1.7.3; 1.4; 1.5; 1.6 & Protocolo integral 1.10.2.1.1 / 1.10.2.2.1 y evidencia específica de la cláusula\\\hline
RT-06.33 & 1.7.1; 1.7.3; 1.4; 1.5; 1.6 & Protocolo integral 1.10.2.1.1 / 1.10.2.2.1 y evidencia específica de la cláusula\\\hline
RT-06.34 & Excluido Deseable BT & No procede; conservar exclusión\\\hline
RT-07.01 & 1.3.2; 1.7.1 & Protocolo integral 1.10.2.1.1 / 1.10.2.2.1 y evidencia específica de la cláusula\\\hline
RT-07.02 & 1.3.2; 1.7.1 & Protocolo integral 1.10.2.1.1 / 1.10.2.2.1 y evidencia específica de la cláusula\\\hline
RT-07.03 & 1.3.2; 1.7.1 & Protocolo integral 1.10.2.1.1 / 1.10.2.2.1 y evidencia específica de la cláusula\\\hline
RT-07.04 & 1.3.2; 1.7.1 & Protocolo integral 1.10.2.1.1 / 1.10.2.2.1 y evidencia específica de la cláusula\\\hline
RT-07.05 & 1.3.2; 1.7.1 & Protocolo integral 1.10.2.1.1 / 1.10.2.2.1 y evidencia específica de la cláusula\\\hline
RT-07.06 & 1.3.2; 1.7.1 & Protocolo integral 1.10.2.1.1 / 1.10.2.2.1 y evidencia específica de la cláusula\\\hline
RT-07.07 & 1.3.2; 1.7.1 & Protocolo integral 1.10.2.1.1 / 1.10.2.2.1 y evidencia específica de la cláusula\\\hline
RT-07.08 & 1.3.2; 1.7.1 & Protocolo integral 1.10.2.1.1 / 1.10.2.2.1 y evidencia específica de la cláusula\\\hline
RT-07.09 & 1.3.2; 1.7.1 & Protocolo integral 1.10.2.1.1 / 1.10.2.2.1 y evidencia específica de la cláusula\\\hline
RT-07.10 & 1.3.2; 1.7.1 & Protocolo integral 1.10.2.1.1 / 1.10.2.2.1 y evidencia específica de la cláusula\\\hline
RT-07.11 & 1.3.2; 1.7.1 & Protocolo integral 1.10.2.1.1 / 1.10.2.2.1 y evidencia específica de la cláusula\\\hline
RT-07.12 & 1.3.2; 1.7.1 & Protocolo integral 1.10.2.1.1 / 1.10.2.2.1 y evidencia específica de la cláusula\\\hline
RT-07.13 & 1.3.2; 1.7.1 & Protocolo integral 1.10.2.1.1 / 1.10.2.2.1 y evidencia específica de la cláusula\\\hline
RT-07.14 & 1.3.2; 1.7.1 & Protocolo integral 1.10.2.1.1 / 1.10.2.2.1 y evidencia específica de la cláusula\\\hline
RT-08.01 & 1.2.3; 1.3.1 & Protocolo integral 1.10.2.1.1 / 1.10.2.2.1 y evidencia específica de la cláusula\\\hline
RT-08.02 & 1.2.3; 1.3.1 & Protocolo integral 1.10.2.1.1 / 1.10.2.2.1 y evidencia específica de la cláusula\\\hline
RT-08.03 & 1.2.3; 1.3.1 & Protocolo integral 1.10.2.1.1 / 1.10.2.2.1 y evidencia específica de la cláusula\\\hline
RT-08.04 & 1.2.3; 1.3.1 & Protocolo integral 1.10.2.1.1 / 1.10.2.2.1 y evidencia específica de la cláusula\\\hline
RT-08.05 & 1.2.3; 1.3.1 & Protocolo integral 1.10.2.1.1 / 1.10.2.2.1 y evidencia específica de la cláusula\\\hline
RT-08.06 & 1.2.3; 1.3.1 & Protocolo integral 1.10.2.1.1 / 1.10.2.2.1 y evidencia específica de la cláusula\\\hline
RT-08.07 & 1.2.3; 1.3.1 & Protocolo integral 1.10.2.1.1 / 1.10.2.2.1 y evidencia específica de la cláusula\\\hline
RT-08.08 & 1.2.3; 1.3.1 & Protocolo integral 1.10.2.1.1 / 1.10.2.2.1 y evidencia específica de la cláusula\\\hline
RT-08.09 & 1.2.3; 1.3.1 & Protocolo integral 1.10.2.1.1 / 1.10.2.2.1 y evidencia específica de la cláusula\\\hline
RT-08.10 & 1.2.3; 1.3.1 & Protocolo integral 1.10.2.1.1 / 1.10.2.2.1 y evidencia específica de la cláusula\\\hline
RT-08.11 & 1.2.3; 1.3.1 & Protocolo integral 1.10.2.1.1 / 1.10.2.2.1 y evidencia específica de la cláusula\\\hline
RT-08.12 & 1.2.3; 1.3.1 & Protocolo integral 1.10.2.1.1 / 1.10.2.2.1 y evidencia específica de la cláusula\\\hline
RT-08.13 & 1.2.3; 1.3.1 & Protocolo integral 1.10.2.1.1 / 1.10.2.2.1 y evidencia específica de la cláusula\\\hline
RT-08.14 & 1.2.3; 1.3.1 & Protocolo integral 1.10.2.1.1 / 1.10.2.2.1 y evidencia específica de la cláusula\\\hline
RT-08.15 & Excluido Deseable BT & No procede; conservar exclusión\\\hline
RT-08.16 & 1.2.3; 1.3.1 & Protocolo integral 1.10.2.1.1 / 1.10.2.2.1 y evidencia específica de la cláusula\\\hline
RT-08.17 & 1.2.3; 1.3.1 & Protocolo integral 1.10.2.1.1 / 1.10.2.2.1 y evidencia específica de la cláusula\\\hline
RT-08.18 & 1.2.3; 1.3.1 & Protocolo integral 1.10.2.1.1 / 1.10.2.2.1 y evidencia específica de la cláusula\\\hline
RT-08.19 & Excluido Deseable BT & No procede; conservar exclusión\\\hline
RT-09.01 & 1.10.2 & Protocolo integral 1.10.2.1.1 / 1.10.2.2.1 y evidencia específica de la cláusula\\\hline
RT-09.02 & 1.10.2 & Protocolo integral 1.10.2.1.1 / 1.10.2.2.1 y evidencia específica de la cláusula\\\hline
RT-09.03 & 1.10.2 & Protocolo integral 1.10.2.1.1 / 1.10.2.2.1 y evidencia específica de la cláusula\\\hline
RT-09.04 & 1.10.2 & Protocolo integral 1.10.2.1.1 / 1.10.2.2.1 y evidencia específica de la cláusula\\\hline
RT-09.05 & 1.10.2 & Protocolo integral 1.10.2.1.1 / 1.10.2.2.1 y evidencia específica de la cláusula\\\hline
RT-09.06 & 1.10.2 & Protocolo integral 1.10.2.1.1 / 1.10.2.2.1 y evidencia específica de la cláusula\\\hline
RT-09.07 & 1.10.2 & Protocolo integral 1.10.2.1.1 / 1.10.2.2.1 y evidencia específica de la cláusula\\\hline
RT-09.08 & 1.10.2 & Protocolo integral 1.10.2.1.1 / 1.10.2.2.1 y evidencia específica de la cláusula\\\hline
RT-09.09 & 1.10.2 & Protocolo integral 1.10.2.1.1 / 1.10.2.2.1 y evidencia específica de la cláusula\\\hline
RT-09.10 & Excluido Deseable BT & No procede; conservar exclusión\\\hline
RT-10.01 & 1.12.3; 1.12.4; 1.12.8 & Protocolo integral 1.10.2.1.1 / 1.10.2.2.1 y evidencia específica de la cláusula\\\hline
RT-10.02 & 1.12.3; 1.12.4; 1.12.8 & Protocolo integral 1.10.2.1.1 / 1.10.2.2.1 y evidencia específica de la cláusula\\\hline
RT-10.03 & 1.12.3; 1.12.4; 1.12.8 & Protocolo integral 1.10.2.1.1 / 1.10.2.2.1 y evidencia específica de la cláusula\\\hline
RT-10.04 & 1.12.3; 1.12.4; 1.12.8 & Protocolo integral 1.10.2.1.1 / 1.10.2.2.1 y evidencia específica de la cláusula\\\hline
RT-10.05 & 1.12.3; 1.12.4; 1.12.8 & Protocolo integral 1.10.2.1.1 / 1.10.2.2.1 y evidencia específica de la cláusula\\\hline
RT-10.06 & 1.12.3; 1.12.4; 1.12.8 & Protocolo integral 1.10.2.1.1 / 1.10.2.2.1 y evidencia específica de la cláusula\\\hline
RT-10.07 & 1.12.3; 1.12.4; 1.12.8 & Protocolo integral 1.10.2.1.1 / 1.10.2.2.1 y evidencia específica de la cláusula\\\hline
RT-10.08 & 1.12.3; 1.12.4; 1.12.8 & Protocolo integral 1.10.2.1.1 / 1.10.2.2.1 y evidencia específica de la cláusula\\\hline
RT-10.09 & 1.12.3; 1.12.4; 1.12.8 & Protocolo integral 1.10.2.1.1 / 1.10.2.2.1 y evidencia específica de la cláusula\\\hline
RT-11.01 & 1.2.4; 1.7.2 & Protocolo integral 1.10.2.1.1 / 1.10.2.2.1 y evidencia específica de la cláusula\\\hline
RT-11.02 & 1.2.4; 1.7.2 & Protocolo integral 1.10.2.1.1 / 1.10.2.2.1 y evidencia específica de la cláusula\\\hline
RT-11.03 & 1.2.4; 1.7.2 & Protocolo integral 1.10.2.1.1 / 1.10.2.2.1 y evidencia específica de la cláusula\\\hline
RT-11.04 & 1.2.4; 1.7.2 & Protocolo integral 1.10.2.1.1 / 1.10.2.2.1 y evidencia específica de la cláusula\\\hline
RT-11.05 & 1.2.4; 1.7.2 & Protocolo integral 1.10.2.1.1 / 1.10.2.2.1 y evidencia específica de la cláusula\\\hline
RT-11.06 & 1.2.4; 1.7.2 & Protocolo integral 1.10.2.1.1 / 1.10.2.2.1 y evidencia específica de la cláusula\\\hline
RT-11.07 & 1.2.4; 1.7.2 & Protocolo integral 1.10.2.1.1 / 1.10.2.2.1 y evidencia específica de la cláusula\\\hline
RT-11.08 & 1.2.4; 1.7.2 & Protocolo integral 1.10.2.1.1 / 1.10.2.2.1 y evidencia específica de la cláusula\\\hline
RT-11.09 & 1.2.4; 1.7.2 & Protocolo integral 1.10.2.1.1 / 1.10.2.2.1 y evidencia específica de la cláusula\\\hline
RT-11.10 & 1.2.4; 1.7.2 & Protocolo integral 1.10.2.1.1 / 1.10.2.2.1 y evidencia específica de la cláusula\\\hline
RT-11.11 & 1.2.4; 1.7.2 & Protocolo integral 1.10.2.1.1 / 1.10.2.2.1 y evidencia específica de la cláusula\\\hline
RT-11.12 & 1.2.4; 1.7.2 & Protocolo integral 1.10.2.1.1 / 1.10.2.2.1 y evidencia específica de la cláusula\\\hline
RT-11.13 & 1.2.4; 1.7.2 & Protocolo integral 1.10.2.1.1 / 1.10.2.2.1 y evidencia específica de la cláusula\\\hline
RT-11.14 & 1.2.4; 1.7.2 & Protocolo integral 1.10.2.1.1 / 1.10.2.2.1 y evidencia específica de la cláusula\\\hline
RT-11.15 & 1.2.4; 1.7.2 & Protocolo integral 1.10.2.1.1 / 1.10.2.2.1 y evidencia específica de la cláusula\\\hline
RT-11.16 & 1.2.4; 1.7.2 & Protocolo integral 1.10.2.1.1 / 1.10.2.2.1 y evidencia específica de la cláusula\\\hline
RT-11.17 & 1.2.4; 1.7.2 & Protocolo integral 1.10.2.1.1 / 1.10.2.2.1 y evidencia específica de la cláusula\\\hline
RT-11.18 & 1.2.4; 1.7.2 & Protocolo integral 1.10.2.1.1 / 1.10.2.2.1 y evidencia específica de la cláusula\\\hline
RT-11.19 & 1.2.4; 1.7.2 & Protocolo integral 1.10.2.1.1 / 1.10.2.2.1 y evidencia específica de la cláusula\\\hline
RT-11.20 & 1.2.4; 1.7.2 & Protocolo integral 1.10.2.1.1 / 1.10.2.2.1 y evidencia específica de la cláusula\\\hline
RT-11.21 & Excluido Deseable BT & No procede; conservar exclusión\\\hline
RT-11.22 & 1.2.4; 1.7.2 & Protocolo integral 1.10.2.1.1 / 1.10.2.2.1 y evidencia específica de la cláusula\\\hline
RT-11.23 & 1.2.4; 1.7.2 & Protocolo integral 1.10.2.1.1 / 1.10.2.2.1 y evidencia específica de la cláusula\\\hline
RT-11.24 & 1.2.4; 1.7.2 & Protocolo integral 1.10.2.1.1 / 1.10.2.2.1 y evidencia específica de la cláusula\\\hline
RT-11.25 & 1.2.4; 1.7.2 & Protocolo integral 1.10.2.1.1 / 1.10.2.2.1 y evidencia específica de la cláusula\\\hline
RT-11.26 & 1.2.4; 1.7.2 & Protocolo integral 1.10.2.1.1 / 1.10.2.2.1 y evidencia específica de la cláusula\\\hline
RT-11.27 & 1.2.4; 1.7.2 & Protocolo integral 1.10.2.1.1 / 1.10.2.2.1 y evidencia específica de la cláusula\\\hline
RT-11.28 & 1.2.4; 1.7.2 & Protocolo integral 1.10.2.1.1 / 1.10.2.2.1 y evidencia específica de la cláusula\\\hline
RT-12.01 & 1.2.4; 1.12.5 & Protocolo integral 1.10.2.1.1 / 1.10.2.2.1 y evidencia específica de la cláusula\\\hline
RT-12.02 & 1.2.4; 1.12.5 & Protocolo integral 1.10.2.1.1 / 1.10.2.2.1 y evidencia específica de la cláusula\\\hline
RT-12.03 & 1.2.4; 1.12.5 & Protocolo integral 1.10.2.1.1 / 1.10.2.2.1 y evidencia específica de la cláusula\\\hline
RT-12.04 & 1.2.4; 1.12.5 & Protocolo integral 1.10.2.1.1 / 1.10.2.2.1 y evidencia específica de la cláusula\\\hline
RT-12.05 & 1.2.4; 1.12.5 & Protocolo integral 1.10.2.1.1 / 1.10.2.2.1 y evidencia específica de la cláusula\\\hline
RT-12.06 & 1.2.4; 1.12.5 & Protocolo integral 1.10.2.1.1 / 1.10.2.2.1 y evidencia específica de la cláusula\\\hline
RT-12.07 & 1.2.4; 1.12.5 & Protocolo integral 1.10.2.1.1 / 1.10.2.2.1 y evidencia específica de la cláusula\\\hline
RT-12.08 & 1.2.4; 1.12.5 & Protocolo integral 1.10.2.1.1 / 1.10.2.2.1 y evidencia específica de la cláusula\\\hline
RT-12.09 & 1.2.4; 1.12.5 & Protocolo integral 1.10.2.1.1 / 1.10.2.2.1 y evidencia específica de la cláusula\\\hline
RT-12.10 & 1.2.4; 1.12.5 & Protocolo integral 1.10.2.1.1 / 1.10.2.2.1 y evidencia específica de la cláusula\\\hline
RT-12.11 & 1.2.4; 1.12.5 & Protocolo integral 1.10.2.1.1 / 1.10.2.2.1 y evidencia específica de la cláusula\\\hline
RT-12.12 & 1.2.4; 1.12.5 & Protocolo integral 1.10.2.1.1 / 1.10.2.2.1 y evidencia específica de la cláusula\\\hline
RT-12.13 & 1.2.4; 1.12.5 & Protocolo integral 1.10.2.1.1 / 1.10.2.2.1 y evidencia específica de la cláusula\\\hline
RT-13.01 & 1.3; 1.1.5 & Protocolo integral 1.10.2.1.1 / 1.10.2.2.1 y evidencia específica de la cláusula\\\hline
RT-13.02 & 1.3; 1.1.5 & Protocolo integral 1.10.2.1.1 / 1.10.2.2.1 y evidencia específica de la cláusula\\\hline
RT-13.03 & 1.3; 1.1.5 & Protocolo integral 1.10.2.1.1 / 1.10.2.2.1 y evidencia específica de la cláusula\\\hline
RT-13.04 & 1.3; 1.1.5 & Protocolo integral 1.10.2.1.1 / 1.10.2.2.1 y evidencia específica de la cláusula\\\hline
RT-13.05 & 1.3; 1.1.5 & Protocolo integral 1.10.2.1.1 / 1.10.2.2.1 y evidencia específica de la cláusula\\\hline
RT-13.06 & 1.3; 1.1.5 & Protocolo integral 1.10.2.1.1 / 1.10.2.2.1 y evidencia específica de la cláusula\\\hline
RT-13.07 & 1.3; 1.1.5 & Protocolo integral 1.10.2.1.1 / 1.10.2.2.1 y evidencia específica de la cláusula\\\hline
RT-13.08 & 1.3; 1.1.5 & Protocolo integral 1.10.2.1.1 / 1.10.2.2.1 y evidencia específica de la cláusula\\\hline
RT-13.09 & 1.3; 1.1.5 & Protocolo integral 1.10.2.1.1 / 1.10.2.2.1 y evidencia específica de la cláusula\\\hline
RT-13.10 & 1.3; 1.1.5 & Protocolo integral 1.10.2.1.1 / 1.10.2.2.1 y evidencia específica de la cláusula\\\hline
RT-13.11 & 1.3; 1.1.5 & Protocolo integral 1.10.2.1.1 / 1.10.2.2.1 y evidencia específica de la cláusula\\\hline
RT-13.12 & 1.3; 1.1.5 & Protocolo integral 1.10.2.1.1 / 1.10.2.2.1 y evidencia específica de la cláusula\\\hline
RT-14.01 & 1.3.2; 1.1.5 & Protocolo integral 1.10.2.1.1 / 1.10.2.2.1 y evidencia específica de la cláusula\\\hline
RT-14.02 & 1.3.2; 1.1.5 & Protocolo integral 1.10.2.1.1 / 1.10.2.2.1 y evidencia específica de la cláusula\\\hline
RT-14.03 & 1.3.2; 1.1.5 & Protocolo integral 1.10.2.1.1 / 1.10.2.2.1 y evidencia específica de la cláusula\\\hline
RT-14.04 & 1.3.2; 1.1.5 & Protocolo integral 1.10.2.1.1 / 1.10.2.2.1 y evidencia específica de la cláusula\\\hline
RT-14.05 & 1.3.2; 1.1.5 & Protocolo integral 1.10.2.1.1 / 1.10.2.2.1 y evidencia específica de la cláusula\\\hline
RT-14.06 & 1.3.2; 1.1.5 & Protocolo integral 1.10.2.1.1 / 1.10.2.2.1 y evidencia específica de la cláusula\\\hline
RT-14.07 & 1.3.2; 1.1.5 & Protocolo integral 1.10.2.1.1 / 1.10.2.2.1 y evidencia específica de la cláusula\\\hline
RT-14.08 & 1.3.2; 1.1.5 & Protocolo integral 1.10.2.1.1 / 1.10.2.2.1 y evidencia específica de la cláusula\\\hline
RT-14.09 & Excluido Deseable BT & No procede; conservar exclusión\\\hline
RT-15.01 & 1.2.2.5; 1.2.4.2; 1.10.2; 1.12.4 & Protocolo integral 1.10.2.1.1 / 1.10.2.2.1 y evidencia específica de la cláusula\\\hline
RT-15.02 & 1.2.2.5; 1.2.4.2; 1.10.2; 1.12.4 & Protocolo integral 1.10.2.1.1 / 1.10.2.2.1 y evidencia específica de la cláusula\\\hline
RT-15.03 & 1.2.2.5; 1.2.4.2; 1.10.2; 1.12.4 & Protocolo integral 1.10.2.1.1 / 1.10.2.2.1 y evidencia específica de la cláusula\\\hline
RT-15.04 & 1.2.2.5; 1.2.4.2; 1.10.2; 1.12.4 & Protocolo integral 1.10.2.1.1 / 1.10.2.2.1 y evidencia específica de la cláusula\\\hline
RT-15.05 & Excluido Deseable BT & No procede; conservar exclusión\\\hline
RT-15.06 & Excluido Deseable BT & No procede; conservar exclusión\\\hline
RT-15.07 & 1.2.2.5; 1.2.4.2; 1.10.2; 1.12.4 & Protocolo integral 1.10.2.1.1 / 1.10.2.2.1 y evidencia específica de la cláusula\\\hline
RT-15.08 & 1.2.2.5; 1.2.4.2; 1.10.2; 1.12.4 & Protocolo integral 1.10.2.1.1 / 1.10.2.2.1 y evidencia específica de la cláusula\\\hline
RT-15.09 & 1.2.2.5; 1.2.4.2; 1.10.2; 1.12.4 & Protocolo integral 1.10.2.1.1 / 1.10.2.2.1 y evidencia específica de la cláusula\\\hline
RT-16.01 & 1.7.3; 1.6.1; 1.5.3; 1.5.1.5 & Protocolo integral 1.10.2.1.1 / 1.10.2.2.1 y evidencia específica de la cláusula\\\hline
RT-16.02 & 1.7.3; 1.6.1; 1.5.3; 1.5.1.5 & Protocolo integral 1.10.2.1.1 / 1.10.2.2.1 y evidencia específica de la cláusula\\\hline
RT-16.03 & 1.7.3; 1.6.1; 1.5.3; 1.5.1.5 & Protocolo integral 1.10.2.1.1 / 1.10.2.2.1 y evidencia específica de la cláusula\\\hline
RT-16.04 & 1.7.3; 1.6.1; 1.5.3; 1.5.1.5 & Protocolo integral 1.10.2.1.1 / 1.10.2.2.1 y evidencia específica de la cláusula\\\hline
RT-16.05 & Excluido Deseable BT & No procede; conservar exclusión\\\hline
RT-16.06 & 1.7.3; 1.6.1; 1.5.3; 1.5.1.5 & Protocolo integral 1.10.2.1.1 / 1.10.2.2.1 y evidencia específica de la cláusula\\\hline
RT-16.07 & 1.7.3; 1.6.1; 1.5.3; 1.5.1.5 & Protocolo integral 1.10.2.1.1 / 1.10.2.2.1 y evidencia específica de la cláusula\\\hline
RT-16.08 & 1.7.3; 1.6.1; 1.5.3; 1.5.1.5 & Protocolo integral 1.10.2.1.1 / 1.10.2.2.1 y evidencia específica de la cláusula\\\hline
RT-16.09 & 1.7.3; 1.6.1; 1.5.3; 1.5.1.5 & Protocolo integral 1.10.2.1.1 / 1.10.2.2.1 y evidencia específica de la cláusula\\\hline
RT-16.10 & 1.7.3; 1.6.1; 1.5.3; 1.5.1.5 & Protocolo integral 1.10.2.1.1 / 1.10.2.2.1 y evidencia específica de la cláusula\\\hline
RT-16.11 & 1.7.3; 1.6.1; 1.5.3; 1.5.1.5 & Protocolo integral 1.10.2.1.1 / 1.10.2.2.1 y evidencia específica de la cláusula\\\hline
RT-16.12 & 1.7.3; 1.6.1; 1.5.3; 1.5.1.5 & Protocolo integral 1.10.2.1.1 / 1.10.2.2.1 y evidencia específica de la cláusula\\\hline
RT-16.13 & 1.7.3; 1.6.1; 1.5.3; 1.5.1.5 & Protocolo integral 1.10.2.1.1 / 1.10.2.2.1 y evidencia específica de la cláusula\\\hline
RT-16.14 & 1.7.3; 1.6.1; 1.5.3; 1.5.1.5 & Protocolo integral 1.10.2.1.1 / 1.10.2.2.1 y evidencia específica de la cláusula\\\hline
RT-16.15 & 1.7.3; 1.6.1; 1.5.3; 1.5.1.5 & Protocolo integral 1.10.2.1.1 / 1.10.2.2.1 y evidencia específica de la cláusula\\\hline
RT-16.16 & 1.7.3; 1.6.1; 1.5.3; 1.5.1.5 & Protocolo integral 1.10.2.1.1 / 1.10.2.2.1 y evidencia específica de la cláusula\\\hline
RT-16.17 & 1.7.3; 1.6.1; 1.5.3; 1.5.1.5 & Protocolo integral 1.10.2.1.1 / 1.10.2.2.1 y evidencia específica de la cláusula\\\hline
RT-16.18 & 1.7.3; 1.6.1; 1.5.3; 1.5.1.5 & Protocolo integral 1.10.2.1.1 / 1.10.2.2.1 y evidencia específica de la cláusula\\\hline
RT-16.19 & 1.7.3; 1.6.1; 1.5.3; 1.5.1.5 & Protocolo integral 1.10.2.1.1 / 1.10.2.2.1 y evidencia específica de la cláusula\\\hline
RT-16.20 & 1.7.3; 1.6.1; 1.5.3; 1.5.1.5 & Protocolo integral 1.10.2.1.1 / 1.10.2.2.1 y evidencia específica de la cláusula\\\hline
RT-16.21 & 1.7.3; 1.6.1; 1.5.3; 1.5.1.5 & Protocolo integral 1.10.2.1.1 / 1.10.2.2.1 y evidencia específica de la cláusula\\\hline
RT-16.22 & 1.7.3; 1.6.1; 1.5.3; 1.5.1.5 & Protocolo integral 1.10.2.1.1 / 1.10.2.2.1 y evidencia específica de la cláusula\\\hline
RT-16.23 & 1.7.3; 1.6.1; 1.5.3; 1.5.1.5 & Protocolo integral 1.10.2.1.1 / 1.10.2.2.1 y evidencia específica de la cláusula\\\hline
RT-16.24 & 1.7.3; 1.6.1; 1.5.3; 1.5.1.5 & Protocolo integral 1.10.2.1.1 / 1.10.2.2.1 y evidencia específica de la cláusula\\\hline
RT-16.25 & 1.7.3; 1.6.1; 1.5.3; 1.5.1.5 & Protocolo integral 1.10.2.1.1 / 1.10.2.2.1 y evidencia específica de la cláusula\\\hline
RT-16.26 & Excluido Deseable BT & No procede; conservar exclusión\\\hline
RT-16.27 & 1.7.3; 1.6.1; 1.5.3; 1.5.1.5 & Protocolo integral 1.10.2.1.1 / 1.10.2.2.1 y evidencia específica de la cláusula\\\hline
RT-16.28 & 1.7.3; 1.6.1; 1.5.3; 1.5.1.5 & Protocolo integral 1.10.2.1.1 / 1.10.2.2.1 y evidencia específica de la cláusula\\\hline
RT-16.29 & 1.7.3; 1.6.1; 1.5.3; 1.5.1.5 & Protocolo integral 1.10.2.1.1 / 1.10.2.2.1 y evidencia específica de la cláusula\\\hline
RT-16.30 & 1.7.3; 1.6.1; 1.5.3; 1.5.1.5 & Protocolo integral 1.10.2.1.1 / 1.10.2.2.1 y evidencia específica de la cláusula\\\hline
RT-16.31 & 1.7.3; 1.6.1; 1.5.3; 1.5.1.5 & Protocolo integral 1.10.2.1.1 / 1.10.2.2.1 y evidencia específica de la cláusula\\\hline
RT-16.32 & 1.7.3; 1.6.1; 1.5.3; 1.5.1.5 & Protocolo integral 1.10.2.1.1 / 1.10.2.2.1 y evidencia específica de la cláusula\\\hline
RT-16.33 & 1.7.3; 1.6.1; 1.5.3; 1.5.1.5 & Protocolo integral 1.10.2.1.1 / 1.10.2.2.1 y evidencia específica de la cláusula\\\hline
RT-16.34 & 1.7.3; 1.6.1; 1.5.3; 1.5.1.5 & Protocolo integral 1.10.2.1.1 / 1.10.2.2.1 y evidencia específica de la cláusula\\\hline
RT-17.01 & 1.7.2; 1.7.3; 1.2.3 & Protocolo integral 1.10.2.1.1 / 1.10.2.2.1 y evidencia específica de la cláusula\\\hline
RT-17.02 & 1.7.2; 1.7.3; 1.2.3 & Protocolo integral 1.10.2.1.1 / 1.10.2.2.1 y evidencia específica de la cláusula\\\hline
RT-17.03 & 1.7.2; 1.7.3; 1.2.3 & Protocolo integral 1.10.2.1.1 / 1.10.2.2.1 y evidencia específica de la cláusula\\\hline
RT-17.04 & 1.7.2; 1.7.3; 1.2.3 & Protocolo integral 1.10.2.1.1 / 1.10.2.2.1 y evidencia específica de la cláusula\\\hline
RT-17.05 & 1.7.2; 1.7.3; 1.2.3 & Protocolo integral 1.10.2.1.1 / 1.10.2.2.1 y evidencia específica de la cláusula\\\hline
RT-17.06 & 1.7.2; 1.7.3; 1.2.3 & Protocolo integral 1.10.2.1.1 / 1.10.2.2.1 y evidencia específica de la cláusula\\\hline
RT-17.07 & 1.7.2; 1.7.3; 1.2.3 & Protocolo integral 1.10.2.1.1 / 1.10.2.2.1 y evidencia específica de la cláusula\\\hline
RT-17.08 & Excluido Deseable BT & No procede; conservar exclusión\\\hline
RT-18.01 & 1.9.3; 1.3.2; 1.1.5.3 & Protocolo integral 1.10.2.1.1 / 1.10.2.2.1 y evidencia específica de la cláusula\\\hline
RT-18.02 & 1.9.3; 1.3.2; 1.1.5.3 & Protocolo integral 1.10.2.1.1 / 1.10.2.2.1 y evidencia específica de la cláusula\\\hline
RT-18.03 & 1.9.3; 1.3.2; 1.1.5.3 & Protocolo integral 1.10.2.1.1 / 1.10.2.2.1 y evidencia específica de la cláusula\\\hline
RT-18.04 & 1.9.3; 1.3.2; 1.1.5.3 & Protocolo integral 1.10.2.1.1 / 1.10.2.2.1 y evidencia específica de la cláusula\\\hline
RT-18.05 & 1.9.3; 1.3.2; 1.1.5.3 & Protocolo integral 1.10.2.1.1 / 1.10.2.2.1 y evidencia específica de la cláusula\\\hline
RT-18.06 & 1.9.3; 1.3.2; 1.1.5.3 & Protocolo integral 1.10.2.1.1 / 1.10.2.2.1 y evidencia específica de la cláusula\\\hline
RT-18.07 & 1.9.3; 1.3.2; 1.1.5.3 & Protocolo integral 1.10.2.1.1 / 1.10.2.2.1 y evidencia específica de la cláusula\\\hline
RT-18.08 & 1.9.3; 1.3.2; 1.1.5.3 & Protocolo integral 1.10.2.1.1 / 1.10.2.2.1 y evidencia específica de la cláusula\\\hline
RT-18.09 & 1.9.3; 1.3.2; 1.1.5.3 & Protocolo integral 1.10.2.1.1 / 1.10.2.2.1 y evidencia específica de la cláusula\\\hline
RT-18.10 & Excluido Deseable BT & No procede; conservar exclusión\\\hline
RT-19.01 & 1.1 & Protocolo integral 1.10.2.1.1 / 1.10.2.2.1 y evidencia específica de la cláusula\\\hline
RT-19.02 & 1.1 & Protocolo integral 1.10.2.1.1 / 1.10.2.2.1 y evidencia específica de la cláusula\\\hline
RT-19.03 & 1.1 & Protocolo integral 1.10.2.1.1 / 1.10.2.2.1 y evidencia específica de la cláusula\\\hline
RT-19.04 & 1.1 & Protocolo integral 1.10.2.1.1 / 1.10.2.2.1 y evidencia específica de la cláusula\\\hline
RT-19.05 & 1.1 & Protocolo integral 1.10.2.1.1 / 1.10.2.2.1 y evidencia específica de la cláusula\\\hline
RT-19.06 & 1.1 & Protocolo integral 1.10.2.1.1 / 1.10.2.2.1 y evidencia específica de la cláusula\\\hline
RT-19.07 & 1.1 & Protocolo integral 1.10.2.1.1 / 1.10.2.2.1 y evidencia específica de la cláusula\\\hline
RT-19.08 & 1.1 & Protocolo integral 1.10.2.1.1 / 1.10.2.2.1 y evidencia específica de la cláusula\\\hline
RT-19.09 & 1.1 & Protocolo integral 1.10.2.1.1 / 1.10.2.2.1 y evidencia específica de la cláusula\\\hline
RT-19.10 & 1.1 & Protocolo integral 1.10.2.1.1 / 1.10.2.2.1 y evidencia específica de la cláusula\\\hline
RT-20.01 & 1.10; 1.11 & Protocolo integral 1.10.2.1.1 / 1.10.2.2.1 y evidencia específica de la cláusula\\\hline
RT-20.02 & 1.10; 1.11 & Protocolo integral 1.10.2.1.1 / 1.10.2.2.1 y evidencia específica de la cláusula\\\hline
RT-20.03 & 1.10; 1.11 & Protocolo integral 1.10.2.1.1 / 1.10.2.2.1 y evidencia específica de la cláusula\\\hline
RT-20.04 & 1.10; 1.11 & Protocolo integral 1.10.2.1.1 / 1.10.2.2.1 y evidencia específica de la cláusula\\\hline
RT-20.05 & 1.10; 1.11 & Protocolo integral 1.10.2.1.1 / 1.10.2.2.1 y evidencia específica de la cláusula\\\hline
RT-20.06 & 1.10; 1.11 & Protocolo integral 1.10.2.1.1 / 1.10.2.2.1 y evidencia específica de la cláusula\\\hline
RT-20.07 & 1.10; 1.11 & Protocolo integral 1.10.2.1.1 / 1.10.2.2.1 y evidencia específica de la cláusula\\\hline
RT-20.08 & 1.10; 1.11 & Protocolo integral 1.10.2.1.1 / 1.10.2.2.1 y evidencia específica de la cláusula\\\hline
RT-21.01 & 1.12.1; 1.12.3 & Recepción centro propio, turnos NOC y continuidad; cuadrante 13--56\\\hline
RT-21.02 & 1.1.4; 1.12.1 & Designación gerente, dedicación y registros de comité; SD1\\\hline
RT-21.03 & 1.12.1; 1.12.12 & Nominación especialistas y dedicaciones antes de intervención; recursos propios futuros\\\hline
RT-21.04 & 1.12 & Recepción periódica del paquete y control contractual SLA\\\hline
RT-21.05 & 1.12.1; 1.12.3 & Recepción centro de atención, tres agentes simultáneos y relevo bilingüe\\\hline
RT-21.06 & 1.12.1; 1.12.3 & SLA de atención, abandono y resolución medidos; C07 horarios e inglés; canal emergencia NOC separado\\\hline
RT-21.07 & 1.12 & Recepción periódica del paquete y control contractual SLA\\\hline
RT-21.08 & 1.12 & Recepción periódica del paquete y control contractual SLA\\\hline
RT-21.09 & 1.12 & Recepción periódica del paquete y control contractual SLA\\\hline
RT-21.10 & 1.12 & Recepción periódica del paquete y control contractual SLA\\\hline
RT-21.11 & 1.12 & Recepción periódica del paquete y control contractual SLA\\\hline
RT-21.12 & Excluido Deseable BT & No procede; conservar exclusión\\\hline
RT-21.13 & 1.12 & Recepción periódica del paquete y control contractual SLA\\\hline
RT-21.14 & 1.12.1 & Modelo Erlang C, demanda/TMA/sensibilidad y contraste con registros reales; dimensionamiento\_servicio.tex\\\hline
RT-21.15 & 1.12 & Recepción periódica del paquete y control contractual SLA\\\hline
RT-21.16 & 1.12.1; 1.12.3 & Especialista/suplente, embarcación, ventana segura y costos incluidos para Bahía Ilque; sin dejar descubierto otro sitio\\\hline
RT-21.17 & 1.12 & Recepción periódica del paquete y control contractual SLA\\\hline
RT-21.18 & 1.12 & Recepción periódica del paquete y control contractual SLA\\\hline
RT-21.19 & 1.12.12.evolutiva & Bolsa anual 960 HH, composición por perfil, aprobación/liquidación y acumulación de saldo\\\hline
RT-21.20 & 1.12.12.evolutiva; 1.10.2 & Cambio evolutivo recibido con pruebas, regresión y seguridad equivalentes al desarrollo original\\\hline
RT-21.21 & 1.12.12.deuda & Registro cuantificado y 32/240 HH mensuales de capacidad de ingeniería a deuda técnica\\\hline
RT-21.22 & 1.12.5; 1.12.12.normativa & Plan anual de versiones, autorización de ventana, ensayos y reversión\\\hline
RT-22.01 & 1.11.2; 1.11.7 & Cinco perfiles, cohortes, materiales y registro individual por etapa\\\hline
RT-22.02 & 1.11.2; 1.11.4; 1.11.7 & Presencial/sincrónica/autoformación y acompañamiento; fechas naturales de cuadrante\\\hline
RT-22.03 & 1.11.2; 1.12.6 & Material español editable, manuales/FAQ/videos/base de conocimiento y propiedad del CLIENTE\\\hline
RT-22.04 & 1.11.7; 1.12.10 & Turnos y formación previa al congelamiento; nuevos estacionales y monitores, sin formación interna delegada al CLIENTE\\\hline
RT-22.05 & 1.11.3; 1.11.7 & Evaluación/certificación individual de referentes administradores/técnicos designados; no supone área TI existente\\\hline
RT-22.06 & 1.12.6; 1.12.10 & Dos jornadas anuales y alta de personal nuevo, evaluación y recuperación, sin cargo adicional\\\hline
RT-22.07 & 1.12.13 & Seis paquetes mensuales 21--26; 144 HH; referentes/suplentes, casos, documentación y verificación independiente\\\hline
RT-22.08 & 1.12.6; 1.12.3 & Base de conocimiento versionada y métricas de uso/utilidad durante todo el contrato\\\hline
RT-22.09 & Excluido Deseable BT & No procede; conservar exclusión\\\hline
RT-23.01 & 1.1.5 & Control previo de expediente de oferta; contrato: gestión de vigencia 1.1.5\\\hline
RT-23.02 & 1.1.5 & Control previo de expediente de oferta; contrato: gestión de vigencia 1.1.5\\\hline
RT-23.03 & 1.1.5 & Control previo de expediente de oferta; contrato: gestión de vigencia 1.1.5\\\hline
RT-23.04 & 1.1.5 & Control previo de expediente de oferta; contrato: gestión de vigencia 1.1.5\\\hline
RT-23.05 & 1.1.5 & Control previo de expediente de oferta; contrato: gestión de vigencia 1.1.5\\\hline
RT-23.06 & 1.1.5 & Control previo de expediente de oferta; contrato: gestión de vigencia 1.1.5\\\hline
RT-23.07 & Excluido Deseable BT & No procede; conservar exclusión\\\hline
RT-23.08 & Excluido Deseable BT & No procede; conservar exclusión\\\hline
RT-23.09 & Excluido Deseable BT & No procede; conservar exclusión\\\hline
RT-23.10 & Excluido Deseable BT & No procede; conservar exclusión\\\hline
RT-24.01 & 1.1.5 & Control previo de expediente de oferta; contrato: gestión de vigencia 1.1.5\\\hline
RT-24.02 & 1.1.5 & Control previo de expediente de oferta; contrato: gestión de vigencia 1.1.5\\\hline
RT-24.03 & Excluido Deseable BT & No procede; conservar exclusión\\\hline
RT-24.04 & 1.1.5 & Control previo de expediente de oferta; contrato: gestión de vigencia 1.1.5\\\hline
RT-24.05 & 1.1.5 & Control previo de expediente de oferta; contrato: gestión de vigencia 1.1.5\\\hline
RT-24.06 & 1.1.5 & Control previo de expediente de oferta; contrato: gestión de vigencia 1.1.5\\\hline
RT-25.01 & 1.7.3; 1.4; 1.5; 1.6 & Recepción periódica del paquete y control contractual SLA\\\hline
RT-25.02 & 1.7.3; 1.4; 1.5; 1.6 & Recepción periódica del paquete y control contractual SLA\\\hline
RT-25.03 & 1.7.3; 1.4; 1.5; 1.6 & Recepción periódica del paquete y control contractual SLA\\\hline
RT-25.04 & 1.7.3; 1.4; 1.5; 1.6 & Recepción periódica del paquete y control contractual SLA\\\hline
RT-25.05 & 1.7.3; 1.4; 1.5; 1.6 & Recepción periódica del paquete y control contractual SLA\\\hline
RT-25.06 & 1.7.3; 1.4; 1.5; 1.6 & Recepción periódica del paquete y control contractual SLA\\\hline
RT-25.07 & 1.7.3; 1.4; 1.5; 1.6 & Recepción periódica del paquete y control contractual SLA\\\hline
RT-25.08 & 1.7.3; 1.4; 1.5; 1.6 & Recepción periódica del paquete y control contractual SLA\\\hline
RT-25.09 & 1.7.3; 1.4; 1.5; 1.6 & Recepción periódica del paquete y control contractual SLA\\\hline
RT-25.10 & Excluido Deseable BT & No procede; conservar exclusión\\\hline
RT-26.01 & 1.9 & Recepción periódica del paquete y control contractual SLA\\\hline
RT-26.02 & 1.9 & Recepción periódica del paquete y control contractual SLA\\\hline
RT-26.03 & 1.9 & Recepción periódica del paquete y control contractual SLA\\\hline
RT-26.04 & 1.9 & Recepción periódica del paquete y control contractual SLA\\\hline
RT-26.05 & 1.9 & Recepción periódica del paquete y control contractual SLA\\\hline
RT-26.06 & 1.9 & Recepción periódica del paquete y control contractual SLA\\\hline
RT-26.07 & 1.9 & Recepción periódica del paquete y control contractual SLA\\\hline
RT-26.08 & 1.9 & Recepción periódica del paquete y control contractual SLA\\\hline
\end{longtable}\FuenteTabla{T-12 vigente y catálogo de paquetes integrados Synaptix; detalle literal y correspondencia terminal en registros CSV asociados.}

\subsection{Red global, ruta crítica y holguras}
La red seleccionada contiene 2279 nodos y 5075 enlaces. Las asignaciones incluyen los lotes UCP, habilitaciones, ventanas, reservas y actos de recepción requeridos. Los cuadrantes conservan fechas naturales exactas; los índices hábiles son controles de precedencia y no convierten turnos nocturnos o fines de semana en ingeniería diurna.
Con duración $d_i$ y liberación $r_i$, $IT_i=\max(r_i,\max FT_j)$ y $FT_i=IT_i+d_i$. En retroceso, $FL_i=\min IL_s$ y $IL_i=FL_i-d_i$; $HT_i=IL_i-IT_i$ y $HL_i=\min IT_s-FT_i$. El CPM de precedencias y la asignación con capacidad son cálculos distintos. Las holguras de \path{componentes/cpm_implementacion.csv} se refieren al cierre calculado de implementación, mientras la tabla muestra el margen hacia cada obligación contractual.
El CPM hacia el cierre calculado de implementación contiene cuatro ramas críticas simultáneas: observación e informe de IN1, IN3, IN4 e IN5. Convergen en \texttt{PILOTOS-REV} y \texttt{PILOTOS-SUB}; el tramo común continúa por \texttt{1.9.1.6.1}, \texttt{1.9.4.4.1}, H12 y \texttt{FIN-IMPLEMENTACION}. La Tabla~\ref{tab:t15-ramas-cpm} identifica los nodos con holgura cero. Esas holguras se refieren al cierre calculado de esta red; no constituyen autorización para desplazar los hitos intermedios ni el acta/corte E2. El corte E2 tiene además una puerta rígida: cierre natural de observación, acta expresa de marcha blanca y habilitación lógica al inicio de 21, sin holgura. La aceptación final H12 conserva su ventana dentro del mes 21; el acta de marcha blanca es requisito previo del corte y no se sustituye por el acta posterior de cierre del proyecto.
\begingroup\small\begin{longtable}{L{25mm}L{26mm}L{26mm}R{26mm}}
\caption{Programación seleccionada y márgenes}\label{tab:sd7-hitos-globales}\\\hline
Hito/puerta & Programado & Límite & Margen hábil\\\hline\endfirsthead
\hline Hito/puerta & Programado & Límite & Margen hábil\\\hline\endhead
H1 & 2027-01-20 & 2027-01-29 & 7,67\\\hline
H2 & 2027-03-01 & 2027-03-31 & 21,44\\\hline
H3 & 2027-05-27 & 2027-05-31 & 2,33\\\hline
H4 & 2027-09-23 & 2027-09-30 & 5,22\\\hline
H5 & 2027-11-08 & 2027-11-30 & 16,01\\\hline
H6 & 2027-12-01 & 2027-12-31 & 22\\\hline
H7 & 2028-03-29 & 2028-03-31 & 2,44\\\hline
H8 & 2028-01-17 & 2028-01-31 & 10,78\\\hline
H9 & 2028-04-18 & 2028-04-28 & 8,22\\\hline
H10 & 2028-05-11 & 2028-05-31 & 14,67\\\hline
H11 & 2028-06-01 & 2028-06-30 & 20\\\hline
ACTA-MB-E2 & 2028-08-01 & 2028-08-01 & 0\\\hline
H12 & 2028-08-16 & 2028-08-31 & 11,67\\\hline
CORTE-E1 & 2028-03-30 & 2028-03-31 & 1,44\\\hline
CORTE-E2 & 2028-08-01 & 2028-08-01 & 0\\\hline
\end{longtable}\endgroup
La revisión continua es el procedimiento seleccionado para el acta de marcha blanca E2: el calendario demuestra capacidad propia y precedencias bajo el supuesto adoptado de participación de la Contraparte en la revisión continua y firma expresa al cierre de la evidencia. La selección de la línea base queda cerrada bajo ese supuesto; el acto futuro sigue siendo una puerta obligatoria y no se presume obtenido. Una revisión íntegra posterior al cierre natural produce atraso y se conserva como contraste, sin cambiar la fecha ofrecida. Las reservas de diez días marinos se imputan una sola vez antes de habilitar cada marcha blanca; no son holgura del corte ni presupuesto automático de retrabajo.\parencite[arts.~17--18]{bases-admin}
\begingroup\small\begin{longtable}{L{47mm}R{16mm}R{16mm}R{16mm}R{16mm}R{12mm}R{12mm}}
\caption{Ramas críticas respecto del cierre calculado de implementación}\label{tab:t15-ramas-cpm}\\\hline
Nodo & IT & FT & IL & FL & HT & HL\\\hline\endfirsthead
\hline Nodo & IT & FT & IL & FL & HT & HL\\\hline\endhead
\path{OBS-IN1-PILOTO} & 378 & 396 & 378 & 396 & 0 & 0\\\hline
\path{OBS-IN3-PILOTO} & 378 & 396 & 378 & 396 & 0 & 0\\\hline
\path{OBS-IN4-PILOTO} & 378 & 396 & 378 & 396 & 0 & 0\\\hline
\path{OBS-IN5-PILOTO} & 378 & 396 & 378 & 396 & 0 & 0\\\hline
\path{1.9.1.5.6} & 396 & 399,33 & 396 & 399,33 & 0 & 0\\\hline
\path{1.9.3.5.3} & 396 & 399,33 & 396 & 399,33 & 0 & 0\\\hline
\path{1.9.4.3.3} & 396 & 399,33 & 396 & 399,33 & 0 & 0\\\hline
\path{1.9.5.3.4} & 396 & 399,33 & 396 & 399,33 & 0 & 0\\\hline
\path{PILOTOS-REV} & 399,33 & 409,33 & 399,33 & 409,33 & 0 & 0\\\hline
\path{PILOTOS-SUB} & 409,33 & 419,33 & 409,33 & 419,33 & 0 & 0\\\hline
\path{1.9.1.6.1} & 419,33 & 422,67 & 419,33 & 422,67 & 0 & 0\\\hline
\path{1.9.4.4.1} & 422,67 & 423,78 & 422,67 & 423,78 & 0 & 0\\\hline
\path{H12} & 423,78 & 423,78 & 423,78 & 423,78 & 0 & 0\\\hline
\path{FIN-IMPLEMENTACION} & 423,78 & 423,78 & 423,78 & 423,78 & 0 & 0\\\hline
\end{longtable}\endgroup
Índices en días hábiles relativos al escenario de inicio. IT/FT: inicio/fin temprano; IL/FL: inicio/fin tardío; HT/HL: holgura total/libre. Las cuatro ramas parten de OBS-IN1/IN3/IN4/IN5-PILOTO y llegan a PILOTOS-REV mediante 1.9.1.5.6, 1.9.3.5.3, 1.9.4.3.3 y 1.9.5.3.4, respectivamente. La liberación fija de observación se conserva como condición temporal, sin inventar una predecesora que explique su espera.


\subsection{Escenario PERT nuevamente nivelado}
Se recalcula la asignación con $t_E=1,1d$ para ingeniería y once días para ensayos de diez, manteniendo las HH seleccionadas y las capacidades. Los actos de revisión, observación natural, turnos y corte no se multiplican. Es una sensibilidad determinista, sin probabilidad de entrega ni presupuesto implícito de corrección.
\begin{longtable}{L{25mm}L{26mm}L{26mm}R{26mm}}
\caption{PERT con recursos recalculados}\label{tab:sd7-pert-nivelado}\\\hline
Hito/puerta & Directo & PERT recursos & Margen PERT\\\hline\endfirsthead
\hline Hito/puerta & Directo & PERT recursos & Margen PERT\\\hline\endhead
H1 & 2027-01-20 & 2027-01-20 & 7,33\\\hline
H2 & 2027-03-01 & 2027-03-02 & 20,89\\\hline
H3 & 2027-05-27 & 2027-05-26 & 3,44\\\hline
H4 & 2027-09-23 & 2027-09-24 & 4,47\\\hline
H5 & 2027-11-08 & 2027-11-11 & 13,17\\\hline
H6 & 2027-12-01 & 2027-12-01 & 22\\\hline
H7 & 2028-03-29 & 2028-03-29 & 2,39\\\hline
H8 & 2028-01-17 & 2028-01-17 & 10,56\\\hline
H9 & 2028-04-18 & 2028-04-19 & 7,94\\\hline
H10 & 2028-05-11 & 2028-05-15 & 12,83\\\hline
H11 & 2028-06-01 & 2028-06-01 & 20\\\hline
ACTA-MB-E2 & 2028-08-01 & 2028-08-01 & 0\\\hline
H12 & 2028-08-16 & 2028-08-17 & 10,33\\\hline
CORTE-E1 & 2028-03-30 & 2028-03-30 & 1,39\\\hline
CORTE-E2 & 2028-08-01 & 2028-08-01 & 0\\\hline
\end{longtable}
El margen cero del acta/corte E2 permanece bajo PERT y la línea base conserva el supuesto de recepción continua y firma expresa programada. PL-R04 controla su incumplimiento; más recursos propios no reemplazan el acto de aceptación. La capacidad máxima del escenario PERT se contrasta por persona en su registro independiente antes de adoptar un cambio de duración.

\subsection{Fundamento y límites de los escenarios y reservas}
La estimación conserva tres duraciones por actividad y contrasta sus efectos con red y recursos, siguiendo la separación entre duración, riesgo y calendario del manual de programación de NASA.~\parencite[secc.~6.3.1 y 7.4, pp.~38--40 y 59--60]{nasa-cronograma} Para la ponderación PERT se aplica $t_E=(t_O+4t_M+t_P)/6$ y $\sigma=(t_P-t_O)/6$ como sensibilidad de diseño. El escenario de tres valores conserva $t_O=t_M=d$ y $t_P=1,6d$ para ingeniería. La igualdad del optimista y probable representa el mínimo técnico del diseño inicial; el pesimista ensaya una recuperación adicional de hasta $0,6d$. Con ello $t_E=1,1d$, $\sigma=0,1d$ y $\sigma^2=0,01d^2$. Los ensayos conservan $(10,10,16)$, $t_E=11$ y $\sigma=1$ día. Revisiones contractuales, turnos y períodos de observación naturales no se ajustan con el mismo multiplicador.

Estos extremos son hipótesis de sensibilidad, no mediciones, percentiles ni una distribución calibrada. El registro \path{componentes/riesgos_reservas_planificacion.csv} identifica evento, paquetes, responsable, disparador, respuesta y límite de cada reserva. La recuperación de ingeniería deberá respaldarse por corrección y comprobación identificadas; si demanda más HH, se actualizarán paquetes y capacidad. La nivelación previa de $1,1d$ con las mismas HH comprueba disponibilidad temporal a menor intensidad; no financia automáticamente un 10\,\% de retrabajo ni demuestra probabilidad de entrega.

Los diez días hábiles marinos son una reserva de contingencia propuesta para una campaña con dos ventanas alternativas y cancelaciones acumuladas de hasta diez días en el escenario. No se atribuye esa duración a un pronóstico acreditado. Su uso requiere aviso o bloqueo registrado, decisión de Operación y autorización de la ventana alternativa. H3 incorpora su reserva en la cadena; H5 y H10 la preservan antes de habilitar su etapa. No se vuelve a contar como revisión, subsanación ni holgura libre de los cortes finales. Se registrarán duración disponible, consumida y restante; excederla exige escalar y recalcular dentro de los hitos contractuales.

El registro PL-R01--PL-R05 conserva evento, paquete, responsable, disparador y límite; T-16 enlaza sus tratamientos con los riesgos de marejada, integridad, contraparte y capacidad. PL-R01/R-01 consume una sola reserva por campaña y PL-R04 conserva la puerta de recepción E2 sin holgura. La actualización de un tratamiento se propagará a red y asignaciones mediante el control de cambios de SD6. La reserva de gestión se mantiene fuera de la línea base y su utilización requiere decisión de dirección y control de cambios; ninguna reserva autoriza desplazar el plazo contractual. La estimación funcional se concilia mediante el esfuerzo seleccionado y los lotes adicionales. La sensibilidad del ambiente y la firma externa de recepción E2 requieren control propio; no se cubren mediante una reserva meteorológica. No se afirma una confianza normal de dos sigmas sin sustentar distribuciones, correlaciones y efecto de las rutas que cambian bajo recursos.

\subsection{Recepción continua seleccionada de E2}
Synaptix selecciona la recepción continua para cumplir producción E2 y Operación de ambos alcances desde el inicio del mes 21. Para inicio el 1 de diciembre de 2026, la fecha es el 1 de agosto de 2028. La línea base adopta como supuesto de planificación la participación de la Contraparte en la revisión continua y en la sesión final de aceptación, con firma expresa prevista al cierre de la evidencia. Synaptix propone este procedimiento para su ejecución contractual; no acredita un acuerdo ni un acta ya suscritos.~\parencite[arts.~17.2.4, 17.3 y 18; pp.~12--13]{bases-admin}

El expediente parcial se prepara desde el 15 de junio, con índice, versión, conciliación, formación, fuentes y fórmulas SLA, incidentes y observaciones. La red reserva diez días hábiles de revisión parcial y diez de respuesta sobre evidencia ya producida. Calidad entregará registros diarios inmutables y resolverá observaciones tempranas; la revisión parcial no acepta mediciones futuras ni consume el derecho de revisar evidencia nueva.

Las cuatro últimas semanas se observan íntegramente del 4 al 31 de julio, con volumen real, SLA y usuarios. El 31 de julio, de 20:00 a 24:00, se reservan dos ejecutores propios separados, uno de Calidad y otro de Proyecto, cuatro HH cada uno. El cuadrante conserva el intervalo natural, sin reutilizar un turno NOC/SOC/mesa o una jornada de ingeniería. Revisan el expediente vivo y la última evidencia; el cierre automático a las 00:00 conserva el registro final e impide habilitación ante diferencias, incidentes impeditivos o medición incompleta. Ninguna proyección sustituye los últimos datos.

El acta expresa de cierre de marcha blanca es una puerta previa del corte. Synaptix preparará y ensayará en PREPROD la configuración, respaldo, retorno y cambio de autoridad antes de la ventana; el corte es una habilitación lógica preparada, cuya secuencia será cierre de evidencia, comprobación final, acta expresa y activación. El ensayo medirá los tiempos de consolidación y activación: una precedencia lógica no acredita duración física nula ni firma instantánea. Sólo se habilitará E2 como registro oficial si se dispone del acta y de las condiciones copulativas. NOC, SOC y mesa atienden ambos alcances desde el inicio de 21; la estabilización de 28 días comienza el 1 de agosto, con 448 HH y cuadrante propio.

H12 corresponde a la aceptación final del proyecto de implementación dentro del mes 21 y agrega las entregas finales, incluido el cierre de pilotos de innovación y sus revisiones. Se distingue del acta de marcha blanca requerida para el corte del primer día. Esa distinción no autoriza producción sin acta ni aceptación anticipada de un producto.

\paragraph{Preparación y autorización del cierre.} Antes de H10, Proyecto presentará a la Contraparte la agenda, referente y suplente, criterios, plantilla de acta, expediente diario y protocolo de cierre para coordinación expresa. La convocatoria y la respuesta se registrarán, sin presumir disponibilidad ni aceptación por silencio. Antes de H11, Calidad e Integración ensayarán en PREPROD la consolidación del último registro, verificación, presentación del acta, habilitación y retorno; conservarán duración medida de cada paso, versión y responsable. El ensayo no incluirá una firma simulada como evidencia de decisión del CLIENTE ni convertirá tiempo de elaboración en tiempo de aceptación externa.

Durante los 28 días finales, el tablero diario contrastará mínimos de volumen, cobertura, SLA, conciliación, formación e incidentes; Proyecto escalará una previsión de incumplimiento tan pronto sea identificada, sin esperar al último día. A siete días del cierre comprobará agenda y suplentes, observaciones abiertas y preparación del corte. El acta sólo se someterá a firma cuando esté completa la última evidencia. Una duración positiva medida, una observación o un pronunciamiento tardío se reflejará en la programación de la transición: el margen cero no absorbe esa demora. Sin acta expresa y comprobación final no se habilita el registro oficial. Este protocolo se incorpora al expediente de aceptación de T-17; H12 posterior no reemplaza la autorización previa.\parencite[arts.~17--18, pp.~12--13]{bases-admin}

\paragraph{Supuesto adoptado y contingencia.} Para la línea base ofrecida se supone que la Contraparte participará en la sesión final del 31 de julio de 2028, de 20:00 a 24:00, y emitirá su pronunciamiento expreso una vez cerrado el registro de las cuatro semanas. Proyecto coordinará la sesión y presentará el acta; Calidad comprobará la evidencia completa. La firma y activación se programan en la transición al 1 de agosto, conservando ese orden y todos los criterios. Proyecto presentará con la línea base la agenda de recepción continua y el protocolo de cierre; el CLIENTE resolverá su coordinación mediante un pronunciamiento expreso. La viabilidad del instante de corte depende de ese procedimiento y de los tiempos medidos en el ensayo; la reserva cero exige tratar cualquier demora como desviación, no como holgura disponible. La decisión de planificación queda cerrada bajo este supuesto, sin atribuir una aceptación ya obtenida. El CLIENTE conserva los diez días de revisión de evidencia nueva: el supuesto prevé un pronunciamiento anticipado válido, sin renuncia a ese derecho ni aceptación tácita. Si falta el acta, existen observaciones impeditivas o se usa íntegramente la revisión posterior, la puerta bloquea el cambio oficial, se registra el atraso y se extiende marcha blanca a costo de Synaptix, sin desplazar el calendario contractual. PL-R04 controla el incumplimiento del supuesto; el escenario máximo se conserva como sensibilidad de plazo.

\subsection{Comprobación de la fecha objetivo de E2}
La revisión distingue el cierre de evidencia, el pronunciamiento contractual y la ejecución del corte. Para inicio el 1 de diciembre de 2026, la evidencia de las cuatro últimas semanas de marcha blanca termina el 31 de julio de 2028 y la producción/Operación simultáneas tienen como objetivo el 1 de agosto. El artículo 18.3 reconoce hasta diez días hábiles de revisión y otros diez de subsanación; son ventanas máximas, no un mínimo obligatorio de espera. La preparación diaria puede reducir trabajo de revisión, pero no impone al CLIENTE un pronunciamiento inmediato sobre evidencia aún no producida.\parencite[arts.~17.2.4, 17.3 y 18.3]{bases-admin}

\begin{tabularx}{\linewidth}{Y L{27mm} L{27mm}}\hline
Escenario & Fin revisión & Fin subsanación\\\hline
Sin observaciones, revisión máxima & 14/08/2028 & No aplica\\\hline
Con observaciones, ambas ventanas máximas & 14/08/2028 & 29/08/2028\\\hline
\end{tabularx}\par\smallskip
El cálculo usa lunes a viernes y feriados chilenos del escenario; sus datos se conservan en \path{componentes/recepcion_e2_contraste_contractual.csv}. No incluye una nueva revisión ni la duración efectiva del corte. Incluso sin observaciones, el uso completo del plazo de revisión termina después de la fecha objetivo. Agregar recursos a Synaptix, adelantar informes parciales o reducir una reserva no elimina esa dependencia externa.

\paragraph{Resultado y responsabilidad.} El escenario conservador no satisface el objetivo del primer día de 21; el corte del 31 de agosto del contraste de revisión posterior es un escenario de atraso, no la fecha contractual ofrecida. La red seleccionada programa acta de marcha blanca y corte el 1 de agosto; H12 es la aceptación final de implementación dentro de 21. La línea base adopta el supuesto de participación de la Contraparte en la recepción continua, la sesión final del 31 de julio y la firma expresa al cierre de la evidencia, antes del corte. Así se cierra la selección del procedimiento y del calendario ofrecido; la firma y la aceptación son actuaciones futuras sujetas a los criterios de las bases. Proyecto coordinará la ejecución y PL-R04 controlará cualquier desviación. El contraste de ventanas máximas cuantifica el atraso si el supuesto no se cumple; no sustituye el escenario seleccionado ni reduce derechos de revisión. Los servicios de ambos alcances se dimensionan desde el primer día de 21; su dotación no acredita por sí sola que E2 haya sido recibido. No se sustituyen las últimas mediciones por proyecciones ni se desplazan los meses contractuales.\parencite[arts.~17--18]{bases-admin}

\subsection{Calendario de intervención y control de ventanas}
El escenario utiliza inicio el 1 de diciembre de 2026. Las habilitaciones externas tienen fecha requerida y responsable: recinto, equipos, red y telemetría recibidos para el 16 de marzo de 2027; licencias e identidad antes de su configuración; contratos contables y de pagos para enero de 2028. No se afirma que el CLIENTE haya ejecutado las obras. La instalación que afecte al varadero se prepara por anticipado y se ejecuta en la ventana de marzo posterior al día 15; la integración central y recepción de puntos de pantalán continúa en abril, que está abierto para el resto del recinto. H3 incluye cinco ambientes y la recepción de medición, con su reserva marina, antes del cierre de mayo.

Construcción E1 y componentes compartidos se concentran en agosto de 2027 con refuerzos dedicados; las vistas E1 ya no se dejan para septiembre. E2 aprueba modelos, contratos y esquemas en diciembre/enero, construye en enero--febrero y prueba desde el 16 de marzo de 2028. El trabajo remoto en DEV/QA no cambia producción ni habilita usuarios durante el congelamiento; toda intervención productiva, prueba destructiva o formación presencial respeta la ventana del sitio. Los ensayos usan manifiestos completos, dos ejecuciones por etapa y sus revisiones/subsanaciones. Una modificación que invalide evidencia obliga a repetirla sobre la versión recibida, a costo de Synaptix.

Las bases de IN1/IN4/IN5 abarcan 1--28 de mayo y la evaluación 1--28 de junio de 2028; en IN1 las solicitudes entran en los primeros 21 días y se siguen siete. La formación posterior a base se realiza 29--31 de mayo, con dos formadores y Calidad reservados. Las acciones de IN5 se aplican entre ambas ventanas, con autorización y evidencia. Los informes, revisión y subsanación se preparan en julio y las innovaciones se materializan en 21. IN3 conserva 120 pólizas distintas, orden intercalado y control de errores. IN2 tiene fecha concreta de invernada en junio de 2029 y controles de botadura antes de sus maniobras; sus actividades documentales no autorizan una intervención en varadero durante las campañas.

Los controles de DR se anticipan a noviembre, marzo después del día 15 y agosto, con nueve controles en 24/28/33/36/40/45/48/52/56, evitando exceder seis meses. Las pruebas ofensivas de operación se realizan en marzo después del día 15 de 2029, 2030 y 2031. Las jornadas de actualización y la incorporación estacional se preparan fuera del congelamiento; la autoformación no sustituye la evaluación. La salida se acompaña desde abril hasta julio de 2031, al menos noventa días reales, y la eliminación de origen depende de exportaciones, credenciales y acta de transferencia.

Proyecto registrará los catorce eventos anuales y sus bloqueos de uno a tres días, con fechas conocidas un año antes, y las fechas efectivas de las campañas de seis semanas. El escenario no inventa esas fechas. Las actividades de campo sólo se liberan tras comprobar zona, maniobra, evento, viento/marejada, acceso y autorización; ante aviso se cancelan y se consume reserva identificada. Si la fecha real de inicio o un evento anula la ventana, se recalculará la programación completa antes de fijar la línea base, sin trasladar los meses contractuales. La serie de medición comienza en E1 y se conserva antes de la temporada 2029; los residuos y la capacitación ambiental tienen recepción propia y no quedan demostrados sólo por instalar medidores.

\subsection{Hitos de recepción y pagos de implementación}
Los doce hitos E-25 conservan sus meses y condiciones de devengo. La Etapa 1 suma 60\,\% y la Etapa 2 40\,\% del valor de implementación. Una fecha programada no constituye aprobación ni autoriza pago sin su entregable recibido \parencite[Formulario~E-25; p.~74]{bases-admin}.
La habilitación H3 comprende además DR conforme a RT-04.01 de las Bases Técnicas Transversales, aunque la descripción de E-25 enumere DEV, QA, PREPROD y PROD. La EDT incluye una ficha por cada uno de los cinco ambientes \parencite[RT-04.01; pp.~10--11]{bases-transversales}.
La Tabla~\ref{tab:sd7-e25} relaciona los doce hitos con mes, porcentaje y condición de recepción.
\begin{longtable}{L{12mm}R{12mm}R{14mm}L{110mm}}
\caption{Correspondencia contractual con E-25}\label{tab:sd7-e25}\\\hline
 \EncabezadoTabla{Hito} & \EncabezadoTabla{Mes} & \EncabezadoTabla{Pago} & \EncabezadoTabla{Condición de devengo}\\\hline\endfirsthead
\hline \EncabezadoTabla{Hito} & \EncabezadoTabla{Mes} & \EncabezadoTabla{Pago} & \EncabezadoTabla{Condición de devengo}\\\hline\endhead
H1 & 2 & 8\,\% & Alcance y matriz de trazabilidad aprobados \\\hline
H2 & 4 & 7\,\% & Arquitectura, seguridad y modelo de datos aprobados\\\hline
H3 & 6 & 10\,\% & Infraestructura híbrida y cuatro ambientes con observabilidad \\\hline
H4 & 10 & 10\,\% & Software E1 con QA superado y evidencia de cobertura\\\hline
H5 & 12 & 10\,\% & E1 certificada: UAT, carga, resiliencia y seguridad ofensiva \\\hline
H6 & 13 & 5\,\% & Inicio de marcha blanca E1; retorno activo y usuarios capacitados\\\hline
H7 & 16 & 10\,\% & Producción E1 y cierre conforme de marcha blanca \\\hline
H8 & 14 & 5\,\% & Alcance E2 y diseño detallado aprobados\\\hline
H9 & 17 & 10\,\% & Software E2 con QA superado \\\hline
H10 & 18 & 10\,\% & Certificación E2 y cierre de desarrollo\\\hline
H11 & 19 & 5\,\% & Inicio de marcha blanca E2 con E1 en producción \\\hline
H12 & 21 & 10\,\% & Producción E2 y aceptación final de implementación\\\hline
\end{longtable}\FuenteTabla{Bases Administrativas, Formulario E-25. Porcentajes sobre el valor total de implementación.}


La operación comprende 36 períodos, meses 21--56, con pago al mes siguiente vencido sujeto a los descuentos de SLA aplicables. El último período de servicio es el 56 y su pago corresponde al mes siguiente; ese pago no agrega un mes de operación. No se anticiparán ni prorratearán estos importes dentro de implementación. Los porcentajes E-25 se convertirán en montos al consolidar la valoración económica, conservando la separación entre ambas fases.

\subsection{Gantt contractual de 56 meses}
La carta muestra los períodos de trabajo de las doce cuentas, marchas blancas, cobertura inicial y Operación, junto con los doce hitos E-25. IN2 conserva junio de 2029 como primera invernada; las demás innovaciones siguen SD13. Las barras agregadas no autorizan intervenciones durante congelamientos, eventos o marejadas ni sustituyen las precedencias de la red.~\parencite[art.~17, pp.~12--13; Formulario E-25, p.~74]{bases-admin}
\clearpage\begin{landscape}
La Figura~\ref{fig:sd7-gantt-1} presenta la programación agregada de los meses 1--28.
La carta agrega los intervalos calculados por cuenta; las fases contractuales se muestran en filas separadas. Las barras no autorizan cambios durante un congelamiento ni acreditan recepciones ejecutadas.
\begin{figure}[H]\centering\begin{tikzpicture}[x=0.52cm,y=0.55cm,font=\fontsize{9}{11}\selectfont]
\node at (0.5,.6) {1};\draw[SynLine] (0,.2)--(0,-16.5);
\node at (1.5,.6) {2};\draw[SynLine] (1,.2)--(1,-16.5);
\node at (2.5,.6) {3};\draw[SynLine] (2,.2)--(2,-16.5);
\node at (3.5,.6) {4};\draw[SynLine] (3,.2)--(3,-16.5);
\node at (4.5,.6) {5};\draw[SynLine] (4,.2)--(4,-16.5);
\node at (5.5,.6) {6};\draw[SynLine] (5,.2)--(5,-16.5);
\node at (6.5,.6) {7};\draw[SynLine] (6,.2)--(6,-16.5);
\node at (7.5,.6) {8};\draw[SynLine] (7,.2)--(7,-16.5);
\node at (8.5,.6) {9};\draw[SynLine] (8,.2)--(8,-16.5);
\node at (9.5,.6) {10};\draw[SynLine] (9,.2)--(9,-16.5);
\node at (10.5,.6) {11};\draw[SynLine] (10,.2)--(10,-16.5);
\node at (11.5,.6) {12};\draw[SynLine] (11,.2)--(11,-16.5);
\node at (12.5,.6) {13};\draw[SynLine] (12,.2)--(12,-16.5);
\node at (13.5,.6) {14};\draw[SynLine] (13,.2)--(13,-16.5);
\node at (14.5,.6) {15};\draw[SynLine] (14,.2)--(14,-16.5);
\node at (15.5,.6) {16};\draw[SynLine] (15,.2)--(15,-16.5);
\node at (16.5,.6) {17};\draw[SynLine] (16,.2)--(16,-16.5);
\node at (17.5,.6) {18};\draw[SynLine] (17,.2)--(17,-16.5);
\node at (18.5,.6) {19};\draw[SynLine] (18,.2)--(18,-16.5);
\node at (19.5,.6) {20};\draw[SynLine] (19,.2)--(19,-16.5);
\node at (20.5,.6) {21};\draw[SynLine] (20,.2)--(20,-16.5);
\node at (21.5,.6) {22};\draw[SynLine] (21,.2)--(21,-16.5);
\node at (22.5,.6) {23};\draw[SynLine] (22,.2)--(22,-16.5);
\node at (23.5,.6) {24};\draw[SynLine] (23,.2)--(23,-16.5);
\node at (24.5,.6) {25};\draw[SynLine] (24,.2)--(24,-16.5);
\node at (25.5,.6) {26};\draw[SynLine] (25,.2)--(25,-16.5);
\node at (26.5,.6) {27};\draw[SynLine] (26,.2)--(26,-16.5);
\node at (27.5,.6) {28};\draw[SynLine] (27,.2)--(27,-16.5);
\node[anchor=east,align=right,text width=4.3cm] at (-.2,-1) {1.1 Gobierno y contrato};
\fill[SynBlue!55] (0,-1.22) rectangle (1,-0.78);
\fill[SynBlue!55] (1,-1.22) rectangle (2,-0.78);
\fill[SynBlue!55] (2,-1.22) rectangle (3,-0.78);
\fill[SynBlue!55] (3,-1.22) rectangle (4,-0.78);
\fill[SynBlue!55] (4,-1.22) rectangle (5,-0.78);
\fill[SynBlue!55] (5,-1.22) rectangle (6,-0.78);
\fill[SynBlue!55] (6,-1.22) rectangle (7,-0.78);
\fill[SynBlue!55] (7,-1.22) rectangle (8,-0.78);
\fill[SynBlue!55] (8,-1.22) rectangle (9,-0.78);
\fill[SynBlue!55] (9,-1.22) rectangle (10,-0.78);
\fill[SynBlue!55] (10,-1.22) rectangle (11,-0.78);
\fill[SynBlue!55] (11,-1.22) rectangle (12,-0.78);
\fill[SynBlue!55] (12,-1.22) rectangle (13,-0.78);
\fill[SynBlue!55] (13,-1.22) rectangle (14,-0.78);
\fill[SynBlue!55] (14,-1.22) rectangle (15,-0.78);
\fill[SynBlue!55] (15,-1.22) rectangle (16,-0.78);
\fill[SynBlue!55] (16,-1.22) rectangle (17,-0.78);
\fill[SynBlue!55] (17,-1.22) rectangle (18,-0.78);
\fill[SynBlue!55] (18,-1.22) rectangle (19,-0.78);
\fill[SynBlue!55] (19,-1.22) rectangle (20,-0.78);
\node[anchor=east,align=right,text width=4.3cm] at (-.2,-2) {1.2 Plataforma y activos};
\fill[SynBlue!55] (2,-2.22) rectangle (3,-1.78);
\fill[SynBlue!55] (3,-2.22) rectangle (4,-1.78);
\fill[SynBlue!55] (4,-2.22) rectangle (5,-1.78);
\node[anchor=east,align=right,text width=4.3cm] at (-.2,-3) {1.3 Seguridad y auditoría};
\fill[SynBlue!55] (1,-3.22) rectangle (2,-2.78);
\fill[SynBlue!55] (2,-3.22) rectangle (3,-2.78);
\fill[SynBlue!55] (3,-3.22) rectangle (4,-2.78);
\fill[SynBlue!55] (4,-3.22) rectangle (5,-2.78);
\fill[SynBlue!55] (7,-3.22) rectangle (8,-2.78);
\fill[SynBlue!55] (8,-3.22) rectangle (9,-2.78);
\node[anchor=east,align=right,text width=4.3cm] at (-.2,-4) {1.4 Núcleo operacional};
\fill[SynBlue!55] (7,-4.22) rectangle (8,-3.78);
\fill[SynBlue!55] (8,-4.22) rectangle (9,-3.78);
\fill[SynBlue!55] (1,-4.22) rectangle (2,-3.78);
\fill[SynBlue!55] (2,-4.22) rectangle (3,-3.78);
\node[anchor=east,align=right,text width=4.3cm] at (-.2,-5) {1.5 Núcleo de servicios};
\fill[SynBlue!55] (1,-5.22) rectangle (2,-4.78);
\fill[SynBlue!55] (2,-5.22) rectangle (3,-4.78);
\fill[SynBlue!55] (7,-5.22) rectangle (8,-4.78);
\fill[SynBlue!55] (8,-5.22) rectangle (9,-4.78);
\fill[SynBlue!55] (12,-5.22) rectangle (13,-4.78);
\fill[SynBlue!55] (13,-5.22) rectangle (14,-4.78);
\fill[SynBlue!55] (14,-5.22) rectangle (15,-4.78);
\node[anchor=east,align=right,text width=4.3cm] at (-.2,-6) {1.6 Núcleo administrativo};
\fill[SynBlue!55] (1,-6.22) rectangle (2,-5.78);
\fill[SynBlue!55] (2,-6.22) rectangle (3,-5.78);
\fill[SynBlue!55] (7,-6.22) rectangle (8,-5.78);
\fill[SynBlue!55] (8,-6.22) rectangle (9,-5.78);
\fill[SynBlue!55] (12,-6.22) rectangle (13,-5.78);
\fill[SynBlue!55] (13,-6.22) rectangle (14,-5.78);
\fill[SynBlue!55] (14,-6.22) rectangle (15,-5.78);
\fill[SynBlue!55] (15,-6.22) rectangle (16,-5.78);
\node[anchor=east,align=right,text width=4.3cm] at (-.2,-7) {1.7 Canales e integración};
\fill[SynBlue!55] (3,-7.22) rectangle (4,-6.78);
\fill[SynBlue!55] (4,-7.22) rectangle (5,-6.78);
\fill[SynBlue!55] (5,-7.22) rectangle (6,-6.78);
\fill[SynBlue!55] (7,-7.22) rectangle (8,-6.78);
\fill[SynBlue!55] (8,-7.22) rectangle (9,-6.78);
\fill[SynBlue!55] (13,-7.22) rectangle (14,-6.78);
\node[anchor=east,align=right,text width=4.3cm] at (-.2,-8) {1.8 Migración e historia};
\fill[SynBlue!55] (0,-8.22) rectangle (1,-7.78);
\fill[SynBlue!55] (1,-8.22) rectangle (2,-7.78);
\fill[SynBlue!55] (2,-8.22) rectangle (3,-7.78);
\fill[SynBlue!55] (6,-8.22) rectangle (7,-7.78);
\fill[SynBlue!55] (8,-8.22) rectangle (9,-7.78);
\fill[SynBlue!55] (13,-8.22) rectangle (14,-7.78);
\node[anchor=east,align=right,text width=4.3cm] at (-.2,-9) {1.9 Innovaciones y pilotos};
\fill[SynBlue!55] (1,-9.22) rectangle (2,-8.78);
\fill[SynBlue!55] (2,-9.22) rectangle (3,-8.78);
\fill[SynBlue!55] (3,-9.22) rectangle (4,-8.78);
\fill[SynBlue!55] (12,-9.22) rectangle (13,-8.78);
\fill[SynBlue!55] (13,-9.22) rectangle (14,-8.78);
\fill[SynBlue!55] (14,-9.22) rectangle (15,-8.78);
\fill[SynBlue!55] (15,-9.22) rectangle (16,-8.78);
\fill[SynBlue!55] (16,-9.22) rectangle (17,-8.78);
\fill[SynBlue!55] (17,-9.22) rectangle (18,-8.78);
\fill[SynBlue!55] (18,-9.22) rectangle (19,-8.78);
\fill[SynBlue!55] (19,-9.22) rectangle (20,-8.78);
\fill[SynBlue!55] (20,-9.22) rectangle (21,-8.78);
\node[anchor=east,align=right,text width=4.3cm] at (-.2,-10) {1.10 Pruebas y correcciones};
\fill[SynBlue!55] (8,-10.22) rectangle (9,-9.78);
\fill[SynBlue!55] (9,-10.22) rectangle (10,-9.78);
\fill[SynBlue!55] (10,-10.22) rectangle (11,-9.78);
\fill[SynBlue!55] (11,-10.22) rectangle (12,-9.78);
\fill[SynBlue!55] (13,-10.22) rectangle (14,-9.78);
\fill[SynBlue!55] (14,-10.22) rectangle (15,-9.78);
\fill[SynBlue!55] (15,-10.22) rectangle (16,-9.78);
\fill[SynBlue!55] (16,-10.22) rectangle (17,-9.78);
\node[anchor=east,align=right,text width=4.3cm] at (-.2,-11) {1.11 Formación e implantación};
\fill[SynBlue!55] (10,-11.22) rectangle (11,-10.78);
\fill[SynBlue!55] (11,-11.22) rectangle (12,-10.78);
\fill[SynBlue!55] (12,-11.22) rectangle (13,-10.78);
\fill[SynBlue!55] (13,-11.22) rectangle (14,-10.78);
\fill[SynBlue!55] (14,-11.22) rectangle (15,-10.78);
\fill[SynBlue!55] (15,-11.22) rectangle (16,-10.78);
\fill[SynBlue!55] (16,-11.22) rectangle (17,-10.78);
\fill[SynBlue!55] (17,-11.22) rectangle (18,-10.78);
\fill[SynBlue!55] (18,-11.22) rectangle (19,-10.78);
\fill[SynBlue!55] (19,-11.22) rectangle (20,-10.78);
\fill[SynBlue!55] (20,-11.22) rectangle (21,-10.78);
\node[anchor=east,align=right,text width=4.3cm] at (-.2,-12) {1.12 Servicio y salida};
\fill[SynBlue!55] (17,-12.22) rectangle (18,-11.78);
\fill[SynBlue!55] (19,-12.22) rectangle (20,-11.78);
\fill[SynBlue!55] (20,-12.22) rectangle (21,-11.78);
\fill[SynBlue!55] (21,-12.22) rectangle (22,-11.78);
\fill[SynBlue!55] (22,-12.22) rectangle (23,-11.78);
\fill[SynBlue!55] (23,-12.22) rectangle (24,-11.78);
\fill[SynBlue!55] (24,-12.22) rectangle (25,-11.78);
\fill[SynBlue!55] (25,-12.22) rectangle (26,-11.78);
\fill[SynBlue!55] (26,-12.22) rectangle (27,-11.78);
\fill[SynBlue!55] (27,-12.22) rectangle (28,-11.78);
\node[anchor=east,text width=4.3cm,align=right] at (-.2,-13) {Marcha blanca E1};
\fill[SynTeal!65] (12,-13.22) rectangle (15,-12.78);
\node[anchor=east,text width=4.3cm,align=right] at (-.2,-14) {Marcha blanca E2};
\fill[SynTeal!65] (18,-14.22) rectangle (20,-13.78);
\node[anchor=east,text width=4.3cm,align=right] at (-.2,-15) {Operación ambos alcances};
\fill[SynNavyLight!65] (20,-15.22) rectangle (28,-14.78);
\node[anchor=east,text width=4.3cm,align=right] at (-.2,-16) {Cobertura inicial E1/E2};
\fill[SynNavyLight!40] (12,-16.22) rectangle (20,-15.78);
\node[rotate=90,anchor=west,text=SynNavy] at (1.5,-17.2) {H1};
\node[rotate=90,anchor=west,text=SynNavy] at (3.5,-17.2) {H2};
\node[rotate=90,anchor=west,text=SynNavy] at (5.5,-17.2) {H3};
\node[rotate=90,anchor=west,text=SynNavy] at (9.5,-17.2) {H4};
\node[rotate=90,anchor=west,text=SynNavy] at (11.5,-17.2) {H5};
\node[rotate=90,anchor=west,text=SynNavy] at (12.5,-17.2) {H6};
\node[rotate=90,anchor=west,text=SynNavy] at (15.5,-17.2) {H7};
\node[rotate=90,anchor=west,text=SynNavy] at (13.5,-17.2) {H8};
\node[rotate=90,anchor=west,text=SynNavy] at (16.5,-17.2) {H9};
\node[rotate=90,anchor=west,text=SynNavy] at (17.5,-17.2) {H10};
\node[rotate=90,anchor=west,text=SynNavy] at (18.5,-17.2) {H11};
\node[rotate=90,anchor=west,text=SynNavy] at (20.5,-17.2) {H12};
\end{tikzpicture}\caption{Carta Gantt agregada de los meses 1--28}\label{fig:sd7-gantt-1}
\FuenteFigura{Intervalos de red y asignaciones con esfuerzo UCP seleccionado; acta de marcha blanca E2 y corte requeridos al inicio de 21, bajo el supuesto adoptado de recepción continua y firma expresa programada.}
\end{figure}
Las filas permiten contrastar la continuidad de las cuentas con las ventanas contractuales. Los hitos identifican recepciones; la extensión de una barra no modifica las precedencias ni habilita trabajo durante una restricción operacional.
\end{landscape}\clearpage
\clearpage\begin{landscape}
La Figura~\ref{fig:sd7-gantt-2} presenta la programación agregada de los meses 29--56.
La carta agrega los intervalos calculados por cuenta; las fases contractuales se muestran en filas separadas. Las barras no autorizan cambios durante un congelamiento ni acreditan recepciones ejecutadas.
\begin{figure}[H]\centering\begin{tikzpicture}[x=0.52cm,y=0.55cm,font=\fontsize{9}{11}\selectfont]
\node at (0.5,.6) {29};\draw[SynLine] (0,.2)--(0,-16.5);
\node at (1.5,.6) {30};\draw[SynLine] (1,.2)--(1,-16.5);
\node at (2.5,.6) {31};\draw[SynLine] (2,.2)--(2,-16.5);
\node at (3.5,.6) {32};\draw[SynLine] (3,.2)--(3,-16.5);
\node at (4.5,.6) {33};\draw[SynLine] (4,.2)--(4,-16.5);
\node at (5.5,.6) {34};\draw[SynLine] (5,.2)--(5,-16.5);
\node at (6.5,.6) {35};\draw[SynLine] (6,.2)--(6,-16.5);
\node at (7.5,.6) {36};\draw[SynLine] (7,.2)--(7,-16.5);
\node at (8.5,.6) {37};\draw[SynLine] (8,.2)--(8,-16.5);
\node at (9.5,.6) {38};\draw[SynLine] (9,.2)--(9,-16.5);
\node at (10.5,.6) {39};\draw[SynLine] (10,.2)--(10,-16.5);
\node at (11.5,.6) {40};\draw[SynLine] (11,.2)--(11,-16.5);
\node at (12.5,.6) {41};\draw[SynLine] (12,.2)--(12,-16.5);
\node at (13.5,.6) {42};\draw[SynLine] (13,.2)--(13,-16.5);
\node at (14.5,.6) {43};\draw[SynLine] (14,.2)--(14,-16.5);
\node at (15.5,.6) {44};\draw[SynLine] (15,.2)--(15,-16.5);
\node at (16.5,.6) {45};\draw[SynLine] (16,.2)--(16,-16.5);
\node at (17.5,.6) {46};\draw[SynLine] (17,.2)--(17,-16.5);
\node at (18.5,.6) {47};\draw[SynLine] (18,.2)--(18,-16.5);
\node at (19.5,.6) {48};\draw[SynLine] (19,.2)--(19,-16.5);
\node at (20.5,.6) {49};\draw[SynLine] (20,.2)--(20,-16.5);
\node at (21.5,.6) {50};\draw[SynLine] (21,.2)--(21,-16.5);
\node at (22.5,.6) {51};\draw[SynLine] (22,.2)--(22,-16.5);
\node at (23.5,.6) {52};\draw[SynLine] (23,.2)--(23,-16.5);
\node at (24.5,.6) {53};\draw[SynLine] (24,.2)--(24,-16.5);
\node at (25.5,.6) {54};\draw[SynLine] (25,.2)--(25,-16.5);
\node at (26.5,.6) {55};\draw[SynLine] (26,.2)--(26,-16.5);
\node at (27.5,.6) {56};\draw[SynLine] (27,.2)--(27,-16.5);
\node[anchor=east,align=right,text width=4.3cm] at (-.2,-1) {1.1 Gobierno y contrato};
\node[anchor=east,align=right,text width=4.3cm] at (-.2,-2) {1.2 Plataforma y activos};
\node[anchor=east,align=right,text width=4.3cm] at (-.2,-3) {1.3 Seguridad y auditoría};
\node[anchor=east,align=right,text width=4.3cm] at (-.2,-4) {1.4 Núcleo operacional};
\node[anchor=east,align=right,text width=4.3cm] at (-.2,-5) {1.5 Núcleo de servicios};
\node[anchor=east,align=right,text width=4.3cm] at (-.2,-6) {1.6 Núcleo administrativo};
\node[anchor=east,align=right,text width=4.3cm] at (-.2,-7) {1.7 Canales e integración};
\node[anchor=east,align=right,text width=4.3cm] at (-.2,-8) {1.8 Migración e historia};
\node[anchor=east,align=right,text width=4.3cm] at (-.2,-9) {1.9 Innovaciones y pilotos};
\fill[SynBlue!55] (1,-9.22) rectangle (2,-8.78);
\fill[SynBlue!55] (2,-9.22) rectangle (3,-8.78);
\fill[SynBlue!55] (3,-9.22) rectangle (4,-8.78);
\fill[SynBlue!55] (4,-9.22) rectangle (5,-8.78);
\fill[SynBlue!55] (5,-9.22) rectangle (6,-8.78);
\fill[SynBlue!55] (6,-9.22) rectangle (7,-8.78);
\fill[SynBlue!55] (7,-9.22) rectangle (8,-8.78);
\fill[SynBlue!55] (8,-9.22) rectangle (9,-8.78);
\node[anchor=east,align=right,text width=4.3cm] at (-.2,-10) {1.10 Pruebas y correcciones};
\node[anchor=east,align=right,text width=4.3cm] at (-.2,-11) {1.11 Formación e implantación};
\node[anchor=east,align=right,text width=4.3cm] at (-.2,-12) {1.12 Servicio y salida};
\fill[SynBlue!55] (0,-12.22) rectangle (1,-11.78);
\fill[SynBlue!55] (1,-12.22) rectangle (2,-11.78);
\fill[SynBlue!55] (2,-12.22) rectangle (3,-11.78);
\fill[SynBlue!55] (3,-12.22) rectangle (4,-11.78);
\fill[SynBlue!55] (4,-12.22) rectangle (5,-11.78);
\fill[SynBlue!55] (5,-12.22) rectangle (6,-11.78);
\fill[SynBlue!55] (6,-12.22) rectangle (7,-11.78);
\fill[SynBlue!55] (7,-12.22) rectangle (8,-11.78);
\fill[SynBlue!55] (8,-12.22) rectangle (9,-11.78);
\fill[SynBlue!55] (9,-12.22) rectangle (10,-11.78);
\fill[SynBlue!55] (10,-12.22) rectangle (11,-11.78);
\fill[SynBlue!55] (11,-12.22) rectangle (12,-11.78);
\fill[SynBlue!55] (12,-12.22) rectangle (13,-11.78);
\fill[SynBlue!55] (13,-12.22) rectangle (14,-11.78);
\fill[SynBlue!55] (14,-12.22) rectangle (15,-11.78);
\fill[SynBlue!55] (15,-12.22) rectangle (16,-11.78);
\fill[SynBlue!55] (16,-12.22) rectangle (17,-11.78);
\fill[SynBlue!55] (17,-12.22) rectangle (18,-11.78);
\fill[SynBlue!55] (18,-12.22) rectangle (19,-11.78);
\fill[SynBlue!55] (19,-12.22) rectangle (20,-11.78);
\fill[SynBlue!55] (20,-12.22) rectangle (21,-11.78);
\fill[SynBlue!55] (21,-12.22) rectangle (22,-11.78);
\fill[SynBlue!55] (22,-12.22) rectangle (23,-11.78);
\fill[SynBlue!55] (23,-12.22) rectangle (24,-11.78);
\fill[SynBlue!55] (24,-12.22) rectangle (25,-11.78);
\fill[SynBlue!55] (25,-12.22) rectangle (26,-11.78);
\fill[SynBlue!55] (26,-12.22) rectangle (27,-11.78);
\fill[SynBlue!55] (27,-12.22) rectangle (28,-11.78);
\node[anchor=east,text width=4.3cm,align=right] at (-.2,-13) {Marcha blanca E1};
\node[anchor=east,text width=4.3cm,align=right] at (-.2,-14) {Marcha blanca E2};
\node[anchor=east,text width=4.3cm,align=right] at (-.2,-15) {Operación ambos alcances};
\fill[SynNavyLight!65] (0,-15.22) rectangle (28,-14.78);
\node[anchor=east,text width=4.3cm,align=right] at (-.2,-16) {Cobertura inicial E1/E2};
\end{tikzpicture}\caption{Carta Gantt agregada de los meses 29--56}\label{fig:sd7-gantt-2}
\FuenteFigura{Intervalos de red y asignaciones con esfuerzo UCP seleccionado; acta de marcha blanca E2 y corte requeridos al inicio de 21, bajo el supuesto adoptado de recepción continua y firma expresa programada.}
\end{figure}
Las filas permiten contrastar la continuidad de las cuentas con las ventanas contractuales. Los hitos identifican recepciones; la extensión de una barra no modifica las precedencias ni habilita trabajo durante una restricción operacional.
\end{landscape}\clearpage



\CierreSubdocumento
\end{document}
